\documentclass[10pt]{ltjbook}
\usepackage{luatexja}
\usepackage[left=1cm,right=1cm,bottom=1.5cm,top=1.5cm]{geometry}
\usepackage[table]{xcolor}
\usepackage{amsmath}
\usepackage{ruby}
\usepackage{tikz,pgfplots,xfrac,makecell,aurical,array,booktabs,multicol,multirow,calc}
\usepackage{environ}
\usepackage{tabularx}
\usepackage{hhline}
\usepackage{textcomp}
\usepackage{wasysym}
\usepackage{multicol}
\usepackage{microtype}
\usetikzlibrary{patterns}

\newcolumntype{A}{>{\centering\let\newline\\\arraybackslash}X}
\newcolumntype{B}{>{\raggedright\let\newline\\\arraybackslash}X}

\newcommand{\specialc}[2][c]{\begin{tabular}[#1]{@{}c@{}}#2\end{tabular}}
\newcommand{\speciall}[2][c]{\begin{tabular}[#1]{@{}l@{}}#2\end{tabular}}
\newcommand{\specialr}[2][c]{\begin{tabular}[#1]{@{}r@{}}#2\end{tabular}}


\usepackage{graphicx}
\graphicspath{
    {/home/henri/stuff/japanese/kanji/chap1/}{/home/henri/stuff/japanese/kanji/chap2/}{/home/henri/stuff/japanese/kanji/chap3/}{/home/henri/stuff/japanese/kanji/chap4/}{/home/henri/stuff/japanese/kanji/chap5/}{/home/henri/stuff/japanese/kanji/chap6/}{/home/henri/stuff/japanese/kanji/chap7/}{/home/henri/stuff/japanese/kanji/chap8/}{/home/henri/stuff/japanese/kanji/chap9/}{/home/henri/stuff/japanese/kanji/chap10/}{/home/henri/stuff/japanese/kanji/chap11/}{/home/henri/stuff/japanese/kanji/chap12/} }

\pgfplotsset{compat=newest,compat/show suggested version=false}
\usetikzlibrary{pgfplots.groupplots,matrix,positioning,fit,calc,shapes,patterns,automata,plotmarks}

\setlength{\parindent}{0cm}

\renewcommand{\tabularxcolumn}[1]{>{\small}m{#1}}

\newcommand{\rye}[2]{\multicolumn{#1}{|l|}{\begin{tabularx}{0.95\columnwidth}{@{}B@{}}#2\end{tabularx}}\\\hline}

\usepackage{tocloft}
\newcommand{\listofkanji}{List of Kanji}
\newlistof{glyph}{lok}{\listofkanji}
\NewEnviron{glyph}[2]% environment name
{	\refstepcounter{glyph}%
	\addcontentsline{lok}{glyph}{\protect\numberline{\theglyph}{\;}{\;}#1}
	\begin{tabularx}{\columnwidth}{@{}|@{}c@{}|@{}c@{}|B@{}|@{}}
	\multicolumn{3}{@{}|l|@{}}{\cellcolor{black}{\color{white}\textbf{\large{\theglyph{:}\,{#1}}}}}\\\hline
	\raisebox{-0.80\totalheight}{\includegraphics[width=2cm,height=2cm]{#2.pdf}} & \raisebox{-0.80\totalheight}{\includegraphics[width=2cm,height=2cm]{so_#2.pdf}} & \textbf{\underline{Notes:}} \\\hline
	\multicolumn{3}{@{}|l|@{}}{\begin{tabularx}{0.95\columnwidth}{@{}B@{}}
	\BODY
	\end{tabularx}}\\\hline
	\end{tabularx}
}

\NewEnviron{comps}
{
%\keepXColumns
\begin{tabularx}{\columnwidth}{@{}|l|A|@{}}
\hhline{~|-|}\multicolumn{1}{c}{} & \multicolumn{1}{|c|}{\cellcolor{lightgray}$\in$ \textbf{Composition}} \\
\hline
\BODY
\end{tabularx}
}

\NewEnviron{somme}%
{%
%\keepXColumns
\begin{tabularx}{\columnwidth}{@{}|c|B|A|@{}}%
\hline%
\multicolumn{3}{|c|}{\cellcolor{lightgray}\bfseries{+ Words and Compounds}}\\\hline%
\multicolumn{1}{|l|}{\cellcolor{lightgray}Kanji} & \cellcolor{lightgray}English & \cellcolor{lightgray}Pronounciation\\\hline
\BODY
\end{tabularx}%
}

\NewEnviron{prons}%
{%
%\keepXColumns
\begin{tabularx}{\columnwidth}{@{}|l|B|B|@{}}%
\hhline{~|--|}\multicolumn{1}{c}{} & \multicolumn{2}{|c|}{\cellcolor{lightgray}\bfseries{{\textmusicalnote} Pronounciations}}\\
\hhline{~|-|-|}\multicolumn{1}{c}{} & \multicolumn{1}{|l|}{\cellcolor{lightgray}ONYOMI} & \multicolumn{1}{l|}{\cellcolor{lightgray}kunyomi}\\\hline
\BODY
\end{tabularx}%
}

\NewEnviron{remglyph}
{
%\keepXColumns
\begin{tabularx}{\columnwidth}{@{}|B|@{}}%
\hhline{|-|}\multicolumn{1}{|c|}{\cellcolor{lightgray}$\bell\in$ \textbf{Meaning $\leftrightarrow$ Glyph Components}}\\%
\hhline{|-|}\BODY%
\end{tabularx}%
}

\NewEnviron{remprons}
{
%\keepXColumns
\begin{tabularx}{\columnwidth}{@{}|B|@{}}
\hhline{|-|}\multicolumn{1}{|c|}{\cellcolor{lightgray}\textbf{{\bell\textmusicalnote} Meaning $\leftrightarrow$ Pronounciations}}\\
\hhline{|-|}\BODY
\end{tabularx}%
}

\NewEnviron{irrr}%
{%
\begin{tabularx}{\columnwidth}{@{}|c|B|A|@{}}%
\hline%
\multicolumn{3}{|c|}{\cellcolor{lightgray}\bfseries{++ Irregular Reading(s)}}\\\hline%
\cellcolor{lightgray}Kanji & \cellcolor{lightgray}English & \cellcolor{lightgray}Pronounciation\\\hline
\BODY
\end{tabularx}%
}

\NewEnviron{sentence}
{
%\keepXColumns
\begin{tabularx}{\columnwidth}{|B|}
\hhline{|-|}\multicolumn{1}{|c|}{\cellcolor{lightgray}\textbf{$\lbrace\cdot\cdot\cdot\rbrace$ Sample Sentence}}\\
\hhline{|-|}\rule{0pt}{\baselineskip}
\BODY
\end{tabularx}
}

\newcommand{\middel}{\hspace{-0.9pt}\ensuremath{\cdot}\hspace{-0.9pt}}
\newcommand{\diff}{\textsuperscript{*}\hspace{-0.7pt}}
\newcommand{\ekwiv}{\ensuremath{\equiv}}

\newcommand{\hooglig}[1]{\colorbox{lightgray!60!white}{#1}}

\newcommand{\comon}[1]{\uppercase{\textbf{$\prec$#1$\succ$}}}
\newcommand{\comkun}[1]{\textbf{$\prec$#1$\succ$}}
\newcommand{\ucomon}[1]{\uppercase{$\prec$#1$\succ$}}
\newcommand{\ucomkun}[1]{$\prec$#1$\succ$}

\newcommand{\remcom}[1]{$\prec$#1$\succ$}

\newcommand{\comprtwocone}{\cellcolor{lightgray}Glyphs}
\newcommand{\comprthreecone}{\cellcolor{lightgray}English (\#)}
\newcommand{\pronsrtwocone}{\cellcolor{lightgray}\textbf{Cm.}}
\newcommand{\pronsrthreecone}{\cellcolor{lightgray}L.C.}

\setlength{\parskip}{4pt}

\makeatletter
   \def\vhrulefill#1{\leavevmode\leaders\hrule\@height#1\hfill \kern\z@}
\makeatother
%-------------------------------------------------------
%-------------------------------------------------------
%-------------------------------------------------------
\begin{document}
\listofglyph
\clearpage
\phantomsection

%\begin{multicols}{2}
%\chapter{Kanji 1--20}
%\vspace*{\fill}
\begin{glyph}{Mountain (山)}{mountain}\label{glyph:mountain}
Mountain.\\\hline
Obvious from the glyph at first glance.
\end{glyph}

\begin{comps}
\comprtwocone & 山 \\\hline
\comprthreecone & Mountain \\\hline
\end{comps}

\begin{remglyph}
Part of the Kanji radicals (\#{46})\\\hline
\end{remglyph}

\begin{prons}
\pronsrtwocone & \textbf{サン (SAN)} & \textbf{やま (yama)} \\\hline
\pronsrthreecone & --- & --- \\\hline
\end{prons}

\begin{remprons}
On the \hooglig{mountain}, my \comon{son} played a \comkun{yamaha} piano.\\\hline
\end{remprons}

\begin{somme}
    山 & \textit{mountain} & やま \mbox{(yama)} \\\hline
    \notyet{一山} & one + mountain = \textit{a pile of something} &  ひと{\middel}やま \mbox{(hito{\middel}yama)}\\\hline
    \notyet{高山} & tall + mountain = \textit{alpine} & コウ{\middel}ザン \mbox{(K\={O}{\middel}ZAN)}\\\hline
    \notyet{火山} & fire + mountain = \textit{volcano} & カ{\middel}ザン \mbox{(KA{\middel}ZAN)}\\\hline
    \notyet{山村} & mountain + village = \textit{mountain village} & サン{\middel}ソン \mbox{(SAN{\middel}SON)}\\\hline
\end{somme}

\begin{sentence}
\ruby{\underline{山}}{やま}に\ruby{\underline{木}}{き}が\underline{ありません}。\\\hline
yama ni ki ga arimasen.\\\hline
(mountain) (tree) (are not)\\\hline
=\textit{There are no trees on the mountain.}\\\hline
\end{sentence}
\end{multicols}
\vhrulefill{5pt}
%--------------------------------------------------------------------
%--------------------------------------------------------------------
\begin{multicols}{2}
\begin{glyph}{Person (人)}{person}\label{glyph:person}
Person; Human being. \\\hline
When appearing at the end of country names, this kanji denotes an individual's nationality.
\end{glyph}

\begin{comps}
\comprtwocone & 人 \\\hline
\comprthreecone & Person \\\hline
\end{comps}

\begin{remglyph}
Part of the Kanji radicals (\#{9}).\\\hline
\end{remglyph}

\begin{prons}
\pronsrtwocone & \textbf{ジン、ニン (JIN, NIN)} & \textbf{ひと (hito)} \\\hline
\rye{3}{ジン tends to signify that a person belongs to a certain subgroup of
	humanity.  Look for ジン in the first position and at the end of words signifying countries.\\\hline
	ニン tends to indicate a person engaged in an activity specified by the kanji preceding it.\\\hline
	The 訓読み becomes always voiced when it appears outside of the first position.}
\pronsrthreecone & --- & --- \\\hline
\end{prons}

\begin{remprons}
That \hooglig{person} is a \comon{nincompoop} because he \comon{jinxes} with
    \comkun{hit-or-miss} attempts.\\\hline
\end{remprons}

\begin{somme}
    人 & \textit{person} & ひと \mbox{(hito)} \\\hline
    \notyet{白人} & white + person = \textit{caucasian} & ハク{\middel}ジン \mbox{(HAKU{\middel}JIN)} \\\hline
    \notyet{人口} & person + mouth = \textit{population} & ジン{\middel}コウ \mbox{(JIN{\middel}K\={O})} \\\hline
    \notyet{三人} & three + person = \textit{three people} & サン{\middel}ニン \mbox{(SAN{\middel{NIN}})}\\\hline
    \notyet{四人} & four + person = \textit{four people} & よ{\middel}{ニン} \mbox{(yo{\middel}{NIN})}\\\hline
    \notyet{村人} & village + person = \textit{villager} & むら{\middel}{びと} \mbox{(mura{\middel}bito)}\\\hline
    \notyet{外国人} & outside + country + person = \textit{foreigner} & ガイ{\middel}コク{\middel}ジン \mbox{(GAI{\middel}KOKU{\middel}JIN)}\\\hline
\end{somme}

\begin{irrr}
\rye{3}{The irregular readings in the first two examples belong to 人.  In the final compound, both readings are irregular.}
    \notyet{一人} & one + person = \textit{one person} & ひと{\middel}り \mbox{(hito{\middel}ri)}\\\hline
    \notyet{二人} & two + person = \textit{two people} & ふた{\middel}り \mbox{(futa{\middel}ri)}\\\hline
    \notyet{大人} & large + person = \textit{adult} & おとな \mbox{(otona)}\\\hline
\end{irrr}

\begin{sentence}
\underline{あの} \ruby{\underline{人}}{ひと} の\ruby{\underline{手}}{て}\ruby{は}{わ}\ruby{\underline{小}}{ちい}\underline{さい}です。\\\hline
ano hito no te wa chii{\middel}sai desu.\\\hline
(that) (person) (hand) (small).\\\hline
=\textit{That person's hands are small.}\\\hline
\end{sentence}
\end{multicols}
\vhrulefill{5pt}
%--------------------------------------------------------------------
%--------------------------------------------------------------------
\begin{multicols}{2}
\begin{glyph}{One (一)}{one}\label{glyph:one}
One.
\end{glyph}

\begin{comps}
\comprtwocone & 一 \\\hline
\comprthreecone & One \\\hline
\end{comps}

\begin{remglyph}
Part of the Kanji radicals (\#{1}).\\\hline
\end{remglyph}

\begin{prons}
\pronsrtwocone & \textbf{イチ、イツ (ICHI, ITSU)} & \textbf{ひと (hito)} \\\hline
\pronsrthreecone & --- & --- \\\hline
\rye{3}{When this kanji appears as the first character in a compound it is almost always read with its on-yomi.  \\\hline
    Its pronounciation イチ or イツ will depend on whether the initial sound of the following kanji is voiced or voiceless.}
\end{prons}

\begin{remprons}
\hooglig{One} \comkun{hit-or-miss} attempt to \comon{eat soup} that reduces \comon{itchiness}.\\\hline
\end{remprons}

\begin{somme}
    一 & \textit{one} & イチ \mbox{(ICHI)}\\\hline
    一つ & \textit{one (general counter)} & ひと{\middel}つ \mbox{(hito{\middel}tsu)}\\\hline
    \notyet{一日} & one + sun (day) = \textit{one day} & イチ{\middel}ニチ \mbox{(ICHI{\middel}NICHI)}\\\hline
    \notyet{一月} & one + moon (month) = \textit{January} & イチ{\middel}ガツ \mbox{(ICHI{\middel}GATSU)}\\\hline
    \notyet{同一} & same + one = \textit{identical} & ドウ{\middel}イツ \mbox{(D\={O}{\middel}ITSU)}\\\hline
    \notyet{一回} & one + rotate = \textit{once} & イッ{\middel}カイ \mbox{(IK{\middel}KAI)}\\\hline
    \notyet{一時} & one + time = \textit{one o'clock} & イチ{\middel}ジ \mbox{(ICHI{\middel}JI)}\\\hline
\end{somme}

\begin{irrr}
    \notyet{一日} & one + sun (day) = \textit{first day of the month} & ついたち (tsuitachi)\\\hline
\rye{3}{The same compound is shown above with a different pronunciation and meaning.  Context will determine which is the appropriate reading to apply.  For example, when 一日 is seen at the top of a newspaper page, it clearly refers to the first day of the month.}
\end{irrr}

\begin{sentence}
\ruby{\underline{一時}}{イチジ}に\ruby{\underline{会}}{あ}\underline{いましょう}。\\\hline
ICHI{\middel}JI ni a{\middel}imash\={o}.\\\hline
(one o'clock) (let's meet)\\\hline
=\textit{Let's meet at one o'clock.}\\\hline
\end{sentence}
\end{multicols}
\vspace*{\fill}
%%--------------------------------------------------------------------
%%--------------------------------------------------------------------
\pagebreak
\begin{multicols}{2}
\begin{glyph}{Two (二)}{two}\label{glyph:two}
Two.  Including the ideas of ``double'' and ``bi-'', etc.
\end{glyph}

\begin{comps}
\comprtwocone & 一 + 一 \\\hline
\comprthreecone & One + One \\\hline
\end{comps}

\begin{remglyph}
Part of the Kanji radicals (\#{7}).\\\hline
\end{remglyph}

\begin{prons}
\pronsrtwocone & \textbf{ニ (NI)} & \textbf{ふた (futa)} \\\hline
\pronsrthreecone & --- & --- \\\hline
\end{prons}

\begin{remprons}
    \hooglig{Two} \comon{nickels} to tell you \comkun{who tans}?\\\hline
\end{remprons}

\begin{somme}
    ニ & \textit{two} & ニ \mbox{(NI)}\\\hline
    二つ & \textit{two (general counter)} & ふた{\middel}つ \mbox{(futa{\middel}tsu)}\\\hline
    \notyet{二月} & two + moon (month) = \textit{February} & ニガ{\middel}ツ \mbox{(NI{\middel}GATSU)}\\\hline
    \notyet{二十} & two + ten = \textit{twenty} & ニ{\middel}ジュウ \mbox{(NI{\middel}J\={U})}\\\hline
    \notyet{二百} & two + hundred = \textit{two hundred} & ニ{\middel}ヒャク \mbox{(NI{\middel}HYAKU)}\\\hline
    \notyet{二時} & two + time = \textit{two o'clock} & ニ{\middel}ジ \mbox{(NI{\middel}JI)}\\\hline
    \notyet{二週間} & two + week + interval = \textit{two weeks} & ニ{\middel}シュウ{\middel}カン \mbox{(NI{\middel}SH\={U}{\middel}KAN)}\\\hline
\end{somme}

\begin{irrr}
    \notyet{二日} & two + sun (day) = \textit{second day of the month} & ふつか \mbox{(futsuka)}\\\hline
    \notyet{二十日} & two + ten + sun (day) = \textit{twentieth day of the month} & はつか \mbox{(hatsuka)}\\\hline
    \notyet{二十歳} & two + ten + annual = \textit{twenty years old} & はたち \mbox{(hatachi)}\\\hline
\end{irrr}

\begin{sentence}
\ruby{\underline{二月}}{ニガツ}の\underline{オーストラリア}\ruby{は}{わ}\ruby{\underline{美}}{うつく}\underline{しい}。\\\hline
NI{\middel}GATSU no \={o}sutoraria wa utsuku{\middel}shii.\\\hline
(February) (Australia) (beautiful)\\\hline
=\textit{Australia is beautiful in February.}\\\hline
\end{sentence}
\end{multicols}
\vhrulefill{5pt}
%%--------------------------------------------------------------------
%%--------------------------------------------------------------------
\begin{multicols}{2}
\begin{glyph}{Three (三)}{three}\label{glyph:three}
Three/Triple/Tri-, etc.
\end{glyph}

\begin{comps}
\comprtwocone & 一 + 一 + 一 \\\hline
\comprthreecone & One + One + One \\\hline
\end{comps}

\begin{remglyph}
---\\\hline
\end{remglyph}

\begin{prons}
\pronsrtwocone & \textbf{サン (SAN)} & --- \\\hline
\pronsrthreecone & --- & みっ、み (mi-, mi) \\\hline
\end{prons}

\begin{remprons}
My \hooglig{triplet} \comon{sons} \ucomkun{mirror} one another.\\\hline
\end{remprons}

\begin{somme}
    三 & \textit{three} & サン \mbox{(SAN)}\\\hline
    三つ & \textit{three (general counter)} & みっ{\middel}つ \mbox{(mit{\middel}tsu)}\\\hline
    \notyet{三人} & three + person = \textit{three people} & サン{\middel}ニン \mbox{(SAN{\middel}NIN)}\\\hline
    \notyet{三月} & three + moon (month) = \textit{March} & サン{\middel}ガツ \mbox{(SAN{\middel}GATSU)}\\\hline
    \notyet{三時} & three + time = \textit{three o'clock} & サン{\middel}ジ \mbox{(SAN{\middel}JI)}\\\hline
    \notyet{三十} & three + ten = \textit{thirty} & サン{\middel}ジュウ \mbox{(SAN{\middel}J\={U})}\\\hline
    \notyet{三千} & three + thousand = \textit{three thousand} & サン{\middel}ゼン \mbox{(SAN{\middel}ZEN)}\\\hline
    \notyet{三日} & three + sun (day) = \textit{third day of the month} & みっ{\middel}か \mbox{(mik{\middel}ka)}\\\hline
\end{somme}

\begin{sentence}
\ruby{\underline{三人}}{サンニン}が\ruby{\underline{早}}{はや}\underline{く} \ruby{\underline{来}}{き}\underline{ました}。\\\hline
SAN{\middel}NIN ga haya{\middel}ku ki{\middel}mashita.\\\hline
(three people) (early) (came).\\\hline
=\textit{Three people came early.}\\\hline
\end{sentence}
\end{multicols}
%%--------------------------------------------------------------------
%%--------------------------------------------------------------------
\pagebreak
\begin{multicols}{2}
\begin{glyph}{Sun (日)}{sun}\label{glyph:sun}
Sun, and Day.\\\hline
It is the initial kanji in the compound for Japan.
\end{glyph}

\begin{comps}
\comprtwocone & 日 \\\hline
\comprthreecone & Sun \\\hline
\end{comps}

\begin{remglyph}
Part of the Kanji radicals (\#{72})\\\hline
\end{remglyph}

\begin{prons}
\pronsrtwocone & \textbf{ニチ、ジツ (NICHI, JITSU)} & \textbf{ひ (hi)} \\\hline
\pronsrthreecone & --- & か (ka) \\\hline
\rye{3}{ニチ will be encountered overwhelmingly in the first position (ひ only makes a few appearances).\\\hline
ジツ is typically encountered in the second position.\\\hline
ひ is typically encountered in the third position.\\\hline
か is used only to identify the days of the month from 2--10, 14, 17, 20, 24, and 27.}
\end{prons}

\begin{remprons}
The \hooglig{sun} inspired \comon{Nichiren} \comon{Ju-Jitsu}; many other
    \ucomkun{customs} are \comkun{heaped} under the \hooglig{sun}.\\\hline
\end{remprons}

\begin{somme}
    日 & \textit{sun; day} & ひ \mbox{(hi)}\\\hline
    一日 & one + day = \textit{one day} & イチ{\middel}ニチ \mbox{(ICHI{\middel}NICHI)} \\\hline
    \notyet{日本} & sun + main = \textit{Japan} & ニ{\middel}ホン、ニッ{\middel}ポン \mbox{(NI{\middel}HON, NIP{\middel}PON)}\\\hline
    \notyet{休日} & rest + day = \textit{holiday} & キュウ{\middel}ジツ \mbox{(KY\={U}{\middel}JITSU)}\\\hline
    \notyet{毎日} & every + day = \textit{daily} & マイ{\middel}ニチ \mbox{(MAI{\middel}NICHI)}\\\hline
    \notyet{朝日} & morning + sun = \textit{morning sun} & あさ{\middel}ひ \mbox{(asa{\middel}hi)}\\\hline
    \notyet{日曜日} & sun + day of the week + day = \textit{Sunday} & ニチ{\middel}ヨウ{\middel}ビ \mbox{(NICHI{\middel}Y\={O}{\middel}bi)}\\\hline
\end{somme}

\begin{irrr}
    一日 & one + day = \textit{first day of the month} & ついたち \mbox{(tsuitachi)}\\\hline
    \notyet{昨日} & past + day = \textit{yesterday} & きのう \mbox{(kin\={o})}\\\hline
    \notyet{今日} & now + day = \textit{today} & きょう \mbox{(ky\={o})}\\\hline
    \notyet{明日} & bright + day = \textit{tomorrow} & あした、あす \mbox{(ashita, asu)}\\\hline
\end{irrr}

\begin{sentence}
\ruby{\underline{毎日}}{マイニチ}{\;}\ruby{\underline{日本語}}{ニホンゴ}の\ruby{\underline{本}}{ホン}\ruby{を}{お}\ruby{\underline{読}}{よ}\underline{みます}。\\\hline
MAI{\middel}NICHI NI{\middel}HON{\middel}GO no HON o yo{\middel}mimasu.\\\hline
(daily) (Japanese) (book) (read).\\\hline
=\textit{(I) read Japanese books everyday.}\\\hline
\end{sentence}
\end{multicols}
\vhrulefill{5pt}
%%--------------------------------------------------------------------
%%--------------------------------------------------------------------
\begin{multicols}{2}
\begin{glyph}{White (白)}{white}\label{glyph:white}
White.  As well as shades of cleanliness and purity.
\end{glyph}

\begin{comps}
\comprtwocone & 白 \\\hline
\comprthreecone & White \\\hline
\end{comps}

\begin{remglyph}
    Part of the Kanji radicals (\#{106}).\\\hline
\end{remglyph}

\begin{prons}
\pronsrtwocone & \textbf{ハク (HAKU)} & \textbf{しろ (shiro)} \\\hline
\pronsrthreecone & ビャク (BYAKU) & しら (shira) \\\hline
\end{prons}

\begin{remprons}
    In \hooglig{white} robes, the \comkun{sheerly obstinate} \ucomon{Byakuya}
    challenged \comon{Haku}; it was \ucomkun{sheer ugliness}.\\\hline
\end{remprons}

\begin{somme}
    白 & \textit{white (noun)} & しろ \mbox{(shiro)}\\\hline
    白い & \textit{white (adjective)} & しろ{\middel}い \mbox{(shiro{\middel}i)}\\\hline
    白人 & white + person = \textit{caucasian} & ハク{\middel}ジン \mbox{(HAKU{\middel}JIN)}\\\hline
    \notyet{明白} & bright + white = \textit{obvious} & メイ{\middel}ハク \mbox{(MEI{\middel}HAKU)}\\\hline
    \notyet{白鳥} & white + bird = \textit{swan} & ハク{\middel}チョウ \mbox{(HAKU{\middel}CH\={O})}\\\hline
    \notyet{白米} & white + rice = \textit{polished rice} & ハク{\middel}マイ \mbox{(HAKU{\middel}MAI)}\\\hline
\end{somme}

\begin{sentence}
\underline{この} \ruby{\underline{白}}{しろ}\underline{い} \ruby{\underline{馬}}{うま}の\ruby{\underline{名前}}{なまえ}\ruby{は}{わ}\ruby{\underline{``雪''}}{ゆき}です。\\\hline
kono shiro{\middel}i uma no na{\middel}mae wa yuki desu.\\\hline
(this) (white) (horse) (name) (``snow'').\\\hline
=\textit{This white horse's name is ``Snow''}\\\hline
\end{sentence}
\end{multicols}
%%--------------------------------------------------------------------
%%--------------------------------------------------------------------
\pagebreak
\vspace*{\fill}
\begin{multicols}{2}
\begin{glyph}{Mouth (口)}{mouth}\label{glyph:mouth}
Mouth, and things oral.  This kanji also relates to openings in general, from cavers and harbors to taps and bottles, for instance.
\end{glyph}

\begin{comps}
\comprtwocone & 口 \\\hline
\comprthreecone & Mouth \\\hline
\end{comps}

\begin{remglyph}
Part of the Kanji radicals (\#{30}).\\\hline
\end{remglyph}

\begin{prons}
\pronsrtwocone & \textbf{コウ (K\={O})} & \textbf{くち (kuchi)} \\\hline
\pronsrthreecone & ク (KU) & --- \\\hline
\end{prons}

\begin{remprons}
His \hooglig{mouth} \comon{caused} the \ucomon{cook} to \comkun{coochi-coo} at
    him.\\\hline
\end{remprons}

\begin{somme}
    口 & \textit{mouth} & くち \mbox{(kuchi)}\\\hline
    人口 & person + mouth = \textit{population} & ジン{\middel}コウ \mbox{(JIN{\middel}K\={O})}\\\hline
    \notyet{早口} & early (fast) + mouth = \textit{rapid speaking} & はや{\middel}くち \mbox{(haya{\middel}kuchi)}\\\hline
    \notyet{口語} & mouth + words = \textit{colloquial language} & コウ{\middel}ゴ \mbox{(K\={O}{\middel}GO)}\\\hline
    \notyet{出口} & exit + mouth = \textit{exit} & で{\middel}ぐち \mbox{(de{\middel}guchi)}\\\hline
    \notyet{口先} & mouth + precede = \textit{lip service} & くち{\middel}さき \mbox{(kuchi{\middel}saki)}\\\hline
\end{somme}

\begin{sentence}
\underline{この} \ruby{\underline{国}}{くに}の\ruby{\underline{人口}}{ジンコウ}\ruby{は}{わ}\ruby{\underline{少}}{すく}\underline{ない}。\\\hline
kono kuni no JIN{\middel}K\={O} wa suku{\middel}nai.\\\hline
(this) (country) (population) (few).\\\hline
=\textit{This country has a small population.}\\\hline
\end{sentence}
\end{multicols}
\vhrulefill{5pt}
%%--------------------------------------------------------------------
%%--------------------------------------------------------------------
\begin{multicols}{2}
\begin{glyph}{Rotate (回)}{rotate}\label{glyph:rotate}
The general sense of this kanji relates to the ideas of rotation and things going around in both space and time.\\\hline
It can also refer to the vicinity of things (neighbourhoods, surroundings, etc.).\\\hline
This kanji is well-known to Japanese baseball fans as the character used to denote innings.
\end{glyph}

\begin{comps}
\comprtwocone & 囗 + 口 \\\hline
\comprthreecone & Enclosure + Mouth \\\hline
\end{comps}

\begin{remglyph}
\hooglig{Rotate} \remcom{dracula} in the \remcom{enclosure}.\\\hline
\end{remglyph}

\begin{prons}
\pronsrtwocone & \textbf{カイ (KAI)} & \textbf{まわ (mawa)} \\\hline
\pronsrthreecone & --- & --- \\\hline
\end{prons}

\begin{remprons}
The \hooglig{rotating} \comon{kites} that \comkun{ma wacks}.\\\hline
\end{remprons}

\begin{somme}
\rye{3}{Note the difference between intransitive (intr.) and transitive (tr.) verbs.  An intransitive verb does not require an object to complete the action, whereas a transitive verb does.}
    回る (intr.) & \textit{to go around} & まわ{\middel}る \mbox{(mawa{\middel}ru)}\\\hline
    回す (tr.) & \textit{to send around} & まわ{\middel}す \mbox{(mawa{\middel}su)}\\\hline
    一回 & one + rotate = \textit{once} & イッ{\middel}カイ \mbox{(IK{\middel}KAI)}\\\hline
    二回 & two + rotate = \textit{twice} & ニ{\middel}カイ \mbox{(NI{\middel}KAI)}\\\hline
    \notyet{引き回す} & pull + rotate = \textit{to pull around} & ひ{\middel}き　まわ{\middel}す \mbox{(hi{\middel}ki mawa{\middel}su)}\\\hline
    \notyet{言い回し} & say + rotate = \textit{turn of expression} & い{\middel}い　まわ{\middel}し \mbox{(i{\middel}i mawa{\middel}shi)}\\\hline
\end{somme}

\begin{sentence}
\ruby{\underline{私}}{わたし}\ruby{は}{わ}\ruby{\underline{二回}}{ニカイ}{\;}\underline{カナダ}\ruby{へ}{え}\ruby{\underline{行}}{い}\underline{きました}。\\\hline
watashi wa NI{\middel}KAI Canada e i{\middel}kimashita.\\\hline
(I) (twice) (Canada) (went)\\\hline
I went to Canada twice.\\\hline
\end{sentence}
\end{multicols}
\vspace*{\fill}
%%--------------------------------------------------------------------
%%--------------------------------------------------------------------
\pagebreak
\vspace*{\fill}
\begin{multicols}{2}
\begin{glyph}{Four (四)}{four}\label{glyph:four}
Four, Quad-, etc.
\end{glyph}

\begin{comps}
\comprtwocone & 囗 + 儿 \\\hline
    \comprthreecone & Enclosure + Person's legs
    (\ref{glyph:person}) \\\hline
\end{comps}

\begin{remglyph}
\hooglig{Four} \remcom{people} danced to the \hooglig{four} corners of the
    \remcom{prison cell}.\\\hline
\end{remglyph}

\begin{prons}
\pronsrtwocone & \textbf{シ (SHI)} & \textbf{よん、よ (yon, yo)} \\\hline
\pronsrthreecone & --- & よっ (yo-) \\\hline
\end{prons}

\begin{remprons}
\hooglig{Four} \comon{shields} \comkun{yomped} \comkun{yonder} to the
    \hooglig{four} corners of the world.\\\hline
\end{remprons}

\begin{somme}
    四 & \textit{four} & シ \mbox{(SHI)} \\\hline
    四人 & four + person = \textit{four people} & よ{\middel}ニン \mbox{(yo{\middel}NIN)} \\\hline
    四日 & four + day = \textit{the fourth day of the month} & よっか \mbox{(yok{\middel}ka)} \\\hline
    四つ & \textit{four (general counter)} & よっつ \mbox{(yot{\middel}tsu)} \\\hline
    \notyet{四月} & four + moon (month) = \textit{April} & シ{\middel}ガツ \mbox{(SHI{\middel}GATSU)} \\\hline
    \notyet{四時} & four + time = \textit{four o'clock} & よ{\middel}ジ \mbox{(yo{\middel}JI)} \\\hline
    \notyet{四十} & four + ten = \textit{forty} & よん{\middel}ジュウ \mbox{(yon{\middel}J\={U})} \\\hline
    \notyet{四百} & four + hundred = \textit{four hundred} & よん{\middel}ヒャク \mbox{(yon{\middel}HYAKU)} \\\hline
\end{somme}

\begin{sentence}
\ruby{\underline{四月}}{シガツ}に\ruby{は}{わ}\ruby{\underline{雨}}{あめ}が\ruby{\underline{多}}{おお}\underline{い}。\\\hline
    SHI{\middel}GATSU ni wa ame ga \={o}{\middel}i.\\\hline
(April) (rain) (many).\\\hline
=\textit{There's a lot of rain in April.}\\\hline
\end{sentence}
\end{multicols}
\vhrulefill{5pt}
%%--------------------------------------------------------------------
%%--------------------------------------------------------------------
\begin{multicols}{2}
\begin{glyph}{Moon (月)}{moon}\label{glyph:moon}
Moon, Month.
\end{glyph}

\begin{comps}
\comprtwocone & 月 \\\hline
\comprthreecone & Moon \\\hline
\end{comps}

\begin{remglyph}
Part of the Kanji radicals (\#{74}).\\\hline
\end{remglyph}

\begin{prons}
\pronsrtwocone & \textbf{ゲツ、ガツ (GETSU, GATSU)} & \textbf{つき (tsuki)} \\\hline
\pronsrthreecone & --- & --- \\\hline
\end{prons}

\begin{remprons}
Every \hooglig{month} everyone in the \comkun{stew kitchen} gets a Hyundai
    \comon{Getz}.\\\hline
\end{remprons}

\begin{somme}
    月 & \textit{moon; month} & つき \mbox{(tsu{\middel}ki)}\\\hline
    一月 & one + month = \textit{January} & イチ{\middel}ガツ \mbox{(ICHI{\middel}GATSU)}\\\hline
    二月 & two + month = \textit{February} & ニ{\middel}ガツ \mbox{(NI{\middel}GATSU)}\\\hline
    三月 & three + month = \textit{March} & サン{\middel}ガツ \mbox{(SAN{\middel}GATSU)}\\\hline
    四月 & four + month = \textit{April} & シ{\middel}ガツ \mbox{(SHI{\middel}GATSU)}\\\hline
    \notyet{先月} & precede + month = \textit{last month} & セン{\middel}ゲツ \mbox{(SEN{\middel}GETSU)}\\\hline
    \notyet{月見} & moon + see = \textit{moon viewing} & つき{\middel}み \mbox{(tsuki{\middel}mi)}\\\hline
\end{somme}

\begin{sentence}
\ruby{\underline{九}}{ク}\ruby{\underline{月}}{ガツ}に\ruby{\underline{月}}{つき}\ruby{\underline{見}}{み}に\ruby{\underline{行}}{い}きました。\\\hline
KU{\middel}GATSU ni tsuki{\middel}mi ni i{\middel}kimashita.\\\hline
(September) (moon viewing) (went)\\\hline
I went moon viewing in September.\\\hline
\end{sentence}
\end{multicols}
\vspace*{\fill}
%%--------------------------------------------------------------------
%%--------------------------------------------------------------------
\pagebreak
\vspace*{\fill}
\begin{multicols}{2}
\begin{glyph}{Bright (明)}{bright}\label{glyph:bright}
Bright; Clear; Light.
\end{glyph}

\begin{comps}
\comprtwocone & 日 + 月 \\\hline
\comprthreecone & Sun + Moon \\\hline
\end{comps}

\begin{remglyph}
    The \hooglig{bright} \remcom{sun} and \remcom{moon}.\\\hline
\end{remglyph}

\begin{prons}
\pronsrtwocone & \textbf{メイ (MEI)} & \textbf{あ、あか (a, aka)}\\\hline
\pronsrthreecone & ミョウ、ミン (MY\={O}, MIN) & あき (aki)\\\hline
\end{prons}

\begin{remprons}
The \hooglig{bright} \ucomon{myou} dolls, \comkun{a.k.a.} \ucomkun{akimbo} dolls, generated a \ucomon{mingled} \comon{mayhem} after the \comkun{update}.\\\hline
\end{remprons}

\begin{somme}
    明るい & \textit{bright} & あか{\middel}るい \mbox{(aka{\middel}rui)}\\\hline
    明かり & \textit{clearness} & あ{\middel}かり \mbox{(a{\middel}kari)}\\\hline
    明白 & bright + white = \textit{obvious} & メイ{\middel}ハク \mbox{(MEI{\middel}HAKU)}\\\hline
    \notyet{不明} & not + bright = \textit{unclear} & フ{\middel}メイ \mbox{(FU{\middel}MEI)}\\\hline
\end{somme}

\begin{sentence}
\ruby{\underline{月}}{つき}\ruby{は}{わ}\ruby{\underline{明}}{あか}\underline{るい}です。\\\hline
tsuki wa aka{\middel}rui desu.\\\hline
(moon) (bright)\\\hline
The moon is bright.\\\hline
\end{sentence}
\end{multicols}
\vhrulefill{5pt}
%%--------------------------------------------------------------------
%%--------------------------------------------------------------------
\begin{multicols}{2}
\begin{glyph}{Tree (木)}{tree}\label{glyph:tree}
Tree; Wood.
\end{glyph}

\begin{comps}
\comprtwocone & 木 \\\hline
\comprthreecone & Tree \\\hline
\end{comps}

\begin{remglyph}
Part of the Kanji radicals (\#{75}).\\\hline
\end{remglyph}

\begin{prons}
\pronsrtwocone & \textbf{モク、ボク (MOKU, BOKU)} & \textbf{き (ki)}\\\hline
\rye{3}{Look for モク to show up in the first position.\\\hline
	Look for ボク to show up in the second position.\\\hline
	き occurs less often, but appears with equal frequency in the first or second position.}
\pronsrthreecone & --- & こ (ko)\\\hline
\end{prons}

\begin{remprons}
From the \hooglig{tree} I was drinking \comon{mock} \comon{bock} to \comkun{keep} the \ucomkun{cobbles} of everyday life away.\\\hline
\end{remprons}

\begin{somme}
    木 & \textit{tree} & き \mbox{(ki)}\\\hline
    \notyet{木曜日} & tree + day of the week + sun (day) = \textit{Thursday} & モク{\middel}ヨウ{\middel}ビ \mbox{(MOKU{\middel}Y\={O}{\middel}BI)}\\\hline
    \notyet{木星} & tree + star = \textit{Jupiter (planet)} & モク{\middel}セイ \mbox{(MOKU{\middel}SEI)}\\\hline
    \notyet{大木} & large + tree = \textit{large tree} & タイ{\middel}ボク \mbox{(TAI{\middel}BOKU)}\\\hline
    \notyet{土木} & earth +tree = \textit{civil engineering} & ド{\middel}ボク \mbox{(DO{\middel}BOKU)}\\\hline
\end{somme}

\begin{irrr}
    木綿{\notcov} & tree + cottom = \textit{cotton (cloth)} & も{\middel}メン \mbox{(mo{\middel}MEN)}\\\hline
\rye{3}{The second kanji incidentally is not covered in this book.  Such
    characters will be included only when they form words in which the
    pronounciation of the character we are considering has an irregular reading.
    Irregular readings containing these kanji will be identified with a
    star{\notcov}. }
\end{irrr}

\begin{sentence}
\underline{あの} \ruby{\underline{高}}{たか}\underline{い} \ruby{\underline{木}}{き} の\ruby{\underline{名前}}{なまえ}が\ruby{\underline{分}}{わ}\underline{かります}か。\\\hline
ano taka{\middel}i ki no na{\middel}mae ga wa{\middel}karimasu ka.\\\hline
(that) (tall) (tree) (name) (understand).\\\hline
=\textit{Do you know the name of that tall tree?}\\\hline
\end{sentence}
\end{multicols}
\vspace*{\fill}
%%--------------------------------------------------------------------
%%--------------------------------------------------------------------
\pagebreak
\vspace*{\fill}
\begin{multicols}{2}
\begin{glyph}{Five (五)}{five}\label{glyph:five}
Five/Pent-.
\end{glyph}

\begin{comps}
\comprtwocone & 五 \\\hline
\comprthreecone & Five \\\hline
\end{comps}

\begin{remglyph}
    \textit{This is how most people feel when their work ends at five o'clock.}\\\hline
\end{remglyph}

\begin{prons}
\pronsrtwocone & \textbf{ ゴ (GO)} & ---\\\hline
\pronsrthreecone & --- & いつ (itsu)\\\hline
\end{prons}

\begin{remprons}
At \hooglig{5} o'clock we don't \ucomkun{eat soup}, we \comon{gobble}
    it.\\\hline
\end{remprons}

\begin{somme}
    五 & \textit{five} & ゴ \mbox{(GO)}\\\hline
    五人 & five + person = \textit{five people} & ゴ{\middel}ニン \mbox{(GO{\middel}NIN)}\\\hline
    五月 & five + moon (month) = \textit{May} & ゴ{\middel}ガツ \mbox{(GO{\middel}GATSU)}\\\hline
    五日 & five + sun (day) = \textit{fifth day of the month} & いつ{\middel}か \mbox{(itsu{\middel}ka)}\\\hline
    五つ & \textit{five (general counter)} & いつ{\middel}つ \mbox{(itsu{\middel}tsu)}\\\hline
    \notyet{五十} & five + ten = \textit{fifty} & ゴ{\middel}ジュウ \mbox{(GO{\middel}J\={U})}\\\hline
    \notyet{五時} & five + time = \textit{five o'clock} & ゴ{\middel}ジ \mbox{(GO{\middel}JI)}\\\hline
\end{somme}

\begin{sentence}
\ruby{\underline{五十人}}{ゴジュウニン}で\ruby{\underline{山}}{やま}え\ruby{\underline{行}}{い}\underline{きました}。\\\hline
GO{\middel}J\={U}{\middel}NIN de yama e i{\middel}kimashita.\\\hline
(fifty people) (mountain) (went)\\\hline
=\textit{Fifty (of us) went to the mountain.}\\\hline
\end{sentence}
\end{multicols}
\vhrulefill{5pt}
%%--------------------------------------------------------------------
%%--------------------------------------------------------------------
\begin{multicols}{2}
\begin{glyph}{Eye (目)}{eye}\label{glyph:eye}
Eye.\\\hline
In some instances it is read with its kun-yomi, and indicates the \textbf{``-th'' suffix} in words such as ``fourth'' and ``seventh''; it also expresses the related endings of ``first'', ``second'', ``third'', etc.
\end{glyph}

\begin{comps}
\comprtwocone & 目 \\\hline
\comprthreecone & Eye \\\hline
\end{comps}

\begin{remglyph}
Part of the Kanji radicals (\#{109}).\\\hline
\end{remglyph}

\begin{prons}
\pronsrtwocone & \textbf{モク (MOKU)} & \textbf{め (me)}\\\hline
\rye{3}{め is the more common reading of the two.}
\pronsrthreecone & ボク (BOKU) & ま (ma)\\\hline
\end{prons}

\begin{remprons}
I \hooglig{eyed} the \ucomkun{muddy} \comkun{mess} caused by the \comon{mock} \ucomon{bock}.\\\hline
\end{remprons}

\begin{somme}
    目 & \textit{eye} & め \mbox{(me)}\\\hline
    \notyet{目玉} & eye + jewel = \textit{eyeball} & め{\middel}だま \mbox{(me{\middel}dama)}\\\hline
    \notyet{八回目} & eight + rotate + eye = \textit{the eighth time} & ハチ{\middel}カイ{\middel}め \mbox{(HACHI{\middel}KAI{\middel}me)}\\\hline
    \notyet{目立つ} & eye + stand = \textit{to stand out} & め{\middel}だつ \mbox{(me{\middel}datsu)}\\\hline
    \notyet{注目} & pour + eye = \textit{notice} & チュウ{\middel}モク \mbox{(CH\={U}{\middel}MOKU)}\\\hline
\end{somme}

\begin{sentence}
\ruby{\underline{王子}}{オウジ}の\ruby{\underline{目}}{め} \ruby{は}{わ}\ruby{\underline{青い}}{あおい}です。\\\hline
\={O}{\middel}JI no me wa ao{\middel}i desu.\\\hline
(prince) (eye) (blue)\\\hline
=\textit{The prince's eyes are blue.}\\\hline
\end{sentence}
\end{multicols}
\vspace*{\fill}
%%--------------------------------------------------------------------
%%--------------------------------------------------------------------
\pagebreak
\begin{multicols}{2}
\begin{glyph}{Woman (女)}{woman}\label{glyph:woman}
Woman; Female.
\end{glyph}

\begin{comps}
\comprtwocone & 女 \\\hline
\comprthreecone & Woman \\\hline
\end{comps}

\begin{remglyph}
Part of the Kanji radicals (\#{38}).\\\hline
\end{remglyph}

\begin{prons}
\pronsrtwocone & \textbf{ジョ (JO)} & \textbf{おんな (onna)}\\\hline
\pronsrthreecone & ニョ、ニョウ (NYO, NY\={O}) & め (me)\\\hline
\end{prons}

\begin{remprons}
Many \hooglig{woman}, whose \comon{jobs} are in \ucomkun{messy} \ucomon{Neon} areas (an appropriate \ucomon{neologism}), are \comkun{on narcotic} drugs.\\\hline
\end{remprons}

\begin{somme}
    女 & \textit{woman} & おんな \mbox{(onna)}\\\hline
    \notyet{女王} & woman + king = \textit{queen} & ジョ{\middel}オウ \mbox{(JO{\middel}\={O})}\\\hline
    \notyet{少女} & few + woman = \textit{girl} & ショウ{\middel}ジョ \mbox{(SH\={O}{\middel}JO)}\\\hline
    \notyet{女心} & woman + heart = \textit{a woman's heart} & おんな{\middel}ごころ \mbox{(onna{\middel}gokoro)}\\\hline
    \notyet{雪女} & snow + heart = \textit{snow fairy} & ゆき{\middel}おんな \mbox{(yuki{\middel}onna)}\\\hline
    \notyet{男女} & man + woman = \textit{men and women} & ダン{\middel}ジョ \mbox{(DAN{\middel}JO)}\\\hline
\end{somme}

\begin{sentence}
\ruby{\underline{女}}{おんな}の\ruby{\underline{人}}{ひと}は\ruby{\underline{木}}{き}の\ruby{\underline{下}}{した}で\ruby{\underline{止}}{と}\underline{まりました}。\\\hline
onna no hito wa ki no shita de to{\middel}marimashita.\\\hline
(woman) (person) (tree) (under) (stopped).\\\hline
=\textit{The woman stopped under the tree.}\\\hline
\end{sentence}
\end{multicols}
\vhrulefill{5pt}
%%--------------------------------------------------------------------
%%--------------------------------------------------------------------
\begin{multicols}{2}
\begin{glyph}{Large (大)}{large}\label{glyph:large}
Large; Big.\\\hline
The meaning of ``important'' and ``great'' (not necessarily always in a positive sense) can also be implied in words containing this kanji.
\end{glyph}

\begin{comps}
\comprtwocone & 大 \\\hline
\comprthreecone & Large\\\hline
\end{comps}

\begin{remglyph}
 Part of the Kanji radicals (\#{37}).\\\hline
\end{remglyph}

\begin{prons}
\pronsrtwocone & \textbf{ダイ、タイ (DAI, TAI)} & \textbf{おお (\={o})}\\\hline
\rye{3}{ダイ and タイ occur with equal frequency when this kanji appears in the
	first position.  おお much less often.\\\hline When it occupies other
	positions in a compound, however, ダイ is by far the more common reading.}
\pronsrthreecone & --- & ---\\\hline
\end{prons}

\begin{remprons}
The \hooglig{large}, \comon{tight} fabric had an \comkun{awful} \comon{dye} to it.\\\hline
\end{remprons}

\begin{somme}
    大きい & \textit{large} & おお{\middel}きい \mbox{(\={o}{\middel}kii)}\\\hline
    大木 & large + tree = \textit{large tree} & タイ{\middel}ボク \mbox{(TAI{\middel}BOKU)} \\\hline
    \notyet{大好き} & large + like = \textit{to really like} & ダイ　す{\middel}き \mbox{(DAI su{\middel}ki)}\\\hline
    \notyet{大体} & large + body = \textit{in general} & ダイ{\middel}タイ \mbox{(DAI{\middel}TAI)}\\\hline
    \notyet{大雨} & large + rain = \textit{heavy rain} & おお{\middel}あめ \mbox{(\={o}{\middel}ame)}\\\hline
    \notyet{大切} & large + cut = \textit{important} & ダイ{\middel}セツ \mbox{(TAI{\middel}SETSU)}\\\hline
    \notyet{大学} & large + study = \textit{university} & ダイ{\middel}ガク \mbox{(DAI{\middel}GAKU)}\\\hline
\end{somme}

\begin{irrr}
    大人  & large + person = \textit{adult}  & おとな \mbox{(otona)}  \\\hline
\end{irrr}

\begin{sentence}
\underline{あの} \ruby{\underline{森}}{もり} にわ\ruby{\underline{大木}}{タイボク}が\ruby{\underline{多}}{おお}\underline{い}。\\\hline
    ano mori ni wa TAI{\middel}BOKU ga \={o}{\middel}i.\\\hline
(that) (forest) (large trees) (many)\\\hline
=\textit{There are many large trees in that forest}.\\\hline
\end{sentence}
\end{multicols}
%%--------------------------------------------------------------------
%%--------------------------------------------------------------------
\pagebreak
\vspace*{\fill}
\begin{multicols}{2}
\begin{glyph}{Middle (中)}{middle}\label{glyph:middle}
A versatile character that expresses the ideas of ``middle'', ``medium'', ``average'', etc.\\\hline
This kanji can also serves as an abbreviation for China.\\\hline
日 is used as the equivalent with respect to Japan.
\end{glyph}

\begin{comps}
\comprtwocone & 口 + 丨 \\\hline
\comprthreecone & Mouth + Pole (Vertical Stroke) \\\hline
\end{comps}

\begin{remglyph}
    The \hooglig{middle} \remcom{pole} will kill \remcom{dracula}.\\\hline
\end{remglyph}

\begin{prons}
\pronsrtwocone & \textbf{チュウ (CH\={U})} & \textbf{なか (naka)}\\\hline
\pronsrthreecone & ジュウ (J\={U}) & ---\\\hline
\end{prons}

\begin{remprons}
In the \hooglig{middle} of the \ucomon{juice} container is \comon{chewable}
    \comkun{knick-knacks}.\\\hline
\end{remprons}

\begin{somme}
    中 & \textit{middle} & なか \mbox{(naka)}\\\hline
    \notyet{中心} & middle + heart = \textit{center} & チュウ{\middel}シン \mbox{(CH\={U}{\middel}SHIN)}\\\hline
    \notyet{中米} & middle + rice = \textit{Central America} & チュウ{\middel}ベイ \mbox{(CH\={U}{\middel}BEI)}\\\hline
    \notyet{中国} & middle + country = \textit{China} & チュウ{\middel}ゴク \mbox{(CH\={U}{\middel}GOKU)}\\\hline
    \notyet{中立} & middle + stand = \textit{neutrality} & チュウ{\middel}リツ \mbox{(CH\={U}{\middel}RITSU)}\\\hline
    \notyet{中古} & middle + old = \textit{secondhand} & クウ{\middel}コ \mbox{(CH\={U}{\middel}KO)}\\\hline
    \notyet{水中} & water + middle = \textit{underwater/in the water} & スイ{\middel}チュウ \mbox{(SUI{\middel}CH\={U})}\\\hline
\end{somme}

\begin{sentence}
\ruby{\underline{車}}{くるま}の\ruby{\underline{中}}{なか}に\ruby{\underline{何}}{なに}が\underline{あります}か。\\\hline
kuruma no naka ni nani ga arimasu ka.\\\hline
(car) (middle) (what) (is)\\\hline
=\textit{What is in your car?}\\\hline
\end{sentence}
\end{multicols}
\vhrulefill{5pt}
%%--------------------------------------------------------------------
%%--------------------------------------------------------------------
\begin{multicols}{2}
\begin{glyph}{Eight (八)}{eight}\label{glyph:eight}
Eight.
\end{glyph}

\begin{comps}
\comprtwocone & 八 \\\hline
\comprthreecone & Eight\\\hline
\end{comps}

\begin{remglyph}
Part of the Kanji radicals (\#{12}).\\\hline
\end{remglyph}

\begin{prons}
\pronsrtwocone & \textbf{ハチ (HACHI)} & ---\\\hline
\pronsrthreecone &  & よう、やっ、や (y\={o}, ya-, ya)\\\hline
\end{prons}

\begin{remprons}
\hooglig{\comon{Hachi}} \ucomkun{yachted} in the days of \ucomkun{yore}.\\\hline
\end{remprons}

\begin{somme}
    八 & \textit{eight} & ハチ \mbox{(HACHI)}\\\hline
    八月 & eight + moon (month) = \textit{August} & ハチ{\middel}ガツ \mbox{(HACHI{\middel}GATSU)}\\\hline
    八人 & eight + person = \textit{eight people} & ハチ{\middel}ニン \mbox{(HACHI{\middel}NIN)}\\\hline
    八日 & eight + sun (day) = \textit{the eighth day of the month} & よう{\middel}か \mbox{(y\={o}{\middel}ka)}\\\hline
    八つ & \textit{eight (general counter)} & やっ{\middel}つ \mbox{(yat{\middel}tsu)}\\\hline
    \notyet{八円} & eight + circle = \textit{eight yen} & ハチ{\middel}エン \mbox{(HACHI{\middel}EN)}\\\hline
    \notyet{八時} & eight + time = \textit{eight o'clock} & ハチ{\middel}ジ \mbox{(HACHI{\middel}JI)}\\\hline
\end{somme}

\begin{sentence}
\ruby{\underline{八月}}{ハチガツ}に\ruby{\underline{秋田}}{あきた}え\ruby{\underline{行}}{い}\underline{きましょう}。\\\hline
HACHI{\middel}GATSU ni Aki{\middel}ta e i{\middel}kimash\={o}.\\\hline
(August) (Akita) (let's go)\\\hline
=\textit{Let's go to Akita in August.}\\\hline
\end{sentence}
\end{multicols}
\vspace*{\fill}
%%--------------------------------------------------------------------
%%--------------------------------------------------------------------
\pagebreak
\vspace*{\fill}
\begin{multicols}{2}
\begin{glyph}{Small (小)}{small}\label{glyph:small}
Small; Little.
\end{glyph}

\begin{comps}
\comprtwocone & 小 \\\hline
\comprthreecone & Small\\\hline
\end{comps}

\begin{remglyph}
Part of the Kanji radicals (\#{42}).\\\hline
\end{remglyph}

\begin{prons}
\pronsrtwocone & \textbf{ショウ (SH\={O})} & \textbf{こ (ko)}\\\hline
\pronsrthreecone & --- & お (o)\\\hline
\end{prons}

\begin{remprons}
    \hooglig{Small} cheese \comkun{cobbles} along the \comon{shore} will
    \ucomkun{oblige}.\\\hline
\end{remprons}

\begin{somme}
    小さい & \textit{small} & ちい{\middel}さい \mbox{(chii{\middel}sai)}\\\hline
    \notyet{小春} & small + spring = \textit{Indian Summer} & こ{\middel}はる \mbox{(ko{\middel}haru)}\\\hline
    \notyet{小国} & small + country = \textit{small country} & ショウ{\middel}コク \mbox{(SH\={O}{\middel}KOKU)}\\\hline
    \notyet{小高い} & small + tall = \textit{slightly elevated} & こ　だ{\middel}かい \mbox{(ko da{\middel}kai)}\\\hline
    \notyet{最小} & most + small = \textit{smallest} & サイ{\middel}ショウ \mbox{(SAI{\middel}SH\={O})}\\\hline
\end{somme}

\begin{sentence}
\underline{あの} \ruby{\underline{車}}{くるま} \ruby{は}{わ}\underline{とても} \ruby{\underline{小}}{ちい}\underline{さい} \underline{ですね}。\\\hline
ano kuruma wa totemo chii{\middel}sai desu ne.\\\hline
(that) (car) (really) (small) (isn't it).\\\hline
=\textit{That car is really small, isn't it?}\\\hline
\end{sentence}
\end{multicols}
\vhrulefill{5pt}
\vspace*{\fill}
%%--------------------------------------------------------------------
%%--------------------------------------------------------------------
\pagebreak

%
%\begin{multicols}{2}
%\chapter{Kanji 21--40}
%\vspace*{\fill}
\begin{glyph}{Shellfish (貝)}{shellfish}\label{glyph:shellfish}
Shellfish; Seashell.
\end{glyph}

\begin{comps}
\comprtwocone & 目 + 八 \\\hline
\comprthreecone & Eye (\ref{glyph:eye}) + Eight (\ref{glyph:eight})\\\hline
\end{comps}

\begin{remglyph}
\remcom{Cyclops} roasts his \hooglig{shellfish} over a \remcom{volcano}.\\\hline
\end{remglyph}

\begin{prons}
\pronsrtwocone & --- & \textbf{かい (kai)}\\\hline
\pronsrthreecone & バイ (BAI) & ---\\\hline
\end{prons}

\begin{remprons}
\comkun{Kindly} \ucomon{buy} the \hooglig{shellfish}.\\\hline
\end{remprons}

\begin{somme}
貝 & \textit{shellfish} & かい (kai)\\\hline
生貝 & life + shellfish = \textit{raw shellfish} & なま{\middel}がい (nama{\middel}gai)\\\hline
赤貝 & red + shellfish = \textit{ark shell} & あか{\middel}がい (aka{\middel}gai)\\\hline
\end{somme}

\begin{sentence}
\underline{この} \underline{レストラン}の\ruby{\underline{赤貝}}{あかがい}\ruby{は}{わ}\ruby{\underline{有名}}{ユウメイ}{\;}\underline{でぅ}。\\\hline
kono resutoran no aka{\middel}gai wa Y\={U}{\middel}MEI desu.\\\hline
(this) (restaurant) (ark shell) (famous) (is)\\\hline
=\textit{The ark shell in this restaurant is famous.}\\\hline
\end{sentence}
\end{multicols}
\vspace*{-\baselineskip}
\vhrulefill{5pt}
%%--------------------------------------------------------------------
%%--------------------------------------------------------------------
\begin{multicols}{2}
\begin{glyph}{Six (六)}{six}\label{glyph:six}
Six.
\end{glyph}

\begin{comps}
\comprtwocone &  亠 + 八\\\hline
\comprthreecone & Head/Above + Eight (\ref{glyph:eight}) \\\hline
\end{comps}

\begin{remglyph}
The \hooglig{snail} protected itself \remcom{above} the \remcom{volcano}.\\\hline
\end{remglyph}

\begin{prons}
\pronsrtwocone & \textbf{ロク (ROKU)} & ---\\\hline
\pronsrthreecone & リク (RIKU) & むい、むっ、む (mui, mu-, mu)\\\hline
\end{prons}

\begin{remprons}
The \hooglig{6} \ucomkun{mooing} \ucomkun{muslims} were \comon{locked} up to prevent any further \ucomon{licking}.\\\hline
\end{remprons}

\begin{somme}
六 & \textit{six} & ロク (ROKU)\\\hline
六人 & six + person = \textit{six people} & ロク{\middel}ニン (ROKU{\middel}NIN)\\\hline
六月 & six + moon (month) = \textit{June} & ロク{\middel}ガツ (ROKU{\middel}GATSU)\\\hline
六円 & six + circle (yen) = \textit{six yen} & ロク{\middel}エン (ROKU{\middel}EN)\\\hline
六時 & six + time = \textit{six o'clock} & ロク{\middel}ジ (ROKU{\middel}JI)\\\hline
六日 & six + sun (day) = \textit{the sixth day of the month} & むい{\middel}か (mui{\middel}ka)\\\hline
六つ & \textit{six (general counter)} & むっ{\middel}つ (mut{\middel}tsu)\\\hline
\end{somme}

\begin{sentence}
\ruby{\underline{六人}}{ロクニン}\ruby{は}{わ}\ruby{\underline{日本}}{ニホン}\ruby{へ}{え}\ruby{\underline{行}}{い}\underline{きました}。\\\hline
ROKU{\middel}NIN wa NI{\middel}HON e i{\middel}kimashita.\\\hline
(six people) (Japan) (went)\\\hline
=\textit{The six people went to Japan.}\\\hline
\end{sentence}
\end{multicols}
\vspace*{-\baselineskip}
\vhrulefill{5pt}
\vspace*{\fill}
%%--------------------------------------------------------------------
%%--------------------------------------------------------------------
\pagebreak
\vspace*{\fill}
\begin{multicols}{2}
\begin{glyph}{King (王)}{king}\label{glyph:king}
King; royalty.
\end{glyph}

\begin{comps}
\comprtwocone & 王\\\hline
\comprthreecone & King (\ref{glyph:king}) \\\hline
\end{comps}

\begin{remglyph}
---\\\hline
\end{remglyph}

\begin{prons}
\pronsrtwocone & \textbf{オウ (\={O})} & ---\\\hline
\pronsrthreecone & --- & ---\\\hline
\end{prons}

\begin{remprons}
\comon{Awesome} \hooglig{king}.\\\hline
\end{remprons}

\begin{somme}
女王 & woman + king = \textit{queen} & ジョ{\middel}オウ (JO{\middel}\={O})\\\hline
王国 & king + country = \textit{kingdom} & オウ{\middel}コク (\={O}{\middel}KOKU)\\\hline
王朝 & king + morning = \textit{dynasty} & オウ{\middel}チョウ (\={O}{\middel}CH\={O})\\\hline
王子 & king + child = \textit{prince} & オウ{\middel}ジ (\={O}{\middel}JI)\\\hline
海王星 & sea + king + star = \textit{Neptune} & カイ{\middel}オウ{\middel}セイ (KAI{\middel}\={O}{\middel}SEI)\\\hline
\end{somme}

\begin{sentence}
\ruby{\underline{女王}}{ジョオウ}が\ruby{\underline{日本}}{日本}\ruby{へ}{え}\ruby{\underline{行}}{い}\underline{きました}。\\\hline
JO{\middel}\={O} ga NI{\middel}HON e i{\middel}kimashita.\\\hline
(Queen) (Japan) (went)\\\hline
=\textit{The Queen went to Japan.}\\\hline
\end{sentence}
\end{multicols}
\vspace*{-\baselineskip}
\vhrulefill{5pt}
%%--------------------------------------------------------------------
%%--------------------------------------------------------------------
\begin{multicols}{2}
\begin{glyph}{Jewel (玉)}{jewel}\label{glyph:jewel}
Jewel.  Also incorporates the sense of ball-like objects.
\end{glyph}

\begin{comps}
\comprtwocone & 王 + 丶\\\hline
\comprthreecone & King (\ref{glyph:king}) + Dot \\\hline
\end{comps}

\begin{remglyph}
A \remcom{king} carelessly treats his \hooglig{jewels} as if they are random \remcom{dots}.\\\hline
\end{remglyph}

\begin{prons}
\pronsrtwocone & --- & \textbf{たま (tama)}\\\hline
\pronsrthreecone & ギョク (GYOKU) & ---\\\hline
\end{prons}

\begin{remprons}
\comkun{Tamales} is a \hooglig{jewel} of a dish.\\\hline
\end{remprons}

\begin{somme}
玉 & \textit{jewel} & たま (tama)\\\hline
玉ねぎ & \textit{onion} & たま{\middel}ねぎ (tama{\middel}negi)\\\hline
目玉 & eye + jewel = \textit{eyeball} & め{\middel}だま (me{\middel}dama)\\\hline
水玉 & water + jewel = \textit{water droplet} & みず{\middel}たま (mizu{\middel}tama)\\\hline
火玉 & fire + jewel = \textit{fireball} & ひ{\middel}だま (hi{\middel}dama)\\\hline
雪玉 & snow + jewel = \textit{snowball} & ゆき{\middel}だま (yuki{\middel}dama)\\\hline
\end{somme}

\begin{sentence}
\ruby{\underline{玉}}{たま}\underline{ねぎ}\ruby{を}{お}\ruby{\underline{切}}{き}\underline{って} \ruby{\underline{下}}{くだ}\underline{さい}。\\\hline
tama{\middel}negi o ki{\middel}tte kuda{\middel}sai.\\\hline
(onion(s)) (cut) (please).\\\hline
=\textit{Please cut the onion(s)}\\\hline
\end{sentence}
\end{multicols}
\vspace*{-\baselineskip}
\vhrulefill{5pt}
\vspace*{\fill}
%%--------------------------------------------------------------------
%%--------------------------------------------------------------------
\pagebreak
\vspace*{\fill}
\begin{multicols}{2}
\begin{glyph}{Heart (心)}{heart}\label{glyph:heart}
Heart.\\\hline
The sense of feelings.\\\hline
The concept of something's core.
\end{glyph}

\begin{comps}
\comprtwocone & 丶 + 亅 + 丶 + 丶\\\hline
\comprthreecone & Dot + Hook + Dot + Dot \\\hline
\end{comps}

\begin{remglyph}
You can win a \hooglig{heart} by \remcom{doting} or by a \remcom{hook} using brute force.\\\hline
\end{remglyph}

\begin{prons}
\pronsrtwocone & \textbf{シン (SHIN)} & \textbf{こころ (kokoro)}\\\hline
\pronsrthreecone & --- & ---\\\hline
\rye{3}{The kun-yomi for this kanji is always voiced when not in the first position.}
\end{prons}

\begin{remprons}
Her \hooglig{heart} had a pink, romantic \comon{sheen} to it when she gave me the \comkun{Cocoa Roast}.\\\hline
\end{remprons}

\begin{somme}
心 & \textit{heart} & こころ (kokoro)\\\hline
女心 & woman + heart = \textit{a woman's heart} & おんな{\middel}ごころ (onna{\middel}gokoro)\\\hline
中心 & middle + heart = \textit{center} & チュウ{\middel}シン (CH\={U}{\middel}SHIN)\\\hline
下心 & lower + heart = \textit{ulterior motive} & した{\middel}ごころ (shita{\middel}gokoro)\\\hline
安心 & ease + heart = \textit{peace of mind} & アン{\middel}シン (AN{\middel}SHIN)\\\hline
愛国心 & love + country + heart = \textit{patriotism} & アイ{\middel}コク{\middel}シン (AI{\middel}KOKU{\middel}SHIN)\\\hline
\end{somme}

\begin{sentence}
\ruby{\underline{女王}}{ジョオウ}{\;}\underline{によると} \ruby{\underline{愛国心}}{アイコクシン}\ruby{は}{わ}\ruby{\underline{大切}}{タイセツ}だそうです。\\\hline
JO{\middel}\={O} ni yoru to AI{\middel}KOKU{\middel}SHIN wa TAI{\middel}SETSU da s\={o} desu.\\\hline
(queen) (according to) (patriotism) (important)\\\hline
=\textit{According to the queen, patriotism is important.}\\\hline
\end{sentence}
\end{multicols}
\vspace*{-\baselineskip}
\vhrulefill{5pt}
%%--------------------------------------------------------------------
%%--------------------------------------------------------------------
\begin{multicols}{2}
\begin{glyph}{Country (国)}{country}\label{glyph:country}
Country; Nation; State.
\end{glyph}

\begin{comps}
\comprtwocone & 囗 + 玉\\\hline
\comprthreecone & Prison/Enclosure + Jewel (\ref{glyph:jewel}) \\\hline
\end{comps}

\begin{remglyph}
Within these \remcom{borders} is a \remcom{jewel} of a \hooglig{country}.\\\hline
\end{remglyph}

\begin{prons}
\pronsrtwocone & \textbf{コク (KOKU)} & \textbf{くに (kuni)}\\\hline
\pronsrthreecone & --- & ---\\\hline
\end{prons}

\begin{remprons}
That \hooglig{country's} \comkun{cuneiform} \comon{cocked} my interest.\\\hline
\end{remprons}

\begin{somme}
国 & \textit{country} & コク (KOKU)\\\hline
王国 & king + country = \textit{Kingdom} & オウ{\middel}コク (\={O}{\middel}KOKU)\\\hline
全国 & complete + country = \textit{the whole country/nationwide} & ゼン{\middel}コク (ZEN{\middel}KOKU)\\\hline
入国 & enter + country = \textit{to enter a country} & ニュウ{\middel}コク (NY\={U}{\middel}KOKU)\\\hline
国内 & country + inside = \textit{domestic} & コク{\middel}ナイ (KOKU{\middel}NAI)\\\hline
外国人 & outside + country + person = \textit{foreigner} & ガイ{\middel}コク{\middel}ジン (GAI{\middel}KOKU{\middel}JIN)\\\hline
愛国心 & love + country + heart = \textit{patriotism} & アイ{\middel}コク{\middel}シン (AI{\middel}KOKU{\middel}SHIN)\\\hline
\end{somme}

\begin{sentence}
\underline{あの} \ruby{\underline{国}}{くに}には\ruby{\underline{外国人}}{ガイコクジン}が\ruby{\underline{多}}{おお}\underline{い}。\\\hline
ano kuni ni wa GAI{\middel}KOKU{\middel}JIN ga \={o}{\middel}i.\\\hline
(that) (country) (foreigners) (many)\\\hline
=\textit{There are many foreigners in that country.}\\\hline
\end{sentence}
\end{multicols}
\vspace*{-\baselineskip}
\vhrulefill{5pt}
\vspace*{\fill}
%%--------------------------------------------------------------------
%%--------------------------------------------------------------------
\vspace*{\fill}
\pagebreak
\begin{multicols}{2}
\begin{glyph}{Complete (全)}{complete}\label{glyph:complete}
Complete; Wholeness; Completion.
\end{glyph}

\begin{comps}
\comprtwocone & 𠆢 + 王 \\\hline
\comprthreecone & Person (\ref{glyph:person}) + King (\ref{glyph:king}) \\\hline
\end{comps}

\begin{remglyph}
A \remcom{king} is a \hooglig{complete} \remcom{person}.\\\hline
\end{remglyph}

\begin{prons}
\pronsrtwocone & \textbf{ゼン (ZEN)} & --- \\\hline
\pronsrthreecone & --- & まった、まっと (matta, matto) \\\hline
\end{prons}

\begin{remprons}
I had \hooglig{complete} \comon{Zen} until the \ucomkun{mutt} attacked me with a \ucomkun{mattock}.\\\hline
\end{remprons}

\begin{somme}
全国 & complete + country = \textit{the whole country/nationwide} & ゼン{\middel}コク (ZEN{\middel}KOKU)\\\hline
安全 & ease + complete = \textit{safety} & アン{\middel}ゼン (AN{\middel}ZEN)\\\hline
全校 & complete + school = \textit{the whole school/all schools} & ゼン{\middel}コウ (ZEN{\middel}K\={O})\\\hline
全力 & complete + strength = \textit{all one's power} & ゼン{\middel}リョク (ZEN{\middel}RYOKU)\\\hline
全体 & complete + body = \textit{the whole} & ゼン{\middel}タイ (ZEN{\middel}TAI)\\\hline
全部 & complete + part = \textit{all} & ゼン{\middel}ブ (ZEN{\middel}BU)\\\hline
\end{somme}

\begin{sentence}
\ruby{\underline{田中}}{たなか}\underline{さん}\ruby{は}{わ}\ruby{\underline{全国}}{ゼンコク}\ruby{を}{お}\ruby{\underline{回}}{まわ}\underline{る}。\\\hline
Ta{\middel}naka-san wa ZEN{\middel}KOKU o mawa{\middel}ru.\\\hline
(Tanaka-san) (whole country) (go around)\\\hline
=\textit{Tanaka-san is going around the whole country.}\\\hline
\end{sentence}
\end{multicols}
\vspace*{-\baselineskip}
\vhrulefill{5pt}
%%--------------------------------------------------------------------
%%--------------------------------------------------------------------
\begin{multicols}{2}
\begin{glyph}{Ten (十)}{ten}\label{glyph:ten}
Ten.\\\hline
The first two examples of Entry 176 show another use for this character, an interesting application based entirely on its shape.
\end{glyph}

\begin{comps}
\comprtwocone & 十 \\\hline
\comprthreecone & Ten (\ref{glyph:ten}) \\\hline
\end{comps}

\begin{remglyph}
\remcom{Scarecrow}.\\\hline
\end{remglyph}

\begin{prons}
\pronsrtwocone & \textbf{ジュウ (J\={U})} & --- \\\hline
\pronsrthreecone & ジツ (JITSU) & とお、と (t\={o}, to) \\\hline
\end{prons}

\begin{remprons}
As the Brave \hooglig{10} were \ucomkun{tormented} by the \ucomkun{torrents} of rain, they had to get their creative \comon{juices} flowing to keep on practising \ucomon{ju-jitsu}.\\\hline
\end{remprons}

\begin{somme}
十 & \textit{ten} & ジュウ (J\={U})\\\hline
十月 & ten + moon (month) = \textit{October} & ジュウ{\middel}ガツ (J\={U}{\middel}GATSU)\\\hline
十一月 & ten + one + moon (month) = \textit{November} & ジュウ{\middel}イチ{\middel}ガツ (J\={U}{\middel}ICHI{\middel}GATSU)\\\hline
十二月 & ten + two + moon (month) = \textit{December} & ジュウ{\middel}ニ{\middel}ガツ (J\={U}{\middel}NI{\middel}GATSU)\\\hline
十六 & ten + six = \textit{sixteen} & ジュウ{\middel}ロク (J\={U}{\middel}ROKU)\\\hline
十八 & ten + eight = \textit{eighteen} & ジュウ{\middel}ハチ (J\={U}{\middel}HACHI)\\\hline
十日 & ten + sun (day) = \textit{the tenth day of the month} & とお{\middel}か (t\={o}{\middel}ka)\\\hline
\end{somme}

\begin{irrr}
二十日 & two + ten + sun (day) = \textit{the twentieth day of the month} & はつか (hatsuka)\\\hline
二十歳 & two + ten + annual = \textit{twenty years old} & はたち (hatachi)\\\hline
\end{irrr}

\begin{sentence}
\ruby{\underline{十月}}{ジュウガツ}に\ruby{は}{わ}\underline{あの} \ruby{\underline{木}}{き}\ruby{は}{わ}\ruby{\underline{美}}{うつく}\underline{しく} \underline{なります}。\\\hline
J\={U}{\middel}GATSU ni wa ano ki wa utsuku{\middel}shiku narimasu.\\\hline
(October) (that) (tree) (beautiful) (becomes)\\\hline
=\textit{That tree becomes beautiful in October.}\\\hline
\end{sentence}
\end{multicols}
\vspace*{-\baselineskip}
\vhrulefill{5pt}
\vspace*{\fill}
%%--------------------------------------------------------------------
%%--------------------------------------------------------------------
\pagebreak
\vspace*{\fill}
\begin{multicols}{2}
\begin{glyph}{Early (早)}{early}\label{glyph:early}
Early; Fast.
\end{glyph}

\begin{comps}
\comprtwocone & 日 + 十 \\\hline
\comprthreecone & Sun (\ref{glyph:sun}) + Ten (\ref{glyph:ten}) \\\hline
\end{comps}

\begin{remglyph}
\hooglig{Early} in the \remcom{morning-sun} the \remcom{scarecrow} gets into shape.\\\hline
\end{remglyph}

\begin{prons}
\pronsrtwocone & \textbf{ソウ (S\={O})} & \textbf{はや (haya)} \\\hline
\pronsrthreecone & サツ (SATSU) & --- \\\hline
\end{prons}

\begin{remprons}
The \ucomon{sutler soup} eased the \comon{sorely} displeasing \comkun{hay allergy} I had \hooglig{earlier}.\\\hline
\end{remprons}

\begin{somme}
早い & \textit{early; fast} & はや{\middel}い (haya{\middel}i)\\\hline
早まる & \textit{to be in a hurry} & はや{\middel}まる (haya{\middel}maru)\\\hline
早める & \textit{to hurry something up} & はや{\middel}める (haya{\middel}meru)\\\hline
早口 & early + mouth = \textit{rapid speaking} & はや{\middel}くち (haya{\middel}kuchi)\\\hline
早春 & early + spring = \textit{early spring} & ソウ{\middel}シュン (S\={O}{\middel}SHUN)\\\hline
早朝 & early + morning = \textit{early morning} & ソウ{\middel}チョウ (S\={O}{\middel}CH\={O})\\\hline
\end{somme}

\begin{sentence}
\underline{どうして} \underline{そんな}に\ruby{\underline{早}}{はや}\underline{く} \ruby{\underline{来}}{き}\underline{た}のですか。\\\hline
d\={o}shite sonna ni haya{\middel}ku ki{\middel}ta no desu ka.\\\hline
(why) (so) (early) (come)\\\hline
=\textit{Why did you come so early?}\\\hline
\end{sentence}
\end{multicols}
\vspace*{-\baselineskip}
\vhrulefill{5pt}
%%--------------------------------------------------------------------
%%--------------------------------------------------------------------
\begin{multicols}{2}
\begin{glyph}{Upper (上)}{upper}\label{glyph:upper}
Upper; On; Over.\\\hline
It can refer to anything from goods of high quality to superiors at work.
\end{glyph}

\begin{comps}
\comprtwocone & 卜 + 一 \\\hline
\comprthreecone & Divination + One (\ref{glyph:one}) \\\hline
\end{comps}

\begin{remglyph}
\remcom{One} \hooglig{upper} \remcom{divination}.\\\hline
\end{remglyph}

\begin{prons}
\pronsrtwocone & \textbf{ジョウ (J\={O})} & \textbf{かみ、うえ、のぼ、あ、うわ (kami, ue, nobo, a, uwa)} \\\hline
\pronsrthreecone & ショウ (SH\={O}) & --- \\\hline
\rye{3}{J\={O} is by far the most frequently used pronounciation.}
\end{prons}

\begin{remprons}
\hooglig{Upper} \comon{jaw}.\\\hline
\end{remprons}

\begin{somme}
上 & \textit{upper} & かみ (kami)\\\hline
上 & \textit{over; on} & うえ (ue)\\\hline
上る (intr) & \textit{to climb} & のぼ{\middel}る (nobo{\middel}ru)\\\hline
上がる (intr) & \textit{to rise} & あ{\middel}がる (a{\middel}garu)\\\hline
上げる (tr/intr) & \textit{to lift} & あ{\middel}げる (a{\middel}geru)\\\hline
上手 & upper + hand = \textit{skillful} & じょう{\middel}ズ (J\={O}{\middel}ZU)\\\hline
\rye{3}{The verb あ{\middel}げる is almost always used in a transitive sense (that is, it lifts some objects in either a physical or symbolic way), and is best thought of as being paired with the intransitive あ{\middel}がる.}
\end{somme}

\begin{sentence}
\ruby{\underline{肉}}{ニク}\ruby{は}{わ}\underline{テーベル}の\ruby{\underline{上}}{うえ}に\underline{あります}。\\\hline
NIKU wa T\={E}BERU no ue ni arimasu.\\\hline
(meat) (table) (on) (is)\\\hline
=\textit{The meat is on the table.}\\\hline
\end{sentence}
\end{multicols}
\vspace*{-\baselineskip}
\vhrulefill{5pt}
\vspace*{\fill}
%%--------------------------------------------------------------------
%%--------------------------------------------------------------------
\pagebreak
\begin{multicols}{2}
\begin{glyph}{Lower (下)}{lower}\label{glyph:lower}
Lower; Under; Below.
\end{glyph}

\begin{comps}
\comprtwocone & 一 + 卜 \\\hline
\comprthreecone & One (\ref{glyph:one}) + Divination \\\hline
\end{comps}

\begin{remglyph}
\remcom{One} \hooglig{lower} \remcom{divination}.\\\hline
\end{remglyph}

\begin{prons}
\pronsrtwocone & \textbf{カ、ゲ (KA, GE)} & \textbf{した、しま、お、さ、くだ (shita, shimo, o, sa, kuda)}\\\hline
\pronsrthreecone & --- & もと (moto) \\\hline
\rye{3}{お、さ、くだ above are all verb stems, and will thus be accompanied by hiragana hinting at the correct pronounciation.\\\hline Regarding compounds, カ and ゲ will be encountered far more than the others, although there is no obvious pattern as to when each of these is used.}
\end{prons}

\begin{remprons}
---\\\hline
\end{remprons}

\begin{somme}
下 & \textit{low; below; under} & した (shita)\\\hline
下りる (intr) & \textit{to come down} & お{\middel}りる (o{\middel}riru)\\\hline
下ろす (tr) & \textit{to take down} & お{\middel}ろす (o{\middel}rosu)\\\hline
下がる (intr) & \textit{to hang down (on one's own)} & さ{\middel}がる (sa{\middel}garu)\\\hline
下げる (tr) & \textit{to lower (something)} & さ{\middel}げる (sa{\middel}geru)\\\hline
下さる & \textit{to oblige} & くだ{\middel}さる (kuda{\middel}saru)\\\hline
上下 & upper + lower = \textit{high and low} & ジョウ{\middel}ゲ (J\={O}{\middel}GE)\\\hline
天下 & heaven + lower = \textit{the whole land} & テン{\middel}カ (TEN{\middel}KA)\\\hline
\end{somme}

\begin{irrr}
下手 & lower + hand = \textit{to be poor at something} & へ{\middel}た (he{\middel}ta)\\\hline
\end{irrr}

\begin{sentence}
\ruby{\underline{木}}{き}{\;}\underline{から} \ruby{\underline{下}}{お}\underline{りて} \ruby{\underline{下}}{くだ}\underline{さい}。\\\hline
ki kara o{\middel}rite kuda{\middel}sai.\\\hline
(tree) (from) (come down) (please)\\\hline
=\textit{Please come down from the tree.}\\\hline
\end{sentence}
\end{multicols}
\vspace*{-\baselineskip}
\vhrulefill{5pt}
%%--------------------------------------------------------------------
%%--------------------------------------------------------------------
\begin{multicols}{2}
\begin{glyph}{Rice (米)}{rice}\label{glyph:rice}
(Uncooked) rice.\\\hline
This character is also used to symbolize the America's.
\end{glyph}

\begin{comps}
\comprtwocone & 米 (\ref{glyph:rice}) \\\hline
\comprthreecone & Bunny Ears + Tree (\ref{glyph:tree}) \\\hline
\end{comps}

\begin{remglyph}
---\\\hline
\end{remglyph}

\begin{prons}
\pronsrtwocone & \textbf{マイ (MAI)} & \textbf{こめ (kome)}\\\hline
\pronsrthreecone & ベイ (BEI) & --- \\\hline
\rye{3}{ベイ is used as the reading in compounds relating to the Americas.\\\hline
マイ is the primary choice when it refers to rice.}
\end{prons}

\begin{remprons}
We'll make our \comkun{comeback} by \comon{basically} \comon{migrating} to regions with plentiful \hooglig{rice}.\\\hline
\end{remprons}

\begin{somme}
米 & \textit{(uncooked) rice} & こめ (kome)\\\hline
白米 & white + rice = \textit{polished rice} & ハク{\middel}マイ (HAKU{\middel}MAI)\\\hline
米国 & rice + country = \textit{The United States} & ベイ{\middel}コク (BEI{\middel}KOKU)\\\hline
日米 & sun + rice = \textit{Japan-U.S} & ニチ{\middel}ベイ (NICHI{\middel}BEI)\\\hline
中米 & middle + rice = \textit{Central America} & チュウ{\middel}ベイ (CH\={U}{\middel}BEI)\\\hline
北米 & North + rice = \textit{North America} & ホク{\middel}ベイ (HOKU{\middel}BEI)\\\hline
南米 & South + rice = \textit{South America} & ナン{\middel}ベイ (NAN{\middel}BEI)\\\hline
\end{somme}

\begin{sentence}
\ruby{\underline{米国}}{ベイコク}の\ruby{\underline{牛肉}}{ギュウニク}\ruby{は}{わ}\ruby{\underline{高}}{たか}\underline{い}。\\\hline
BEI{\middel}KOKU no GY\={U}{\middel}NIKU wa taka{\middel}i.\\\hline
(United States) (beef) (expensive)\\\hline
=\textit{U.S. beef is expensive.}\\\hline
\end{sentence}
\end{multicols}
\vspace*{-\baselineskip}
\vhrulefill{5pt}
%%--------------------------------------------------------------------
%%--------------------------------------------------------------------
\pagebreak
\vspace*{\fill}
\begin{multicols}{2}
\begin{glyph}{Self (自)}{self}\label{glyph:self}
Self.
\end{glyph}

\begin{comps}
\comprtwocone & 丶 + 目 \\\hline
\comprthreecone & Dot + Eye (\ref{glyph:eye}) \\\hline
\end{comps}

\begin{remglyph}
\remcom{Cyclops} loses \hooglig{self}-control when he sees a \remcom{dot}.\\\hline
\end{remglyph}

\begin{prons}
\pronsrtwocone & \textbf{ジ (JI)} & ---\\\hline
\pronsrthreecone & シ (SHI) & みずか (mizuka) \\\hline
\rye{3}{シ only appears in the compound 自然: self + nature = \textit{Nature}, シ{\middel}ゼン.}
\end{prons}

\begin{remprons}
I told \hooglig{myself}: ``Make \ucomkun{me zoom completely} on the \comon{jigged} \ucomon{shields}''.\\\hline
\end{remprons}

\begin{somme}
自国 & self + country = \textit{one's homeland} & ジ{\middel}コク (JI{\middel}KOKU)\\\hline
自立 & self + stand = \textit{independence} & ジ{\middel}リツ (JI{\middel}RITSU)\\\hline
自分 & self + part = \textit{oneself} & ジ{\middel}ブン (JI{\middel}BUN)\\\hline
自体 & self + body = \textit{in itself} & ジ{\middel}タイ (JI{\middel}TAI)\\\hline
自然 & self + body = \textit{nature} & シ{\middel}ゼン (SHI{\middel}ZEN)\\\hline
\end{somme}

\begin{sentence}
\ruby{\underline{自分}}{ジブン}の\ruby{\underline{車}}{くるま}が\ruby{\underline{買}}{か}\underline{いたい}。\\\hline
JI{\middel}BUN no kuruma ga ka{\middel}itai.\\\hline
(oneself) (car) (want to buy)\\\hline
=\textit{I want to buy my own car.}\\\hline
\end{sentence}
\end{multicols}
\vspace*{-\baselineskip}
\vhrulefill{5pt}
%%--------------------------------------------------------------------
%%--------------------------------------------------------------------
\begin{multicols}{2}
\begin{glyph}{Inside (内)}{inside}\label{glyph:inside}
Inside; Interior; Within.\\\hline
This kanji can also convey the idae of ``home''.
\end{glyph}

\begin{comps}
\comprtwocone & 冂 + 人 \\\hline
\comprthreecone & Wilderness + Person (\ref{glyph:person}) \\\hline
\end{comps}

\begin{remglyph}
The \hooglig{inside} of a \remcom{person} is a \remcom{wilderness}.\\\hline
\end{remglyph}

\begin{prons}
\pronsrtwocone & \textbf{ナイ (NAI)} & \textbf{うち (uchi)}\\\hline
\pronsrthreecone & ダイ (DAI) & --- \\\hline
\end{prons}

\begin{remprons}
\hooglig{In} the \comkun{Uchica} community are \comon{nicely} \ucomon{dyed} fabrics.\\\hline
\end{remprons}

\begin{somme}
内 & \textit{inside} & うち (うち)\\\hline
国内 & country + inside = \textit{domestic} & コク{\middel}ナイ (KOKU{\middel}NAI)\\\hline
市内 & city + inside = \textit{within the city} & シ{\middel}ナイ (SHI{\middel}NAI)\\\hline
内海 & inside + sea = \textit{inland sea} & ナイ{\middel}カイ (NAI{\middel}KAI)\\\hline
車内 & car + inside = \textit{inside the car} & シャ{\middel}ナイ (SHA{\middel}NAI)\\\hline
町内 & town + inside = \textit{in the town} & チョウ{\middel}ナイ (CH\={O}{\middel}NAI)\\\hline
\end{somme}

\begin{sentence}
\underline{その} \ruby{\underline{人}}{ひと}\ruby{は}{わ}\ruby{\underline{国内}}{コクナイ}で\ruby{\underline{有名}}{ユウメイ}です。\\\hline
sono hito wa KOKU{\middel}NAI de Y\={U}{\middel}MEI desu.\\\hline
(that) (person) (domestic) (famous)\\\hline
=\textit{That person is famous domestically.}\\\hline
\end{sentence}
\end{multicols}
\vspace*{-\baselineskip}
\vhrulefill{5pt}
\vspace*{\fill}
%%--------------------------------------------------------------------
%%--------------------------------------------------------------------
\pagebreak
\vspace*{\fill}
\begin{multicols}{2}
\begin{glyph}{Right (右)}{migi}\label{glyph:right}
Right; Right-hand side.
\end{glyph}

\begin{comps}
\comprtwocone & 一 + ノ + 口  \\\hline
\comprthreecone & One (\ref{glyph:one}) + No (katakana)/Slash + Mouth (\ref{glyph:mouth}) \\\hline
\end{comps}

\begin{remglyph}
\hooglig{Rightly} \remcom{slash} that \remcom{one's} \remcom{mouth}.\\\hline
\end{remglyph}

\begin{prons}
\pronsrtwocone & --- & \textbf{みぎ (migi)}\\\hline
\pronsrthreecone & ユウ、ウ (Y\={U}, U) & --- \\\hline
\end{prons}

\begin{remprons}
On \ucomon{Youtube} was an \ucomon{ooky} clip of a guy \comkun{midging} with his \hooglig{right} hand.\\\hline
\end{remprons}

\begin{somme}
右 & \textit{right} & みぎ (migi)\\\hline
右回り & right + rotate = \textit{clockwise} & みぎ　まわ{\middel}り (migi mawa{\middel}ri)\\\hline
右上 & right + upper = \textit{upper right} & みぎ{\middel}うえ (migi{\middel}ue)\\\hline
右下 & right + lower = \textit{lower right} & みぎ{\middel}した (migi{\middel}shita)\\\hline
右手 & right + hand = \textit{right hand} & みぎ{\middel}て (migi{\middel}te)\\\hline
右左 & right + leg = \textit{right leg} & みぎ{\middel}あし (migi{\middel}ashi)\\\hline
\end{somme}

\begin{sentence}
\ruby{\underline{山口}}{やまぐち}\underline{さん}\ruby{は}{わ}\ruby{\underline{右}}{みぎ}の\ruby{\underline{目}}{め}が\underline{かゆい}。\\\hline
Yama{\middel}guchi-san wa migi no me ga kayui.\\\hline
(Yamaguchi-san) (right) (eye) (itchy)\\\hline
=\textit{Yamaguchi-san's right eye is itchy.}\\\hline
\end{sentence}
\end{multicols}
\vspace*{-\baselineskip}
\vhrulefill{5pt}
%%--------------------------------------------------------------------
%%--------------------------------------------------------------------
\begin{multicols}{2}
\begin{glyph}{Have (有)}{have}\label{glyph:have}
Have; Possess.
\end{glyph}

\begin{comps}
\comprtwocone & 一 + ノ + 月  \\\hline
\comprthreecone & One (\ref{glyph:one}) + No (katakana)/Slash + Moon (\ref{glyph:moon}) \\\hline
\end{comps}

\begin{remglyph}
\hooglig{Have} \remcom{one} \remcom{slash} at the \remcom{moon}.\\\hline
\end{remglyph}

\begin{prons}
\pronsrtwocone & \textbf{ユウ (Y\={U})} & \textbf{あ (a)}\\\hline
\pronsrthreecone & ウ (U) & --- \\\hline
\end{prons}

\begin{remprons}
Many people on \comon{Youtube} \hooglig{have} \ucomon{ooky} \comkun{updates}.\\\hline
\end{remprons}

\begin{somme}
有る & \textit{to have} & あ{\middel}る (a{\middel}ru)\\\hline
有力 & have + strength = \textit{powerful} & ユウ{\middel}リョク (Y\={U}{\middel}RYOKU)\\\hline
有名 & have + name = \textit{famous} & ユウ{\middel}メイ (Y\={U}{\middel}MEI)\\\hline
私有 & private + have = \textit{privately owned} & シ{\middel}ユウ (SHI{\middel}Y\={U})\\\hline
公有 & public + have = \textit{publicly owned} & コウ{\middel}ユウ (K\={O}{\middel}Y\={U})\\\hline
国有化 & country + have + change = \textit{nationalization} & コク{\middel}ユウ{\middel}カ (KOKU{\middel}Y\={U}{\middel}KA)\\\hline
\end{somme}

\begin{sentence}
\underline{エジプト}の\underline{ピラミッド}\ruby{は}{わ}\ruby{\underline{有名}}{なまえ}です。\\\hline
Ejiputo no piramiddo wa Y\={U}{\middel}MEI desu.\\\hline
(Egypt) (pyramids) (famous)\\\hline
=\textit{The Egyptian pyramids are famous.}\\\hline
\end{sentence}
\end{multicols}
\vspace*{-\baselineskip}
\vhrulefill{5pt}
\vspace*{\fill}
%%--------------------------------------------------------------------
%%--------------------------------------------------------------------
\pagebreak
\vspace*{\fill}
\begin{multicols}{2}
\begin{glyph}{Meat (肉)}{meat}\label{glyph:meat}
Meat; Flesh.\\\hline
This kanji also has a sexual connotation through its appearance in words such as ``lust'' and ``sensuality''.
\end{glyph}

\begin{comps}
\comprtwocone & 内 + 人  \\\hline
\comprthreecone & Inside (\ref{glyph:inside}) + Person (\ref{glyph:person}) \\\hline
\end{comps}

\begin{remglyph}
There's \hooglig{meat} \remcom{inside} a \remcom{person}.\\\hline
\end{remglyph}

\begin{prons}
\pronsrtwocone & \textbf{ニク (NIKU)} & ---\\\hline
\pronsrthreecone & --- & --- \\\hline
\end{prons}

\begin{remprons}
The \hooglig{flesh} lusts after \comon{nicotine}.\\\hline
\end{remprons}

\begin{somme}
肉 & \textit{meat} & ニク (NIKU)\\\hline
人肉 & person + meat = \textit{human flesh} & ジン{\middel}ニク (JIN{\middel}NIKU)\\\hline
牛肉 & cow + meat = \textit{beef} & ギュウ{\middel}ニク (GY\={U}{\middel}NIKU)\\\hline
肉親 & meat + parent = \textit{blood relative} & ニク{\middel}シン (NIKU{\middel}SHIN)\\\hline
肉体 & meat + body = \textit{the body} & ニク{\middel}タイ (NIKU{\middel}TAI)\\\hline
\end{somme}

\begin{sentence}
\underline{あの} \ruby{\underline{四人}}{よニン}\ruby{は}{わ}\ruby{\underline{牛肉}}{ギュウニク}が\ruby{\underline{大好}}{ダイす}\underline{き}です。\\\hline
ano yo{\middel}NIN wa GY\={U}{\middel}NIKU ga DAIsu{\middel}ki desu.\\\hline
(those) (four people) (beef) (really like)\\\hline
=\textit{Those four people really like beef.}\\\hline
\end{sentence}
\end{multicols}
\vspace*{-\baselineskip}
\vhrulefill{5pt}
%%--------------------------------------------------------------------
%%--------------------------------------------------------------------
\begin{multicols}{2}
\begin{glyph}{Half (半)}{half}\label{glyph:half}
Half.\\\hline
This character also incorporates a sense of ``partial'', forming many words that English conveys with the prefix ``semi-''.
\end{glyph}

\begin{comps}
\comprtwocone & 丶丶 + 二 + 丨\\\hline
\comprthreecone & Dots + Two (\ref{glyph:two}) + Pole \\\hline
\end{comps}

\begin{remglyph}
A \remcom{pole} can \hooglig{half} a the \remcom{two dots}.\\\hline
\end{remglyph}

\begin{prons}
\pronsrtwocone & \textbf{ハン (HAN)} & ---\\\hline
\pronsrthreecone & --- & なか (naka) \\\hline
\end{prons}

\begin{remprons}
I have a \comon{hunch} he has a \ucomkun{knack} at \hooglig{halving} wood.\\\hline
\end{remprons}

\begin{somme}
上半 & upper + half = \textit{upper half} & ジョウ{\middel}ハン (J\={O}{\middel}HAN)\\\hline
下半 & lower + half = \textit{lower half} & カ{\middel}ハン (KA{\middel}HAN)\\\hline
半日 & half + sun (day) = \textit{half a day} & ハン{\middel}ニチ (HAN{\middel}NICHI)\\\hline
半分 & half + part = \textit{half} & ハン{\middel}ブン (HAN{\middel}BUN)\\\hline
九時半 & nine + time + half = \textit{nine-thirty} & キュウ{\middel}ジ{\middel}ハン (KY\={U}{\middel}JI{\middel}HAN)\\\hline
半円 & half + circle = \textit{semicircle} & ハン{\middel}エン (HAN{\middel}EN)\\\hline
\end{somme}

\begin{sentence}
\ruby{\underline{木口}}{きぐち}\underline{さん}\ruby{は}{わ}\ruby{\underline{半日}}{ハンニチ}{\;}\underline{ぐらい} \ruby{\underline{新聞}}{シンブン}\ruby{を}{お}\ruby{\underline{読}}{よ}\underline{みました}。\\\hline
Ki{\middel}guchi-san wa HAN{\middel}NICHI gurai SHIN{\middel}BUN o yo{\middel}mimashita.\\\hline
(Kiguchi-san) (half a day) (about) (newspaper) (read)\\\hline
=\textit{Kiguchi-san read the newspaper for about half a day.}\\\hline
\end{sentence}
\end{multicols}
\vspace*{-\baselineskip}
\vhrulefill{5pt}
\vspace*{\fill}
%%--------------------------------------------------------------------
%%--------------------------------------------------------------------
\pagebreak
\vspace*{\fill}
\begin{multicols}{2}
\begin{glyph}{Nine (九)}{nine}\label{glyph:nine}
Nine.
\end{glyph}

\begin{comps}
\comprtwocone & 九　\\\hline
\comprthreecone & Nine (\ref{glyph:nine}) \\\hline
\end{comps}

\begin{remglyph}
---\\\hline
\end{remglyph}

\begin{prons}
\pronsrtwocone & \textbf{キュウ、ク (KY\={U}, KU)} & ---\\\hline
\pronsrthreecone & --- & ここの (kokono) \\\hline
\rye{3}{キュウ is the general reading for this kanji.\\\hline
Only likely to encounter ク in the third and fourth examples below.}
\end{prons}

\begin{remprons}
\comon{Cooking} \ucomkun{coconut} at \hooglig{nine o'clock} will \comon{kill you}.\\\hline
\end{remprons}

\begin{somme}
九 & \textit{nine} & キュウ (KY\={U})\\\hline
九十 & nine + ten = \textit{ninety} & キュウ{\middel}ジュウ (KY\={U}{\middel}J\={U})\\\hline
九月 & nine + moon (month) = \textit{September} & ク{\middel}ガツ (KU{\middel}GATSU)\\\hline
九時 & nine + time = \textit{nine o'clock} & キュウ{\middel}ジ (KU{\middel}JI)\\\hline
\end{somme}

\begin{sentence}
\ruby{\underline{中山}}{なかやま}\underline{さん}\ruby{は}{わ}\ruby{\underline{九月}}{クガツ}に\ruby{\underline{日本}}{ニホン}\ruby{へ}{え}\ruby{\underline{行}}{い}\underline{きました}。\\\hline
Naka{\middel}yama-san wa KU{\middel}GATSU ni NI{\middel}HON e i{\middel}kimashita.\\\hline
(Nakayama-san) (September) (Japan) (went)\\\hline
=\textit{Nakayama-san went to Japan in September.}\\\hline
\end{sentence}
\end{multicols}
\vspace*{-\baselineskip}
\vhrulefill{5pt}
%%--------------------------------------------------------------------
%%--------------------------------------------------------------------
\begin{multicols}{2}
\begin{glyph}{Spring (春)}{spring}\label{glyph:spring}
Spring.\\\hline
This kanji also carries a hint of sexuality and amorousness---this is apparent by its presence in the Japanese words for puberty, etc.
\end{glyph}

\begin{comps}
\comprtwocone & 三 + 人 + 日　\\\hline
\comprthreecone & Three (\ref{glyph:three}) + Person (\ref{glyph:person}) + Sun (\ref{glyph:sun}) \\\hline
\end{comps}

\begin{remglyph}
\hooglig{Spring} is the \remcom{third} \remcom{personage} of the \remcom{sun}.\\\hline
\end{remglyph}

\begin{prons}
\pronsrtwocone & \textbf{シュン (SHUN)} & \textbf{はる (haru)}\\\hline
\pronsrthreecone & --- & --- \\\hline
\end{prons}

\begin{remprons}
\comon{Shun} the \hooglig{spring} \comkun{haruspex}.\\\hline
\end{remprons}

\begin{somme}
春 & \textit{spring} & はる (haru)\\\hline
小春 & small + spring = \textit{Indian Summer} & こ{\middel}はる (ko{\middel}haru)\\\hline
青春 & Blue + spring = \textit{youth} & セイ{\middel}シュン (SEI{\middel}SHUN)\\\hline
売春 & sell + spring = \textit{prostitution} & バイ{\middel}シュン (BAI{\middel}SHUN)\\\hline
春分 & spring + part = \textit{vernal equinox} & シュン{\middel}ブン (SHUN{\middel}BUN)\\\hline
来春 & come + spring = \textit{next spring} & ライ{\middel}シュン (RAI{\middel}SHUN)\\\hline
\end{somme}

\begin{sentence}
\ruby{\underline{日本}}{ニホン}の\ruby{\underline{春}}{はる}\ruby{は}{わ}\ruby{\underline{美}}{うつく}\underline{しい}。\\\hline
NI{\middel}HON no haru wa utsuku{\middel}shii.\\\hline
(Japan) (spring) (beautiful)\\\hline
=\textit{Japan's spring is beautiful.}\\\hline
\end{sentence}
\end{multicols}
\vspace*{-\baselineskip}
\vhrulefill{5pt}
\vspace*{\fill}
%%--------------------------------------------------------------------
%%--------------------------------------------------------------------
\pagebreak
%
%\begin{multicols}{2}
%\chapter{Kanji 41--60}
%\vspace*{\fill}
\begin{glyph}{Capital (京)}{capital}\label{glyph:capital}
Capital; Large metropolis.
\end{glyph}

\begin{comps}
\comprtwocone & 亠 + 口 + 小  \\\hline
\comprthreecone & Head/Above + Mouth (\ref{glyph:mouth}) + Small (\ref{glyph:small}) \\\hline
\end{comps}

\begin{remglyph}
\remcom{Small} \remcom{vampire} \remcom{heads} in the \hooglig{capital}.\\\hline
\end{remglyph}

\begin{prons}
\pronsrtwocone & \textbf{キョウ (KY\={O})} & ---\\\hline
\pronsrthreecone & ケイ (KEI) & --- \\\hline
\end{prons}

\begin{remprons}
The \comon{Kyoto} \hooglig{capital} \ucomon{caters} for all your needs.\\\hline
\end{remprons}

\begin{somme}
上京 & upper + capital = \textit{come/go to Tokyo} & ジョウ{\middel}キョウ (J\={O}{\middel}KY\={O})\\\hline
東京 & east + capital = \textit{Tokyo} & トウ{\middel}キョウ (T\={O}{\middel}KY\={O})\\\hline
\end{somme}

\begin{sentence}
\ruby{\underline{中山}}{なかやま}\underline{さん}\ruby{は}{わ}\ruby{\underline{上京}}{ジョウキョウ}しました。\\\hline
Naka{\middel}yama-san wa J\={O}{\middel}KY\={O} shimashita.\\\hline
(Nakayama-san) (went to Tokyo)\\\hline
=\textit{Nakayama-san went to Tokyo.}\\\hline
\end{sentence}
\end{multicols}
\vspace*{-\baselineskip}
\vhrulefill{5pt}
%%--------------------------------------------------------------------
%%--------------------------------------------------------------------
\begin{multicols}{2}
\begin{glyph}{See (見)}{see}\label{glyph:see}
See; Watch.\\\hline
The idea of ``show'' can also be expressed.
\end{glyph}

\begin{comps}
\comprtwocone & 目 + 儿  \\\hline
    \comprthreecone & Eye (\ref{glyph:eye}) + man/a person (at the bottom of a
    character)\\\hline
\end{comps}

\begin{remglyph}
    \remcom{People} with \remcom{eyes} can \hooglig{see}.\\\hline
\end{remglyph}

\begin{prons}
\pronsrtwocone & \textbf{ケン (KEN)} & \textbf{み (mi)}\\\hline
\pronsrthreecone & --- & --- \\\hline
\end{prons}

\begin{remprons}
I \hooglig{see} the \comon{kennels} have \comkun{mirrors}.\\\hline
\end{remprons}

\begin{somme}
見える (intr) & \textit(to be visible) & み{\middel}える (mi{\middel}eru)\\\hline
見る (tr) & \textit{to see; to watch)} & み{\middel}る (mi{\middel}ru)\\\hline
見せる (tr) & \textit{to show (something)} & み{\middel}せる (mi{\middel}seru)\\\hline
見上げる & see + upper = \textit{to look up (to/at)} & み　あ{\middel}げる (mi a{\middel}geru)\\\hline
月見 & moon + see = \textit{moon viewing} & つき{\middel}み (tsuki{\middel}mi)\\\hline
見失う & see + lose = \textit{to lose sight of} & み{\middel}うしな{\middel}う (mi{\middel}ushina{\middel}u)\\\hline
外見 & outside + see = \textit{external appearance} & ガイ{\middel}ケン (GAI{\middel}KEN)\\\hline
\end{somme}

\begin{sentence}
\ruby{\underline{山}}{やま}の\ruby{\underline{上}}{うえ}に \ruby{\underline{小}}{ちい}\underline{さい}{\;}\ruby{\underline{木}}{き}が\ruby{\underline{見}}{み}\underline{えます}。\\\hline
yama no ue ni chii{\middel}sai ki ga mi{\middel}emasu.\\\hline
(mountain) (upper) (small) (tree) (is visible)\\\hline
=\textit{A small tree is visible on top of the mountain.}\\\hline
\end{sentence}
\end{multicols}
\vspace*{-\baselineskip}
\vhrulefill{5pt}
\vspace*{\fill}
%%--------------------------------------------------------------------
%%--------------------------------------------------------------------
\pagebreak
\vspace*{\fill}
\begin{multicols}{2}
\begin{glyph}{Stop (止)}{stop}\label{glyph:stop}
Stop.
\end{glyph}

\begin{comps}
\comprtwocone & ⺊ + 丨 + 一\\\hline
\comprthreecone & Divination + Pole + One (\ref{glyph:one}) \\\hline
\end{comps}

\begin{remglyph}
    \hooglig{Stop} using that \remcom{divination pole} at \remcom{once}.\\\hline
\end{remglyph}

\begin{prons}
\pronsrtwocone & \textbf{シ (SHI)} & \textbf{と (to)}\\\hline
\pronsrthreecone & --- & --- \\\hline
\end{prons}

\begin{remprons}
\hooglig{Stop} the \comkun{torrent} with a \comon{shield}.\\\hline
\end{remprons}

\begin{somme}
止まる (intr) & \textit{to stop} & と{\middel}まる (to{\middel}maru)\\\hline
止める (tr) & \textit{to stop (an object)} & と{\middel}める (to{\middel}meru)\\\hline
中止 & middle + stop = \textit{discontinue} & チュウ{\middel}シ (CH\={U}{\middel}SHI)\\\hline
立ち止まる & stand + stop = \textit{discontinue} & た{\middel}ち　ど{\middel}まる (ta{\middel}chi do{\middel}maru)\\\hline
車止め & car + stop = \textit{closed to vehicles} & くるま　ど{\middel}め (kuruma do{\middel}me)\\\hline
\end{somme}

\begin{sentence}
\underline{こちら}に\ruby{\underline{止}}{と}\underline{まって} \ruby{\underline{下}}{くだ}\underline{さい}。\\\hline
kochira ni to{\middel}matte kuda{\middel}sai.\\\hline
(here) (stop) (please)\\\hline
=\textit{Please stop here.}\\\hline
\end{sentence}
\end{multicols}
\vspace*{-\baselineskip}
\vhrulefill{5pt}
%%--------------------------------------------------------------------
%%--------------------------------------------------------------------
\begin{multicols}{2}
\begin{glyph}{Seven (七)}{seven}\label{glyph:seven}
Seven.
\end{glyph}

\begin{comps}
\comprtwocone & 一 + 亅\\\hline
\comprthreecone & One + Hook \\\hline
\end{comps}

\begin{remglyph}
    \remcom{Captain Hook} sailed the \hooglig{seven} seas \remcom{once} in his entire life.\\\hline
\end{remglyph}

\begin{prons}
\pronsrtwocone & \textbf{シチ (SHICHI)} & \textbf{なな (nana)}\\\hline
\pronsrthreecone & --- & なの (nano) \\\hline
\end{prons}

\begin{remprons}
\comon{She cheated} the \comkun{nana} with \hooglig{seven} \ucomkun{nanoparticles}.\\\hline
\end{remprons}

\begin{somme}
七 & \textit{seven} & シチ (SHICHI)\\\hline
七つ & \textit{seven (general counter)} & なな{\middel}つ (nana{\middel}tsu)\\\hline
七月 & seven + moon (month) \textit{July} & シチ{\middel}ガツ (SHICHI{\middel}GATSU)\\\hline
十七 & ten + seven = \textit{seventeen} & ジュウ{\middel}なな (J\={U}{\middel}nana)\\\hline
七十 & seven + ten = \textit{seventy} & なな{\middel}ジュウ (nana{\middel}J\={U})\\\hline
七時 & seven + time = \textit{seven o'clock} & シチ{\middel}ジ (SHICHI{\middel}JI)\\\hline
\end{somme}

\begin{sentence}
\ruby{\underline{高木}}{たかぎ}\underline{さん}が\ruby{\underline{七月}}{シチガツ}に\ruby{\underline{中米}}{チュウベイ}から\ruby{\underline{来}}{き}\underline{ます}。\\\hline
Taka{\middel}gi-san ga SHICHI{\middel}GATSU ni CH\={U}{\middel}BEI kara ki{\middel}masu.\\\hline
(Takagi-san) (July) (Central America) (will come)\\\hline
=\textit{Takagi-san will come from Central America in July.}\\\hline
\end{sentence}
\end{multicols}
\vspace*{-\baselineskip}
\vhrulefill{5pt}
\vspace*{\fill}
%%--------------------------------------------------------------------
%%--------------------------------------------------------------------
\pagebreak
\vspace*{\fill}
\begin{multicols}{2}
\begin{glyph}{Few (少)}{few}\label{glyph:few}
Few; Little (in terms of quantity).
\end{glyph}

\begin{comps}
\comprtwocone & 小 + ノ \\\hline
    \comprthreecone & Small (\ref{glyph:small}) + NO (katakana)\\\hline
\end{comps}

\begin{remglyph}
A \hooglig{few} occurences of \remcom{small} \remcom{noses}.\\\hline
\end{remglyph}

\begin{prons}
\pronsrtwocone & \textbf{ショウ (SH\={O})} & \textbf{すく、すこ (suku, suko)}\\\hline
\pronsrthreecone & --- & --- \\\hline
\rye{3}{ショウ will appear in compounds.\\\hline
すく and すこ are only found in the first and second examples below.}
\end{prons}

\begin{remprons}
He brought a \hooglig{few} \comkun{scones} with his \comkun{scooter} to the \comon{shore}.\\\hline
\end{remprons}

\begin{somme}
少ない & \textit{few} & すく{\middel}ない (suku{\middel}nai)\\\hline
少し & \textit{a little; a few} & すこ{\middel}し (suko{\middel}shi)\\\hline
少々 & few + few = \textit{a little} & ショウ{\middel}ショウ (SH\={O}{\middel}SH\={O})\\\hline
少女 & few + woman = \textit{girl} & ショウ{\middel}ジョ (SH\={O}{\middel}JO)\\\hline
少年 & few + year = \textit{boy} & ショウ{\middel}ネン (SH\={O}{\middel}NEN)\\\hline
最小 & most + few = \textit{fewest} & サイ{\middel}ショウ (SAI{\middel}SH\={O})\\\hline
\end{somme}

\begin{sentence}
\underline{この} \ruby{\underline{家}}{いえ}に\ruby{は}{わ}\ruby{\underline{物}}{もの}が\ruby{\underline{少}}{すく}\underline{ない}。\\\hline
kono ie ni wa mono ga suku{\middel}nai.\\\hline
(this) (house) (thing) (few)\\\hline
=\textit{There are few things in this house.}\\\hline
\end{sentence}
\end{multicols}
\vspace*{-\baselineskip}
\vhrulefill{5pt}
%%--------------------------------------------------------------------
%%--------------------------------------------------------------------
\begin{multicols}{2}
\begin{glyph}{South (南)}{south}\label{glyph:south}
South.
\end{glyph}

\begin{comps}
\comprtwocone & 十 + 冂 + 䒑 + 十\\\hline
\comprthreecone & Ten (\ref{glyph:ten}) + Wilderness + Grass + Ten (\ref{glyph:ten})\\\hline
\end{comps}

\begin{remglyph}
    The \hooglig{south} \remcom{wilderness} do not even have \remcom{ten}
    \remcom{grass} blades.\\\hline
\end{remglyph}

\begin{prons}
\pronsrtwocone & \textbf{ナン (NAN)} & \textbf{みなみ (minami)}\\\hline
\pronsrthreecone & --- & --- \\\hline
\end{prons}

\begin{remprons}
The \hooglig{southern} \comon{nun} at a \comkun{mean of a meal}.\\\hline
\end{remprons}

\begin{somme}
南 & \textit{south} & みなみ (minami)\\\hline
南口 & south + mouth = \textit{southern exit/entrance} & みなみ{\middel}ぐち (minami{\middel}guchi)\\\hline
南米 & south + rice = \textit{Southern America} & ナン{\middel}ベイ (NAN{\middel}BEI)\\\hline
最南 & most + south = \textit{southernmost} & サイ{\middel}ナン (SAI{\middel}NAN)\\\hline
南西 & south + west = \textit{southwest} & ナン{\middel}セイ (NAN{\middel}SEI)\\\hline
南東 & south + east = \textit{south east} & ナン{\middel}トウ (NAN{\middel}T\={O})\\\hline
\end{somme}

\begin{sentence}
\underline{バスターミナル}の\ruby{\underline{南口}}{みなみぐし}で\ruby{\underline{会}}{あ}\underline{いましょう}。\\\hline
basut\={a}minaru no minami{\middel}guchi de a{\middel}imash\={o}.\\\hline
(bus terminal) (south entrance) (let's meet)\\\hline
=\textit{Let's meet at the south entrance of the bus terminal.}\\\hline
\end{sentence}
\end{multicols}
\vspace*{-\baselineskip}
\vhrulefill{5pt}
\vspace*{\fill}
%%--------------------------------------------------------------------
%%--------------------------------------------------------------------
\pagebreak
\vspace*{\fill}
\begin{multicols}{2}
\begin{glyph}{Craft (工)}{craft}\label{glyph:craft}
Carft; Construction; Manufacturing; Handicrafts; Industrial products.\\\hline
When found at the end of a compound (as in the final example below), this kanji can take on the meaning of ``worker'', and usually implies a person involved in some part of manual labor.
\end{glyph}

\begin{comps}
\comprtwocone & 工\\\hline
\comprthreecone & Craft (\ref{glyph:craft})\\\hline
\end{comps}

\begin{remglyph}
Anvil.\\\hline
\end{remglyph}

\begin{prons}
\pronsrtwocone & \textbf{コウ (K\={O})} & ---\\\hline
\pronsrthreecone & ク (KU) & --- \\\hline
\end{prons}

\begin{remprons}
His \ucomon{cooking} brilliance \comon{caused} me to refine my own \hooglig{craft}.\\\hline
\end{remprons}

\begin{somme}
人工 & person + craft = \textit{artificial}  & ジン{\middel}コウ (JIN{\middel}K\={O})\\\hline
人工林 & person + craft + grove = \textit{planted forest} & ジン{\middel}コウ{\middel}リン (JIN{\middel}K\={O}{\middel}RIN)\\\hline
工学 & craft + study = \textit{engineering} & コウ{\middel}ガク (K\={O}{\middel}GAKU)\\\hline
工学士 & craft + study + gentleman = \textit{Bachelor of Engineering} & コウ{\middel}ガク{\middel}シ (K\={O}{\middel}GAKU{\middel}SHI)\\\hline
刀工 & sword + craft = \textit{swordmaker} & トウ{\middel}コウ (T\={O}{\middel}K\={O})\\\hline
\end{somme}

\begin{sentence}
\underline{これ}\ruby{は}{わ}\ruby{\underline{人工}}{ジンコウ}の\ruby{\underline{島}}{しま}{\;}\underline{ですね}。\\\hline
kore wa JIN{\middel}K\={O} no shima desu ne.\\\hline
(this) (artificial) (island) (isn't it)\\\hline
=\textit{This is an artificial island, isn't it?}\\\hline
\end{sentence}
\end{multicols}
\vspace*{-\baselineskip}
\vhrulefill{5pt}
%%--------------------------------------------------------------------
%%--------------------------------------------------------------------
\begin{multicols}{2}
\begin{glyph}{Left (左)}{hidari}\label{glyph:left}
Left; Left-hand side.
\end{glyph}

\begin{comps}
\comprtwocone & 一 + ノ + 工\\\hline
\comprthreecone & One (\ref{glyph:one}) + NO ノ Katakana + Craft (\ref{glyph:craft})\\\hline
\end{comps}

\begin{remglyph}
The \remcom{superhero} \hooglig{left} the \remcom{anvil} on the monster's head.\\\hline
\end{remglyph}

\begin{prons}
\pronsrtwocone & --- & \textbf{ひだり (hidari)}\\\hline
\pronsrthreecone & サ (SA) & --- \\\hline
\end{prons}

\begin{remprons}
We \hooglig{left} the \comon{subject} of where the \comkun{ring was hid}.\\\hline
\end{remprons}

\begin{somme}
左 & \textit{left} & ひだり (hidari)\\\hline
左回り & left + rotate = \textit{counter-clockwise} & ひだり　まわ{\middel}り (hidari mawa{\middel}ri)\\\hline
左上 & left + upper = \textit{upper left} & ひだり{\middel}うえ (hidari{\middel}ue)\\\hline
左下 & left + lower = \textit{lower left} & ひだり{\middel}した (hidari{\middel}shita)\\\hline
左手 & left + hand = \textit{left hand} & ひだり{\middel}て (hidari{\middel}te)\\\hline
左足 & left + leg = \textit{left leg} & ひだり{\middel}あし (hidari{\middel}ashi)\\\hline
\end{somme}

\begin{sentence}
\underline{あの} \ruby{\underline{人}}{ひと}の\ruby{\underline{左手}}{ひだりて}\ruby{を}{お}\ruby{\underline{見}}{み}\underline{て} \ruby{\underline{下}}{くだ}\underline{さい}。\\\hline
ano hito no hidari{\middel}te o mi{\middel}te kuda{\middel}sai.\\\hline
(that) (person) (left hand) (see) (please)\\\hline
=\textit{Please watch that person's left hand.}\\\hline
\end{sentence}
\end{multicols}
\vspace*{-\baselineskip}
\vhrulefill{5pt}
\vspace*{\fill}
%%--------------------------------------------------------------------
%%--------------------------------------------------------------------
\pagebreak
\begin{multicols}{2}
\begin{glyph}{Tall (高)}{tall}\label{glyph:tall}
Tall; High; Best.\\\hline
This kanji is also used to express the idea of expensive.
\end{glyph}

\begin{comps}
\comprtwocone & 亠 + 口 + 冂 + 口 \\\hline
\comprthreecone & above/head + Mouth (\ref{glyph:mouth}) + wilderness + Mouth(\ref{glyph:mouth})\\\hline
\end{comps}

\begin{remglyph}
    The \hooglig{tall} \remcom{vampire's} \remcom{head} was cast into the
    \remcom{wilderness}.\\\hline
\end{remglyph}

\begin{prons}
\pronsrtwocone & \textbf{コウ (K\={O})} & \textbf{たか (taka)}\\\hline
\pronsrthreecone & --- & --- \\\hline
\end{prons}

\begin{remprons}
I \comkun{tucked} the \hooglig{tall} stash under the bed \comon{'cause} I was afraid.\\\hline
\end{remprons}

\begin{somme}
高い & \textit{tall; expensive} & たか{\middel}い (taka{\middel}i)\\\hline
高まる (intr) & \textit{to rise} & たか{\middel}まる (taka{\middel}maru)\\\hline
高める (tr) & \textit{to raise} & たか{\middel}める (taka{\middel}meru)\\\hline
高山 & tall + mountain = \textit{alpine} & コウ{\middel}ザン (K\={O}{\middel}ZAN)\\\hline
高校 & tall + school = \textit{high school} & コウ{\middel}コウ (K\={O}{\middel}K\={O})\\\hline
円高 & circle (yen) + tall = \textit{strong yen} & エン{\middel}だか (EN{\middel}daka)\\\hline
最高 & most + tall = \textit{highest; the best} & サイ{\middel}コウ (SAI{\middel}K\={O})\\\hline
\end{somme}

\begin{sentence}
\underline{あの} \ruby{\underline{高}}{たか}\underline{い} \ruby{\underline{山}}{やま}\ruby{は}{わ}\ruby{\underline{美}}{うつく}\underline{しい}{\;}\underline{ですね}。\\\hline
ano takai yama wa utsuku{\middel}shii desu ne.\\\hline
(that) (tall) (mountain) (beautiful) (isn't it)\\\hline
=\textit{That tall mountain is beautiful, isn't it.}\\\hline
\end{sentence}
\end{multicols}
\vspace*{-\baselineskip}
\vhrulefill{5pt}
%%--------------------------------------------------------------------
%%--------------------------------------------------------------------
\begin{multicols}{2}
\begin{glyph}{Buy (買)}{buy}\label{glyph:buy}
Buy.
\end{glyph}

\begin{comps}
\comprtwocone & 罒 + 貝 \\\hline
\comprthreecone & net + Shellfish (\ref{glyph:shellfish})\\\hline
\end{comps}

\begin{remglyph}
\hooglig{Buy} the \remcom{shellfish} freshly caught in the \remcom{net}.\\\hline
\end{remglyph}

\begin{prons}
\pronsrtwocone & --- & \textbf{か (ka)}\\\hline
\pronsrthreecone & バイ (BAI) & --- \\\hline
\end{prons}

\begin{remprons}
A \comkun{customary} \hooglig{\ucomon{buy}}.\\\hline
\end{remprons}

\begin{somme}
買う & \textit{to buy} & か{\middel}う (ka{\middel}u)\\\hline
買い上げる & buy + upper = \textit{to buy (up/out)} & か{\middel}い　あ{\middel}げる (ka{\middel}i a{\middel}geru)\\\hline
買上げ & buy + upper = \textit{a purchase} & かい　あ{\middel}げ (kai a{\middel}ge)\\\hline
買い物 & buy + thing = \textit{shopping} & か{\middel}い　もの (ka{\middel}i mono)\\\hline
買い入れる & buy + enter = \textit{to stock up on} & か{\middel}い　い{\middel}れる (ka{\middel}i i{\middel}reru)\\\hline
買い手 & buy + hand = \textit{buyer} & か{\middel}い　て (ka{\middel}i te)\\\hline
\rye{3}{By comparing the second and third examples above, we se that the noun in example 3 has ``lost'' the い that would be present if the kanji were in its -ます verb form, かいます (買います).  It is in a sense implied, based on the hiragana ending after 上.  Unfortunately, there are no hard and fast rules for this aspect of the language, and it can be frustrating when such words are encountered.  Example 3 can be seen variously as 買い上げ、買上げ、or even 買上 in different dictionaries.\\\hline
Nonetheless, not many kanji present with such difficulties, and once a `primary' word (usually the -ます form of the verb) has been learned, its related nounds will be easily recognized.}
\end{somme}

\begin{sentence}
\ruby{\underline{山中}}{やまなか}\underline{さん}\ruby{は}{わ}\ruby{\underline{高}}{たか}\underline{い} \ruby{\underline{物}}{もの}\ruby{を}{お}\ruby{\underline{買}}{か}\underline{う} \underline{つもり}です。\\\hline
Yama{\middel}naka-san wa taka{\middel}i mono o ka{\middel}u tsumori desu.\\\hline
(Yamanaka-san) (expensive) (thing) (buy) (plans)\\\hline
=\textit{Yamanaka-san plans to buy something expensive.}\\\hline
\end{sentence}
\end{multicols}
\vspace*{-\baselineskip}
\vhrulefill{5pt}
%%--------------------------------------------------------------------
%%--------------------------------------------------------------------
\pagebreak
\vspace*{\fill}
\begin{multicols}{2}
\begin{glyph}{Hundred (百)}{hundred}\label{glyph:hundred}
Hundred.\\\hline
In a few instances, the idea of `many' is conveyed.
\end{glyph}

\begin{comps}
\comprtwocone & 一 + 白 \\\hline
\comprthreecone & One (\ref{glyph:one}) + White (\ref{glyph:white})\\\hline
\end{comps}

\begin{remglyph}
\hooglig{Hundred} percent of \remcom{top buns} are \remcom{white}.\\\hline
\end{remglyph}

\begin{prons}
\pronsrtwocone & \textbf{ヒャク (HYAKU)} & ---\\\hline
\pronsrthreecone & --- & --- \\\hline
\end{prons}

\begin{remprons}
\hooglig{Hundred} percent \comon{yucky}.\\\hline
\end{remprons}

\begin{somme}
百 & \textit{one hundred} & ヒャク (HYAKU)\\\hline
二百 & two + hundred = \textit{two hundred} & ニ{\middel}ヒャク (NI{\middel}HYAKU)\\\hline
三百 & three + hundred = \textit{three hundred} & サン{\middel}ビャク (SAN{\middel}BYAKU)\\\hline
四百 & four + hundred = \textit{four hundred} & よん{\middel}ヒャク (yon{\middel}HYAKU)\\\hline
五百 & five + hundred = \textit{five hundred} & ゴ{\middel}ヒャク (GO{\middel}HYAKU)\\\hline
六百 & six + hundred = \textit{six hundred} & ロク{\middel}ヒャク (ROKU{\middel}HYAKU)\\\hline
七百 & seven + hundred = \textit{seven hundred} & なな{\middel}ヒャク (nana{\middel}HYAKU)\\\hline
八百 & eight + hundred = \textit{eight hundred} & ハチ{\middel}ヒャク (HACHI{\middel}HYAKU)\\\hline
九百 & nine + hundred = \textit{nine hundred} & キュウ{\middel}ヒャク (KY\={U}{\middel}HYAKU)\\\hline
\end{somme}

\begin{sentence}
\ruby{\underline{肉}}{ニク}\ruby{を}{お}\ruby{\underline{五百}}{ゴヒャク}{\;}\underline{グラム}{\;}\ruby{\underline{買}}{か}\underline{いました}。\\\hline
NIKU o GO{\middel}HYAKU guramu ka{\middel}imashita.\\\hline
(meat) (five hundred) (grams) (bought)\\\hline
=\textit{(I) bought five hundred grams of meat.}\\\hline
\end{sentence}
\end{multicols}
\vspace*{-\baselineskip}
\vhrulefill{5pt}
%%--------------------------------------------------------------------
%%--------------------------------------------------------------------
\begin{multicols}{2}
\begin{glyph}{Circle (円)}{circle}\label{glyph:circle}
Circle; Round.\\\hline
This kanji is also used to denote the Yen, Japan's Currency.  The connection originated from round one-yen coins.
\end{glyph}

\begin{comps}
\comprtwocone & 冂 + 亠 \\\hline
\comprthreecone & Wilderness + above/head\\\hline
\end{comps}

\begin{remglyph}
    \remcom{Above} the \remcom{wilderness} is a \hooglig{circular} sun.\\\hline
\end{remglyph}

\begin{prons}
\pronsrtwocone & \textbf{エン (EN)} & \textbf{まる (maru)}\\\hline
\pronsrthreecone & --- & --- \\\hline
\end{prons}

\begin{remprons}
The \hooglig{round} \comkun{Maru} cat made quite an \comon{entrance}.\\\hline
\end{remprons}

\begin{somme}
円い & \textit{circular} & まる{\middel}い (maru{\middel}i)\\\hline
円高 & yen + tall = \textit{strong yen} & エン{\middel}だか (EN{\middel}daka)\\\hline
二円 & two + yen = \textit{two yen} & ニ{\middel}エン (NI{\middel}EN)\\\hline
半円 & haf + circle = \textit{semi-circle} & ハン{\middel}エン (HAN{\middel}EN)\\\hline
円安 & yen + ease = \textit{weak yen} & エン{\middel}やす (EN{\middel}yasu)\\\hline
円周 & circle + around = \textit{circumference} & エン{\middel}シュウ (EN{\middel}SH\={U})\\\hline
\end{somme}

\begin{sentence}
\underline{あの} \ruby{\underline{子}}{こ}\ruby{は}{わ}\ruby{\underline{百円}}{ヒャクエン}\ruby{を}{お}\underline{なくした}。\\\hline
ano ko wa HYAKU{\middel}EN o nakushita.\\\hline
(that) (child) (one hundred yen) (lost)\\\hline
=\textit{That child lost a hundred yen.}\\\hline
\end{sentence}
\end{multicols}
\vspace*{-\baselineskip}
\vhrulefill{5pt}
\vspace*{\fill}
%%--------------------------------------------------------------------
%%--------------------------------------------------------------------
\pagebreak
\vspace*{\fill}
\begin{multicols}{2}
\begin{glyph}{Basis (元)}{basis}\label{glyph:basis}
This character expresses the idea of a basis or the original state of something.\\\hline
When used with its kun-yomi before words such as ``president'', it can be translated as ``ex'' or ``former''.
\end{glyph}

\begin{comps}
\comprtwocone & 二 + 儿 \\\hline
    \comprthreecone & Two (\ref{glyph:two}) + Man/A person (at the bottom of a
    character)\\\hline
\end{comps}

\begin{remglyph}
    \remcom{Two} \remcom{people} are the \hooglig{basis} of any
    relationship.\\\hline
\end{remglyph}

\begin{prons}
\pronsrtwocone & \textbf{ゲン、ガン (GEN, GAN)} & \textbf{もと (moto)}\\\hline
\pronsrthreecone & --- & --- \\\hline
\rye{3}{ガン only occurs in the first position with a few common words.}
\end{prons}

\begin{remprons}
\comon{Genghis Khan} holding the \comkun{motorists} at \comon{gunpoint} had no \hooglig{basis}.\\\hline
\end{remprons}

\begin{somme}
元々 & basis + basis = \textit{from the first} & もと{\middel}もと (moto{\middel}moto)\\\hline
元日 & basis + sun (day) = \textit{New Year's Day} & ガン{\middel}ジツ (GAN{\middel}JITSU)\\\hline
手元 & hand + basis = \textit{at hand} & て{\middel}もと (te{\middel}moto)\\\hline
元気 & basis + spirit = \textit{high spirits} & ゲン{\middel}キ (GEN{\middel}KI)\\\hline
火元 & fire + basis = \textit{origin of a fire} & ひ{\middel}もと (hi{\middel}moto)\\\hline
元来 & basis + come = \textit{by nature} & ガン{\middel}ライ (GAN{\middel}RAI)\\\hline
\end{somme}

\begin{sentence}
\ruby{\underline{元日}}{ガンジツ}な\underline{の}で\ruby{\underline{人}}{ひと}が\ruby{\underline{少}}{すく}\underline{ない}。\\\hline
GAN{\middel}JITSU na no de hito ga suku{\middel}nai.\\\hline
(New Year's Day) (because) (person) (few).\\\hline
=\textit{It's New Year's Day, so there aren't many people.}\\\hline
\end{sentence}
\end{multicols}
\vspace*{-\baselineskip}
\vhrulefill{5pt}
%%--------------------------------------------------------------------
%%--------------------------------------------------------------------
\begin{multicols}{2}
\begin{glyph}{Neck (首)}{neck}\label{glyph:neck}
Neck in the physical sense.\\\hline
As well as the ideas of chief/leader.
\end{glyph}

\begin{comps}
\comprtwocone & 䒑 + 自 \\\hline
\comprthreecone & grass + Self (\ref{glyph:neck})\\\hline
\end{comps}

\begin{remglyph}
I want to buy a \remcom{grass} \hooglig{necklace} \remcom{myself}.\\\hline
\end{remglyph}

\begin{prons}
\pronsrtwocone & \textbf{シュ (SHU)} & \textbf{くび (kubi)}\\\hline
\pronsrthreecone & --- & --- \\\hline
\end{prons}

\begin{remprons}
\comon{Shooting} \comkun{cubic} \hooglig{necks}.\\\hline
\end{remprons}

\begin{somme}
首 & \textit{neck} & くび (kubi)\\\hline
自首 & self + neck = \textit{surrender} & ジ{\middel}シュ (JI{\middel}SHU)\\\hline
手首 & hand + neck = \textit{wrist} & て{\middel}くび (te{\middel}kubi)\\\hline
足首 & leg + neck = \textit{ankle} & あし{\middel}くび (ashi{\middel}kubi)\\\hline
\end{somme}

\begin{sentence}
\ruby{\underline{高木}}{たかぎ}\underline{さん}\ruby{は}{わ}\ruby{\underline{首}}{くび}に\ruby{\underline{水}}{みず}\ruby{を}{お}\underline{かけました}。\\\hline
Takagi-san wa kubi ni mizu o kakemashita\\\hline
(Takagi-san) (neck) (water) (put)\\\hline
=\textit{Takagi-san put some water on his neck.}\\\hline
\end{sentence}
\end{multicols}
\vspace*{-\baselineskip}
\vhrulefill{5pt}
\vspace*{\fill}
%%--------------------------------------------------------------------
%%--------------------------------------------------------------------
\pagebreak
\vspace*{\fill}
\begin{multicols}{2}
\begin{glyph}{Walk (歩)}{walk}\label{glyph:walk}
Walk; Step; Pace.\\\hline
A secondary meaning of ``rate'' (as used in the financial sense) is found only with ブ, a less common on-yomi.
\end{glyph}

\begin{comps}
\comprtwocone & 止 + 少 \\\hline
\comprthreecone & Stop (\ref{glyph:stop}) + Few (\ref{glyph:small})\\\hline
\end{comps}

\begin{remglyph}
There are \remcom{few} \remcom{stops} in \hooglig{walking}.\\\hline
\end{remglyph}

\begin{prons}
\pronsrtwocone & \textbf{ホ (HO)} & \textbf{ある (aru)}\\\hline
\pronsrthreecone & ブ、フ (BU, FU) & あゆ (ayu) \\\hline
\end{prons}

\begin{remprons}
I \hooglig{walked} to the \comon{hospital} to receive \ucomkun{ayurvedic} treatment since I was \ucomon{bullied} to eat \comkun{arugula} which turned out to be a \ucomon{foofy}.\\\hline
\end{remprons}

\begin{somme}
歩く & \textit{to walk} & ある{\middel}く (aru{\middel}ku)\\\hline
一歩 & one + walk = \textit{one step} & イッ{\middel}ポ ()\\\hline
歩き回る & walk + rotate = \textit{to ramble about} & ある{\middel}き　まわ{\middel}る (aru{\middel}ki mawa{\middel}ru)\\\hline
歩道 & walk + road = \textit{sidewalk} & ホ{\middel}ドウ (HO{\middel}D\={O})\\\hline
歩行 & walk + go = \textit{walking} & ホ{\middel}コウ (HO{\middel}K\={O})\\\hline
歩行者 & walk + go + individual = \textit{pedestrian} & ホ{\middel}コウ{\middel}シャ (HO{\middel}K\={O}{\middel}SHA)\\\hline
\end{somme}

\begin{sentence}
\ruby{\underline{家}}{いえ}{\;}\underline{から}{\;}\ruby{\underline{山}}{やま}{\;}\underline{まで} \ruby{\underline{歩}}{ある}\underline{きました}。\\\hline
ie kara yama made aru{\middel}kimashita.\\\hline
(house) (from) (mountain) (until) (walked)\\\hline
=\textit{(We) walked from the house to the mountain.}\\\hline
\end{sentence}
\end{multicols}
\vspace*{-\baselineskip}
\vhrulefill{5pt}
%%--------------------------------------------------------------------
%%--------------------------------------------------------------------
\begin{multicols}{2}
\begin{glyph}{Child (子)}{child}\label{glyph:child}
Child.\\\hline
This character can also be used to suggest something small.
\end{glyph}

\begin{comps}
\comprtwocone & 子 \\\hline
\comprthreecone & Child (\ref{glyph:child})\\\hline
\end{comps}

\begin{remglyph}
---\\\hline
\end{remglyph}

\begin{prons}
\pronsrtwocone & \textbf{シ (SHI)} & \textbf{こ (ko)}\\\hline
\pronsrthreecone & ス (SU) & --- \\\hline
\rye{3}{Both of these readings can at times become voiced in the second position.\\\hline
Japanese parents often choose this kanji for a suffix when naming their children.  It is read シ in the case of males, and こ when applied to females.}
\end{prons}

\begin{remprons}
It \ucomon{suited} the \hooglig{child} to be \comon{shielded} on the \comkun{cobbled} road.\\\hline
\end{remprons}

\begin{somme}
子 & \textit{child} & こ (ko)\\\hline
女子 & woman + child = \textit{women's} & ジョ{\middel}シ (JO{\middel}SHI)\\\hline
男子 & man + child = \textit{men's} & ダン{\middel}シ (DAN{\middel}SHI)\\\hline
王子 & king + child = \textit{prince} & オウ{\middel}ジ (\={O}{\middel}JI)\\\hline
子牛 & child + cow = \textit{calf} & こ{\middel}うし (ko{\middel}ushi)\\\hline
電子 & electric + child = \textit{electron} & デン{\middel}シ (DEN{\middel}SHI)\\\hline
\end{somme}

\begin{sentence}
\ruby{\underline{五月}}{ゴガツ}に\ruby{\underline{女王}}{ジョオウ}と\ruby{\underline{王子}}{オウジ}が\ruby{\underline{米国}}{ベイコク}\ruby{へ}{え}\ruby{\underline{行}}{い}\underline{きました}。\\\hline
GO{\middel}GATSU ni JO{\middel}\={O} to \={O}{\middel}JI ga BEI{\middel}KOKU e i{\middel}kimashita.\\\hline
(May) (queen) (prince) (America) (went)\\\hline
=\textit{In May, the queen and prince went to America.}\\\hline
\end{sentence}
\end{multicols}
\vspace*{-\baselineskip}
\vhrulefill{5pt}
\vspace*{\fill}
%%--------------------------------------------------------------------
%%--------------------------------------------------------------------
\pagebreak
\vspace*{\fill}
\begin{multicols}{2}
\begin{glyph}{Like (好)}{like}\label{glyph:like}
Like; Good; Favorable.\\\hline
Nothing but positive feelings are associated with this character.
\end{glyph}

\begin{comps}
\comprtwocone & 女 + 子 \\\hline
    \comprthreecone & Woman (\ref{glyph:woman}) +  Child (\ref{glyph:child})\\\hline
\end{comps}

\begin{remglyph}
Woman likes child.\\\hline
\end{remglyph}

\begin{prons}
\pronsrtwocone & \textbf{コウ (K\={O})} & \textbf{す (su)}\\\hline
\pronsrthreecone & --- & この (kono) \\\hline
\end{prons}

\begin{remprons}
I \hooglig{like} \comkun{sudoku} because it \ucomkun{connotes} \comon{causality}.\\\hline
\end{remprons}

\begin{somme}
好き & \textit{like} & す{\middel}き (su{\middel}ki)\\\hline
大好き & large + like = \textit{to really like} & ダイ　す{\middel}き (DAI su{\middel}ki)\\\hline
人好き & person + like = \textit{amiability} & ひと　ず{\middel}き (hito zu{\middel}ki)\\\hline
話し好き & speak + like = \textit{talkative} & はな{\middel}し　ず{\middel}き (hana{\middel}shi zu{\middel}ki)\\\hline
友好 & friend + like = \textit{friendship} & ユウ{\middel}コウ (Y\={U}{\middel}K\={O})\\\hline
同好 & same + like = \textit{same tastes} & ドウ{\middel}コウ (D\={O}{\middel}K\={O})\\\hline
\end{somme}

\begin{sentence}
\ruby{\underline{山口}}{やまぐち}\underline{さん}\ruby{は}{わ}\underline{あそこ}の\ruby{\underline{女}}{おんな}の\ruby{\underline{人}}{ひと}が\ruby{\underline{大好}}{ダイす}\underline{き}です。\\\hline
Yama{\middel}guchi-san wa asoko no onna no hito ga DAI su{\middel}ki desu.\\\hline
(Yamaguchi-san) (over there) (woman) (person) (really likes)\\\hline
=\textit{Yamaguchi-san really likes that woman over there.}\\\hline
\end{sentence}
\end{multicols}
\vspace*{-\baselineskip}
\vhrulefill{5pt}
%%--------------------------------------------------------------------
%%--------------------------------------------------------------------
\begin{multicols}{2}
\begin{glyph}{Old (古)}{old}\label{glyph:old}
Old; Ancient; Antique.
\end{glyph}

\begin{comps}
\comprtwocone & 十 + 口 \\\hline
\comprthreecone & Ten (\ref{glyph:ten}) + Mouth (\ref{glyph:mouth})\\\hline
\end{comps}

\begin{remglyph}
Both the \remcom{scarecrow} and \remcom{vampire} are \hooglig{old} fictional characters.\\\hline
\end{remglyph}

\begin{prons}
\pronsrtwocone & \textbf{コ (KO)} & \textbf{ふる (furu)}\\\hline
\pronsrthreecone & --- & --- \\\hline
\end{prons}

\begin{remprons}
The \hooglig{old} \comkun{furuncular} \comon{cop}.\\\hline
\end{remprons}

\begin{somme}
古い & \textit{old} & ふる{\middel}い (furu{\middel}i)\\\hline
中古 & middle + old = \textit{secondhand} & チュウ{\middel}コ (CH\={U}{\middel}KO)\\\hline
最古 & most + old = \textit{oldest} & サイ{\middel}コ (SAI{\middel}KO)\\\hline
古語 & old + words = \textit{archaic word} & コ{\middel}ゴ (KO{\middel}GO)\\\hline
古来 & old + come = \textit{time honored} & コ{\middel}ライ (KO{\middel}RAI)\\\hline
古里 & old + hamlet = \textit{old hamlet} & ふる{\middel}さと (furu{\middel}sato)\\\hline
\end{somme}

\begin{sentence}
\ruby{\underline{山下}}{やました}\underline{さん}\ruby{は}{さ}\ruby{\underline{中古}}{チュウコ}の\ruby{\underline{品}}{しな}\ruby{を}{お} \underline{よく} \ruby{\underline{買}}{か}\underline{います}。\\\hline
Yama{\middel}shita-san wa CH\={U}{\middel}KO no shina o yoku ka{\middel}imasu.\\\hline
(Yamashita-san) (secondhand) (goods) (often) (buys)\\\hline
=\textit{Yamashita-san often buys secondhand goods.}\\\hline
\end{sentence}
\end{multicols}
\vspace*{-\baselineskip}
\vhrulefill{5pt}
\vspace*{\fill}
%%--------------------------------------------------------------------
%%--------------------------------------------------------------------
\pagebreak
\vspace*{\fill}
\begin{multicols}{2}
\begin{glyph}{Fire (火)}{fire}\label{glyph:fire}
Fire.
\end{glyph}

\begin{comps}
\comprtwocone & 火 \\\hline
\comprthreecone & Fire (\ref{glyph:fire})\\\hline
\end{comps}

\begin{remglyph}
---\\\hline
\end{remglyph}

\begin{prons}
\pronsrtwocone & \textbf{カ (KA)} & \textbf{ひ (hi)}\\\hline
\pronsrthreecone & --- & ほ (ho) \\\hline
\end{prons}

\begin{remprons}
A \ucomkun{hobby} of \comon{cussing} \comkun{heaps} of \hooglig{fire}.\\\hline
\end{remprons}

\begin{somme}
火 & \textit{fire} & ひ (hi)\\\hline
火元 & fire + basis = \textit{the origin of a fire} & ひ{\middel}もと (hi{\middel}moto)\\\hline
大火 & large + fire = \textit{conflagration} & タイ{\middel}カ (TAI{\middel}KA)\\\hline
火山 & fire + mountain = \textit{volcano} & カ{\middel}ザン (KA{\middel}SEI)\\\hline
花火 & flower + fire = \textit{fireworks} & はな{\middel}び (hana{\middel}bi)\\\hline
火曜日 & fire + day of the week + sun (day) = \textit{Tuesday} & カ{\middel}ヨウ{\middel}び (KA{\middel}Y\={O}{\middel}bi)\\\hline
火星 & fire + star = \textit{Mars} & カ{\middel}セイ (KA{\middel}SEI)\\\hline
\end{somme}

\begin{sentence}
\ruby{\underline{日本}}{ニホン}に\ruby{は}{わ}\ruby{\underline{火山}}{カザン}が\underline{あります}。\\\hline
NI{\middel}HON ni wa KA{\middel}ZAN ga arimasu.\\\hline
(Japan) (volcano) (are)\\\hline
=\textit{There are volcanoes in Japan.}\\\hline
\end{sentence}
\end{multicols}
\vspace*{-\baselineskip}
\vhrulefill{5pt}
%%--------------------------------------------------------------------
%%--------------------------------------------------------------------
\begin{multicols}{2}
\begin{glyph}{Winter (冬)}{winter}\label{glyph:winter}
Winter.
\end{glyph}

\begin{comps}
\comprtwocone & 夂 + 丶 + 丶 \\\hline
\comprthreecone & Overtaking + Dot + Dot\\\hline
\end{comps}

\begin{remglyph}
    \hooglig{Winter} \remcom{overtaking} two \remcom{dots}.\\\hline
\end{remglyph}

\begin{prons}
\pronsrtwocone & \textbf{トウ (T\={O})} & \textbf{ふゆ (fuyu)}\\\hline
\pronsrthreecone & --- & --- \\\hline
\end{prons}

\begin{remprons}
\comkun{Few usage} during a \comon{torturous} \hooglig{winter}.\\\hline
\end{remprons}

\begin{somme}
冬 & \textit{winter} & ふゆ (fuyu)\\\hline
冬休み & winter + rest = \textit{winter vacation} & ふゆ　やす{\middel}み (fuyu yasu{\middel}mi)\\\hline
立冬 & stand + winter = \textit{first day of winter} & リッ{\middel}トウ (RIT{\middel}T\={O})\\\hline
冬物 & winter + thing = \textit{winter clothing} & ふゆ{\middel}もの (fuyu{\middel}mono)\\\hline
冬空 & winter + empty = \textit{winter sky} & ふゆ{\middel}ぞら (fuyu{\middel}zora)\\\hline
\end{somme}

\begin{sentence}
\ruby{\underline{冬}}{ふゆ}に\underline{なる}と\underline{スイス}の\ruby{\underline{山}}{やま}\ruby{は}{わ}\ruby{\underline{美}}{うつく}\underline{しい}。\\\hline
fuyu ni naru to SUISU no yama wa utsuku{\middel}shii.\\\hline
(winter) (becomes) (Switzerland) (mountain) (beautiful)\\\hline
=\textit{The mountains of Switzerland are beautiful in winter.}\\\hline
\end{sentence}
\end{multicols}
\vspace*{-\baselineskip}
\vhrulefill{5pt}
\vspace*{\fill}
%%--------------------------------------------------------------------
%%--------------------------------------------------------------------
\pagebreak

%
%\begin{multicols}{2}
%\chapter{Kanji 61--80}
%\vspace*{\fill}
\begin{glyph}{Water (水)}{water}\label{glyph:water}
Water.
\end{glyph}

\begin{comps}
\comprtwocone & 水 \\\hline
\comprthreecone & Water (\ref{glyph:water}) \\\hline
\end{comps}

\begin{remglyph}
---\\\hline
\end{remglyph}

\begin{prons}
\pronsrtwocone & \textbf{スイ (SUI)} & \textbf{みず (mizu)}\\\hline
\pronsrthreecone & --- & --- \\\hline
\end{prons}

\begin{remprons}
\comkun{Make me zoom} on to the \hooglig{watery} \comon{chop suey}.\\\hline
\end{remprons}

\begin{somme}
水 & \textit{water} & みず (mizu)\\\hline
水中 & water + middle = \textit{underwater; in the water} & スイ{\middel}チュウ (SUI{\middel}CH\={U})\\\hline
下水 & lower + water = \textit{sewerage} & ゲ{\middel}スイ (GE{\middel}SUI)\\\hline
水玉 & water + jewel = \textit{water droplet} & みず{\middel}だま (mizu{\middel}dama)\\\hline
    \notyet{水星} & water + star = \textit{Mercury} & スイ{\middel}セイ (SUI{\middel}SEI)\\\hline
    \notyet{水曜日} & water + day of the week + sun (day) = \textit{Mercury} & スイ{\middel}ヨウ{\middel}び (SUI{\middel}Y\={O}{\middel}bi)\\\hline
    \notyet{水死} & water + death = \textit{drowning} & スイ{\middel}シ (SUI{\middel}SHI)\\\hline
\end{somme}

\begin{sentence}
\underline{あの} \underline{ダム}に\ruby{は}{}\ruby{\underline{水}}{みず}が\underline{ありますえん}。\\\hline
ano damu ni wa mizu ga arimasen.\\\hline
(that) (dam) (water) (isn't)\\\hline
=\textit{There is no water at that dam.}\\\hline
\end{sentence}
\end{multicols}
\vspace*{-\baselineskip}
\vhrulefill{5pt}
%%--------------------------------------------------------------------
%%--------------------------------------------------------------------
\begin{multicols}{2}
\begin{glyph}{Ease (安)}{ease}\label{glyph:ease}
Ease; stability; peacefulness.\\\hline
An important secondary meaning is ``inexpensive''.  In this regard, this kanji is the opposite of 高.
\end{glyph}

\begin{comps}
\comprtwocone & 宀 + 女 \\\hline
\comprthreecone & roof/cover + Woman (\ref{glyph:woman}) \\\hline
\end{comps}

\begin{remglyph}
Every \remcom{woman} needs a \remcom{roof} over her head for a sense of \hooglig{ease and stability}.\\\hline
\end{remglyph}

\begin{prons}
\pronsrtwocone & \textbf{アン (AN)} & \textbf{やす (yasu)}\\\hline
\pronsrthreecone & --- & --- \\\hline
\end{prons}

\begin{remprons}
\hooglig{Ease} your \comon{uncle} with \comkun{ya suit}.\\\hline
\end{remprons}

\begin{somme}
安い & \textit{easy; cheap} & やす{\middel}い (yasu{\middel}i)\\\hline
円安 & circle (yen) + ease = \textit{weak yen} & エン{\middel}だか (EN{\middel}daka)\\\hline
安心 & ease + heart = \textit{peace of mind} & アン{\middel}シン (AN{\middel}SHIN)\\\hline
安全 & ease + complete = \textit{safety} & アン{\middel}ゼン (AN{\middel}ZEN)\\\hline
    \notyet{不安} & not + ease = \textit{uneasiness} & フ{\middel}アン (FU{\middel}AN)\\\hline
    \notyet{公安} & public + ease = \textit{public peace} & コウ{\middel}アン (K\={O}{\middel}AN)\\\hline
\end{somme}

\begin{sentence}
    \ruby{\underline{安心}}{アンシン}{\;}\underline{して} \ruby{\underline{下}}{くだ}\underline{さい}。\\\hline
AN{\middel}SHIN shite kuda{\middel}sai.\\\hline
(peace of mind) (do) (please)\\\hline
=\textit{Please don't worry.}\\\hline
\end{sentence}
\end{multicols}
\vspace*{-\baselineskip}
\vhrulefill{5pt}
\vspace*{\fill}
%%--------------------------------------------------------------------
%%--------------------------------------------------------------------
\pagebreak
\vspace*{\fill}
\begin{multicols}{2}
\begin{glyph}{Strength (力)}{strength}\label{glyph:strength}
Strength; Power; Force.
\end{glyph}

\begin{comps}
\comprtwocone & 力 \\\hline
\comprthreecone & Strength (\ref{glyph:strength}) \\\hline
\end{comps}

\begin{remglyph}
---\\\hline
\end{remglyph}

\begin{prons}
\pronsrtwocone & \textbf{リョク、リキ (RYOKU, RIKI)} & \textbf{ちから (chikara)}\\\hline
\pronsrthreecone & --- & --- \\\hline
\rye{3}{リョク will be encountered far more often than リキ.  The kun-yomi is a common word on its own.}
\end{prons}

\begin{remprons}
By \hooglig{strengthening} the \comon{reeky}, \comon{lean yorkie}, its \comkun{cheek arrangement} will be re-established.\\\hline
\end{remprons}

\begin{somme}
力 & \textit{strenght} & ちから (chikara)\\\hline
全力 & complete + strength = \textit{all one's power} & ゼン{\middel}リョク (ZEN{\middel}RYOKU)\\\hline
有力 & have + strength = \textit{powerful} & ユウ{\middel}リョク (Y\={U}{\middel}RYOKU)\\\hline
水力 & water + strength = \textit{hydro power} & スイ{\middel}リョク (SUI{\middel}RYOKU)\\\hline
火力 & fire + strength = \textit{thermal power} & スイ{\middel}リョク (KA{\middel}RYOKU)\\\hline
    \notyet{電力} & electric + strength = \textit{electric power} & デン{\middel}リョク (DEN{\middel}RYOKU)\\\hline
    \notyet{力士} & strenght + gentleman = \textit{sumo wrestler} & リキ{\middel}シ (RIKI{\middel}SHI)\\\hline
\end{somme}

\begin{sentence}
    \ruby{\underline{中山}}{なかやま}\underline{さん}\ruby{は}{わ}\ruby{\underline{車}}{くるま}の\underline{ビジネス}で\underline{とても}{\;}\ruby{\underline{有力}}{ユウリョク}な\ruby{\underline{人}}{ひと}{\;}\underline{です}。\\\hline
Naka{\middel}yama-san wa kuruma no bijinesu de totemo Y\={U}{\middel}RYOKU na hito desu.\\\hline
(Nakayama-san) (car) (business) (very) (powerful) (person) (is)\\\hline
=\textit{In the automotive business, Nakayama-san is a very powerful person.}\\\hline
\end{sentence}
\end{multicols}
\vspace*{-\baselineskip}
\vhrulefill{5pt}
%%--------------------------------------------------------------------
%%--------------------------------------------------------------------
\begin{multicols}{2}
\begin{glyph}{Summer (夏)}{summer}\label{glyph:summer}
Summer.
\end{glyph}

\begin{comps}
\comprtwocone & 一 + 自 + 夂 \\\hline
\comprthreecone & One (\ref{glyph:one}) + Self (\ref{glyph:self}) + overtaking \\\hline
\end{comps}

\begin{remglyph}
    \hooglig{Summer} \remcom{overtakes} \remcom{one} \remcom{self}.\\\hline
\end{remglyph}

\begin{prons}
\pronsrtwocone & \textbf{カ (KA)} & \textbf{なつ (natsu)}\\\hline
\pronsrthreecone & ゲ (GE) & --- \\\hline
\end{prons}

\begin{remprons}
In \hooglig{summer} I \ucomon{get} \comkun{nuts} such that I start \comon{cussing}.\\\hline
\end{remprons}

\begin{somme}
夏 & \textit{summer} & なつ (natsu)\\\hline
    \notyet{夏休み} & summer + rest = \textit{summer vacation} & なつ　やす{\middel}み (natsu yasu{\middel}mi)\\\hline
    \notyet{夏物} & summer + thing = \textit{summer cloting} & なつ{\middel}もの (natsu{\middel}mono)\\\hline
    \notyet{立夏} & stand + summer = \textit{first day of summer} & リッ{\middel}カ (RIK{\middel}KA)\\\hline
\end{somme}

\begin{sentence}
\underline{この} \ruby{\underline{夏}}{なつ}に\ruby{は}{わ}\ruby{\underline{東京}}{トウキョウ}\ruby{へ}{え}\ruby{\underline{行}}{い}\underline{きました}。\\\hline
kono natsu ni wa T\={O}{\middel}KY\={O} e i{\middel}kimashite.\\\hline
(this) (summer) (Tokyo) (go)\\\hline
=\textit{(We'll) be going to Tokyo this summer.}\\\hline
\end{sentence}
\end{multicols}
\vspace*{-\baselineskip}
\vhrulefill{5pt}
\vspace*{\fill}
%%--------------------------------------------------------------------
%%--------------------------------------------------------------------
\pagebreak
\vspace*{\fill}
\begin{multicols}{2}
\begin{glyph}{Thousand (千)}{thousand}\label{glyph:thousand}
Thousand.\\\hline
Like 百, this kanji sometimes conveys a general idea of `many'.
\end{glyph}

\begin{comps}
\comprtwocone & ノ + 十 \\\hline
\comprthreecone & NO katakana/bend + Ten (\ref{glyph:thousand}) \\\hline
\end{comps}

\begin{remglyph}
    \remcom{No} \remcom{ten} \hooglig{thousand}.\\\hline
\end{remglyph}

\begin{prons}
\pronsrtwocone & \textbf{セン (SEN)} & ---\\\hline
\pronsrthreecone &  & いち (ichi) \\\hline
\end{prons}

\begin{remprons}
The \ucomkun{chief} is \comon{central} to the morale of \hooglig{thousands} of soldiers.\\\hline
\end{remprons}

\begin{somme}
千 & \textit{one thousand} & セン (SEN)\\\hline
二千 & two + thousand = \textit{two thousand} & ニ{\middel}セン (NI{\middel}SEN)\\\hline
三千 & three + thousand = \textit{three thousand} & サン{\middel}ゼン (SAN{\middel}ZEN)\\\hline
四千 & four + thousand = \textit{four thousand} & よん{\middel}セン (yon{\middel}SEN)\\\hline
五千 & five + thousand = \textit{five thousand} & ゴ{\middel}セン (GO{\middel}SEN)\\\hline
六千 & six + thousand = \textit{six thousand} & ロク{\middel}セン (ROKU{\middel}SEN)\\\hline
七千 & seven + thousand = \textit{seven thousand} & なな{\middel}セン (nana{\middel}SEN)\\\hline
八千 & eight + thousand = \textit{eight thousand} & ハチ{\middel}セン (HACHI{\middel}SEN)\\\hline
九千 & nine + thousand = \textit{nine thousand} & キュウ{\middel}セン (KY\={U}{\middel}SEN)\\\hline
\end{somme}

\begin{sentence}
    \ruby{\underline{米}}{こめ}\ruby{は}{わ}\ruby{\underline{七千円}}{ななセンエン}{\;}\underline{なので} \ruby{\underline{安}}{やす}\underline{くない}。\\\hline
kome wa nana{\middel}SEN{\middel}EN na no de yasu{\middel}kunai.\\\hline
(rice) (seven thousand yen) (because) (not cheap)\\\hline
=\textit{At seven thousand yen, rice isn't cheap.}\\\hline
\end{sentence}
\end{multicols}
\vspace*{-\baselineskip}
\vhrulefill{5pt}
%%--------------------------------------------------------------------
%%--------------------------------------------------------------------
\begin{multicols}{2}
\begin{glyph}{Rest (休)}{rest}\label{glyph:rest}
Rest.
\end{glyph}

\begin{comps}
\comprtwocone & 亻 + 木 \\\hline
\comprthreecone & Person (variant) (\ref{glyph:person}) + Tree (\ref{glyph:tree}) \\\hline
\end{comps}

\begin{remglyph}
\remcom{People} typically find \hooglig{rest} under \remcom{trees.}\\\hline
\end{remglyph}

\begin{prons}
\pronsrtwocone & \textbf{キュウ (KY\={U})} & \textbf{やす (yasu)}\\\hline
\pronsrthreecone & --- & --- \\\hline
\end{prons}

\begin{remprons}
If \comkun{ya suit} yourself to too much \hooglig{rest}, it will \comon{kill you}.\\\hline
\end{remprons}

\begin{somme}
休む & \textit{to rest} & やす{\middel}む (yasu{\middel}mu)\\\hline
お休み & \textit{``Good Night!''} & お{\middel}やす{\middel}み (o{\middel}yasu{\middel}mi)\\\hline
休日 & rest + day = \textit{holiday} & キュウ{\middel}ジツ (KY\={U}{\middel}JITSU)\\\hline
夏休み & summer + rest = \textit{summer vacation} & なつ　やす{\middel}み (natsu yasu{\middel}mi)\\\hline
冬休み & winter + rest = \textit{winter vacation} & ふゆ　やす{\middel}み (fuyu yasu{\middel}mi)\\\hline
休火山 & rest + fire + mountain = \textit{dormant volcano} & キュウ{\middel}カ{\middel}ザン (KY\={U}{\middel}KA{\middel}ZAN)\\\hline
\end{somme}

\begin{sentence}
    \ruby{\underline{山口}}{やまぐち}\underline{さん}\ruby{は}{わ}\ruby{\underline{半日}}{ハンニチ}{\;}\ruby{\underline{休}}{やす}\underline{みました}。\\\hline
Yama{\middel}guchi-san wa HAN{\middel}NICHI yasu{\middel}mimashita.\\\hline
(Yamaguchi-san) (half a day) (rested)\\\hline
=\textit{Yamaguchi-san rested for half a day.}\\\hline
\end{sentence}
\end{multicols}
\vspace*{-\baselineskip}
\vhrulefill{5pt}
\vspace*{\fill}
%%--------------------------------------------------------------------
%%--------------------------------------------------------------------
\pagebreak
\vspace*{\fill}
\begin{multicols}{2}
\begin{glyph}{Hand (手)}{te}\label{glyph:hand}
Hand.
\end{glyph}

\begin{comps}
\comprtwocone & 手 \\\hline
\comprthreecone & Hand (\ref{glyph:hand}) \\\hline
\end{comps}

\begin{remglyph}
---\\\hline
\end{remglyph}

\begin{prons}
\pronsrtwocone & \textbf{シュ (SHU)} & \textbf{て (te)}\\\hline
\pronsrthreecone & --- & た (ta) \\\hline
\end{prons}

\begin{remprons}
\hooglig{Hand} me the \comkun{tempting} \ucomkun{tunnel} \comon{shooter}.\\\hline
\end{remprons}

\begin{somme}
手 & \textit{hand} & て (te)\\\hline
右手 & right + hand = \textit{right hand} & みぎ{\middel}て (migi{\middel}te)\\\hline
左手 & left + hand = \textit{left hand} & ひだり{\middel}て (hidari{\middel}te)\\\hline
手元 & hand + basis = \textit{at hand} & て{\middel}もと (te{\middel}moto)\\\hline
手首 & hand + neck = \textit{wrist} & て{\middel}くび (te{\middel}kubi)\\\hline
    \notyet{手話} & hand + speak = \textit{sign language} & シュ{\middel}ワ (SHU{\middel}WA)\\\hline
    \notyet{空手} & empty + hand = \textit{karate} & から{\middel}て (kara{\middel}te)\\\hline
\end{somme}

\begin{irrr}
上手 & upper + hand = \textit{to be good at (something)} & ジョウ{\middel}ず (J\={O}{\middel}zu)\\\hline
下手 & lower + hand = \textit{to be poor at (something)} & へ{\middel}た (he{\middel}ta)\\\hline
\end{irrr}

\begin{sentence}
\underline{あの} \underline{ピアニスト}の\ruby{\underline{手}}{て}\ruby{は}{わ}\ruby{\underline{美}}{うつく}\underline{しい} \underline{ですね}。\\\hline
ano pianisuto no te wa utsuku{\middel}shii desu ne.\\\hline
(that) (pianist) (hands) (beautiful) (aren't they)\\\hline
=\textit{That pianist's hands are beautiful, aren't they.}\\\hline
\end{sentence}
\end{multicols}
\vspace*{-\baselineskip}
\vhrulefill{5pt}
%%--------------------------------------------------------------------
%%--------------------------------------------------------------------
\begin{multicols}{2}
\begin{glyph}{Main (本)}{main}\label{glyph:main}
Main; Origin.  Also used to emphasize the idea of ``this''.\\\hline
When encountered alone, it often indicates a book.\\\hline
It is also a `counter' for long, cylindrical objects such as pencils and bottles.\\\hline
The sense of ``origin'', incidentally, is evident in the compound for Japan itself, the ``Land of the Rising Sun.''
\end{glyph}

\begin{comps}
\comprtwocone & 木 + 一 \\\hline
\comprthreecone & Tree (\ref{glyph:tree}) + One (\ref{glyph:one}) \\\hline
\end{comps}

\begin{remglyph}
One main tree.\\\hline
\end{remglyph}

\begin{prons}
\pronsrtwocone & \textbf{ホン (HON)} & \textbf{もと (moto)}\\\hline
\pronsrthreecone & --- & --- \\\hline
\rye{3}{もと is common in Japanese family and place names.}
\end{prons}

\begin{remprons}
The \comkun{motorists} were \hooglig{mainly} \comon{haunted}.\\\hline
\end{remprons}

\begin{somme}
本 & \textit{book} & ホン (HON)\\\hline
日本 & sun + main = \textit{Japan} & ニ{\middel}ホン、ニッ{\middel}ポン (NI{\middel}HON, NIP{\middel}PON)\\\hline
全日本 & complete + sun + main = \textit{all Japan} & ゼン{\middel}ニ{\middel}ホン (ZEN{\middel}NI{\middel}HON)\\\hline
日本人 & sun + main + person = \textit{a Japanese person} & ニ{\middel}ホン{\middel}ジン (NI{\middel}HON{\middel}JIN)\\\hline
山本さん & mountain + main = \textit{Yamamoto-san} & やま{\middel}もと{\middel}さん (yama{\middel}moto{\middel}san)\\\hline
    \notyet{日本語} & sun + main + words = \textit{Japanese (language)} & ニ{\middel}ホン{\middel}ゴ (NI{\middel}HON{\middel}GO)\\\hline
    \notyet{本土} & main + earth = \textit{mainland} & ホン{\middel}ド (HON{\middel}DO)\\\hline
\end{somme}

\begin{sentence}
\ruby{\underline{山下}}{やました}\underline{さん}\ruby{は}{わ}\ruby{\underline{高}}{たか}\underline{い} \ruby{\underline{本}}{ホン}\ruby{を}{お}\ruby{\underline{買}}{か}\underline{いました}。\\\hline
Yama{\middel}shita-san wa taka{\middel}i HON o ka{\middel}imashita.\\\hline
(Yamashita-san) (expensive) (book) (bought)\\\hline
=\textit{Yamashita-san bought an expensive book.}\\\hline
\end{sentence}
\end{multicols}
\vspace*{-\baselineskip}
\vhrulefill{5pt}
\vspace*{\fill}
%%--------------------------------------------------------------------
%%--------------------------------------------------------------------
\pagebreak
\vspace*{\fill}
\begin{multicols}{2}
\begin{glyph}{Change (化)}{change}\label{glyph:change}
Change; Transform.\\\hline
When used as a suffix, this character expresses a variety of English words ending in ``-ification'' or ``-ization''.
\end{glyph}

\begin{comps}
\comprtwocone & 亻 + 匕 \\\hline
\comprthreecone & Person (variant) (\ref{glyph:person}) + spoon \\\hline
\end{comps}

\begin{remglyph}
That \remcom{person} \hooglig{changed} the \remcom{spoon}.\\\hline
\end{remglyph}

\begin{prons}
\pronsrtwocone & \textbf{カ (KA)} & ---\\\hline
\pronsrthreecone & KE (ケ) & ば (ba) \\\hline
\end{prons}

\begin{remprons}
People \comon{cussed} the \ucomkun{bulging} \ucomon{keg} that \hooglig{changed} the alcoholic taste.\\\hline
\end{remprons}

\begin{somme}
国有化 & country + have + change = \textit{nationalization} & コク{\middel}ユウ{\middel}カ (KOKU{\middel}Y\={U}{\middel}KA)\\\hline
    \notyet{美化} & beautiful + change = \textit{beautification} & ビ{\middel}カ (BI{\middel}KA)\\\hline
    \notyet{化学} & change + study = \textit{chemistry} & キ{\middel}ガク (KA{\middel}GAKU)\\\hline
    \notyet{気化} & spirit + change = \textit{vaporization} & キ{\middel}カ (KI{\middel}KA)\\\hline
    \notyet{電化} & electric + change = \textit{electrification} & デン{\middel}カ (DEN{\middel}KA)\\\hline
    \notyet{同化} & same + change = \textit{assimilation} & ドウ{\middel}カ (D\={O}{\middel}KA)\\\hline
\end{somme}

\begin{sentence}
\ruby{\underline{山本}}{やまもと}\underline{さん}\ruby{は}{わ}\ruby{\underline{化学}}{カガク}が\ruby{\underline{大好}}{ダイす}\underline{き}です。\\\hline
Yama{\middel}moto-san wa KA{\middel}GAKU ga DAI su{\middel}ki desu.\\\hline
(Yamamoto-san) (chemistry) (really likes)\\\hline
=\textit{Yamamoto-san really likes chemistry.}\\\hline
\end{sentence}
\end{multicols}
\vspace*{-\baselineskip}
\vhrulefill{5pt}
%%--------------------------------------------------------------------
%%--------------------------------------------------------------------
\begin{multicols}{2}
\begin{glyph}{River (川)}{river}\label{glyph:river}
River.
\end{glyph}

\begin{comps}
\comprtwocone & 川 \\\hline
\comprthreecone & River (\ref{glyph:river}) \\\hline
\end{comps}

\begin{remglyph}
---\\\hline
\end{remglyph}

\begin{prons}
\pronsrtwocone & --- & \textbf{カワ (kawa)}\\\hline
\pronsrthreecone & セン (SEN) & --- \\\hline
\rye{3}{When appearing as the final character in the names of rivers, this kanji is invariable voiced as がわ.}
\end{prons}

\begin{remprons}
The \ucomon{central} \hooglig{river} attraction was the \comkun{carwash.}\\\hline
\end{remprons}

\begin{somme}
川 & \textit{river} & かわ (kawa)\\\hline
川上 &  river + upper = \textit{upstream} & かわ{\middel}かみ (kawa{\middel}kami)\\\hline
川下 & river + lower = \textit{downstream} & かわ{\middel}した (kawa{\middel}shita)\\\hline
川口 & river + mouth = \textit{river mouth} & かわ{\middel}ぐち (kawa{\middel}guchi)\\\hline
川本さん & river + main = \textit{Kawamoto-san} & かわ{\middel}もと{\middel}さん (kawa{\middel}moto{\middel}san)\\\hline
\end{somme}

\begin{sentence}
\ruby{\underline{山}}{やま}の\ruby{\underline{上}}{うえ}に\ruby{\underline{行}}{い}\underline{く}と\ruby{\underline{川}}{かわ}が\ruby{\underline{見}}{み}\underline{えます}か。\\\hline
yama no ue ni i{\middel}ku to kawa ga mi{\middel}emasu ka.\\\hline
(mountain) (upper) (go) (river) (can see)\\\hline
=\textit{Can you see the river if you go to the top of the mountain.}\\\hline
\end{sentence}
\end{multicols}
\vspace*{-\baselineskip}
\vhrulefill{5pt}
\vspace*{\fill}
%%--------------------------------------------------------------------
%%--------------------------------------------------------------------
\pagebreak
\vspace*{\fill}
\begin{multicols}{2}
\begin{glyph}{Body (体)}{body}\label{glyph:body}
The physical body of a person or object.\\\hline
A secondary meaning of ``substance'' and ``style''.
\end{glyph}

\begin{comps}
\comprtwocone & 亻 + 本 \\\hline
\comprthreecone & Person (variant) (\ref{glyph:person}) + Main (\ref{glyph:main}) \\\hline
\end{comps}

\begin{remglyph}
The \hooglig{body} is the \remcom{main} aspect of any \remcom{person}.\\\hline
\end{remglyph}

\begin{prons}
\pronsrtwocone & \textbf{タイ (TAI)} & \textbf{からだ (karada)}\\\hline
\pronsrthreecone & テイ (TEI) & --- \\\hline
\end{prons}

\begin{remprons}
My \hooglig{body} \comon{tightened} as the \comkun{car adapter} got \ucomon{taped} to me.\\\hline
\end{remprons}

\begin{somme}
体 & \textit{body} & からだ (karada)\\\hline
大体 & large + body = \textit{in general} & ダイ{\middel}タイ (DAI{\middel}TAI)\\\hline
全体 & complete + body = \textit{the whole} & ゼン{\middel}タイ (ZEN{\middel}TAI)\\\hline
肉体 & meat + body = \textit{the body} & ニク{\middel}タイ (NIKU{\middel}TAI)\\\hline
体力 & body + strength = \textit{physical strength} & タイ{\middel}リョク (TAI{\middel}RYOKU)\\\hline
自体 & self + body = \textit{in itself} & ジ{\middel}タイ (JI{\middel}TAI)\\\hline
    \notyet{団体} & group + body = \textit{organization} & ダン{\middel}タイ (DAN{\middel}TAI)\\\hline
\end{somme}

\begin{sentence}
\underline{あの} \ruby{\underline{人}}{ひと}\ruby{は}{わ}\underline{すごく} \ruby{\underline{体力}}{タイリョク}が\ruby{\underline{有}}{あ}\underline{ります}。\\\hline
ano hito wa sugoku TAI{\middel}RYOKU ga a{\middel}rimasu.\\\hline
(that) (person) (incredible) (physical strength) (has)\\\hline
=\textit{That person has incredible physical strength.}\\\hline
\end{sentence}
\end{multicols}
\vspace*{-\baselineskip}
\vhrulefill{5pt}
%%--------------------------------------------------------------------
%%--------------------------------------------------------------------
\begin{multicols}{2}
\begin{glyph}{North (北)}{north}\label{glyph:north}
North.
\end{glyph}

\begin{comps}
\comprtwocone & 爿 + 匕 \\\hline
\comprthreecone & bed + spoon \\\hline
\end{comps}

\begin{remglyph}
    The \hooglig{northern} \remcom{bed} \remcom{spoon}.\\\hline
\end{remglyph}

\begin{prons}
\pronsrtwocone & \textbf{ホク (HOKU)} & \textbf{きた (kita)}\\\hline
\pronsrthreecone & --- & --- \\\hline
\end{prons}

\begin{remprons}
\comon{Hawks} move \hooglig{northward} to the \comkun{citadel}.\\\hline
\end{remprons}

\begin{somme}
北 & \textit{north} & きた (kita)\\\hline
北米 & north + rice = \textit{North America} & ホク{\middel}ベイ (HOKU{\middel}BEI)\\\hline
北口 & north + mouth = \textit{north exit/entrance} & きた{\middel}ぐち (kita{\middel}guci)\\\hline
    \notyet{北東} & north + east = \textit{northeast} & ホク{\middel}トウ (HOKU{\middel}T\={O})\\\hline
    \notyet{北西} & north + west = \textit{northwest} & ホク{\middel}セイ (HOKU{\middel}SEI)\\\hline
    \notyet{北海道} & north + sea + road = \textit{Hokkaido} & ホッ{\middel}カイ{\middel}ドウ (HOK{\middel}KAI{\middel}D\={O})\\\hline
\end{somme}

\begin{sentence}
\ruby{\underline{北日本}}{きたニホン}の\ruby{\underline{冬}}{ふゆ}\ruby{は}{わ}\underline{とても} \ruby{\underline{美}}{うつく}\underline{しい}。\\\hline
kita{\middel}NI{\middel}HON no fuyu wa totemo utsuku{\middel}shii.\\\hline
(northern Japan) (winter) (very) (beautiful)\\\hline
=\textit{Northern Japan's winter is very beautiful.}\\\hline
\end{sentence}
\end{multicols}
\vspace*{-\baselineskip}
\vhrulefill{5pt}
\vspace*{\fill}
%%--------------------------------------------------------------------
%%--------------------------------------------------------------------
\pagebreak
\vspace*{\fill}
\begin{multicols}{2}
\begin{glyph}{Rice field (田)}{rice_field}\label{glyph:rice_field}
Rice field, or a field in general.
\end{glyph}

\begin{comps}
\comprtwocone & 田 \\\hline
\comprthreecone & Rice Field (\ref{glyph:rice_field}) \\\hline
\end{comps}

\begin{remglyph}
---\\\hline
\end{remglyph}

\begin{prons}
\pronsrtwocone & \textbf{デン (DEN)} & \textbf{た (ta)}\\\hline
\pronsrthreecone & --- & --- \\\hline
\rye{3}{This kanji is commonly used in Japanese family names; it is often voiced, as in the fourth and fifth examples below, the latter of which is a mix of on and kun-yomi.}
\end{prons}

\begin{remprons}
In the \hooglig{rice field} was a \comkun{tub} hid inside a \comon{den}.\\\hline
\end{remprons}

\begin{somme}
田 & \textit{rice field} & た (ta)\\\hline
水田 & water + rice field = \textit{rice paddy} & スイ{\middel}デン (SUI{\middel}DEN)\\\hline
田中さん & rice field + middle = \textit{Tanaka-san} & た{\middel}なか{\middel}さん (ta{\middel}naka{\middel}san)\\\hline
山田さん & mountain + rice field = \textit{Yamada-san} & やま{\middel}だ{\middel}さん (yama{\middel}da{\middel}san)\\\hline
本田さん & main + rice field = \textit{Honda-san} & ホン{\middel}だ{\middel}さん (HON{\middel}da{\middel}san)\\\hline
\end{somme}

\begin{irrr}
    田舎{\notcov} & rice field + building = \textit{the countryside} & いなか (inaka)\\\hline
\end{irrr}

\begin{sentence}
\ruby{\underline{水田}}{スイデン}\ruby{を}{お}\ruby{\underline{見}}{み}に\ruby{\underline{行}}{い}\underline{きました}。\\\hline
SUI{\middel}DEN o mi ni i{\middel}kimashita.\\\hline
(rice paddy) (see) (went)\\\hline
=\textit{(She) went to see the rice paddies.}\\\hline
\end{sentence}
\end{multicols}
\vspace*{-\baselineskip}
\vhrulefill{5pt}
%%--------------------------------------------------------------------
%%--------------------------------------------------------------------
\begin{multicols}{2}
\begin{glyph}{Man (男)}{man}\label{glyph:man}
Man; Male.
\end{glyph}

\begin{comps}
\comprtwocone & 田 + 力 \\\hline
\comprthreecone & Rice Field (\ref{glyph:rice_field}) + Strength (\ref{glyph:strength})\\\hline
\end{comps}

\begin{remglyph}
A real \hooglig{man} shows his \remcom{strength} in the \remcom{rice fields}.\\\hline
\end{remglyph}

\begin{prons}
\pronsrtwocone & \textbf{ダン (DAN)} & \textbf{おとこ (otoko)}\\\hline
\pronsrthreecone & ナン (NAN) & --- \\\hline
\end{prons}

\begin{remprons}
\hooglig{Man} perfected the \comkun{auto-complete} function to dunk \ucomon{nuns}.\\\hline
\end{remprons}

\begin{somme}
男 & \textit{man} & おとこ (otoko)\\\hline
男女 & man + woman = \textit{men and woman} & ダン{\middel}ジョ (DAN{\middel}JO)\\\hline
男子 & man + child = \textit{men's} & ダン{\middel}シ (DAN{\middel}SHI)\\\hline
大男 & large + man = \textit{large man} & おお{\middel}おとこ (\={o}{\middel}otoko)\\\hline
    \notyet{雪男} & snow + man = \textit{the abominable snowman} & ゆき{\middel}おとこ (yuki{\middel}otoko)\\\hline
\end{somme}

\begin{sentence}
\underline{あの} \ruby{\underline{男}}{おとこ}の\ruby{\underline{人}}{ひと}\ruby{は}{わ}\ruby{\underline{五千円}}{ゴセンエン}\ruby{を}{お}\underline{なくした}。\\\hline
ano otoko no hito wa GO{\middel}SEN{\middel}EN o nakushita.\\\hline
(that) (man) (person) (five thousand yen) (lost)\\\hline
=\textit{That man lost five thousand yen.}\\\hline
\end{sentence}
\end{multicols}
\vspace*{-\baselineskip}
\vhrulefill{5pt}
\vspace*{\fill}
%%--------------------------------------------------------------------
%%--------------------------------------------------------------------
\pagebreak
\vspace*{\fill}
\begin{multicols}{2}
\begin{glyph}{House (家)}{house}\label{glyph:house}
Some sense of relation to house or family.\\\hline
When used as a suffix, it denotes a person having some degree of skill or interest with respect to the kanji preceding it.
\end{glyph}

\begin{comps}
\comprtwocone & 宀 + 豕 \\\hline
\comprthreecone & roof + pig\\\hline
\end{comps}

\begin{remglyph}
    A \hooglig{house} is a \remcom{roof} over a \remcom{pig}.\\\hline
\end{remglyph}

\begin{prons}
\pronsrtwocone & \textbf{カ (KA)} & \textbf{いえ (ie)}\\\hline
\pronsrthreecone & ケ (KE) & や (ya) \\\hline
\end{prons}

\begin{remprons}
\comkun{i.e.} we started \comon{cussing} once we realised the \ucomon{keg} was already \hooglig{housed} in the \ucomkun{yacht}.\\\hline
\end{remprons}

\begin{somme}
家 & \textit{house} & いえ (ie)\\\hline
人家 & person + house = \textit{dwelling} & ジン{\middel}カ (JIN{\middel}KA)\\\hline
    \notyet{売り家} & sell + house = \textit{house for sale} & う{\middel}り　いえ (u{\middel}ri ie)\\\hline
    \notyet{家具} & house + tool = \textit{furniture} & カ{\middel}グ (KA{\middel}GU)\\\hline
    \notyet{愛鳥家} & love + bird + house = \textit{bird lover} & アイ{\middel}チョウ{\middel}カ (AI{\middel}CH\={O}{\middel}KA)\\\hline
\end{somme}

\begin{sentence}
\ruby{\underline{本田}}{ホンだ}\underline{さん}の\ruby{\underline{家}}{いえ}\ruby{を}{お}\ruby{\underline{見}}{み}\underline{ました}か。\\\hline
HON{\middel}da{\middel}san no ie o mi{\middel}mashita ka.\\\hline
(Honda-san) (house) (see)\\\hline
=\textit{Did you see Honda-san's house?}\\\hline
\end{sentence}
\end{multicols}
\vspace*{-\baselineskip}
\vhrulefill{5pt}
%%--------------------------------------------------------------------
%%--------------------------------------------------------------------
\begin{multicols}{2}
\begin{glyph}{East (東)}{east}\label{glyph:east}
East.
\end{glyph}

\begin{comps}
\comprtwocone & 木 + 日 \\\hline
\comprthreecone & Tree (\ref{glyph:tree}) + Sun (\ref{glyph:tree})\\\hline
\end{comps}

\begin{remglyph}
In the \hooglig{east} the \remcom{sun} rises behind the \remcom{trees}.\\\hline
\end{remglyph}

\begin{prons}
\pronsrtwocone & \textbf{トウ (T\={O})} & \textbf{ひがし (higashi)}\\\hline
\pronsrthreecone & --- & --- \\\hline
\end{prons}

\begin{remprons}
\hooglig{Eastern} \comon{torture} methods will make \comkun{him gush}.\\\hline
\end{remprons}

\begin{somme}
東 & \textit{east} & ひがし (higashi)\\\hline
東口 & east + mouth = \textit{east exit/entrance} & ひがし{\middel}ぐち (higashi{\middel}guchi)\\\hline
中東 & middle + east = \textit{The Middle East} & チュウ{\middel}トウ (NAN{\middel}T\={O})\\\hline
東京 & east + capital = \textit{Tokyo} & キョウ{\middel}トウ (KY\={O}{\middel}T\={O})\\\hline
南東 & south + east = \textit{south east} & ナン{\middel}トウ (NAN{\middel}T\={O})\\\hline
北北東 & north + north + east = \textit{north-northeast} & ホク{\middel}ホク{\middel}トウ (HOKU{\middel}HOKU{\middel}T\={O})\\\hline
\end{somme}

\begin{sentence}
\ruby{\underline{東口}}{ひがしぐち}の\underline{キオスク}で\ruby{\underline{本}}{ホン}\ruby{を}{お}\ruby{\underline{買}}{か}\underline{いました}。\\\hline
higashi{\middel}guchi no kiosuku de HON o ka{\middel}imashita.\\\hline
(east exit) (kiosk) (book) (bought)\\\hline
=\textit{(I) bought a book at the east exit kiosk.}\\\hline
\end{sentence}
\end{multicols}
\vspace*{-\baselineskip}
\vhrulefill{5pt}
\vspace*{\fill}
%--------------------------------------------------------------------
%--------------------------------------------------------------------
\pagebreak
\vspace*{\fill}
\begin{multicols}{2}
\begin{glyph}{Think (思)}{think}\label{glyph:think}
Think.
\end{glyph}

\begin{comps}
\comprtwocone & 田 + 心 \\\hline
\comprthreecone & Rice field (\ref{glyph:rice_field}) + Heart (\ref{glyph:heart})\\\hline
\end{comps}

\begin{remglyph}
I \hooglig{think} the \remcom{rice field} has forever stolen my \remcom{heart}.\\\hline
\end{remglyph}

\begin{prons}
\pronsrtwocone & \textbf{シ (SHI)} & \textbf{おも (omo)} \\\hline
\pronsrthreecone & --- & --- \\\hline
\end{prons}

\begin{remprons}
I \hooglig{think} \comkun{omo} will give it a nice \comon{sheen}.\\\hline
\end{remprons}

\begin{somme}
思う & \textit{to think} & おも{\middel}う (omo{\middel}u)\\\hline
    \notyet{思い出す} & think + exit = \textit{to remember} & おも{\middel}い　だ{\middel}す (omo{\middel}i da{\middel}su)\\\hline
    \notyet{物思い} & thing + think = \textit{pensiveness} & もの　おも{\middel}い (mono omo{\middel}i)\\\hline
    \notyet{思い切る} & think + cut = \textit{one's intent} & おも{\middel}い　き{\middel}る (omo{\middel}i ki{\middel}ru)\\\hline
    \notyet{意思} & mind + think = \textit{(one's) intent} & イ{\middel}シ (I{\middel}SHI)\\\hline
\end{somme}

\begin{sentence}
    \underline{テーベル}の\ruby{\underline{上}}{うえ}に\ruby{\underline{四百円}}{よんヒャク}{\;}\underline{ある}と\ruby{\underline{思}}{おも}\underline{います}。\\\hline
t\={e}beru no ue ni yon{\middel}HYAKU{\middel}EN aru to omo{\middel}imasu.\\\hline
(table) (above) (four hundred yen) (is) (think)\\\hline
=\textit{(I) think there's four hundred yen on the table.}\\\hline
\end{sentence}
\end{multicols}
\vspace*{-\baselineskip}
\vhrulefill{5pt}
%%--------------------------------------------------------------------
%%--------------------------------------------------------------------
\begin{multicols}{2}
\begin{glyph}{Ear (耳)}{ear}\label{glyph:ear}
Ear.
\end{glyph}

\begin{comps}
\comprtwocone & 耳 \\\hline
\comprthreecone & Ear (\ref{glyph:ear})\\\hline
\end{comps}

\begin{remglyph}
---\\\hline
\end{remglyph}

\begin{prons}
\pronsrtwocone & --- & \textbf{みみ (mimi)} \\\hline
\pronsrthreecone & ジ (JI) & --- \\\hline
\end{prons}

\begin{remprons}
\comkun{Mimic} the \hooglig{ear} \ucomon{jig}.\\\hline
\end{remprons}

\begin{somme}
耳 & \textit{ear} & 耳 (mimi)\\\hline
耳元 & ear + basis = \textit{close to one's ear} & みみ{\middel}もと (mimi{\middel}moto)\\\hline
    \notyet{空耳} & empty + ear = \textit{mishearing; feigned deafness} & そら{\middel}みみ (sora{\middel}mimi)\\\hline
\end{somme}

\begin{sentence}
\ruby{\underline{田中}}{たなか}\underline{さん}の\ruby{\underline{耳}}{みみ}が\ruby{\underline{赤}}{あか}\underline{い}。\\\hline
Ta{\middel}naka-san no mimi ga aka{\middel}i.\\\hline
(Tanaka-san) (ears) (red)\\\hline
=\textit{Tanaka-san's ears are red.}\\\hline
\end{sentence}
\end{multicols}
\vspace*{-\baselineskip}
\vhrulefill{5pt}
\vspace*{\fill}
%%--------------------------------------------------------------------
%%--------------------------------------------------------------------
\pagebreak
\vspace*{\fill}
\begin{multicols}{2}
\begin{glyph}{Father (父)}{father}\label{glyph:father}
Father.
\end{glyph}

\begin{comps}
\comprtwocone & 父 \\\hline
\comprthreecone & Father (\ref{glyph:father})\\\hline
\end{comps}

\begin{remglyph}
---\\\hline
\end{remglyph}

\begin{prons}
\pronsrtwocone & \textbf{フ (FU)} & \textbf{ちち (chichi)} \\\hline
\pronsrthreecone & --- & --- \\\hline
\end{prons}

\begin{remprons}
Any \comkun{chichi} \hooglig{father} is a \comon{foofy}.\\\hline
\end{remprons}

\begin{somme}
\rye{3}{The various words for father shown below reflect differing levels of familiarity and politeness; most terms used for addressing people in Japan have multiple forms such as these.}
父 & \textit{father} & ちち (chichi)\\\hline
父の日 & father + sun (day) = \textit{Father's Day} & ちち{\middel}の{\middel}ひ (chichi{\middel}no{\middel}oya)\\\hline
    \notyet{父母} & father + mother = \textit{father and mother} & フ{\middel}ボ (FU{\middel}BO)\\\hline
    \notyet{父親} & fater + parent = \textit{father} & ちち{\middel}おや (chichi{\middel}oya)\\\hline
\end{somme}

\begin{irrr}
お父さん & \textit{father} & おとうさん (ot\={o}san)\\\hline
    伯父さん{\notcov} & related + father = \textit{uncle} & おじさん (ojisan)\\\hline
\end{irrr}

\begin{sentence}
\ruby{\underline{父}}{ちち}\ruby{は}{わ}\underline{ビジネス}で\ruby{\underline{東京}}{トウキョウ}\ruby{へ}{え}\ruby{\underline{行}}{い}\underline{きました}。\\\hline
chichi wa bijinesu de T\={O}{\middel}KY\={O} e i{\middel}kimashita.\\\hline
(Father) (business) (Tokyo) (went)\\\hline
=\textit{Father went to Tokyo on business.}\\\hline
\end{sentence}
\end{multicols}
\vspace*{-\baselineskip}
\vhrulefill{5pt}
%%--------------------------------------------------------------------
%%--------------------------------------------------------------------
\begin{multicols}{2}
\begin{glyph}{Say (言)}{say}\label{glyph:say}
``Say'' and things of a verbal nature.\\\hline
This character appears in more than 60 kanji.
\end{glyph}

\begin{comps}
\comprtwocone & 言 \\\hline
\comprthreecone & Speech/Speaking/Say (\ref{glyph:say})\\\hline
\end{comps}

\begin{remglyph}
Sound waves traveling outward.\\\hline
\end{remglyph}

\begin{prons}
\pronsrtwocone & \textbf{ゲン (GEN)} & \textbf{い (i)} \\\hline
\pronsrthreecone & ゴン (GON) & こと (koto) \\\hline
\end{prons}

\begin{remprons}
They \hooglig{say} \comon{Genghis Khan's} \comkun{interest} in \ucomkun{koto's} is \ucomon{gone} with the wind.\\\hline
\end{remprons}

\begin{somme}
言う & \textit{to say} & い{\middel}う (i{\middel}u)\\\hline
言い回し & say + rotate = \textit{turn of expression} & い{\middel}い　まわ{\middel}し (i{\middel}i mawa{\middel}shi)\\\hline
明言 & bright + say = \textit{declaration} & メイ{\middel}ゲン (MEI{\middel}GEN)\\\hline
    \notyet{言語} & say + words = \textit{language} & ゲン{\middel}ゴ (GEN{\middel}GO)\\\hline
    \notyet{言い出す} & say + exit = \textit{to begin to say} & い{\middel}い　だ{\middel}す (i{\middel}i da{\middel}su)\\\hline
    \notyet{言い切る} & say + cut = \textit{to say definitively} & い{\middel}い　き{\middel}る (i{\middel}i ki{\middel}ru)\\\hline
    言葉{\notcov} & say + leaf = \textit{word} & こと{\middel}ば (koto{\middel}ba)\\\hline
\end{somme}

\begin{sentence}
\ruby{\underline{川本}}{かわもと}\underline{さん} \underline{も} \underline{``はい''}と\ruby{\underline{言}}{い}\underline{いました}。\\\hline
Kawa{\middel}moto-san mo ``hai'' to i{\middel}imashita.\\\hline
(Kawamoto-san) (also) (``yes'') (said)\\\hline
=\textit{Kawamoto-san also said ``yes''.}\\\hline
\end{sentence}
\end{multicols}
\vspace*{-\baselineskip}
\vhrulefill{5pt}
\vspace*{\fill}
%%--------------------------------------------------------------------
%%--------------------------------------------------------------------
\pagebreak

%
%\begin{multicols}{2}
%\chapter{Kanji 81--100}
%\begin{glyph}{Dry (干)}{dry}\label{glyph:dry}
Dry.  Ebb.
\end{glyph}

\begin{comps}
\comprtwocone & 干 \\\hline
\comprthreecone & Dry/Shield \\\hline
\end{comps}

\begin{remglyph}
Part of the 漢字 radicals (\#{51}).\\\hline
\end{remglyph}

\begin{prons}
\pronsrtwocone & --- & \textbf{ひ、ほ (hi, ho)}\\\hline
\pronsrthreecone & カン (KAN) & --- \\\hline
\end{prons}

\begin{remprons}
    \hooglig{Dry} \comkun{heaps} of \hooglig{dried} fruit garnered by
    \ucomon{cunning} \comkun{hobbyists}.\\\hline
\end{remprons}

\begin{somme}
    干る (intr.) & \textit{to dry (on its own)} & ひ{\middel}る \mbox{(hi{\middel}ru)}\\\hline
    干す (tr.) & \textit{to dry (something)} & ほ{\middel}す \mbox{(ho{\middel}su)}\\\hline
    {\diff}日干し & sun + dry = \textit{sun-dried} & ひ　ぼ{\middel}し \mbox{(hi bo{\middel}shi)}\\\hline
    干上がる & dry + upper = \textit{to dry up} & ひ　あ{\middel}がる \mbox{(hi a{\middel}garu)}\\\hline
    \notyet{干物} & dry + thing = \textit{dried fish} & ひ{\middel}もの \mbox{(hi{\middel}mono)}\\\hline
\end{somme}

\begin{sentence}
\ruby{\underline{干物}}{ひもの}\ruby{を}{お}\ruby{\underline{買}}{か}\underline{って} \ruby{\underline{下}}{くだ}\underline{さい}。\\\hline
hi{\middel}mono o kat{\middel}te kuda{\middel}sai.\\\hline
(dried fish) (buy) (please)\\\hline
=\textit{Please buy dried fish.}\\\hline
\end{sentence}
\end{multicols}
\vhrulefill{2pt}
%%--------------------------------------------------------------------
%%--------------------------------------------------------------------
\begin{multicols}{2}
\begin{glyph}{Gold (金)}{gold}\label{glyph:gold}
Depending on the context in which it appears, this common character signifies either ``gold'', ``metal'', or ``money''.\\\hline
It clearly means ``gold'', for instance, when referring to the status of an Olympic champion.\\\hline
When functioning as a component, however, it hints that the 漢字 involved has some connection to metal in general.
\end{glyph}

\begin{comps}
\comprtwocone & 金\\\hline
\comprthreecone &　Gold\\\hline
\end{comps}

\begin{remglyph}
Part of the Kanji radicals (\#{167}).\\\hline
\end{remglyph}

\begin{prons}
\pronsrtwocone & \textbf{キン (KIN)} & \textbf{かね (kane)}\\\hline
\pronsrthreecone & コン (KON) & かな (kana) \\\hline
\end{prons}

\begin{remprons}
The \hooglig{golden} \ucomkun{canary} \comkun{cane} were stolen by \comon{kindred} \ucomon{con} artists.\\\hline
\end{remprons}

\begin{somme}
    金 & \textit{gold; metal; money} & キン \mbox{(KIN)}\\\hline
    金 & \textit{metal; money} & かね \mbox{(kane)}\\\hline
    金山 & gold + mountain = \textit{gold mine} & キン{\middel}ザン \mbox{(KIN{\middel}ZAN)}\\\hline
    \notyet{年金} & year + gold = \textit{annuity/pension} & ネン{\middel}キン \mbox{(NEN{\middel}KIN)}\\\hline
    \notyet{金曜日} & gold + day of the week + sun (day) = \textit{Friday} & キン{\middel}ヨウ{\middel}び \mbox{(KIN{\middel}Y\={O}{\middel}bi)}\\\hline
    \notyet{金星} & gold + star = \textit{Venus (planet)} & キン{\middel}セイ \mbox{(KIN{\middel}SEI)}\\\hline
\end{somme}

\begin{sentence}
\ruby{\underline{内田}}{うちだ}\underline{さん}\ruby{は}{わ}\underline{あまり} \underline{お}\ruby{\underline{金}}{かね}が\underline{ない}。\\\hline
Uchi{\middel}da-san wa amari o{\middel}kane ga nai.\\\hline
(Uchida-san) (not much) (money) (not)\\\hline
=\textit{Uchida-san doesn't have much money.}\\\hline
\end{sentence}
\end{multicols}
%%--------------------------------------------------------------------
%%--------------------------------------------------------------------
\begin{multicols}{2}
\begin{glyph}{Words (語)}{words}\label{glyph:words}
Words.\\\hline
This 漢字 will become familiar as the character is used to indicate languages.
\end{glyph}

\begin{comps}
\comprtwocone & 言 + 五 + 口 \\\hline
\comprthreecone & Say + Five + Mouth \\\hline
\end{comps}

\begin{remglyph}
\hooglig{Wording} the \remcom{five} \remcom{vampire} \remcom{sayings}.\\\hline
\end{remglyph}

\begin{prons}
\pronsrtwocone & \textbf{ゴ (GO)} & \textbf{かた (kata)}\\\hline
\pronsrthreecone & --- & かたり (katari) \\\hline
\end{prons}

\begin{remprons}
\hooglig{Word} \comkun{catalogues} \comon{gobbled} down to \ucomkun{cut a
    ribbon} at the finishline.\\\hline
\end{remprons}

\begin{somme}
\rye{3}{Note the importance of long and short vowel sounds in Japanese by looking closely at the second and third examples below.}
    語る & \textit{to talk} & かた{\middel}る \mbox{(kata{\middel}ru)}\\\hline
    口語 & mouth words = \textit{colloquial language} & コウ{\middel}ゴ \mbox{(K\={O}{\middel}GO)}\\\hline
    古語 & old + words = \textit{archaic words} & コ{\middel}ゴ \mbox{(KO{\middel}GO)}\\\hline
    日本語 & sun + main + words = \textit{Japanese (language)} & ニ{\middel}ホン{\middel}ゴ \mbox{(NI{\middel}HON{\middel}GO)}\\\hline
    言語 & say + words = \textit{language} & ゲン{\middel}ゴ \mbox{(GEN{\middel}GO)}\\\hline
    \notyet{英語} & English + words = \textit{English (language)} & エイ{\middel}ゴ \mbox{(EI{\middel}GO)}\\\hline
    \notyet{語学} & words + study = \textit{linguistics} & ゴ{\middel}ガク \mbox{(GO{\middel}GAKU)}\\\hline
\end{somme}

\begin{sentence}
\underline{あなた}の\ruby{\underline{日本語}}{ニホンゴ}\ruby{は}{わ}\ruby{\underline{上手}}{ジョウず}ですよ。\\\hline
anata no NI{\middel}HON{\middel}GO wa J\={O}{\middel}zu desu yo.\\\hline
(you) (Japanese) (good at)\\\hline
=\textit{Your Japanese is good.}\\\hline
\end{sentence}
\end{multicols}
\vhrulefill{2pt}
%--------------------------------------------------------------------
%--------------------------------------------------------------------
\begin{multicols}{2}
\begin{glyph}{Gentleman (士)}{gentleman}\label{glyph:gentleman}
Gentleman.\\\hline
A man who has attained status in a military or academic field.
\end{glyph}

\begin{comps}
\comprtwocone & 士 \\\hline
\comprthreecone & Gentleman \\\hline
\end{comps}

\begin{remglyph}
Part of the 漢字 radicals (\#{33}).\\\hline
\end{remglyph}

\begin{prons}
\pronsrtwocone & \textbf{シ (SHI)} & ---\\\hline
\pronsrthreecone & --- & --- \\\hline
\end{prons}

\begin{remprons}
\hooglig{Gentlemen} \comon{shield} other people.\\\hline
\end{remprons}

%{\middel}
\begin{somme}
    力士 & strength + gentleman = \textit{sumo wrestler} & リキ{\middel}シ \mbox{(RIKI{\middel}SHI)}\\\hline
    \notyet{士気} & gentleman + spirit = \textit{morale} & シ{\middel}キ \mbox{(SHI{\middel}KI)}\\\hline
    \notyet{工学士} & craft + study + gentleman = \textit{Bachelor of Engineering} & コウ{\middel}ガク{\middel}シ \mbox{(K\={O}{\middel}GAKU{\middel}SHI)}\\\hline
\end{somme}

\begin{sentence}
\underline{あの} \ruby{\underline{力士}}{リキシ}\ruby{は}{わ}\underline{とても} \ruby{\underline{大}}{おお}\underline{きい} \underline{です}よ。\\\hline
ano RIKI{\middel}SHI wa totemo \={o}{\middel}kii desu yo.\\\hline
(that) (sumo wrestler) (very) (large) (is)\\\hline
=\textit{That sumo wrestler is really large.}\\\hline
\end{sentence}
\end{multicols}
\vspace*{\fill}
%--------------------------------------------------------------------
%--------------------------------------------------------------------
\pagebreak
\vspace*{\fill}
\begin{multicols}{2}
\begin{glyph}{Morning (朝)}{morning}\label{glyph:morning}
Morning.\\\hline
A secondary meaning relates to ``royal dynasties''.
\end{glyph}

\begin{comps}
\comprtwocone & 十 + 早 + 月 \\\hline
\comprthreecone & Ten + Early (\ref{glyph:early}) + Moon\\\hline
\end{comps}

\begin{remglyph}
While the \hooglig{morning} is still \remcom{early}, the \remcom{moon} winks at
    the \remcom{scarecrow}.\\\hline
\end{remglyph}

\begin{prons}
\pronsrtwocone & \textbf{チョウ (CH\={O})} & \textbf{あさ (asa)}\\\hline
\pronsrthreecone & --- & --- \\\hline
\end{prons}

\begin{remprons}
\hooglig{Morning} \comon{chores} must be finished \comkun{asap}.\\\hline
\end{remprons}

%{\middel}
% &  ()\\\hline
\begin{somme}
    朝 & \textit{morning} & あさ \mbox{(asa)}\\\hline
    朝日 & morning + sun = \textit{morning sun} & あさ{\middel}ひ \mbox{(asa{\middel}hi)}\\\hline
    王朝 & king + morning = \textit{dynasty} & オウ{\middel}チョウ \mbox{(\={O}{\middel}CH\={O})}\\\hline
    早朝 & early + morning = \textit{early morning} & ソウ{\middel}チョウ \mbox{(S\={O}{\middel}CH\={O})}\\\hline
    \notyet{毎朝} & every + morning = \textit{every morning} & マイ{\middel}あさ \mbox{(MAI{\middel}asa)}\\\hline
\end{somme}

\begin{irrr}
    \notyet{今朝} & now + morning = \textit{this morning} & けさ (kesa)\\\hline
\end{irrr}

\begin{sentence}
\ruby{\underline{上田}}{うえだ}\underline{さん}\ruby{は}{わ}\ruby{\underline{朝}}{あさ}が\ruby{\underline{大好}}{ダイす}\underline{き}です。\\\hline
Ue{\middel}da-san wa asa ga DAI{\middel}su{\middel}ki desu.\\\hline
(Ueda-san) (morning) (really likes)\\\hline
=\textit{Ueda-san really likes mornings.}\\\hline
\end{sentence}
\end{multicols}
\vhrulefill{2pt}
%--------------------------------------------------------------------
%--------------------------------------------------------------------
\begin{multicols}{2}
\begin{glyph}{Blue (青)}{blue}\label{glyph:blue}
Blue.\\\hline
This character also applies to a few objects that Westerners invariably consider green.  E.g. \textit{blue traffic light}, \textit{blue vegetables}, amongst other things.\\\hline
The sense of ``youth'' and ``immaturity'' is tied up with the idea of ``green''.
\end{glyph}

\begin{comps}
\comprtwocone & 青 \\\hline
\comprthreecone & Blue \\\hline
\end{comps}

\begin{remglyph}
Part of the 漢字 radicals (\#{174}).\\\hline
\end{remglyph}

\begin{prons}
\pronsrtwocone & \textbf{セイ (SEI)} & \textbf{あお (ao)}\\\hline
\pronsrthreecone & ショウ (SH\={O}) &  \\\hline
\end{prons}

\begin{remprons}
\hooglig{Blue} \ucomon{shore} \comon{sailing} really pumps the \comkun{aorta}.\\\hline
\end{remprons}

%{\middel}
% &  ()\\\hline
\begin{somme}
    青 & \textit{blue; green (noun)} & あお \mbox{(ao)}\\\hline
    青い & \textit{blue; green (adjective)} & あお{\middel}い \mbox{(ao{\middel}i)}\\\hline
    青春 & blue + spring = \textit{youth} & セイ{\middel}シュン \mbox{(SEI{\middel}SHUN)}\\\hline
    {\weird}\notyet{青物} & blue + thing = \textit{green vegetables} & あお{\middel}もの \mbox{(ao{\middel}mono)}\\\hline
    \notyet{青空} & blue + empty = \textit{blue sky} & あお{\middel}ぞら \mbox{(ao{\middel}zora)}\\\hline
\end{somme}

\begin{sentence}
\ruby{\underline{山田}}{やまだ}\underline{さん}の\ruby{\underline{家}}{いえ}\ruby{は}{わ}\ruby{\underline{青}}{あお}\underline{い}です。\\\hline
Yama{\middel}da-san no ie wa ao{\middel}i desu.\\\hline
(Yamada-san) (house) (blue)\\\hline
=\textit{Yamada-san's house is blue.}\\\hline
\end{sentence}
\end{multicols}
\vspace*{\fill}
%%--------------------------------------------------------------------
%%--------------------------------------------------------------------
\pagebreak
\vspace*{\fill}
\begin{multicols}{2}
\begin{glyph}{Earth (土)}{earth}\label{glyph:earth}
Earth as in ``dirt'', not the planet.
\end{glyph}

\begin{comps}
\comprtwocone & 土 \\\hline
\comprthreecone & Earth \\\hline
\end{comps}

\begin{remglyph}
Part of the 漢字 radicals (\#{32}).\\\hline
\end{remglyph}

\begin{prons}
\pronsrtwocone & \textbf{ド (DO)} & \textbf{つち (tsuchi)}\\\hline
\pronsrthreecone & ト (TO) & --- \\\hline
\end{prons}

\begin{remprons}
\hooglig{Earth} \comon{dots} made by the \ucomon{top} \comkun{tsuchikage}.\\\hline
{\warn}Tsuchikage is a fictional アニメ character in ナルト.\\\hline
\end{remprons}

%{\middel}
% &  ()\\\hline
\begin{somme}
    土 & \textit{earth} & つち \mbox{(tsuchi)}\\\hline
    土手 & earth + hand = \textit{embankment} & ド{\middel}て \mbox{(DO{\middel}te)}\\\hline
    {\diff}本土 & main + earth = \textit{mainland} & ホン{\middel}ド \mbox{(HON{\middel}DO)}\\\hline
    土木 & earth + tree = \textit{Civil Engineering} & ド{\middel}ボク \mbox{(DO{\middel}BOKU)}\\\hline
    \notyet{土星} & earth + star = \textit{Saturn} & ド{\middel}セイ \mbox{(DO{\middel}SEI)}\\\hline
    \notyet{土曜日} & earth + day of the week + sun (day) = \textit{Saturday} & ド{\middel}ヨウ{\middel}び \mbox{(DO{\middel}Y\={O}{\middel}bi)}\\\hline
\end{somme}

\begin{irrr}
    \notyet{土産} & earth + produce = \textit{souvenir} & みやげ (miyage)\\\hline
\end{irrr}

\begin{sentence}
\ruby{\underline{川}}{かわ}の\ruby{\underline{土手}}{ドて}\ruby{を}{お}\ruby{\underline{上}}{あ}\underline{がる}と\ruby{\underline{水田}}{スイデン}が\ruby{\underline{見}}{み}\underline{えます}。\\\hline
kawa no DO{\middel}te o a{\middel}garu to SUI{\middel}DEN ga mi{\middel}emasu.\\\hline
(river) (embankment) (go up) (rice paddies) (can see)\\\hline
=\textit{If you go up the river embankment, you can see rice paddies.}\\\hline
\end{sentence}
\end{multicols}
\vhrulefill{2pt}
%%--------------------------------------------------------------------
%%--------------------------------------------------------------------
\begin{multicols}{2}
\begin{glyph}{Hang (掛)}{hang}\label{glyph:hang}
Hang.\\\hline
Note how the last stroke of the bottom ``earth'' component leads into the next line by being drawn with an upward slant.
\end{glyph}

\begin{comps}
\comprtwocone & 扌 + 土 + 土 + ⼘ \\\hline
\comprthreecone & Hand + Earth + Earth + Divination  \\\hline
\end{comps}

\begin{remglyph}
\hooglig{Hanging} \remcom{earth} with your \remcom{hands} is \remcom{divination}.\\\hline
\end{remglyph}

\begin{prons}
\pronsrtwocone & --- & \textbf{か (ka)}\\\hline
\pronsrthreecone & --- & --- \\\hline
\end{prons}

\begin{remprons}
\hooglig{Hanging} clothes is a \comkun{customary} thing to do.\\\hline
\end{remprons}

%{\middel}
% &  ()\\\hline
\begin{somme}
\rye{3}{As the final example below shows, this 漢字 can act similarly to 買 (see entry \ref{glyph:buy}).}
掛かる (intr.) & \textit{to hang (by itself)} & か{\middel}かる \mbox{(ka{\middel}karu)}\\\hline
掛ける (tr.) & \textit{to hang (something)} & か{\middel}ける \mbox{(ka{\middel}keru)}\\\hline
掛金 & hang + gold (money) = \textit{installment} & かけ{\middel}キン \mbox{(kake{\middel}KIN)}\\\hline
\end{somme}

\begin{sentence}
\underline{シャツ}\ruby{を}{お}\underline{あそこ}に\ruby{\underline{掛}}{か}\underline{けて} \ruby{\underline{下}}{くだ}\underline{さい}。\\\hline
SHATSU o asoko ni ka{\middel}kete kuda{\middel}sai.\\\hline
(shirts) (over there) (hang) (please)\\\hline
=\textit{Please hang the shirts over there.}\\\hline
\end{sentence}
\end{multicols}
\vspace*{\fill}
%%--------------------------------------------------------------------
%%--------------------------------------------------------------------
\pagebreak
\begin{multicols}{2}
\begin{glyph}{Ten Thousand (万)}{ten_thousand}\label{glyph:ten_thousand}
Ten thousand.\\\hline
More so than 百 and 千, this character can express the general idea of ``countless amount'' of something.
\end{glyph}

\begin{comps}
\comprtwocone & 丿 + 一 \\\hline
\comprthreecone & NO/Slash + One\\\hline
\end{comps}

\begin{remglyph}
\hooglig{Ten thousand} \remcom{slashes} can be conquered by \remcom{no} \remcom{one}.\\\hline
\end{remglyph}

\begin{prons}
\pronsrtwocone & \textbf{マン (MAN)} & ---\\\hline
\pronsrthreecone & バン (BAN) &  \\\hline
\end{prons}

\begin{remprons}
\hooglig{Ten thousand} \comon{mundane} \ucomon{buns}.\\\hline
\end{remprons}

%{\middel}
% &  ()\\\hline
\begin{somme}
一万 & one + ten thousand = \textit{ten thousand} & イチ{\middel}マン \mbox{(ICHI{\middel}MAN)}\\\hline
二万 & two  + ten thousand = \textit{twenty thousand} & ニ{\middel}マン \mbox{(NI{\middel}MAN)}\\\hline
三万 & three + ten thousand = \textit{thirty thousand} & サン{\middel}マン \mbox{(SAN{\middel}MAN)}\\\hline
四万 & four + ten thousand = \textit{forty thousand} & よん{\middel}マン \mbox{(yon{\middel}MAN)}\\\hline
五万 & five + ten thousand = \textit{fifty thousand} & ゴ{\middel}マン \mbox{(GO{\middel}MAN)}\\\hline
六万 & six + ten thousand = \textit{sixty thousand} & ロク{\middel}マン \mbox{(ROKU{\middel}MAN)}\\\hline
七万 & seven + ten thousand = \textit{seventy thousand} & なな{\middel}マン \mbox{(nana{\middel}MAN)}\\\hline
八万 & eight + ten thousand = \textit{eighty thousand} & ハチ{\middel}マン \mbox{(HACHI{\middel}MAN)}\\\hline
九万 & nine + ten thousand = \textit{ninety thousand} & キュウ{\middel}マン \mbox{(KY\={U}{\middel}MAN)}\\\hline
\end{somme}

\begin{sentence}
\underline{あの}
\ruby{\underline{国}}{くに}の\ruby{\underline{人口}}{ジンコウ}\ruby{は}{わ}\ruby{\underline{八万人}}{ハチマンニン}{\;}\underline{です}。\\\hline
ano kuni no JIN{\middel}K\={O} wa HACHI{\middel}MAN{\middel}NIN desu.\\\hline
(that) (country) (population) (eighty thousand people) (is)\\\hline
=\textit{That country's population is eighty thousand.}\\\hline
\end{sentence}
\end{multicols}
\vhrulefill{2pt}
%%--------------------------------------------------------------------
%%--------------------------------------------------------------------
\begin{multicols}{2}
\begin{glyph}{Rain (雨)}{rain}\label{glyph:rain}
Rain.
\end{glyph}

\begin{comps}
\comprtwocone & 雨 \\\hline
\comprthreecone & Rain \\\hline
\end{comps}

\begin{remglyph}
Part of the 漢字 radicals (\#{173}).\\\hline
\end{remglyph}

\begin{prons}
\pronsrtwocone & --- & \textbf{あめ、あま (ame, ama)}\\\hline
\rye{3}{The 訓読み, あま, only appears in the first position.}
  \pronsrthreecone & ウ (U) & --- \\\hline
\end{prons}

\begin{remprons}
\hooglig{Rain} destroyed the \comon{ooky} \comkun{amalgamation}, and therefore \comkun{amends} were not necessary anymore.\\\hline
\end{remprons}

%{\middel}
% &  ()\\\hline
\begin{somme}
雨 & \textit{rain} & あめ \mbox{(ame)}\\\hline
雨上がり & rain + upper = \textit{after the rain} & あめ　あ{\middel}がり \mbox{(ame a{\middel}gari)}\\\hline
大雨 & large + rain = \textit{heavy rain} & おお{\middel}あめ \mbox{(\={o}{\middel}ame)}\\\hline
\notyet{雨空} & rain + sky = \textit{rainy sky} & あま{\middel}ぞら \mbox{(ama{\middel}zora)}\\\hline
\end{somme}

\begin{irrr}
\rye{3}{The following word can be read with its normal 音読み, but is often heard with the irregular reading pronunciation below.}
    梅雨{\notcov} & plum tree + rain = \textit{the rainy season} & つゆ (tsuyu)\\\hline
\end{irrr}

\begin{sentence}
\ruby{\underline{九月}}{クガツ}に\ruby{は}{わ}\ruby{\underline{日本}}{ニホン}で\ruby{\underline{大雨}}{おおあめ}の\ruby{\underline{日}}{ひ}が\ruby{\underline{多}}{おお}\underline{い}です。\\\hline
KU{\middel}GATSU ni wa NI{\middel}HON de \={o}{\middel}ame no hi ga \={o}{\middel}i desu.\\\hline
(September) (Japan) (heavy rain) (day) (many)\\\hline
=\textit{In September, Japan has many days with heavy rain.}\\\hline
\end{sentence}
\end{multicols}
%%--------------------------------------------------------------------
%%--------------------------------------------------------------------
\pagebreak
\vspace*{\fill}
\begin{multicols}{2}
\begin{glyph}{Electric (電)}{electric}\label{glyph:electric}
Electric; electricity.
\end{glyph}

\begin{comps}
\comprtwocone & 雨 + 日 + 乚 \\\hline
\comprthreecone & Rain + Sun + Second/Latter \\\hline
\end{comps}

\begin{remglyph}
\hooglig{Electricity} is a \remcom{latter} result of \remcom{rain} under the \remcom{sun}.\\\hline
\end{remglyph}

\begin{prons}
\pronsrtwocone & \textbf{デン (DEN)} & ---\\\hline
\pronsrthreecone & --- & --- \\\hline
\end{prons}

\begin{remprons}
\hooglig{Electric} \comon{den}.\\\hline
\end{remprons}

%{\middel}
% &  ()\\\hline
\begin{somme}
    電化 & electric + change = \textit{electrification} & デン{\middel}カ \mbox{(DEN{\middel}KA)}\\\hline
    電力 & electric + strength = \textit{electric power} & ダン{\middel}リョク \mbox{(DEN{\middel}RYOKU)}\\\hline
    電子 & electric + child = \textit{electron} & デン{\middel}シ \mbox{(DEN{\middel}SHI)}\\\hline
    \notyet{電気} & electric + spirit = \textit{electricity} & デン{\middel}キ \mbox{(DEN{\middel}KI)}\\\hline
    \notyet{電話} & electric + speak = \textit{telephone} & デン{\middel}ワ \mbox{(DEN{\middel}WA)}\\\hline
    \notyet{電車} & electric + car = \textit{(electric) train} & デン{\middel}シャ \mbox{(DEN{\middel}SHA)}\\\hline
\end{somme}

\begin{sentence}
\ruby{\underline{田中}}{たなか}\underline{さん}\ruby{は}{わ}\ruby{\underline{電車}}{デンシャ}\ruby{を}{お}\ruby{\underline{下}}{お}\underline{りました}。\\\hline
Ta{\middel}naka-san wa DEN{\middel}SHA o o{\middel}rimashita.\\\hline
(Tanaka-san) (train) (descended)\\\hline
=\textit{Tanaka-san got off the train.}\\\hline
\end{sentence}
\end{multicols}
\vhrulefill{2pt}
%%--------------------------------------------------------------------
%%--------------------------------------------------------------------
\begin{multicols}{2}
\begin{glyph}{Sell (売)}{sell}\label{glyph:sell}
Sell.
\end{glyph}

\begin{comps}
\comprtwocone & 士 + 冖 + 儿\\\hline
    \comprthreecone & Gentleman + Cover/Crown + Man/Person\\\hline
\end{comps}

\begin{remglyph}
By \hooglig{selling} products the \remcom{gentleman} is \remcom{covering} for
    another \remcom{person}.\\\hline
\end{remglyph}

\begin{prons}
\pronsrtwocone & \textbf{バイ (BAI)} & \textbf{う (u)}\\\hline
\pronsrthreecone & --- & --- \\\hline
\end{prons}

\begin{remprons}
    \hooglig{Selling} and \comon{buying} is the \comkun{oomph} behind any business.\\\hline
\end{remprons}

%{\middel}
% &  ()\\\hline
\begin{somme}
\rye{3}{As the final example below shows, this is another kanji that can function similar to 買.}
    売る & \textit{to sell} & う{\middel}る \mbox{(u{\middel}ru)}\\\hline
    売り歩く & sell + walk = \textit{to peddle} & う{\middel}り　ある{\middel}く \mbox{(u{\middel}ri aru{\middel}ku)}\\\hline
    売春 & sell + spring = \textit{prostitution} & バイ{\middel}シュン \mbox{(BAI{\middel}SHUN)}\\\hline
    売り家 & sell + house = \textit{``house for sale''} & う{\middel}り　いえ \mbox{(u{\middel}ri ie)}\\\hline
    \notyet{売り切れる} & sell + cut = \textit{to be sold out} & う{\middel}り　き{\middel}れる \mbox{(u{\middel}ri ki{\middel}reru)}\\\hline
    \notyet{売切れ} & sell + cut = \textit{``sold out''} & う{\middel}り　き{\middel}れ \mbox{(u{\middel}ri ki{\middel}re)}\\\hline
\end{somme}

\begin{sentence}
\ruby{\underline{田口}}{たぐち}\underline{さん}\ruby{は}{わ}\ruby{\underline{日本語}}{ニホンゴ}の\ruby{\underline{本}}{ホン}\ruby{を}{お}\ruby{\underline{売}}{}\underline{って}います。\\\hline
Ta{\middel}guchi-san wa NI{\middel}HON{\middel}GO no HON o u{\middel}tte imasu.\\\hline
(Taguchi-san) (Japanese) (books) (sells)\\\hline
=\textit{Taguchi-san sells Japanese books.}\\\hline
\end{sentence}
\end{multicols}
\vspace*{\fill}
%%--------------------------------------------------------------------
%%--------------------------------------------------------------------
\pagebreak
\begin{multicols}{2}
\begin{glyph}{Individual (者)}{individual}\label{glyph:individual}
Think of individual in the sense of a human being.\\\hline
The characters in a compound that precede this kanji define the person.\\\hline
This character never comes first.
\end{glyph}

\begin{comps}
\comprtwocone & 耂老⺹ + 日 \\\hline
\comprthreecone & Old + Sun \\\hline
\end{comps}

\begin{remglyph}
    That \hooglig{individual} is as \remcom{old} as the \remcom{sun}.\\\hline
\end{remglyph}

\begin{prons}
\pronsrtwocone & \textbf{シャ (SHA)} & ---\\\hline
\pronsrthreecone & --- & もの (mono) \\\hline
\end{prons}

\begin{remprons}
    The \hooglig{individual} \comon{shut} up his \ucomkun{monologue}.\\\hline
\end{remprons}

%{\middel}
% &  ()\\\hline
\begin{somme}
    有力者 & have + strength + individual = \textit{powerful person} & ユウ{\middel}リョク{\middel}シャ \mbox{(Y\={U}{\middel}RYOKU{\middel}SHA)}\\\hline
    \notyet{前者} & before + individual = \textit{the former} & ゼン{\middel}シャ \mbox{(ZEN{\middel}SHA)}\\\hline
    \notyet{後者} & after + individual = \textit{the latter} & コウ{\middel}シャ \mbox{(K\={O}{\middel}SHA)}\\\hline
    \notyet{学者} & study + individual = \textit{scholar} & ガク{\middel}シャ \mbox{(GAKU{\middel}SHA)}\\\hline
\end{somme}

\begin{sentence}
\underline{あの} \ruby{\underline{国}}{くに}で\ruby{\underline{王子}}{オウジ}\ruby{は}{わ}\ruby{\underline{有力者}}{ユウリョクシャ}です。\\\hline
ano kuni de \={O}{\middel}JI wa Y\={U}{\middel}RYOKU{\middel}SHA desu.\\\hline
(that) (country) (prince) (powerful person)\\\hline
=\textit{The prince is a powerful person in that country.}\\\hline
\end{sentence}
\end{multicols}
\vhrulefill{2pt}
%%--------------------------------------------------------------------
%%--------------------------------------------------------------------
\begin{multicols}{2}
\begin{glyph}{Enter (入)}{enter}\label{glyph:enter}
Enter; Putting in.
\end{glyph}

\begin{comps}
\comprtwocone & 入 \\\hline
\comprthreecone & Enter \\\hline
\end{comps}

\begin{remglyph}
Part of the Kanji radicals.\\\hline
\end{remglyph}

\begin{prons}
\pronsrtwocone & \textbf{ニュウ (NY\={U})} & \textbf{はい、い (hai, i)}\\\hline
\pronsrthreecone & --- & --- \\\hline
\end{prons}

\begin{remprons}
\hooglig{Enter} \comon{new} \comkun{interesting} \comkun{heights}.\\\hline
\end{remprons}

%{\middel}
% &  ()\\\hline
\begin{somme}
    入る (intr.) & \textit{to enter} & はい{\middel}る \mbox{(hai{\middel}ru)}\\\hline
    入る (intr.) & \textit{to enter} & い{\middel}る \mbox{(i{\middel}ru)}\\\hline
    入れる (tr.) & \textit{to put (something) into} & い{\middel}れる \mbox{(i{\middel}reru)}\\\hline
    入口 & enter + mouth = \textit{entrance} & いり{\middel}ぐち \mbox{(iri{\middel}guchi)}\\\hline
    入国 & enter + country = \textit{to enter a country} & ニュウ{\middel}コク \mbox{(NY\={U}{\middel}KOKU)}\\\hline
    入金 & enter + gold = \textit{payment} & ニュウ{\middel}キン \mbox{(NY\={U}{\middel}KIN)}\\\hline
    買い入れる & buy + enter = \textit{to stock up on} & か{\middel}い　い{\middel}れる \mbox{(ka{\middel}i i{\middel}reru)}\\\hline
\rye{3}{はい{\middel}る expresses the intransitive meaning of the verb ``enter'' (that is, to enter some type of place).  It is rarely used in compounds.\\\hline
い{\middel}る is most commonly found in compounds, primarily when the other kanji is read with kun-yomi (ニュウ performs this function when the other kanji is read with 音読み).\\\hline
い{\middel}れる is simply the transitive companion of はい{\middel}る.\\\hline
As can be seen in the fourth example, this kanji can behave on rare occasions like 買.}
\end{somme}

\begin{sentence}
\underline{お}\ruby{\underline{金}}{かね}\ruby{を}{お}\ruby{\underline{入}}{い}\underline{れて} \ruby{\underline{下}}{くだ}\underline{さい}。\\\hline
o{\middel}kane o i{\middel}rete kuda{\middel}sai.\\\hline
(money) (put in) (please)\\\hline
=\textit{Please insert the money.}\\\hline
\end{sentence}
\end{multicols}
%%--------------------------------------------------------------------
%%--------------------------------------------------------------------
\pagebreak
\vspace*{\fill}
\begin{multicols}{2}
\begin{glyph}{Snow (雪)}{snow}\label{glyph:snow}
Snow.
\end{glyph}

\begin{comps}
\comprtwocone & 雨 + 彐彑 \\\hline
\comprthreecone & Rain + Pig's head/snout \\\hline
\end{comps}

\begin{remglyph}
\hooglig{Snow} is produced when \remcom{rain} comes out a \remcom{pig's snout}.\\\hline
\end{remglyph}

\begin{prons}
\pronsrtwocone & --- & \textbf{ゆき (yuki)}\\\hline
\pronsrthreecone & セツ (SETSU) & --- \\\hline
\end{prons}

\begin{remprons}
    \hooglig{``Snowy''}, also known as \comkun{``Yuki''}, \comon{sets} up a
    \hooglig{snowman}.\\\hline
\end{remprons}

%{\middel}
% &  ()\\\hline
\begin{somme}
    雪 & \textit{snow} & ゆき \mbox{(yuki)}\\\hline
    雪玉 & snow + jewel = \textit{snowball} & ゆき{\middel}だま \mbox{(yuki{\middel}dama)}\\\hline
    雪女 & snow + woman = \textit{snow fairy} & ゆき{\middel}おんな \mbox{(yuki{\middel}onna)}\\\hline
    雪男 & snow + man = \textit{the Abominable Snowman} & ゆき{\middel}おとｔこ \mbox{(yuki{\middel}otoko)}\\\hline
    大雪 & large + snow = \textit{heavy snow} & おお{\middel}ゆき \mbox{(\={o}{\middel}yuki)}\\\hline
    小雪 & small + snow = \textit{light snow} & こ{\middel}ゆき \mbox{(ko{\middel}yuki)}\\\hline
\end{somme}

\begin{irrr}
    吹雪{\notcov} & blow + snow = \textit{blizzard} & ふぶき (fubuki)\\\hline
    雪崩{\notcov} & snow + crumble = \textit{avalance} & なだれ (nadare)\\\hline
\end{irrr}

\begin{sentence}
\underline{テレビ}の\underline{ニュース}に\underline{よる}と\ruby{\underline{明日}}{あした}\hphantom{h}\underline{から} \ruby{\underline{大雪}}{おおゆき}です。\\\hline
TEREBI no NY\={U}SU ni yoru to ashita kara \={o}yuki desu.\\\hline
(television) (news) (according to) (tomorrow) (from) (heavy snow)\\\hline
=\textit{According to the TV news, heavy snow will be starting tomorrow.}\\\hline
\end{sentence}
\end{multicols}
\vhrulefill{2pt}
%%--------------------------------------------------------------------
%%--------------------------------------------------------------------
\begin{multicols}{2}
\begin{glyph}{Gate (門)}{gate}\label{glyph:gate}
Gate.
\end{glyph}

\begin{comps}
\comprtwocone & 門 \\\hline
\comprthreecone & Gate \\\hline
\end{comps}

\begin{remglyph}
Part of the Kanji radicals (\#{169}).\\\hline
\end{remglyph}

\begin{prons}
\pronsrtwocone & \textbf{モン (MON)} & \textbf{かど (kado)}\\\hline
\pronsrthreecone & --- & --- \\\hline
\end{prons}

\begin{remprons}
    The \hooglig{gates} of \comon{monasteries} are not there for
    \comkun{cuddling}.\\\hline
\end{remprons}

%{\middel}
% &  ()\\\hline
\begin{somme}
    門 & \textit{gate} & モン \mbox{(MON)}\\\hline
    門口 & gate + mouth = \textit{front door} & かど{\middel}ぐち \mbox{(kado{\middel}guchi)}\\\hline
    水門 & water + gate = \textit{sluice gate} & スイ{\middel}モン \mbox{(SUI{\middel}MON)}\\\hline
    校門 & school + gate = \textit{school gate} & コウ{\middel}モン \mbox{(K\={O}{\middel}MON)}\\\hline
\end{somme}

\begin{sentence}
\underline{あそこ}の\ruby{\underline{高}}{たか}\underline{い} \ruby{\underline{門}}{モン}\ruby{は}{わ}\underline{とても} \ruby{\underline{古}}{ふる}\underline{い}ですね。\\\hline
asoko no taka{\middel}i MON wa totemo furu{\middel}i desu ne.\\\hline
(over there) (tall) (gate) (very) (old)\\\hline
=\textit{That tall gate over there is very old, isn't it?}\\\hline
\end{sentence}
\end{multicols}
\vspace*{\fill}
%%--------------------------------------------------------------------
%%--------------------------------------------------------------------
\pagebreak
\vspace*{\fill}
\begin{multicols}{2}
\begin{glyph}{Interval (間)}{interval}\label{glyph:interval}
This kanji indicates an interval in time or space, and thus incorporates the ideas of ``during'', ``between'', and ``among''.\\\hline
When used with the kun-yomi ま, it can also mean a room.
\end{glyph}

\begin{comps}
\comprtwocone & 門 + 日 \\\hline
\comprthreecone & Gate + Sun \\\hline
\end{comps}

\begin{remglyph}
    After an \hooglig{interval} the \remcom{sun} will have passed across the
    \remcom{gate}.\\\hline
\end{remglyph}

\begin{prons}
\pronsrtwocone & \textbf{カン (KAN)} & \textbf{あいだ、ま (aida, ma)}\\\hline
\pronsrthreecone & ケン (KEN) & --- \\\hline
\end{prons}

\begin{remprons}
    During the \hooglig{interval} of the \comon{cunning} \comkun{aida}
    performance, \comkun{muddy} \ucomon{kennels} were set up.\\\hline
\end{remprons}

%{\middel}
% &  ()\\\hline
\begin{somme}
    間 & \textit{interval} & あいだ \mbox{(aida)}\\\hline
    手間 & hand + interval = \textit{labor; trouble} & て{\middel}ま \mbox{(te{\middel}ma)}\\\hline
    中間 & middle + interval = \textit{midway} & チュウ{\middel}カン \mbox{(CH\={U}{\middel}KAN)}\\\hline
    \notyet{時間} & time + interval = \textit{time} & ジ{\middel}カン \mbox{(JI{\middel}KAN)}\\\hline
    \notyet{空間} & empty + interval = \textit{space} & クウ{\middel}カン \mbox{(K\={U}{\middel}KAN)}\\\hline
\end{somme}

\begin{sentence}
\underline{アメリカ}に\underline{いる} \ruby{\underline{間}}{あいだ} \ruby{\underline{英語}}{エイゴ}\ruby{を}{お}\ruby{\underline{話}}{}\underline{します}。\\\hline
AMERIKA ni iru aida EI{\middel}GO o hana{\middel}shimasu.\\\hline
(America) (be) (interval) (English) (speak)\\\hline
=\textit{I'll speak English while I'm in America.}\\\hline
\end{sentence}
\end{multicols}
\vhrulefill{2pt}
%%--------------------------------------------------------------------
%%--------------------------------------------------------------------
\begin{multicols}{2}
\begin{glyph}{Exit (出)}{exit}\label{glyph:exit}
Exit; Leave; Let out.
\end{glyph}

\begin{comps}
\comprtwocone & 屮 + 凵 \\\hline
\comprthreecone & Sprout + Wide opened mouth/Container \\\hline
\end{comps}

\begin{remglyph}
    \hooglig{Exit} the \remcom{sprout} in the \remcom{container}.\\\hline
\end{remglyph}

\begin{prons}
\pronsrtwocone & \textbf{シュツ (SHUTSU)} & \textbf{で、だ (de, da)}\\\hline
\pronsrthreecone & スイ (SUI) & --- \\\hline
\end{prons}

\begin{remprons}
    I made my \hooglig{exit} after I got \comkun{dented} by the \comkun{dull}
    \comon{shoots} in the \ucomon{chop suey}.\\\hline
\end{remprons}

%{\middel}
% &  ()\\\hline
\begin{somme}
    出る (intr.) & \textit{to leave} & で{\middel}る \mbox{(de{\middel}ru)}\\\hline
    出す (tr.) & \textit{to let out} & だ{\middel}す \mbox{(da{\middel}su)}\\\hline
    出かける & \textit{to leave (for somewhere)} & で{\middel}かける \mbox{(de{\middel}kakeru)}\\\hline
    思い出す & think + exit = \textit{to remember} & おも{\middel}い　だ{\middel}す \mbox{(omo{\middel}i da{\middel}su)}\\\hline
    出口 & exit + mouth = \textit{exit} & で{\middel}ぐち \mbox{(de{\middel}guchi)}\\\hline
    出国 & exit + country = \textit{to leave a country} & シュッ{\middel}コク \mbox{(SHUK{\middel}KOKU)}\\\hline
\end{somme}

\begin{sentence}
\underline{すみません}、\ruby{\underline{出口}}{でぐち}\ruby{は}{わ}\underline{どこ}ですか。\\\hline
sumimasen de{\middel}guchi wa doko desu ka.\\\hline
(excuse me) (exit) (where)\\\hline
=\textit{Excuse me, where's the exit?}\\\hline
\end{sentence}
\end{multicols}
\vspace*{\fill}
%%--------------------------------------------------------------------
%%--------------------------------------------------------------------
\pagebreak
\vspace*{\fill}
\begin{multicols}{2}
\begin{glyph}{Life (生)}{life}\label{glyph:life}
One of the most fascinating characters in the language, this kanji deals with life in all its guises.  Along with the ideas of ``birth'', ``production'', and ``growing'', it also incorporates the notions of ``pure'', ``raw'', and ``fresh'', in everything from beer to silk.\\\hline
When used as a suffix, it will often indicate ``student''.
\end{glyph}

\begin{comps}
\comprtwocone & 生 \\\hline
    \comprthreecone & Life \\\hline
\end{comps}

\begin{remglyph}
    Part of the Kanji radicals (\#{100}).\\\hline
\end{remglyph}

\begin{prons}
\pronsrtwocone & \textbf{セイ (SEI)} & \textbf{い、う、なま (i, u, nama)}\\\hline
\rye{3}{なま only appears in the first position.}
\pronsrthreecone & ショウ (SH\={O}) & お、は、き (o, ha, ki) \\\hline
\rye{3}{In terms of variety of pronunciations, kanji do not come any more complicated than this.\\\hline
い and う are verb stems (as are お and は of the less-common 訓読み).}
\end{prons}

\begin{remprons}
    \hooglig{Life} is full \comkun{interesting} \comkun{oomph} when
    \comon{sailing} in \comkun{Namaqualand}.\\\hline
    \hooglig{Life} at the \ucomon{shore} \ucomkun{obligates} me to
    \ucomkun{keep} \ucomkun{halting}.\\\hline
\end{remprons}

%{\middel}
% &  ()\\\hline
\begin{somme}
    生きる (intr.) & \textit{to live} & い{\middel}きる \mbox{(i{\middel}kiru)}\\\hline
    生かす (tr.) & \textit{to give life to} & い{\middel}かす \mbox{(i{\middel}kasu)}\\\hline
    生まれる (intr.) & \textit{to be born} & う{\middel}まれる \mbox{(u{\middel}mareru)}\\\hline
    生む (tr.) & \textit{to give birth to} & う{\middel}む \mbox{(u{\middel}mu)}\\\hline
    生ビール & life + beer = \textit{draft beer} & なま{\middel}ビール \mbox{(nama{\middel}b\={i}ru)}\\\hline
    人生 & person + life = \textit{(human) life} & ジン{\middel}セイ \mbox{(JIN{\middel}SEI)}\\\hline
    \notyet{学生} & study + life = \textit{student} & ガク{\middel}セイ \mbox{(GAKU{\middel}SEI)}\\\hline
    \notyet{先生} & precede + life = \textit{teacher} & セン{\middel}セイ \mbox{(SEN{\middel}SEI)}\\\hline
\end{somme}

\begin{irrr}
    芝生{\notcov} & lawn + life = \textit{lawn} & しば{\middel}ふ (shiba{\middel}fu)\\\hline
\end{irrr}

\begin{sentence}
\ruby{\underline{生}}{なま}\underline{ビール}\ruby{は}{わ}\ruby{\underline{好}}{す}\underline{き}ですか。\\\hline
nama{\middel}b\={i}ru wa su{\middel}ki desu ka.\\\hline
(draft beer) (like)\\\hline
=\textit{Do you like draft beer?}\\\hline
\end{sentence}
\end{multicols}
%%--------------------------------------------------------------------
%%--------------------------------------------------------------------
\begin{multicols}{2}
\begin{glyph}{Cow (牛)}{ushi}\label{glyph:cow}
Cow; Ox; Bull.
\end{glyph}

\begin{comps}
\comprtwocone & 牛 \\\hline
\comprthreecone & Cow \\\hline
\end{comps}

\begin{remglyph}
Part of the Kanji radicals (\#{93}).\\\hline
\end{remglyph}

\begin{prons}
\pronsrtwocone & \textbf{ギュウ (GY\={U})} & \textbf{うし (ushi)}\\\hline
\pronsrthreecone & --- & --- \\\hline
\end{prons}

\begin{remprons}
    \hooglig{Cows} \comkun{shipped} by \comon{gyou}?\\\hline
\end{remprons}

%{\middel}
% &  ()\\\hline
\begin{somme}
    牛 & \textit{cow} & うし \mbox{(ushi)}\\\hline
    子牛 & child + cow = \textit{calf} & こ{\middel}うし \mbox{(ko{\middel}ushi)}\\\hline
    水牛 & water + cow = \textit{water buffalo} & スイ{\middel}ギュウ \mbox{(SUI{\middel}GY\={U})}\\\hline
    牛肉 & cow + meat = \textit{beef} & ギュウ{\middel}ニク \mbox{(GY\={U}{\middel}NIKU)}\\\hline
    \notyet{牛馬} & cow + horse = \textit{cattle and horses} & ギュウ{\middel}バ \mbox{(GY\={U}{\middel}BA)}\\\hline
\end{somme}

\begin{sentence}
\underline{アフリカ}で\ruby{\underline{水牛}}{スイギュウ}と\underline{キリン}\ruby{を}{お}\ruby{\underline{見}}{み}\underline{ました}。\\\hline
AFURIKA de SUI{\middel}GY\={U} to KIRIN o mi{\middel}mashita.\\\hline
(Africa) (water buffalo) (giraffe) (saw)\\\hline
=\textit{(I) saw a water buffalo and giraffe in Africa.}\\\hline
\end{sentence}
\end{multicols}
\vhrulefill{5pt}
\vspace*{\fill}
%
%\begin{multicols}{2}
%\chapter{Kanji 101--120}
%%%--------------------------------------------------------------------
%1--------------------------------------------------------------------
\vspace*{\fill}
\begin{glyph}{Beautiful (美)}{beautiful}\label{glyph:beautiful}
All things beautiful.
\end{glyph}

\begin{comps}
\comprtwocone & 羊 + 大 \\\hline
    \comprthreecone & Sheep (\ref{glyph:sheep}) + Large (\ref{glyph:large}) \\\hline
\end{comps}

\begin{remglyph}
    \hooglig{Beautiful} \remcom{large} \remcom{sheep}.\\\hline
\end{remglyph}

\begin{prons}
\pronsrtwocone & \textbf{ビ (BI)} & \textbf{うつく (utsuku)}\\\hline
\pronsrthreecone & --- & --- \\\hline
\end{prons}

\begin{remprons}
The \comon{big} \hooglig{beautiful} \comkun{boots coupon}.\\\hline
\end{remprons}

%{\middel}
% &  ()\\\hline
\begin{somme}
美しい & \textit{beautiful} & うつく{\middel}しい (utsuku{\middel}shii)\\\hline
美人 & beautiful + person = \textit{beautiful woman} &　ビ{\middel}ジン (BI{\middel}JIN) \\\hline
美化 & beautiful + change = \textit{beautification} & ビ{\middel}カ (BI{\middel}KA)\\\hline
    \notyet{美学} & beautiful + study = \textit{aesthetics} & ビ{\middel}ガク (BI{\middel}GAKU)\\\hline
\end{somme}

\begin{irrr}
    \notyet{美味しい} & beautiful + taste = \textit{delicious} & おい{\middel}しい (oi{\middel}shii)\\\hline	
\end{irrr}

\begin{sentence}
\ruby{\underline{山}}{やま}の\ruby{\underline{木}}{き}が\ruby{\underline{美しい}}{うつくしい}です。\\\hline
yama no ki ga utsuku{\middel}shii desu.\\\hline
(mountain) (trees) (beautiful)\\\hline
=\textit{The mountain trees are beautiful.}\\\hline
\end{sentence}
\end{multicols}
\vspace*{-\baselineskip}
\vhrulefill{5pt}
%%--------------------------------------------------------------------
%2--------------------------------------------------------------------
\begin{multicols}{2}
\begin{glyph}{Island (島)}{island}\label{glyph:island}
Island.
\end{glyph}

\begin{comps}
\comprtwocone & 鳥 + 山  \\\hline
    \comprthreecone &Bird (\ref{glyph:bird}) + Mountain (\ref{glyph:mountain})  \\\hline
\end{comps}

\begin{remglyph}
    This \remcom{bird} lives only on the \remcom{mountains} of a particular
    \hooglig{island}.\\\hline
\end{remglyph}

\begin{prons}
\pronsrtwocone & \textbf{トウ (T\={O})} & \textbf{しま (shima)}\\\hline
\pronsrthreecone & --- & --- \\\hline
\end{prons}

\begin{remprons}
\comkun{She mugs} \comon{tortured} \hooglig{islanders}.\\\hline
\end{remprons}

%{\middel}
% &  ()\\\hline
\begin{somme}
    島 & \textit{island} &しま (shima) \\\hline
    小島 & small + island = \textit{small island} & こ{\middel}じま (ko{\middel}jima)\\\hline
    島国 & island + country = \textit{island nation} & しま{\middel}ぐに (shima{\middel}guni)\\\hline
    半島 & half + island = \textit{Peninsula} & ハン{\middel}トウ (HAN{\middel}T\={O})\\\hline
    本島 & main + island = \textit{Main island} & ホン{\middel}トウ (HON{\middel}T\={O})\\\hline
\end{somme}
\begin{sentence}
    \underline{オウストラリア}と\ruby{\underline{日本}}{ニホン}\ruby{は}{わ}\ruby{\underline{島国}}{しまぐに}です。\\\hline
    \={O}sutoraria to NI{\middel}HON wa shima{\middel}guni desu.\\\hline
    (Australia) (Japan) (island nation) \\\hline
    =\textit{Australia and Japan are island nations.}\\\hline
\end{sentence}
\end{multicols}
\vspace*{-\baselineskip}
\vhrulefill{5pt}
\vspace*{\fill}
\pagebreak
%%--------------------------------------------------------------------
%3--------------------------------------------------------------------
\vspace*{\fill}
\begin{multicols}{2}
\begin{glyph}{Tiny (寸)}{tiny}\label{glyph:tiny}
Tiny.
\end{glyph}

\begin{comps}
\comprtwocone & 寸 \\\hline
    \comprthreecone & small/inch (\ref{glyph:tiny}) \\\hline
\end{comps}

\begin{remglyph}
    An \remcom{inch} is \hooglig{small}.\\\hline
\end{remglyph}

\begin{prons}
\pronsrtwocone & \textbf{スン (SUN)} & ---\\\hline
\pronsrthreecone & --- & --- \\\hline
\end{prons}

\begin{remprons}
    \hooglig{Tiny} \comon{soon}.\\\hline
\end{remprons}

%{\middel}
% &  ()\\\hline
\begin{somme}
    一寸  & one + tiny = \textit{a tiny bit}  & イッ{\middel}スン (IS{\middel}SUN) \\\hline
    \notyet{寸前}  & tiny + before = \textit{immediately before}  & スン{\middel}ゼン
    (SUN{\middel}ZEN) \\\hline
\end{somme}

\begin{sentence}
    \ruby{\underline{出}}{で}\underline{かける}{\;}\ruby{\underline{寸前}}{スンゼン}に\ruby{\underline{高木}}{たかぎ}\underline{さん}\ruby{を}{お}\ruby{\underline{見}}{み}\underline{ました}。\\\hline
    de{\middel}kakeru SUN{\middel}ZEN ni Taka{\middel}gi-san o mi{\middel}mashita.\\\hline
    (leave) (immediately before) (Takagi-san) (saw)\\\hline
    =\textit{I saw Takagi-san just before I left.} \\\hline
\end{sentence}
\end{multicols}
\vspace*{-\baselineskip}
\vhrulefill{5pt}
%%--------------------------------------------------------------------
%4--------------------------------------------------------------------
\begin{multicols}{2}
\begin{glyph}{Temple (寺)}{temple}\label{glyph:temple}
    (Buddhist) Temple.
\end{glyph}

\begin{comps}
\comprtwocone &  寸 + 土 \\\hline
    \comprthreecone & Tiny (\ref{glyph:tiny}) + Earth (\ref{glyph:earth}) \\\hline
\end{comps}

\begin{remglyph}
    A \hooglig{temple} is \remcom{tiny} compared to the entire \remcom{earth}.\\\hline
\end{remglyph}

\begin{prons}
\pronsrtwocone & \textbf{ジ (JI)} & \textbf{てら (tera)}\\\hline
\pronsrthreecone & ---  & ---  \\\hline
\end{prons}

\begin{remprons}
    The \comon{jig} is up.  You will not \comkun{terraform} this
    \hooglig{temple}.\\\hline
\end{remprons}

%{\middel}
% &  ()\\\hline
\begin{somme}
    寺  & \textit{temple}  & てら (tera) \\\hline
    寺田さん  & temple + main = \textit{Terada-san}  & てら{\middel}だ{\middel}さん (tera{\middel}da-san) \\\hline
    山寺  & mountain + temple = \textit{mountain temple}  & やま{\middel}でら (yama{\middel}dera) \\\hline
    古寺  & old + temple = \textit{old temple}  & ふる{\middel}でら (furu{\middel}dera) コ{\middel}ジ (KO{\middel}JI) \\\hline
\end{somme}

\begin{sentence}
    \ruby{\underline{山寺}}{やまでら}の\ruby{\underline{古}}{ふる}\underline{い}{\;}\ruby{\underline{門}}{モン}の\ruby{\underline{下}}{した}で\ruby{\underline{休}}{やす}\underline{みました}。\\\hline
    yama{\middel}dera no furu{\middel}i MON no shita de yasu{\middel}mimashita.\\\hline
    (mountain temple) (old) (gate) (under) (rested) \\\hline
    =\textit{(We) rested under the old gate of the mountain temple.}\\\hline
\end{sentence}
\end{multicols}
\vspace*{-\baselineskip}
\vhrulefill{5pt}
\vspace*{\fill}
\pagebreak
%%--------------------------------------------------------------------
%5--------------------------------------------------------------------
\begin{multicols}{2}
\begin{glyph}{Precede (先)}{precede}\label{glyph:precede}
All meanings derived from the general sense of ``preceding'', or the idea of an
    object ``ahead'' in a literal or figurative way.\\\hline
    Depending on the context in which it appears, this kanji can signify the
    ``tip'' or ``point'' of something, the general future (or the past, when used as
    in the third compound), or a destination.\\\hline
    It can also refer to people in the sense of ``the other party''.
\end{glyph}

\begin{comps}
\comprtwocone &  牛 + 儿\\\hline
    \comprthreecone & Cow (\ref{glyph:cow}) + Person (\ref{glyph:person}) \\\hline
\end{comps}

\begin{remglyph}
    The \remcom{cow} \hooglig{preceded} the \remcom{person}.\\\hline
\end{remglyph}

\begin{prons}
\pronsrtwocone & \textbf{セン (SEN)} & \textbf{さき (saki)}\\\hline
\pronsrthreecone & --- & --- \\\hline
\end{prons}

\begin{remprons}
    It's \comkun{sucky} that the \hooglig{preceding} rules are so
    \comon{central} to this school.\\\hline
\end{remprons}

%{\middel}
% &  ()\\\hline
\begin{somme}
    先  & \textit{precede}  & さき (saki) \\\hline
    先生  & precede + life = \textit{teacher}  & セン{\middel}セイ (SEN{\middel}SEI) \\\hline
    先月  & precede + moon (month) = \textit{last month}  & セン{\middel}ゲツ (SEN{\middel}GETSU) \\\hline
    口先  & mouth + precede = \textit{lip service}  & くち{\middel}さき (kuchi{\middel}saki)\\\hline
\end{somme}

\begin{sentence}
    \ruby{\underline{先生}}{センセイ}\ruby{は}{わ}\ruby{\underline{冬休}}{ふゆやす}\underline{み}に\ruby{\underline{日本}}{ニホン}{\;}\underline{アルプス}\ruby{へ}{え}\ruby{\underline{行}}{い}\underline{きました}。\\\hline
    SEN{\middel}SEI wa fuyu yasu{\middel}mi ni NI{\middel}HON arupusu e i{\middel}kimashita.\\\hline
    (teacher) (winter vacation) (Japanese Alps) (went)\\\hline
    =\textit{(Our) teacher went to the Japanese Alps for winter
    vacation.}\\\hline
\end{sentence}
\end{multicols}
%%--------------------------------------------------------------------
%6--------------------------------------------------------------------
\begin{multicols}{2}
\begin{glyph}{Open (開)}{open}\label{glyph:open}
    Open.\\\hline
    Secondary meanings include ``starting'' and the idea of ``developing''
    things (in the sense of opening or reclaiming land for example.)
\end{glyph}

\begin{comps}
\comprtwocone & 門 + 廾 + 一 \\\hline
    \comprthreecone & Gate (\ref{glyph:gate}) + two hands + One (\ref{glyph:one}) \\\hline
\end{comps}

\begin{remglyph}
    It takes \remcom{two hands} to \hooglig{open} the \remcom{gate} in
    \remcom{one} go.\\\hline
\end{remglyph}

\begin{prons}
\pronsrtwocone & \textbf{カイ (KAI)} & \textbf{あ; ひら (a; hira)}\\\hline
\pronsrthreecone & --- & --- \\\hline
    \rye{3}{In the general physical of ``to open'', あく and あける function as
    a normal intrasitive/transitive verb pair.\\\hline
    ひらく is more concerned with a metaphorical or non-physical sense of
    ``opening'', and appears primarily in expressions (in the ``opening'', for
    example, of a conference, of new lands, or of cherry blossoms).\\\hline
    Expressing the sense of a shop's front door being open would thus call for
    the verb あく, while conveying the idea of opening a business would require
    ひらく.\\\hline
    In the first position, this kanji invariably takes the onyomi.}
\end{prons}

\begin{remprons}
    \comkun{He rustled} through the \comkun{updates} in hopes of \hooglig{opening} new possibilities to \comon{kite-surfing}.\\\hline
\end{remprons}

%{\middel}
% &  ()\\\hline
\begin{somme}
    開く (intr)  & \textit{to be open}  & あ{\middel}く (a{\middel}ku) \\\hline
    開ける (tr)  & \textit{to open (something)}  & あ{\middel}ける (a{\middel}keru) \\\hline
    開く (intr/tr)  & \textit{to open; to develop}  & ひら{\middel}く (hira{\middel}ku) \\\hline
    開化  & open + change = \textit{becoming civilized}  & カイ{\middel}カ (KAI{\middel}KA) \\\hline
    開国  & open + country = \textit{to open up a country}  & カイ{\middel}コク (KAI{\middel}KOKU) \\\hline
    \notyet{開会}  & open + meet = \textit{to open a meeting}  & カイ{\middel}カイ (KAI{\middel}KAI) \\\hline
    \notyet{公開}  & public + open = \textit{open to the public}  & コ{\middel}ウカイ (K\={O}{\middel}KAI) \\\hline
\end{somme}

\begin{sentence}
    \ruby{\underline{水門}}{スイモン}\ruby{を}{お}\ruby{\underline{開}}{あ}\underline{けて}{\;}\ruby{\underline{下}}{くだ}\underline{さい}。\\\hline
    SUI{\middel}MON o a{\middel}kete kuda{\middel}sai.\\\hline
    (sluice gate) (open) (please)\\\hline
    =\textit{Please open the sluice gate.}\\\hline
\end{sentence}
\end{multicols}
\pagebreak
%%--------------------------------------------------------------------
%7--------------------------------------------------------------------
\vspace*{\fill}
\begin{multicols}{2}
\begin{glyph}{Bird (鳥)}{bird}\label{glyph:bird}
    Bird.
\end{glyph}

\begin{comps}
\comprtwocone & 鳥 \\\hline
    \comprthreecone & Bird (\ref{glyph:bird}) \\\hline
\end{comps}

\begin{remglyph}
\\\hline
\end{remglyph}

\begin{prons}
    \pronsrtwocone & \textbf{チョウ (CH\={O})} & \textbf{とり (tori)}\\\hline
\pronsrthreecone & --- & --- \\\hline
\end{prons}

\begin{remprons}
    \comkun{Toreador} \comon{chores} are for the \hooglig{birds}.\\\hline
\end{remprons}

%{\middel}
% &  ()\\\hline
\begin{somme}
    鳥  & \textit{bird}  &  とり (tori) \\\hline
    白鳥  & white + bird = \textit{swan}  & ハク{\middel}チョウ (HAKU{\middel}CH\={O}) \\\hline
    小鳥  & small + bird = \textit{small bird}  &  こ{\middel}とり (ko{\middel}tori) \\\hline
    \notyet{夜鳥}  & night + bird = \textit{nocturnal bird}  & ヤ{\middel}チョウ (YA{\middel}CH\={O})  \\\hline
    \notyet{愛鳥家}  & love + bird + house = \textit{bird lover}  & アイ{\middel}チョウ{\middel}カ (AI{\middel}CH\={O}{\middel}KA) \\\hline
\end{somme}

\begin{sentence}
    \underline{あの}{\;}\ruby{\underline{青}}{あお}\underline{い}{\;}\ruby{\underline{小鳥}}{ことり}の\ruby{\underline{名前}}{なまえ}が\ruby{\underline{分}}{わ}\underline{かります}か。\\\hline
    ano ao{\middel}i ko{\middel}tori no na{\middel}mae ga wa{\middel}karimasu ka.\\\hline
    (that) (blue) (little bird) (name) (understand)\\\hline
    =\textit{Do you know the name of that little blue bird?}\\\hline
\end{sentence}
\end{multicols}
\vspace*{-\baselineskip}
\vhrulefill{5pt}
%%--------------------------------------------------------------------
%8--------------------------------------------------------------------
\begin{multicols}{2}
\begin{glyph}{Evening (夕)}{evening}\label{glyph:evening}
Evening.
\end{glyph}

\begin{comps}
\comprtwocone & 夕 \\\hline
    \comprthreecone & Evening (\ref{glyph:evening}) \\\hline
\end{comps}

\begin{remglyph}
\\\hline
\end{remglyph}

\begin{prons}
    \pronsrtwocone & --- & \textbf{ゆう (y\={u})}\\\hline
    \pronsrthreecone & セキ (SEKI) & --- \\\hline
\end{prons}

\begin{remprons}
    \comkun{Youse} look \ucomon{sexy} this \hooglig{evening}.\\\hline
\end{remprons}

%{\middel}
% &  ()\\\hline
\begin{somme}
    夕べ  & \textit{evening}  & ゆう{\middel}べ (y\={u}{\middel}be) \\\hline
    夕日  & evening + sun = \textit{setting sun}  & ゆう{\middel}ひ (y\={u}{\middel}hi) \\\hline
    \notyet{夕立}  & evening + stand = \textit{late afternoon (rain) shower}  & ゆう{\middel}だち (y\={u}{\middel}dachi)\\\hline
\end{somme}

\begin{sentence}
    \ruby{\underline{山}}{やま}{\;}\underline{から}{\;}\ruby{\underline{見}}{み}\underline{る}{\;}\ruby{\underline{夕日}}{ゆうひ}が\ruby{\underline{美}}{うつく}\underline{しい}。\\\hline
    yama kara mi{\middel}ru y\={u}{\middel}hi ga utsuku{\middel}shii.\\\hline
    (mountain) (from) (see) (setting sun) (beautiful)\\\hline
    =\textit{The setting sun is beautiful when seen from the mountain.}\\\hline
\end{sentence}
\end{multicols}
\vspace*{-\baselineskip}
\vhrulefill{5pt}
\vspace*{\fill}
\pagebreak
%%--------------------------------------------------------------------
%9--------------------------------------------------------------------
\vspace*{\fill}
\begin{multicols}{2}
\begin{glyph}{Outside (外)}{outside}\label{glyph:outside}
Outside.\\\hline
The idea of ``foreign'' or ``other'' is also present, as the third example makes
clear.\\\hline
The less common readings convey some secondary meanings of removal and
    disconnection.
\end{glyph}

\begin{comps}
\comprtwocone & 夕 + 卜  \\\hline
    \comprthreecone &Evening (\ref{glyph:evening}) + divination  \\\hline
\end{comps}

\begin{remglyph}
    Go \remcom{outside} to do your \remcom{divination} in the \remcom{evening}.\\\hline
\end{remglyph}

\begin{prons}
\pronsrtwocone & \textbf{ガイ (GAI)} & \textbf{そと (soto)}\\\hline
    \pronsrthreecone &ゲ (GE) &ほか; はず (hoka; hazu) \\\hline
\end{prons}

\begin{remprons}
    \hooglig{Outside} is a \ucomkun{hawker} \comon{guy} that \ucomon{gets} all the \ucomkun{huzz} and buzz with his \comkun{sotto} voce remarks.\\\hline
\end{remprons}

%{\middel}
% &  ()\\\hline
\begin{somme}
    外  & \textit{outside}  & そと (soto) \\\hline
    外見  & outside + see = \textit{external appearance}  & ガイ{\middel}ケン (GAI{\middel}KEN) \\\hline
    外国人  & outside + country + person = \textit{foreigner}  & ガイ{\middel}コク{\middel}ジン (GAI{\middel}KOKU{\middel}JIN) \\\hline
    外出  & outside + exit = \textit{going out}  & ガイ{\middel}シュツ (GAI{\middel}SHUTSU) \\\hline
    \notyet{海外}  & sea + outside = \textit{overseas}  & カイ{\middel}ガイ (KAI{\middel}GAI) \\\hline
\end{somme}

\begin{sentence}
    \ruby{\underline{日本}}{ニホン}には\ruby{\underline{外国人}}{ガイコクジン}が\ruby{\underline{多}}{おお}\underline{い}。\\\hline
    NI{\middel}HON ni wa GAI{\middel}KOKU{\middel}JIN ga \={o}{\middel}i.\\\hline
    (Japan) (foreigners) (many)\\\hline
    =\textit{There are many foreigners in Japan.}\\\hline
\end{sentence}
\end{multicols}
\vspace*{-\baselineskip}
\vhrulefill{5pt}
%%--------------------------------------------------------------------
%10-------------------------------------------------------------------
\begin{multicols}{2}
\begin{glyph}{Tongue (舌)}{tongue}\label{glyph:tongue}
Tongue.
\end{glyph}

\begin{comps}
\comprtwocone &千 + 口  \\\hline
    \comprthreecone & Thousand (\ref{glyph:thousand}) + Mouth (\ref{glyph:mouth}) \\\hline
\end{comps}

\begin{remglyph}
    The \hooglig{tongue} in your \remcom{mouth} can lead to \remcom{thousands} of consequences.\\\hline
\end{remglyph}

\begin{prons}
\pronsrtwocone & --- & \textbf{した (shita)}\\\hline
    \pronsrthreecone &ゼツ (ZETSU)  & --- \\\hline
\end{prons}

\begin{remprons}
    \comkun{She tugged} \ucomon{Zetsu's} \hooglig{tongue}.\\\hline
\end{remprons}

%{\middel}
% &  ()\\\hline
\begin{somme}
    舌  & \textit{tongue}  & した (shita) \\\hline
\end{somme}

\begin{sentence}
    \ruby{\underline{内田}}{うちだ}\underline{さん}\ruby{は}{わ}\ruby{\underline{舌}}{した}に\underline{ピアス}\ruby{を}{お}\underline{しました}。\\\hline
    Uchi{\middel}da-san wa shita ni piasu o shimashita.\\\hline
    (Uchida-san) (togue) (piercing) (did)\\\hline
    =\textit{Uchida-san got her tongue pierced.}\\\hline
\end{sentence}
\end{multicols}
\vspace*{-\baselineskip}
\vhrulefill{5pt}
\vspace*{\fill}
\pagebreak
%%--------------------------------------------------------------------
%11-------------------------------------------------------------------
\vspace*{\fill}
\begin{multicols}{2}
\begin{glyph}{Diagram (図)}{diagram}\label{glyph:diagram}
Diagram/Drawing.\\\hline
By extension comes the idea of ``planning''.
\end{glyph}

\begin{comps}
\comprtwocone & 囗 + 斗 \\\hline
\comprthreecone & enlosure + volume measure  \\\hline
\end{comps}

\begin{remglyph}
    The \hooglig{diagram} depicts an \remcom{enclosure} of \remcom{volume
    measurements}.\\\hline
\end{remglyph}

\begin{prons}
\pronsrtwocone & \textbf{ズ; ト (ZU; TO)} & --- \\\hline
    \pronsrthreecone &---  &はか (haka) \\\hline
\end{prons}

\begin{remprons}
    The \comon{top} \comon{zucchini} \hooglig{diagram} inspired the \ucomkun{haka}.\\\hline
\end{remprons}

%{\middel}
% &  ()\\\hline
\begin{somme}
    図工  & diagram + craft = \textit{manual arts}  & ズ{\middel}コウ (ZU{\middel}K\={O}) \\\hline
    心電図  & heart + electric + diagram = \textit{electrocardiogram}  & シン{\middel}デン{\middel}ズ (SHIN{\middel}DEN{\middel}ZU)  \\\hline
    \notyet{海図}  & sea + diagram = \textit{nautical chart}  & カイ{\middel}ズ (KAI{\middel}ZU)  \\\hline
    \notyet{意図}  & mind + diagram = \textit{(one's) aim}  & イ{\middel}ト (I{\middel}TO) \\\hline
\end{somme}

\begin{sentence}
    \ruby{\underline{心電図}}{シンデンズ}{\;}\underline{モニタアー}\ruby{を}{お}{\;}\ruby{\underline{見}}{み}\underline{せて}{\;}\ruby{\underline{下}}{くだ}\underline{さい}。\\\hline
    SHIN{\middel}DEN{\middel}ZU monit\={a} o mi{\middel}sete kuda{\middel}sai.\\\hline
    (electrocardiogram) (monitor) (show) (please) \\\hline
    =\textit{Please show me the electrocardiogram monitor.}\\\hline
\end{sentence}
\end{multicols}
\vspace*{-\baselineskip}
\vhrulefill{5pt}
%%--------------------------------------------------------------------
%12-------------------------------------------------------------------
\begin{multicols}{2}
\begin{glyph}{Stand (立)}{stand}\label{glyph:stand}
Stand/Rise.
\end{glyph}

\begin{comps}
\comprtwocone & 立 \\\hline
    \comprthreecone & Stand (\ref{glyph:stand}) \\\hline
\end{comps}

\begin{remglyph}
\\\hline
\end{remglyph}

\begin{prons}
\pronsrtwocone & \textbf{リツ (RITSU)} & \textbf{た (ta)}\\\hline
    \pronsrthreecone &リュウ (RY\={U}) & --- \\\hline
    \rye{3}{Consistent with other on-yomi ending in ツ, RITSU ``doubles up''
    with unvoiced consonant sounds that follow it, as the fifth example below
    demonstrates.\\\hline
    It can also absorb its accompanying hiragana, as the final example
    shows.}
\end{prons}

\begin{remprons}
    I cannot \hooglig{stand} \comon{leets}.  It is a \ucomon{leukemia} \comkun{tunneling} its way into our society.\\\hline
\end{remprons}

%{\middel}
% &  ()\\\hline
\begin{somme}
    立つ (intr) & \textit{to stand}  & た{\middel}つ (ta{\middel}tsu) \\\hline
    立てる (tr) & \textit{to set up}  & た{\middel}てる (ta{\middel}teru) \\\hline
    自立  & self + stand = \textit{independence}  & ジ{\middel}リツ (JI{\middel}RITSU)  \\\hline
    中立  & middle + stand = \textit{neutrality}  & チュウ{\middel}リツ (CH\={U}{\middel}RITSU) \\\hline
    立春  & stand + spring = \textit{first day of spring}  &  リッ{\middel}シュン
    (RIS{\middel}SHUN) \\\hline
    立ち入る  & stand + enter = \textit{to go into}  & た{\middel}ち い{\middel}る (ta{\middel}chi i{\middel}ru) \\\hline
    夕立  & evening + stand = \textit{late afternoon (rain) shower}  & ゆう{\middel}だち
    (y\={u}{\middel}dachi)\\\hline
\end{somme}

\begin{sentence}
    \ruby{\underline{女}}{おんな}の\ruby{\underline{人}}{ひと}が\ruby{\underline{美}}{うつく}しい\ruby{\underline{月}}{つき}の\ruby{\underline{下}}{した}に\ruby{\underline{立}}{た}\underline{っています}。\\\hline
    onna no hito ga utsuku{\middel}shii tsuki no shita ni tat{\middel}te imasu.\\\hline
    (woman) (person) (beautiful) (moon) (under) (is standing)\\\hline
    =\textit{The woman is standing under a beautiful moon.}\\\hline
\end{sentence}
\end{multicols}
\vspace*{-\baselineskip}
\vhrulefill{5pt}
\vspace*{\fill}
\pagebreak
%%--------------------------------------------------------------------
%13-------------------------------------------------------------------
\vspace*{\fill}
\begin{multicols}{2}
\begin{glyph}{Parent (親)}{parent}\label{glyph:parent}
``Parent'', together with the senses of intimacy and closeness.
\end{glyph}

\begin{comps}
\comprtwocone & 立 + 木 + 見 \\\hline
    \comprthreecone & stand (\ref{glyph:stand}) + tree (\ref{glyph:tree}) + see
    (\ref{glyph:see})\\\hline
\end{comps}

\begin{remglyph}
    ``\remcom{Stand} on a \remcom{tree}, What do you \remcom{see}?, Your
    \hooglig{parents}, your parents!''\\\hline
\end{remglyph}

\begin{prons}
\pronsrtwocone & \textbf{シン (SHIN)} & \textbf{おや (oya)}\\\hline
    \pronsrthreecone &---  &した (shita) \\\hline
\end{prons}

\begin{remprons}
    The \hooglig{parent} \ucomkun{she tugged} \comon{sheened} with animosity as
    she yelled ``\comkun{Oi ya}! come here''.\\\hline
\end{remprons}

%{\middel}
% &  ()\\\hline
\begin{somme}
    親  & \textit{parent}  & おや (oya) \\\hline
    父親  & father + parent = \textit{father}  & ちち{\middel}おや (chichi{\middel}oya) \\\hline
    親子  & parent + child = \textit{parent and child}  &おや{\middel}こ (oya{\middel}ko) \\\hline
    肉親  & meat + parent = \textit{blood relative}  &ニク{\middel}シン  (NIKU{\middel}SHIN) \\\hline
    \notyet{母親}  & mother + child = \textit{mother}  &はは{\middel}おや (haha{\middel}oya) \\\hline
    \notyet{親切}  & parent + cut = \textit{kindhearted}  &シン{\middel}セツ (SHIN{\middel}SETSU)  \\\hline
\end{somme}

\begin{sentence}
    \ruby{\underline{中川}}{なかがわ}\underline{さん}の\ruby{\underline{父親}}{ちちおや}\ruby{は}{わ}\ruby{\underline{山}}{やま}が\ruby{\underline{大好}}{ダイスキ}\underline{き}です。\\\hline
    Naka{\middel}gawa-san no chichi{\middel}oya wa yama ga DAI{\middel}suki desu.\\\hline
    (Nakagawa-san) (father) (mountain) (really likes)\\\hline
    =\textit{Nakagawa-san's father really likes the mountains.}\\\hline
\end{sentence}
\end{multicols}
\vspace*{-\baselineskip}
\vhrulefill{5pt}
%%--------------------------------------------------------------------
%14-------------------------------------------------------------------
\begin{multicols}{2}
\begin{glyph}{Star (星)}{star}\label{glyph:star}
Star.
\end{glyph}

\begin{comps}
\comprtwocone &日 + 生  \\\hline
    \comprthreecone & Sun (\ref{glyph:sun}) + Life (\ref{glyph:life})  \\\hline
\end{comps}

\begin{remglyph}
    When the \remcom{sun} shines over your \remcom{life}, it is impossible to
    not feel like a \hooglig{star}.\\\hline
\end{remglyph}

\begin{prons}
\pronsrtwocone & \textbf{セイ (SEI)} & \textbf{ほし (hoshi)}\\\hline
    \pronsrthreecone &ショウ (SH\={O})  &---  \\\hline
\end{prons}

\begin{remprons}
    \comon{Sailing} according to the \comkun{horsy} \hooglig{stars} will lead us to the \ucomon{shore}.\\\hline
\end{remprons}

%{\middel}
% &  ()\\\hline
\begin{somme}
    星  & \textit{star}  & ほし (hoshi) \\\hline
    水星  & water + star = \textit{Mercury}  & スイ{\middel}セイ (SUI{\middel}SEI) \\\hline
    金星  & gold + star = \textit{Venus}  & キン{\middel}セイ (KIN{\middel}SEI) \\\hline
    火星  & fire + star = \textit{Mars}  & カ{\middel}セイ (KA{\middel}SEI) \\\hline
    木星  & tree + star = \textit{Jupiter}  & モク{\middel}セイ (MOKU{\middel}SEI) \\\hline
    土星  & earth + star = \textit{Saturn}  & ド{\middel}セイ (DO{\middel}SEI) \\\hline
\end{somme}

\begin{sentence}
    \underline{あの}{\;}\ruby{\underline{明}}{あか}るい\ruby{\underline{星}}{ほし}\ruby{は}{わ}\ruby{\underline{金星}}{キンセイ}だと\ruby{\underline{思}}{おも}\underline{い}ますか。\\\hline
    ano aka{\middel}rui hoshi wa KIN{\middel}SEI da to omo{\middel}imasu ka.\\\hline
    (that) (bright) (star) (Venus) (think)\\\hline
    =\textit{Do you think that bright star is Venus?}\\\hline
\end{sentence}
\end{multicols}
\vspace*{-\baselineskip}
\vhrulefill{5pt}
\vspace*{\fill}
\pagebreak
%%--------------------------------------------------------------------
%15-------------------------------------------------------------------
\vspace*{\fill}
\begin{multicols}{2}
\begin{glyph}{Autumn (秋)}{autumn}\label{glyph:autumn}
Autumn.
\end{glyph}

\begin{comps}
\comprtwocone & 禾 + 火  \\\hline
    \comprthreecone & grain ear + Fire (\ref{glyph:fire})  \\\hline
\end{comps}

\begin{remglyph}
    \hooglig{Autumn} being the time of the year when the stubble of \remcom{wheat} was set on \remcom{fire} after harvesting.\\\hline
\end{remglyph}

\begin{prons}
    \pronsrtwocone & \textbf{シュウ (SH\={U})} & \textbf{あき (aki)}\\\hline
\pronsrthreecone & ---  & --- \\\hline
\end{prons}

\begin{remprons}
    \comkun{Akimbo} \hooglig{autumn} \comon{shoes}.\\\hline
\end{remprons}

%{\middel}
% &  ()\\\hline
\begin{somme}
    秋  & \textit{autumn}  & あき (aki) \\\hline
    立秋  & stand + autumn = \textit{First day of Autumn}  & リッ{\middel}シュウ
    (RIS{\middel}SH\={U}) \\\hline
    春夏秋冬  & spring + summer + autumn + winter = \textit{The four seasons
    year around}  & シュン{\middel}カ{\middel}シュウ{\middel}トウ (SHUN{\middel}KA{\middel}SH\={U}{\middel}T\={O}) \\\hline
    \notyet{秋空}  & autumn + empty = \textit{autumn sky}  & あき{\middel}ぞら (aki{\middel}zora)  \\\hline
    \notyet{秋分}  & autumn + part = \textit{autumnal equinox}  & シュウ{\middel}ブン (SH\={U}{\middel}BUN) \\\hline
\end{somme}

\begin{sentence}
    \ruby{\underline{秋}}{あき}に\ruby{は}{わ}\ruby{\underline{日本}}{ニホン}の\ruby{\underline{山}}{やま}が\ruby{\underline{美}}{うつく}\underline{しい}。\\\hline
    aki ni wa NI{\middel}HON no yama ga utsuku{\middel}shii.\\\hline
    (autumn) (Japan) (mountain) (beautiful)\\\hline
    =\textit{Japan's mountains are beautiful in autumn.}\\\hline
\end{sentence}
\end{multicols}
\vspace*{-\baselineskip}
\vhrulefill{5pt}
%%--------------------------------------------------------------------
%16-------------------------------------------------------------------
\begin{multicols}{2}
\begin{glyph}{West (西)}{west}\label{glyph:west}
West.
\end{glyph}

\begin{comps}
\comprtwocone & 西 \\\hline
    \comprthreecone & West (\ref{glyph:west}) \\\hline
\end{comps}

\begin{remglyph}
\\\hline
\end{remglyph}

\begin{prons}
\pronsrtwocone & \textbf{セイ (SEI)} & \textbf{にし (nishi)}\\\hline
    \pronsrthreecone &サイ (SAI) & --- \\\hline
\end{prons}

\begin{remprons}
    We found our \comkun{niche} of \ucomon{psychedellic} drugs by \comon{sailing} \hooglig{west}.\\\hline
\end{remprons}

%{\middel}
% &  ()\\\hline
\begin{somme}
    西  & \textit{west}  & にし (nishi) \\\hline
    西口  & west + mouth = \textit{western exit}  & にし{\middel}ぐち (nishi{\middel}guchi) \\\hline
    北北西  & north + north + west = \textit{north-northwest}  & ホク{\middel}ホク{\middel}セイ
    (HOKU{\middel}HOKU{\middel}SEI) \\\hline
    南西  & south + west = \textit{southwest}  & ナン{\middel}セイ (NAN{\middel}SEI) \\\hline
    南南西  & south + south + west = \textit{south-southwest}  & ナン{\middel}ナン{\middel}セイ
    (NAN{\middel}NAN{\middel}SEI) \\\hline
    北西  & north + west = \textit{northwest}  & ホク{\middel}セイ (HOKU{\middel}SEI)  \\\hline
\end{somme}

\begin{sentence}
    \ruby{\underline{西口}}{にしぐち}から\ruby{\underline{出}}{で}\underline{て}{\;}\ruby{\underline{父}}{ちち}に\ruby{\underline{会}}{あ}\underline{いました}。\\\hline
    nishi{\middel}guchi kara de{\middel}te chichi ni a{\middel}imashita.\\\hline
    (west exit) (exit) (father) (met)\\\hline
    =\textit{(She) left by the west exit and met her father.}\\\hline
\end{sentence}
\end{multicols}
\vspace*{-\baselineskip}
\vhrulefill{5pt}
\vspace*{\fill}
\pagebreak
%%--------------------------------------------------------------------
%17-------------------------------------------------------------------
\vspace*{\fill}
\begin{multicols}{2}
\begin{glyph}{Sword (刀)}{sword}\label{glyph:sword}
Sword.
\end{glyph}

\begin{comps}
\comprtwocone & 刀 \\\hline
    \comprthreecone & Sword (\ref{glyph:sword}) \\\hline
\end{comps}

\begin{remglyph}
\\\hline
\end{remglyph}

\begin{prons}
    \pronsrtwocone & \textbf{トウ (T\={O})} & \textbf{かたな (katana)}\\\hline
\pronsrthreecone & --- & --- \\\hline
\end{prons}

\begin{remprons}
    He was \comon{tortured} with a \comkun{katana} \hooglig{sword}.\\\hline
\end{remprons}

%{\middel}
% &  ()\\\hline
\begin{somme}
    刀  & \textit{sword}  & かたな (katana) \\\hline
    刀工  & sword + craft = \textit{sword maker}  & トウ{\middel}コウ (T\={O}{\middel}K\={O}) \\\hline
    日本刀  & sun + main + sword = \textit{japanese sword}  & ニ{\middel}ホン{\middel}トウ (NI{\middel}HON{\middel}T\={O}) \\\hline
\end{somme}

\begin{sentence}
    \ruby{\underline{先月}}{センゲツ}{\;}\ruby{\underline{東京}}{トウコウ}で\ruby{\underline{日本刀}}{ニホントウ}\ruby{を}{お}\ruby{\underline{買}}{か}\underline{いました}。\\\hline
    SEN{\middel}GETSU T\={O}{\middel}KY\={O} de NI{\middel}HON{\middel}T\={O} o ka{\middel}imashita.\\\hline
    (last month) (Tokyo) (Japanese Sword) (bought)\\\hline
    =\textit{Last month I bought a Japanese sword in Tokyo.}\\\hline
\end{sentence}
\end{multicols}
\vspace*{-\baselineskip}
\vhrulefill{5pt}
%%--------------------------------------------------------------------
%18-------------------------------------------------------------------
\begin{multicols}{2}
\begin{glyph}{Cut (切)}{cut}\label{glyph:cut}
Cut.\\\hline
When the kunyomi is use as a suffix, it can impart a sense of completely
    finishing something, or acting in a decisive manner.  The final example
    provides an example.
\end{glyph}

\begin{comps}
\comprtwocone &匕 + 刀  \\\hline
    \comprthreecone & spoon + Sword (\ref{glyph:sword}) \\\hline
\end{comps}

\begin{remglyph}
    That \remcom{sword} can \hooglig{cut} a \remcom{spoon}.\\\hline
\end{remglyph}

\begin{prons}
\pronsrtwocone & \textbf{セツ (SETSU)} & \textbf{き (ki)}\\\hline
    \pronsrthreecone &サイ (SAI) & --- \\\hline
\end{prons}

\begin{remprons}
    I \comkun{keep} \hooglig{cutting} \comon{sets} into my \ucomon{psyche}.\\\hline
\end{remprons}

%{\middel}
% &  ()\\\hline
\begin{somme}
    \rye{3}{As the fourth example illustrates, this is another kanji that
    behaves like 買.}
    切れる  & \textit{to break by itself}  & き{\middel}れる (ki{\middel}reru) \\\hline
    切る  & \textit{to cut (something)}  & き{\middel}る (ki{\middel}ru) \\\hline
    切り下げる  & cut + lower = \textit{to devalue}  & き{\middel}り さ{\middel}げる (ki{\middel}ri sa{\middel}geru) \\\hline
    切下げ  & cut + lower = \textit{devaluation}  & き{\middel}り さ{\middel}げ (ki{\middel}ri sa{\middel}ge) \\\hline
    親切  & parent + cut = \textit{kindhearted}  & シン{\middel}セツ (SHIN{\middel}SETSU) \\\hline
    大切  & large + cut = \textit{important}  & タイ{\middel}セツ (TAI{\middel}SETSU) \\\hline
    思い切る  & think + cut = \textit{to make up one's mind}  & おも{\middel}い き{\middel}る (omo{\middel}i ki{\middel}ru) \\\hline
\end{somme}

\begin{irrr}
    \rye{3}{As the kunyomi of 切 ``doubles up'' with the kanji following it in
    this next pair of words, these compounds are best learned as irregular
    readings.}
    切手  & cut + hand = \textit{postage stamp}  & きっ{\middel}て (kit{\middel}te) \\\hline
    切符{\notcov}  & cut + symbol = \textit{ticket}  & きっ{\middel}プ (kip{\middel}pu)  \\\hline
\end{irrr}

\begin{sentence}
    \ruby{\underline{水力}}{スイリョク}と\ruby{\underline{火力}}{カリョク}\ruby{は}{わ}\ruby{\underline{大切}}{タイセツ}ですね。\\\hline
    SUI{\middel}RYOKU to KA{\middel}RYOKU wa TAI{\middel}SETSU desu ne.\\\hline
    (hydro power) (thermal power) (important)\\\hline
    =\textit{Hydro and thermal power are important, aren't they?}\\\hline
\end{sentence}
\end{multicols}
\vspace*{-\baselineskip}
\vhrulefill{5pt}
\vspace*{\fill}
\pagebreak
%%--------------------------------------------------------------------
%19-------------------------------------------------------------------
\vspace*{\fill}
\begin{multicols}{2}
\begin{glyph}{Day of the week (曜)}{day_of_the_week}\label{glyph:day_of_the_week}
Day of the week.
\end{glyph}

\begin{comps}
\comprtwocone &日 + 彐 + 彐 + 隹  \\\hline
    \comprthreecone &Sun (\ref{glyph:sun}) + pig heads + small bird  \\\hline
\end{comps}

\begin{remglyph}
    Every \hooglig{day of the week}, the \remcom{small birds} sit on the \remcom{pig heads}.\\\hline
\end{remglyph}

\begin{prons}
    \pronsrtwocone & \textbf{ヨウ (Y\={O})} & ---\\\hline
\pronsrthreecone & --- & --- \\\hline
\end{prons}

\begin{remprons}
    We \comon{yawn} at least once \hooglig{everday of the week}.\\\hline
\end{remprons}

%{\middel}
% &  ()\\\hline
\begin{somme}
    日曜日  & sun + day of the week + sun (day) = \textit{Sunday}  & ニチ{\middel}ヨウ{\middel}び
    (NICHI{\middel}Y\={O}{\middel}bi)\\\hline
    月曜日  & mooon + day of the week + sun (day) = \textit{Monday}  &
    ゲツ{\middel}ヨウ{\middel}び (GETSU{\middel}Y\={O}{\middel}bi) \\\hline
    火曜日  & fire + day of the week + sun (day) = \textit{Tuesday}  & カ{\middel}ヨウ{\middel}び
    (KA{\middel}Y\={O}{\middel}bi)\\\hline
    水曜日  & water + day of the week + sun (day) = \textit{Wednesday}  &
    スイ{\middel}ヨウ{\middel}び (SUI{\middel}Y\={O}{\middel}bi) \\\hline
    木曜日  & tree + day of the week + sun (day) = \textit{Thursday}  &
    モク{\middel}ヨウ{\middel}び (MOKU{\middel}Y\={O}{\middel}bi) \\\hline
    金曜日  & gold + day of the week + sun (day) = \textit{Friday}  &
    キニ{\middel}ヨウ{\middel}び
    (KIN{\middel}Y\={O}{\middel}bi)\\\hline
    土曜日  & earth + day of the week + sun (day) = \textit{Saturday}  &
    ド{\middel}ヨウ{\middel}び (DO{\middel}Y\={O}{\middel}bi) \\\hline
\end{somme}

\begin{sentence}
    \ruby{\underline{土曜日}}{ドヨウび}に\underline{パーティー}が\underline{あります}。\\\hline
    DO{\middel}Y\={O}{\middel}bi ni p\={a}{\middel}t\={i} ga arimasu.\\\hline
    (Saturday) (party) (is)\\\hline
    =\textit{There's a party on Saturday}\\\hline
\end{sentence}
\end{multicols}
\vspace*{-\baselineskip}
\vhrulefill{5pt}
%%--------------------------------------------------------------------
%20-------------------------------------------------------------------
\begin{multicols}{2}
\begin{glyph}{Around (周)}{around}\label{glyph:around}
Around, as in the sense of a lap or circuit.
\end{glyph}

\begin{comps}
\comprtwocone & 冂 +  土 + 口 \\\hline
    \comprthreecone & wilderness + Earth (\ref{glyph:earth}) + Mouth (\ref{glyph:mouth}) \\\hline
\end{comps}

\begin{remglyph}
    \hooglig{Around} the \remcom{earthen} \remcom{wilderness}, your
    \remcom{mouths} cannot be satisfied.\\\hline
\end{remglyph}

\begin{prons}
    \pronsrtwocone & \textbf{シュウ (SH\={U})} & --- \\\hline
\pronsrthreecone & --- & --- \\\hline
\end{prons}

\begin{remprons}
    \comon{Shoe} \hooglig{lap}.\\\hline
\end{remprons}

%{\middel}
% &  ()\\\hline
\begin{somme}
    円周  & circle + around = \textit{circumference}  & エン{\middel}シュウ (EN{\middel}SH\={U}) \\\hline
    周回  & around + rotate = \textit{going around}  & シュウ{\middel}カイ (SH\={U}{\middel}KAI) \\\hline
    一周  & once + around = \textit{once around}  & イッ{\middel}シュウ (IS{\middel}SH\={U}) \\\hline
\end{somme}

\begin{sentence}
    \ruby{\underline{寺田}}{てらだ}\underline{さん}\ruby{は}{わ}\underline{ニュウジランド}\ruby{を}{お}\ruby{\underline{一周}}{イッシュウ}{\;}\underline{しました}。\\\hline
    Tera{\middel}da-san wa Ny\={u}j\={i}rando o IS{\middel}SH\={U} shimashita.\\\hline
    (Terada-san) (New Zealand) (once around) (did)\\\hline
    =\textit{Terada-san did a circuit of New Zealand.}\\\hline
\end{sentence}
\end{multicols}
\vspace*{-\baselineskip}
\vhrulefill{5pt}
%%--------------------------------------------------------------------
%%--------------------------------------------------------------------
\vspace*{\fill}
\pagebreak

%
%\begin{multicols}{2}
%\chapter{Kanji 121--140}
%%%--------------------------------------------------------------------
%1--------------------------------------------------------------------
\begin{glyph}{Read (読)}{read}\label{glyph:read}
Read.
\end{glyph}
\gindex{Read}{読}\strindex{14}{読}

\begin{comps}
\comprtwocone & 言 + 売 \\\hline
\comprthreecone & Say/Speaking/Speech + Sell (\ref{glyph:sell}) \\\hline
\end{comps}

\begin{remglyph}
\hooglig{Reading} material, they \remcom{say}, \remcom{sells} a lot.\\\hline
\end{remglyph}

\begin{prons}
\pronsrtwocone & \textbf{ドク (DOKU)} & \textbf{よ (yo)}\\\hline
\pronsrthreecone & トク (TOKU)  & --- \\\hline
\end{prons}

\begin{remprons}
\hooglig{Reading} about the \ucomon{talking} \comkun{yomp} moving \comon{dockward}.\\\hline
\end{remprons}

%{\middel}
% &  ()\\\hline
\begin{somme}
読む  & \textit{to read} & よ{\middel}む \mbox{(yo{\middel}mu)} \\\hline
読み切れ  & read + cut = \textit{to finish reading} & よ{\middel}み き{\middel}る \mbox{(yo{\middel}mi ki{\middel}re)} \\\hline
読者  & read + inidividual = \textit{reader}  & ドク{\middel}シャ \mbox{(DOKU{\middel}SHA)} \\\hline
\notyet{読み物}  & read + thing = \textit{reading matter}   & よみ{\middel} もの \mbox{(yo{\middel}mi mono)} \\\hline
\notyet{多読家} & many + read + house = \textit{well-read person} & タ{\middel}ドク{\middel}シャ \mbox{(TA{\middel}DOKU{\middel}KA)} \\\hline
\end{somme}

\begin{decomp}{6}
此の&古い&本&を&読んでいます&か。\\\hline
この&ふるい&ホン&を&よんでいます&か\\\hline
this&old&book&を&are reading&か\\\hline
\multicolumn{6}{|r|}{\textit{Are you reading this old book?}}\\\hline
\end{decomp}

\end{multicols}
%%--------------------------------------------------------------------
%2--------------------------------------------------------------------
\begin{multicols}{2}
\begin{glyph}{Road (道)}{road}\label{glyph:road}
Road.  Way.\\\hline
This interesting 漢字 can signify both a \textit{physical road} and a spiritual \textit{path}.  In this latter sense, it is often translated as \textit{the way}.
\end{glyph}
\gindex{Road}{道}\strindex{12}{道}

\begin{comps}
\comprtwocone & 首 + 辵辶⻌⻍  \\\hline
\comprthreecone & Neck/Head + Road/Walk \\\hline
\end{comps}

\begin{remglyph}
\hooglig{Road} travelling with my \remcom{neck}-support.\\\hline
\end{remglyph}

\begin{prons}
\pronsrtwocone & \textbf{ドウ (D\={O})} & \textbf{みち (michi)}\\\hline
\pronsrthreecone & トウ (T\={O})  & ---  \\\hline
\end{prons}

\begin{remprons}
\hooglig{Roads} in \comkun{Michigan} are so \ucomon{torturous} that you'd rather open the \comon{door}, step out, and start walking.\\\hline
\end{remprons}

%{\middel}
% &  ()\\\hline
\begin{somme}
道 & \textit{road; way}  & みち \mbox{(michi)}  \\\hline
小道 & small + road = \textit{path}  & こ{\middel}みち \mbox{(ko{\middel}michi)}  \\\hline
水道 & water + door = \textit{water supply system}  & スイ{\middel}ドウ \mbox{(SUI{\middel}D\={O})}  \\\hline
下水道 & lower + water + road = \textit{sewer system}  & ゲ{\middel}スイ{\middel}ドウ \mbox{(GE{\middel}SUI{\middel}D\={O})} \\\hline
歩道 & walk + road = \textit{sidewalk}  & ホ{\middel}ドウ \mbox{(HO{\middel}D\={O})}  \\\hline
国道 & country + road = \textit{national highway}  & コク{\middel}ドウ \mbox{(KOKU{\middel}T\={O})}  \\\hline
\notyet{神道} & God + road = \textit{Japanese religion}  & シン{\middel}トウ \mbox{(SHIN{\middel}T\={O})} \\\hline
\end{somme}

\begin{decomp}{5}
右&の&道&を&行け。\\\hline
みち&の&みち&を&いけ\\\hline
right&の&road&を&to go\\\hline
\multicolumn{5}{|r|}{\textit{Take the right road.}}\\\hline
\end{decomp}

\end{multicols}
%%--------------------------------------------------------------------
%3--------------------------------------------------------------------
\begin{multicols}{2}
\begin{glyph}{Heaven (天)}{heaven}\label{glyph:heaven}
Heaven.  The Heavens.\\\hline
This character can indicate both the spiritual sense of \textit{heaven} as well as the more literal notions of \textit{sky} and \textit{air}.
\end{glyph}
\gindex{Heaven}{天}\strindex{4}{天}

\begin{comps}
\comprtwocone & 大 + 一 \\\hline
\comprthreecone & Large + One \\\hline
\end{comps}

\begin{remglyph}
\hooglig{Heaven} is \remcom{one} \remcom{large} place for the saints.\\\hline
\end{remglyph}

\begin{prons}
\pronsrtwocone & \textbf{テン (TEN)} & ---\\\hline
\pronsrthreecone & ---  & あめ、あま (ame, ama) \\\hline
\end{prons}

\begin{remprons}
\hooglig{Heaven} is not \comon{tense} because there is no need for \ucomkun{amends} or \ucomkun{amalgamations}.\\\hline
\end{remprons}

%{\middel}
% &  ()\\\hline
\begin{somme}
\rye{3}{Note how the 音読み for \textit{country} is voiced in the first example.  In the fourth compound, the 音読み of 王 changes from オウ to ノウ in order to make the pronunciation smoother.  This type of phonetic change is extremely rare.}
天国  & heaven + country = \textit{heaven}  & テン{\middel}コク \mbox{(TEN{\middel}GOKU)} \\\hline
天下  & heaven + lower = \textit{the whole land}  & テン{\middel}カ \mbox{(TEN{\middel}KA)} \\\hline
天体  & heaven + body = \textit{heavenly body}  & テン{\middel}タイ \mbox{(TEN{\middel}TAI)}  \\\hline
天王星  & heaven + king + star = \textit{Uranus}  & テン{\middel}ノウ{\middel}セイ \mbox{(TEN{\middel}N\={O}{\middel}SEI)}  \\\hline
\notyet{天気}  & heaven + spirit = \textit{weather}  & テン{\middel}キ \mbox{(TEN{\middel}KI)}  \\\hline
\end{somme}

\begin{decomp}{8}
明日&の&天気&は&雨&だ&そう&です。\\\hline
あした&の&テンキ&は&あめ&だ&そう&です\\\hline
tomorrow&の&weather&は&rain&だ&seems&be / is\\\hline
\multicolumn{8}{|r|}{\textit{It seems the weather will be rainy tomorrow.}}\\\hline
\end{decomp}

\end{multicols}
%%--------------------------------------------------------------------
%4--------------------------------------------------------------------
\begin{multicols}{2}
\begin{glyph}{Part (分)}{part}\label{glyph:part}
This important character suggests the ideas of \textit{dividing} and \textit{portion}.  The sense of \textit{division} is always present in some way.\\\hline
The first compound below can be though of as \textit{parting} the unknown to arrive at understanding.
\end{glyph}
\gindex{Part}{分}\strindex{4}{分}

\begin{comps}
\comprtwocone & 八 + 刀 \\\hline
\comprthreecone & Eight + Sword \\\hline
\end{comps}

\begin{remglyph}
\hooglig{Part} the \remcom{volcano} with this \remcom{sword}.\\\hline
\end{remglyph}

\begin{prons}
\pronsrtwocone & \textbf{ブン (BUN)} & \textbf{わ (wa)}\\\hline
\pronsrthreecone & フン、ブ (FUN, BU) & ---  \\\hline
\rye{3}{フン relates to minutes (of time or arc), as in the example 五分.
ブ deals with percentages and rates; 五分 in this sense would be pronounced ゴブ, and mean either five or fifty percent.\\\hline
Note also that 大分, one of the many Japanese words expressing the idea of \textit{much} or \textit{very}, can be pronounced ダイブン or ダイブ.}
\end{prons}

\begin{remprons}
The \hooglig{part} where I \comkun{wander} to the \comon{boon} such that I \ucomon{foon}.\\\hline
\end{remprons}

%{\middel}
% &  ()\\\hline
\begin{somme}
分かる (intr.) & \textit{to understand}  & わ{\middel}かる \mbox{(wa{\middel}karu)}  \\\hline
分ける (tr.) & \textit{to part; to divide}  & わ{\middel}ける \mbox{(wa{\middel}keru)}  \\\hline
自分 & self + part = \textit{oneself}  & ジ{\middel}ブン \mbox{(JI{\middel}BUN)}  \\\hline
半分 & half + part = \textit{half}  & ハン{\middel}ブン \mbox{(HAN{\middel}BUN)}  \\\hline
春分 & spring + part = \textit{vernal equinox}  & シュン{\middel}ブン \mbox{(SHUN{\middel}BUN)}  \\\hline
秋分 & autumn + part = \textit{autumnal equinox}  & シュウ{\middel}ブン \mbox{(SH\={U}{\middel}BUN)}  \\\hline
五分 & five + part = \textit{five minutes}  & ゴフ{\middel}ン \mbox{(GO{\middel}FUN)}  \\\hline
五分 & five + part = \textit{five/fifty percent}  & ゴ{\middel}ブ \mbox{(GO{\middel}BU)}  \\\hline
大分 & large + part = \textit{much/very}  & ダイ{\middel}ブン、ダイ{\middel}ブ \mbox{(DAI{\middel}BU(N))}  \\\hline
\end{somme}

\begin{decomp}{4}
日本語&が&分かります&か。\\\hline
ニホンゴ&が&わかります&か\\\hline
Japanese&が&understand&か\\\hline
\multicolumn{4}{|r|}{\textit{Do you understand Japanese.}}\\\hline
\end{decomp}

\end{multicols}
\pagebreak
%%--------------------------------------------------------------------
%5--------------------------------------------------------------------

\begin{multicols}{2}
\begin{glyph}{Thing (物)}{thing}\label{glyph:thing}
Thing.  Object.\\\hline
Note that the writing order for \textit{cow} differs when it is used as a component; slanting the final stroke upward makes for a smoother transition to the next part of the character.
\end{glyph}
\gindex{Thing}{物}\strindex{8}{物}

\begin{comps}
\comprtwocone & 牛 + 勿 \\\hline
\comprthreecone & Cow + (\ldots) Not \\\hline
\end{comps}

\begin{remglyph}
That \hooglig{thing} can \remcom{not} be a \remcom{cow}.\\\hline
\end{remglyph}

\begin{prons}
\pronsrtwocone & \textbf{ブツ (BUTSU)} & \textbf{もの (mono)}\\\hline
\pronsrthreecone & モツ (MOTSU)  & --- \\\hline
\end{prons}

\begin{remprons}
The \hooglig{thing} about this \comkun{monologue} is \ucomon{motsurella} \comon{boots}.\\\hline
\end{remprons}

%{\middel}
% &  ()\\\hline
\begin{somme}
物 & \textit{thing}  & もの \mbox{(mono)} \\\hline
買い物 & buy + thing = \textit{shopping}  & か{\middel}い　もの \mbox{(ka{\middel}i mono)} \\\hline
読み物 & read + thing = \textit{reading matter}  & よ{\middel}み　もの \mbox{(yo{\middel}mi mono)} \\\hline
人物 & person + thing = \textit{one's character}  & ジン{\middel}ブツ \mbox{(JIN{\middel}BUTSU)} \\\hline
見物 & see + thing = \textit{sightseeing}  & ケン{\middel}ブツス \mbox{(KEN{\middel}BUTSU)} \\\hline
物体 & thing + body = \textit{an object}  & ブッ{\middel}タイ \mbox{(BUT{\middel}TAI)} \\\hline
\end{somme}

\begin{decomp}{6}
本田さん&は&どんな&人物&です&か。\\\hline
ホンださん&は&どんな&ジンブツ&です&か\\\hline
Honda-san&は&what kind of&character&be / is&か\\\hline
\multicolumn{6}{|r|}{\textit{What kind of person is Honda-san?}}\\\hline
\end{decomp}

\end{multicols}

%%--------------------------------------------------------------------
%6--------------------------------------------------------------------
\begin{multicols}{2}
\begin{glyph}{Axe (斤)}{axe}\label{glyph:axe}
Axe.\\\hline
Though rarely encountered on its own, this character serves as a component in several important 漢字.
\end{glyph}
\gindex{Axe}{斤}\strindex{4}{斤}

\begin{comps}
\comprtwocone & 斤 \\\hline
\comprthreecone & Axe/Kind of Weight Measure \\\hline
\end{comps}

\begin{remglyph}
Part of the 漢字 radicals (\#{69}).\\\hline
\end{remglyph}

\begin{prons}
\pronsrtwocone & \textbf{キン (KIN)} & ---\\\hline
\pronsrthreecone & --- & --- \\\hline
\end{prons}

\begin{remprons}
\hooglig{Axe} the \comon{kindred}.\\\hline
\end{remprons}

%{\middel}
% &  ()\\\hline
\begin{somme}
斤 & \textit{axe}  & キン \mbox{(KIN)}  \\\hline
\end{somme}

\begin{decomp}{4}
斤&が&読んでます&か。\\\hline
キン&が&よんでます&か\\\hline
axe&が&read&か\\\hline
\multicolumn{4}{|r|}{\textit{Can you read ``斤''?}}\\\hline
\end{decomp}

\end{multicols}

\pagebreak
%%--------------------------------------------------------------------
%7--------------------------------------------------------------------

\begin{multicols}{2}
\begin{glyph}{New (新)}{new}\label{glyph:new}
New.
\end{glyph}
\gindex{New}{新}\strindex{13}{新}

\begin{comps}
\comprtwocone & 立 + 木 + 斤 \\\hline
\comprthreecone & Stand + Tree + Axe\\\hline
\end{comps}

\begin{remglyph}
\hooglig{New} \remcom{trees} cannot \remcom{stand} \remcom{axes}.\\\hline
\end{remglyph}

\begin{prons}
\pronsrtwocone & \textbf{シン (SHIN)} & \textbf{あたら (atara)}\\\hline
\pronsrthreecone & ---  & あら、にい (ara, n\={i}) \\\hline
\end{prons}

\begin{remprons}
\hooglig{New} \ucomkun{arabesques} will make you \comon{sheen} with \ucomkun{neat} \comkun{ataraxia}.\\\hline
{\warn}\textit{Arabesque} is a ballet position.  {\warn}\textit{Ataraxia} is a state of serene calmness.\\\hline
\end{remprons}

%{\middel}
% &  ()\\\hline
\begin{somme}
新しい & \textit{new}  & あたら{\middel}しい \mbox{(atara{\middel}shii)} \\\hline
一新 & one + new = \textit{complete renewal}  & イッ{\middel}シン \mbox{(IS{\middel}SHIN)} \\\hline
新月 & new + moon = \textit{new moon}  & シン{\middel}ゲツ \mbox{(SHIN{\middel}GETSU)} \\\hline
新星 & new + star = \textit{nova}  & シン{\middel}セイ \mbox{(SHIN{\middel}SEI)}  \\\hline
\notyet{最新} & most + new = \textit{newest}  & サイ{\middel}シン \mbox{(SAI{\middel}SHIN)} \\\hline
\end{somme}

\begin{decomp}{4}
新しい&日本刀&を&買いました。\\\hline
あたらしい&ニホントウ&を&かいました\\\hline
new&Japanese sword&を&bought\\\hline
\multicolumn{4}{|r|}{\textit{(I) bought a new Japanese sword.}}\\\hline
\end{decomp}

\end{multicols}

%%--------------------------------------------------------------------
%8--------------------------------------------------------------------
\begin{multicols}{2}
\begin{glyph}{Flower (花)}{flower}\label{glyph:flower}
Flower.
\end{glyph}
\gindex{Flower}{花}\strindex{7}{花}

\begin{comps}
\comprtwocone & 艸 + 化 \\\hline
\comprthreecone & Grass + Change (\ref{glyph:change}) \\\hline
\end{comps}

\begin{remglyph}
\hooglig{Flowers} are \remcom{grass} \remcom{changelings}.\\\hline
\end{remglyph}

\begin{prons}
\pronsrtwocone & \textbf{カ (KA)} & \textbf{はな (hana)}\\\hline
\pronsrthreecone & --- & --- \\\hline
\end{prons}

\begin{remprons}
\hooglig{Flowers} that \comon{cuss} are \comkun{hung-up}.\\\hline
\end{remprons}

%{\middel}
% &  ()\\\hline
\begin{somme}
花 & \textit{flower}  & はな \mbox{(hana)} \\\hline
生花 & life + flower = \textit{flower arrangement}  & い{\middel}け　ばな \mbox{(i{\middel}ke bana)} \\\hline
花見 & flower + see = \textit{cherry blossom viewing}  & はな{\middel}み \mbox{(hana{\middel}mi)} \\\hline
花火 & flower + fire = \textit{fireworks}  & はな{\middel}び \mbox{(hana{\middel}bi)} \\\hline
火花 & fire + flower = \textit{sparks}  & ひ{\middel}ばな \mbox{(hi{\middel}bana)} \\\hline
開花 & open + flower = \textit{bloom}  & カイ{\middel}カ \mbox{(KAI{\middel}KA)}  \\\hline
\end{somme}

\begin{decomp}{5}
お&花見&に&行きません&か。\\\hline
お&はなみ&に&いきません&か\\\hline
お&flower viewing&に&to go&か\\\hline
\multicolumn{5}{|r|}{\textit{Why don't we go and see the cherry blossoms?}}\\\hline
\end{decomp}

\end{multicols}

\pagebreak
%%--------------------------------------------------------------------
%9--------------------------------------------------------------------

\begin{multicols}{2}
\begin{glyph}{Name (名)}{name}\label{glyph:name}
Name.\\\hline
By extension, this 漢字 also includes the ideas of \textit{reputation} and \textit{fame}.
\end{glyph}
\gindex{Name}{名}\strindex{6}{名}

\begin{comps}
\comprtwocone & 夕 + 口\\\hline
\comprthreecone & Evening + Mouth\\\hline
\end{comps}

\begin{remglyph}
\hooglig{Names} of \remcom{evening} \remcom{vampires}.\\\hline
\end{remglyph}

\begin{prons}
\pronsrtwocone & \textbf{メイ (MEI)} & \textbf{な (na)}\\\hline
\pronsrthreecone & ミョウ (MY\={O}) & --- \\\hline
\end{prons}

\begin{remprons}
\hooglig{Naming} the \ucomon{Myou} dolls is a \comkun{numbing} \comon{mayem}.\\\hline
\end{remprons}

%{\middel}
% &  ()\\\hline
\begin{somme}
名高い  & name + tall = \textit{renowned}  & な　だ{\middel}かい \mbox{(na da{\middel}kai)} \\\hline
有名  & have + name = \textit{famous}  & ユウ{\middel}メイ \mbox{(Y\={U}{\middel}MEI)} \\\hline
名山  & name + mountain = \textit{famous mountain}  & メイ{\middel}ザン \mbox{(MEI{\middel}ZAN)} \\\hline
名物  & name + thing = \textit{famous regional product}  & メイ{\middel}ブツ \mbox{(MEI{\middel}BUTSU)}  \\\hline
\notyet{名前}  & name + before = \textit{name}  & な{\middel}まえ \mbox{(na{\middel}mae)}  \\\hline
\end{somme}

\begin{decomp}{5}
彼女&は&有名&歌手&だ。\\\hline
かのジョ&は&ユウメイ&カシュ&だ\\\hline
she&は&famous&singer&be / is\\\hline
\multicolumn{5}{|r|}{\textit{She is a noted singer.}}\\\hline
\end{decomp}

\end{multicols}

%%--------------------------------------------------------------------
%10-------------------------------------------------------------------
\begin{multicols}{2}
\begin{glyph}{Mother (母)}{mother}\label{glyph:mother}
Mother.
\end{glyph}
\gindex{Mother}{母}\strindex{5}{母}

\begin{comps}
\comprtwocone & 母 \\\hline
\comprthreecone & Mother \\\hline
\end{comps}

\begin{remglyph}
Part of the 漢字 radicals (\#{80}).\\\hline
\end{remglyph}

\begin{prons}
\pronsrtwocone & \textbf{ボ (BO)} & \textbf{はは (haha)}\\\hline
\pronsrthreecone & --- & --- \\\hline
\end{prons}

\begin{remprons}
\hooglig{Mother} \comon{bottles} up all her \comkun{laughter}.\\\hline
\end{remprons}

%{\middel}
% &  ()\\\hline
\begin{somme}
母 & \textit{mother (responsible occupation)}  &  はは \mbox{(haha)}\\\hline
母親 & mother + parent = \textit{mother (formal)}  &  はは{\middel}おや \mbox{(haha{\middel}oya)}\\\hline
母国語 & mother + country + words = \textit{mother tongue}  & ボ{\middel}コク{\middel}ゴ \mbox{(BO{\middel}KOKU{\middel}GO)} \\\hline
母の日 & mother + sun (day) = \textit{Mother's day}  &  はは{\middel}の{\middel}ひ \mbox{(haha{\middel}no{\middel}hi)} \\\hline
父母 & father + mother = \textit{father and mother}  &  フ{\middel}ボ \mbox{(FU{\middel}BO)} \\\hline
\end{somme}

\begin{irrr}
お母さん  & \textit{mother}  &  おかあさん (ok\={a}san) \\\hline
伯母さん{\notcov}  & related + mother = \textit{aunt}  & おばさん (obasan) \\\hline
\end{irrr}

\end{multicols}
\begin{decomp}{7}
新しい&先生&の&母国語&は&英語&です。\\\hline
あたらしい&センセイ&の&ボコクゴ&は&エイゴ&です\\\hline
new&teacher&の&mother tongue&は&English&be / is\\\hline
\multicolumn{7}{|r|}{\textit{The new teacher's mother tongue is English.}}\\\hline
\end{decomp}

\pagebreak
%%--------------------------------------------------------------------
%11-------------------------------------------------------------------

\begin{multicols}{2}
\begin{glyph}{Private (私)}{private}\label{glyph:private}
Private.\\\hline
This character is also used for the first person singular, \textit{I}.
\end{glyph}
\gindex{Private}{私}\strindex{7}{私}

\begin{comps}
\comprtwocone & 禾 + 厶 \\\hline
\comprthreecone & Two-branch tree/Grain + Private \\\hline
\end{comps}

\begin{remglyph}
\hooglig{\remcom{Private}} \remcom{grain}.\\\hline
\end{remglyph}

\begin{prons}
\pronsrtwocone & \textbf{シ (SHI)} & \textbf{わたし (watashi)}\\\hline
\pronsrthreecone & --- & わたくし (watakushi) \\\hline
\end{prons}

\begin{remprons}
I \hooglig{privately} rejoiced--``\comkun{what a shield}, \ucomkun{what a
cushion}''.\\\hline
\end{remprons}

%{\middel}
% &  ()\\\hline
\begin{somme}
私 & \textit{I}  & わたし \mbox{(watashi)} \\\hline
私有 & private + have = \textit{privately owned}  & シ{\middel}ユウ \mbox{(SHI{\middel}Y\={U})} \\\hline
私道 & private + road = \textit{private road}  & シ{\middel}ドウ \mbox{(SHI{\middel}D\={O})} \\\hline
私物 & private + thing = \textit{private property}  & シ{\middel}ブツ \mbox{(SHI{\middel}BUTSU)} \\\hline
\notyet{私学}  & private + study = \textit{private school}  & シ{\middel}ガク \mbox{(SHI{\middel}GAKU)}  \\\hline
\end{somme}

\begin{decomp}{3}
私&を&愛してる？\\\hline
わたし&を&アイしてる\\\hline
I&を&to love\\\hline
\multicolumn{3}{|r|}{\textit{Do you love me?}}\\\hline
\end{decomp}

\end{multicols}

%%--------------------------------------------------------------------
%12-------------------------------------------------------------------
\begin{multicols}{2}
\begin{glyph}{Night (夜)}{night}\label{glyph:night}
Night.
\end{glyph}
\gindex{Night}{夜}\strindex{8}{夜}

\begin{comps}
\comprtwocone & 亠 + 人亻𠆢 + 夕 \\\hline
\comprthreecone & Top/Lid + Person + Evening\\\hline
\end{comps}

\begin{remglyph}
\hooglig{Night} \remcom{people} wear \remcom{tops} when it is \remcom{evening}.\\\hline
\end{remglyph}

\begin{prons}
\pronsrtwocone & \textbf{ヤ (YA)} & \textbf{よる、よ (yoru, yo)}\\\hline
\rye{3}{よる is only used when this 漢字 appears alone.\\\hline
よ and ヤ are distributed evenly throughout compounds.}
\pronsrthreecone & --- & --- \\\hline
\end{prons}

\begin{remprons}
A \hooglig{night}-time \comkun{yomp} to the \comon{yacht} with \comkun{Yoruichi}.\\\hline
{\warn}\textit{Yoruichi} 「夜一」 is a character in the famous \textit{Bleach} \mbox{アニメ}.\\\hline
\end{remprons}

%{\middel}
% &  ()\\\hline
\begin{somme}
夜  & \textit{night}  & よる \mbox{(yoru)} \\\hline
夜中  & night + middle = \textit{middle of the night}  & よ{\middel}なか \mbox{(yo{\middel}naka)} \\\hline
夜明け  & night + bright = \textit{dawn}  & よ　あ{\middel}け \mbox{(yo a{\middel}ke)} \\\hline
月夜  & moon + night = \textit{moonlit night}  & つき{\middel}よ \mbox{(tsuki{\middel}yo)} \\\hline
夜間  & night + interval = \textit{night time}  & ヤ{\middel}カン \mbox{(YA{\middel}KAN)} \\\hline
\end{somme}

\begin{decomp}{5}
夜明け&に&家&を&出ました。\\\hline
よあけ&に&いえ&を&でました\\\hline
dawn&に&house&を&exit\\\hline
\multicolumn{5}{|r|}{\textit{I left the house at dawn.}}\\\hline
\end{decomp}

\end{multicols}

\pagebreak
%%--------------------------------------------------------------------
%13-------------------------------------------------------------------

\begin{multicols}{2}
\begin{glyph}{Before (前)}{before}\label{glyph:before}
Before.  Front.\\\hline
It is worth comparing the compounds below to those for 先 (Entry \ref{glyph:precede}, page \pageref{glyph:precede}) to get a sense of the subtle differences between these two characters.
\end{glyph}
\gindex{Before}{前}\strindex{9}{前}

\begin{comps}
\comprtwocone & 艸艹 + 月 + 刀 \\\hline
\comprthreecone & Grass + Moon + Sword\\\hline
\end{comps}

\begin{remglyph}
\hooglig{Before} the \remcom{moon} appears I cut the \remcom{grass} with my \remcom{sword}.\\\hline
\end{remglyph}

\begin{prons}
\pronsrtwocone & \textbf{ゼン (ZEN)} & \textbf{まえ (mae)}\\\hline
\pronsrthreecone & --- & --- \\\hline
\end{prons}

\begin{remprons}
\hooglig{Before} the \comkun{maelstrom} there was \comon{zen}.\\\hline
\end{remprons}

%{\middel}
% &  ()\\\hline
\begin{somme}
前  & \textit{before}  & まえ \mbox{(mae)} \\\hline
名前  & name + before = \textit{name}  & な{\middel}まえ \mbox{(na{\middel}mae)} \\\hline
手前  & hand + before = \textit{this side of}  & て{\middel}まえ \mbox{(te{\middel}mae)}  \\\hline
前半  & before + half = \textit{the first half}  & ゼン{\middel}ハン \mbox{(ZEN{\middel}HAN)} \\\hline
前者  & before + individual = \textit{the former}  & ゼン{\middel}シャ \mbox{(ZEN{\middel}SHA)} \\\hline
\notyet{午前}  & noon + before = \textit{A.M.; Morning}  & ゴゼン \mbox{(GO{\middel}ZEN)} \\\hline
\end{somme}

\end{multicols}
\begin{decomp}{8}
門&の&前&に&美しい&花&が&あります。\\\hline
モン&の&まえ&に&うつくしい&はな&が&あります\\\hline
gate&の&before&に&beautiful&flower&が&are\\\hline
\multicolumn{8}{|r|}{\textit{There are beautiful flowers in front of the gate.}}\\\hline
\end{decomp}

%%--------------------------------------------------------------------
%14-------------------------------------------------------------------
\begin{multicols}{2}
\begin{glyph}{Many (多)}{many}\label{glyph:many}
Many.  Much.
\end{glyph}
\gindex{Many}{多}\strindex{6}{多}

\begin{comps}
\comprtwocone & 夕 + 夕 \\\hline
\comprthreecone & Evening + Evening\\\hline
\end{comps}

\begin{remglyph}
\hooglig{Many} sins are done during \remcom{evenings}.\\\hline
\end{remglyph}

\begin{prons}
\pronsrtwocone & \textbf{タ (TA)} & \textbf{おお (\={o})}\\\hline
\pronsrthreecone & --- & --- \\\hline
\end{prons}

\begin{remprons}
\hooglig{Many} \comkun{awful} \comon{tub} experiences.\\\hline
\end{remprons}

%{\middel}
% &  ()\\\hline
\begin{somme}
多い & \textit{many}  & おお{\middel}い \mbox{(\={o}{\middel}i)} \\\hline
多少 & many + few = \textit{more or less}  & タ{\middel}ショウ \mbox{(TA{\middel}SH\={O})} \\\hline
多分 & many + part = \textit{perhaps}  & タ{\middel}ブン \mbox{(TA{\middel}BUN)} \\\hline
多大 & many + large = \textit{great quantity}  & タ{\middel}ダイ \mbox{(TA{\middel}DAI)} \\\hline
多読家 & many + read + house = \textit{well-read person}  & タ{\middel}ドク{\middel}カ \mbox{(TA{\middel}DOKU{\middel}KA)}  \\\hline
\end{somme}

\begin{decomp}{6}
彼&は&多分&有名&に&成らない。\\\hline
かれ&は&タブン&ユウメイ&に&ならない\\\hline
he&は&perhaps&famous&に&not become\\\hline
\multicolumn{6}{|r|}{\textit{Perhaps he'll never become famous.}}\\\hline
\end{decomp}

\end{multicols}

\pagebreak
%%--------------------------------------------------------------------
%15-------------------------------------------------------------------

\begin{multicols}{2}
\begin{glyph}{Hold (持)}{hold}\label{glyph:hold}
Hold.  Possess.
\end{glyph}
\gindex{Hold}{持}\strindex{9}{持}

\begin{comps}
\comprtwocone & 手扌 + 寺 \\\hline
\comprthreecone & Hand + Temple (\ref{glyph:temple}) \\\hline
\end{comps}

\begin{remglyph}
\hooglig{Hold} the \remcom{temple} in your \remcom{hand}.\\\hline
\end{remglyph}

\begin{prons}
\pronsrtwocone & --- & \textbf{も (mo)}\\\hline
\pronsrthreecone & ジ (JI) & --- \\\hline
\end{prons}

\begin{remprons}
\hooglig{Hold} on to your \comkun{morales}, or the \ucomon{jig} is up.\\\hline
\end{remprons}

%{\middel}
% &  ()\\\hline
\begin{somme}
\rye{3}{As the fourth example below shows, this 漢字 can act like 買.}
持つ  & \textit{to hold}  & もつ \mbox{(motsu)}  \\\hline
金持ち  & gold + hold = \textit{rich person}  & かね　も{\middel}ち \mbox{(kane mo{\middel}chi)} \\\hline
持ち上げる  & hold + upper = \textit{to lift up}  & も{\middel}ち　あ{\middel}げる \mbox{(mo{\middel}chi a{\middel}geru)} \\\hline
持物  & hold + thing = \textit{one's belongings}  & もち{\middel}もの \mbox{(mochi{\middel}mono)} \\\hline
\notyet{気持ち}  & spirit + hold = \textit{feeling}  & キ　も{\middel}ち \mbox{(KI mo{\middel}chi)} \\\hline
\end{somme}

\begin{decomp}{4}
気持ち&が&悪く&なった。\\\hline
キもち&が&わるく&なった\\\hline
feeling&が&bad&to \textit{become}\\\hline
\multicolumn{4}{|r|}{\textit{A strange feeling came over me.}}\\\hline
\end{decomp}

\end{multicols}

%%--------------------------------------------------------------------
%16-------------------------------------------------------------------
\begin{multicols}{2}
\begin{glyph}{Write (書)}{write}\label{glyph:write}
Write.\\\hline
This character has to do with all matters related to writing, and can suggest such things as books, notes, documents, and calligraphy, amongst others.
\end{glyph}
\gindex{Write}{書}\strindex{10}{書}

\begin{comps}
\comprtwocone & 聿 + 日 \\\hline
\comprthreecone & Brush + Sun \\\hline
\end{comps}

\begin{remglyph}
\hooglig{Write} with this \remcom{brush} the kanji for \remcom{``Sun''}.\\\hline
\end{remglyph}

\begin{prons}
\pronsrtwocone & \textbf{ショ (SHO)} & \textbf{か (ka)}\\\hline
\pronsrthreecone & --- & --- \\\hline
\end{prons}

\begin{remprons}
\hooglig{Write} down the \comkun{customers'} \comon{shopping} experiences.\\\hline
\end{remprons}

%{\middel}
% &  ()\\\hline
\begin{somme}
\rye{3}{The third and the fifth examples below show that this 漢字 can act like 買.  Be aware that the 訓読み is always voiced when it appears in the second position, as it is in the final compound.}
書く & \textit{to write}  & か{\middel}く \mbox{(ka{\middel}ku)} \\\hline
書き入れる & write + enter = \textit{to write in}  & か{\middel}き　い{\middel}れる \mbox{(ka{\middel}ki i{\middel}reru)} \\\hline
書入れ & write + enter = \textit{an entry}  & か{\middel}き　い{\middel}れ \mbox{(ka{\middel}ki i{\middel}re)} \\\hline
書道 & write + road = \textit{calligraphy}  & ショ{\middel}ドウ \mbox{(SHO{\middel}D\={O})} \\\hline
{\diff}前書 & before + write = \textit{preface}  & まえ{\middel}がき \mbox{(mae{\middel}gaki)}  \\\hline
\end{somme}

\begin{decomp}{4}
鉛筆&で&書いて&下さい。\\\hline
エンピツ&で&かいて&ください\\\hline
pencil&で&to write&please\\\hline
\multicolumn{4}{|r|}{\textit{Please write with a pencil.}}\\\hline
\end{decomp}

\begin{decomp}{3}
手紙&書いて&ね。\\\hline
てがみ&かいて&ね\\\hline
letter&to write&\textit{request confirmation}\\\hline
\multicolumn{3}{|r|}{\textit{Write me sometime, OK?}}\\\hline
\end{decomp}

\end{multicols}

\pagebreak
%%--------------------------------------------------------------------
%17-------------------------------------------------------------------

\begin{multicols}{2}
\begin{glyph}{Sound (音)}{sound}\label{glyph:sound}
Sound.
\end{glyph}
\gindex{Sound}{音}\strindex{9}{音}

\begin{comps}
\comprtwocone & 音 \\\hline
\comprthreecone & Sound \\\hline
\end{comps}

\begin{remglyph}
Part of the 漢字 radicals (\#{180}).\\\hline
\end{remglyph}

\begin{prons}
\pronsrtwocone & \textbf{オン (ON)} & \textbf{おと (oto)}\\\hline
\pronsrthreecone & イン (IN) & ね (ne) \\\hline
\end{prons}

\begin{remprons}
\hooglig{Sounded} like an \ucomon{inkling} of a \ucomkun{necromancer} \comkun{automatically} going \comon{online}.\\\hline
\end{remprons}

%{\middel}
% &  ()\\\hline
\begin{somme}
\rye{3}{The second example is very familiar to you by now.  Note also how the meaning of the same compound changes in the example following this, when both 漢字 are read with their 音読み.}
音  & \textit{sound}  & おと \mbox{(oto)} \\\hline
音読み  & sound + read = \textit{the on-yomi}  & オン　よ{\middel}み \mbox{(ON yo{\middel}mi)} \\\hline
音読  & sound + read = \textit{reading aloud}  & オン{\middel}ドク \mbox{(ON{\middel}DOKU)} \\\hline
物音  & thing + sound = \textit{a noise}  & もの{\middel}おと \mbox{(mono{\middel}oto)} \\\hline
\notyet{足音}  & leg + sound = \textit{sound of footsteps}  & あし{\middel}おと \mbox{(ashi{\middel}oto)}  \\\hline
\end{somme}

\end{multicols}
\begin{decomp}{8}
あれ&は&水&の&音&だ&と&思います。\\\hline
あれ&は&みず&の&おと&だ&と&おもいます\\\hline
that&は&water&の&sound&だ&と&to think\\\hline
\multicolumn{8}{|r|}{\textit{I think that's the sound of water.}}\\\hline
\end{decomp}

%%--------------------------------------------------------------------
%18-------------------------------------------------------------------
\begin{multicols}{2}
\begin{glyph}{Meet (会)}{meet}\label{glyph:meet}
Meet.  Meeting.\\\hline
When found at the end of compounds, this 漢字 often expresses English words such as \textit{society}, \textit{association}, and \textit{party}.
\end{glyph}
\gindex{Meet}{会}\strindex{6}{会}

\begin{comps}
\comprtwocone & 人 + 二 + 厶 \\\hline
\comprthreecone & Person + Two + Private \\\hline
\end{comps}

\begin{remglyph}
\hooglig{Meet} \remcom{two} \remcom{people} \remcom{privately}.\\\hline
\end{remglyph}

\begin{prons}
\pronsrtwocone & \textbf{カイ (KAI)} & \textbf{あ (a)}\\\hline
\pronsrthreecone &エ (E)  &  \\\hline
\end{prons}

\begin{remprons}
\hooglig{Meet} at the \comkun{updated} \comon{kites} anchored by \ucomon{eggs}.\\\hline
\end{remprons}

%{\middel}
% &  ()\\\hline
\begin{somme}
会う  & \textit{to meet}  & あ{\middel}う \mbox{(a{\middel}u)} \\\hline
会見  & meet + see = \textit{interview}  & カイ{\middel}ケン \mbox{(KAI{\middel}KEN)} \\\hline
{\diff}大会  & large + meet = \textit{convention; tournament}  & タイ{\middel}カイ \mbox{(TAI{\middel}KAI)} \\\hline
開会  & open + meet = \textit{to open a meeting}  & カイ{\middel}カイ \mbox{(KAI{\middel}KAI)} \\\hline
入会金  & enter + meet + gold = \textit{enrollment fee}  & ニュウ{\middel}カイ{\middel}キン \mbox{(NY\={U}{\middel}KAI{\middel}KIN)}  \\\hline
\end{somme}

\begin{decomp}{6}
米国大会&は&九月&に&あります&か。\\\hline
ベイコクタイカイ&は&クガツ&に&あります&か\\\hline
U.S. tournament&は&September&に&is&か\\\hline
\multicolumn{6}{|r|}{\textit{Is the U.S. tournament in September?}}\\\hline
\end{decomp}

\end{multicols}

\pagebreak
%%--------------------------------------------------------------------
%19-------------------------------------------------------------------

\begin{multicols}{2}
\begin{glyph}{Hear (聞)}{hear}\label{glyph:hear}
Hear.\\\hline
Ask.\\\hline
Like 早 (Entry \ref{glyph:early}, page \pageref{glyph:early}), this is another 漢字 with two unrelated and distinct meanings.
\end{glyph}
\gindex{Hear}{聞}\strindex{14}{聞}

\begin{comps}
\comprtwocone & 門 + 耳\\\hline
\comprthreecone & Gate + Ear\\\hline
\end{comps}

\begin{remglyph}
\hooglig{Hearing} my \hooglig{question} starts \remcom{gate} of the \remcom{ear}.\\\hline
\end{remglyph}

\begin{prons}
\pronsrtwocone & \textbf{ブン (BUN)} & \textbf{き (ki)}\\\hline
\pronsrthreecone &モン (MON)  & --- \\\hline
\end{prons}

\begin{remprons}
\hooglig{Hear} how the \ucomon{monastery} \comkun{keeps} \comon{boons}.\\\hline
\end{remprons}

%{\middel}
% &  ()\\\hline
\begin{somme}
聞こえる (intr.)  & \textit{to be heard}  & き{\middel}こえる \mbox{(ki{\middel}koeru)}  \\\hline
聞く (tr.)  & \textit{to hear; to ask}  & き{\middel}く \mbox{(ki{\middel}ku)} \\\hline
聞き出す  & hear + exit = \textit{to find out about}  & き{\middel}き　だ{\middel}す \mbox{(ki{\middel}ki da{\middel}su)} \\\hline
聞き入れる  & hear + enter = \textit{to comply with}  & き{\middel}き　い{\middel}れる \mbox{(ki{\middel}ki i{\middel}reru)}  \\\hline
新聞  & new + hear = \textit{newspaper}  & シン{\middel}ブン \mbox{(SHIN{\middel}BUN)} \\\hline
見聞  & see + hear = \textit{information}  & ケン{\middel}ブン \mbox{(KEN{\middel}BUN)}  \\\hline
\end{somme}

\begin{decomp}{5}
夜中&に&物音&を&聞きました。\\\hline
よなか&に&ものおと&を&ききました\\\hline
middle of night&に&sound&を&heard\\\hline
\multicolumn{5}{|r|}{\textit{(She) heard a sound in the middle of the night.}}\\\hline
\end{decomp}

\end{multicols}

%%--------------------------------------------------------------------
%20-------------------------------------------------------------------
\begin{multicols}{2}
\begin{glyph}{Center (央)}{center}\label{glyph:center}
Center.\\\hline
Other than the first example (and examples derived from it), this 漢字 is rarely used.
\end{glyph}
\gindex{Center}{央}\strindex{5}{央}

\begin{comps}
\comprtwocone & 冖 + 大 \\\hline
\comprthreecone & Cover/Crown + Large\\\hline
\end{comps}

\begin{remglyph}
\hooglig{Center} the \remcom{large} \remcom{crown}.\\\hline
\end{remglyph}

\begin{prons}
\pronsrtwocone & \textbf{オウ (\={O})} & --- \\\hline
\pronsrthreecone & --- & --- \\\hline
\end{prons}

\begin{remprons}
At the \hooglig{center} of my being is an \comon{awesome} God.\\\hline
\end{remprons}

%{\middel}
% &  ()\\\hline
\begin{somme}
中央  & middle + center = \textit{center}  & チュウ{\middel}オウ \mbox{(CH\={U}{\middel}\={O})}  \\\hline
中央口  & middle + center + mouth = \textit{central exit/entrance} &  チュウ{\middel}オウ{\middel}ぐち \mbox{(CH\={U}{\middel}\={O}{\middel}guchi)}  \\\hline
\end{somme}
\end{multicols}

\begin{decomp}{5}
中央口&で&秋山さん&を&見ました。\\\hline
チュウオウぐち&で&あきやまさん&を&みました\\\hline
central entrance&で&Akiyama-san&を&saw\\\hline
\multicolumn{5}{|r|}{\textit{(He) saw Akiyama-san at the central entrance.}}\\\hline
\end{decomp}
%%--------------------------------------------------------------------
%%--------------------------------------------------------------------

\pagebreak
%
%\begin{multicols}{2}
%\chapter{Kanji 141--160}
%%%--------------------------------------------------------------------
%1--------------------------------------------------------------------
\vspace*{\fill}
\begin{glyph}{Near (近)}{near}\label{glyph:near}
Near.\\\hline
    The sense of \textit{``recent''} is also incorporated.
\end{glyph}

\begin{comps}
\comprtwocone & 斤 + 辵辶⻌⻍ \\\hline
    \comprthreecone & Axe + Walk/Road  \\\hline
\end{comps}

\begin{remglyph}
    \hooglig{Nearly} got \remcom{axed} while \remcom{walking}.\\\hline
\end{remglyph}

\begin{prons}
\pronsrtwocone & \textbf{キン (KIN)} & \textbf{ちか (chika)}\\\hline
\pronsrthreecone & --- & --- \\\hline
\end{prons}

\begin{remprons}
    \hooglig{Near} \comon{kindred} to us are the \comkun{chickadees}.\\\hline
\end{remprons}

%{\middel}
% &  ()\\\hline
\begin{somme}
    近い  &  \textit{near}   &  ちか{\middel}い \mbox{(chika{\middel}i)}\\\hline
    近づく  &  \textit{to approach} &  ちか{\middel}ずく \mbox{(chika{\middel}zuku)}\\\hline
    近道  &  near + road = \textit{shortcut} &  ちか{\middel}みち \mbox{(chika{\middel}michi)} \\\hline
    近親  &  near + parent = \textit{close relative} &  キン{\middel}シン \mbox{(KIN{\middel}SHIN)} \\\hline
    \notyet{最近}  &  most + near = \textit{recently} &  サイ{\middel}キン \mbox{(SAI{\middel}KIN)} \\\hline
\end{somme}

\begin{sentence}
    \underline{\ruby{新}{あたら}しい}{\;}\underline{\ruby{家}{いえ}}\ruby{は}{わ}\underline{\ruby{川}{かわ}}の\underline{\ruby{近}{ちか}く}に\underline{あります}。\\\hline
atara{\middel}shii ie wa kawa no chika{\middel}ku ni arimasu.\\\hline
    (new) (house) (river) (near) (is)\\\hline
    =\textit{The new house is near a river.}\\\hline
\end{sentence}
\end{multicols}
\vhrulefill{5pt}
%%--------------------------------------------------------------------
%2--------------------------------------------------------------------
\begin{multicols}{2}
\begin{glyph}{Lose (失)}{lose}\label{glyph:lose}
Lose, in the sense of ``mislay''.\\\hline
This character often conveys a hint of ``failure and error''.
\end{glyph}

\begin{comps}
\comprtwocone & 丿 + 二 + 人  \\\hline
    \comprthreecone & Slash + Two + Person\\\hline
\end{comps}

\begin{remglyph}
    \hooglig{Lost} \remcom{two} \remcom{persons} \remcom{slashing} each other.\\\hline
\end{remglyph}

\begin{prons}
\pronsrtwocone & \textbf{シツ (SHITSU)} & \textbf{うしな (ushina)}\\\hline
\pronsrthreecone & --- & --- \\\hline
\end{prons}

\begin{remprons}
    \hooglig{Lose} your \comon{sheets} until \comkun{you shine}.\\\hline
\end{remprons}

%{\middel}
% &  ()\\\hline
\begin{somme}
    失う  & \textit{to lose}  &  うしな{\middel}う \mbox{(ushina{\middel}u)}\\\hline
    見失う  &  see + lose = \textit{to lose sight of} & み　うしな{\middel}う \mbox{(mi ushina{\middel}u)} \\\hline
    失言  &  lose + say = \textit{slip of the tongue} & シス{\middel}ゲン \mbox{(SHITSU{\middel}GEN)} \\\hline
\end{somme}

\begin{sentence}
    \underline{\ruby{大会}{タイカイ}}で\underline{\ruby{失言}{シツゲン}}しました。\\\hline
TAI{\middel}KAI de SHITSU{\middel}GEN shimashita.\\\hline
    (convention) (slip of the tongue) (did)\\\hline
    =\textit{(I) made a slip of the tongue at the convention.} \\\hline
\end{sentence}
\end{multicols}
\vspace*{\fill}
\pagebreak
%%--------------------------------------------------------------------
%3--------------------------------------------------------------------
\vspace*{\fill}
\begin{multicols}{2}
\begin{glyph}{Bow (弓)}{bow}\label{glyph:bow}
    Bow (the partner to an arrow).
\end{glyph}

\begin{comps}
\comprtwocone & 弓 \\\hline
    \comprthreecone & Bow \\\hline
\end{comps}

\begin{remglyph}
Part of the Kanji radicals (\#{57}).\\\hline
\end{remglyph}

\begin{prons}
    \pronsrtwocone & \textbf{キュウ (KY\={U})} & \textbf{ゆみ (yumi)}\\\hline
\pronsrthreecone & --- & --- \\\hline
\end{prons}

\begin{remprons}
    This \hooglig{bow} will \comon{kill you} \comkun{yummily}.\\\hline
\end{remprons}

%{\middel}
% &  ()\\\hline
\begin{somme}
    弓  & \textit{bow}  &  ゆみ \mbox{(yumi)} \\\hline
    弓道  & bow + road = \textit{(Japanese) archery}  & キュウ{\middel}ドウ \mbox{(KY\={U}{\middel}D\={O})}  \\\hline
\end{somme}

\begin{sentence}
    \underline{\ruby{前田}{まえだ}さん}\ruby{は}{わ}\underline{\ruby{古}{ふる}い}{\;}\underline{\ruby{弓}{ゆみ}}と\underline{\ruby{美}{うつく}しい}\underline{\ruby{日本刀}{ニホントウ}}\ruby{を}{お}\underline{\ruby{買}{か}いました}。\\\hline
    Mae{\middel}da-san wa furu{\middel}i yumi to utsuku{\middel}shii NI{\middel}HON{\middel}T\={O} o ka{\middel}imashita.\\\hline
    (Maeda-san) (old) (bow) (beautiful) (Japanese sword) (bought)\\\hline
    =\textit{Maeda-san bough an old bow and a beautiful Japanese sword.}\\\hline
\end{sentence}
\end{multicols}
\vhrulefill{5pt}
%%--------------------------------------------------------------------
%4--------------------------------------------------------------------
\begin{multicols}{2}
\begin{glyph}{Horse (馬)}{horse}\label{glyph:horse}
Horse.
\end{glyph}

\begin{comps}
\comprtwocone & 馬 \\\hline
    \comprthreecone & Horse \\\hline
\end{comps}

\begin{remglyph}
Part of the Kanji radicals (\#{187}).\\\hline
\end{remglyph}

\begin{prons}
\pronsrtwocone & \textbf{バ (BA)} & \textbf{うま (uma)}\\\hline
\pronsrthreecone & --- & --- \\\hline
\end{prons}

\begin{remprons}
    \hooglig{Horse}-\ucomon{bum} \ucomkun{mud} is \comkun{umami}.\\\hline
\end{remprons}

%{\middel}
% &  ()\\\hline
\begin{somme}
    馬  & \textit{horse}  &  うま \mbox{(uma)} \\\hline
    子馬  & child + horse = \textit{colt; pony}  & こ{\middel}うま \mbox{(ko{\middel}uma)}  \\\hline
    牛馬  & cow + horse = \textit{cows and horses}  & ギュウ{\middel}バ \mbox{(GY\={U}{\middel}BA)}  \\\hline
    \notyet{馬具}  & horse + tool = \textit{harness}  &  バ{\middel}グ \mbox{(BA{\middel}GU)} \\\hline
\end{somme}

\begin{sentence}
    \underline{\ruby{馬}{うま}}と\underline{\ruby{牛}{うし}}の\underline{どちら}が\underline{\ruby{好}{す}き}ですか。\\\hline
uma to ushi no dochira ga su{\middel}ki desu ka.\\\hline
    (horse) (cow) (which) (like)\\\hline
    =\textit{Which do you like, horses or cows?}\\\hline
\end{sentence}
\end{multicols}
\vspace*{\fill}
\pagebreak
%%--------------------------------------------------------------------
%5--------------------------------------------------------------------
\vspace*{\fill}
\begin{multicols}{2}
\begin{glyph}{Week (週)}{week}\label{glyph:week}
Week.
\end{glyph}

\begin{comps}
\comprtwocone & 周 + 辵辶⻌ \\\hline
    \comprthreecone & Around (\ref{glyph:around}) + Walk/Road \\\hline
\end{comps}

\begin{remglyph}
    \hooglig{Weeks} of \remcom{walking} \remcom{around}.\\\hline
\end{remglyph}

\begin{prons}
    \pronsrtwocone & \textbf{シュウ (SH\={U})} & ---\\\hline
\pronsrthreecone & --- & --- \\\hline
\end{prons}

\begin{remprons}
    \hooglig{Weekly} \comon{shoes}.\\\hline
\end{remprons}

%{\middel}
% &  ()\\\hline
\begin{somme}
    週間  & week + interval = \textit{week}  & シュウ{\middel}カン \mbox{(SH\={U}{\middel}KAN)}  \\\hline
    一週間  & one + week + interval = \textit{one week}  & イッ{\middel}シュウ{\middel}カン \mbox{(IS{\middel}SH\={U}{\middel}KAN)}  \\\hline
    二週間  & two + week + interval = \textit{two weeks}  & ニ{\middel}シュウ{\middel}カン \mbox{(NI{\middel}SH\={U}{\middel}KAN)}  \\\hline
    週日  & week + sun = \textit{weekday}  & シュウ{\middel}ジツ \mbox{(SH\={U}{\middel}JITSU)}  \\\hline
    先週  & precede + week = \textit{last week}  & セン{\middel}シュウ \mbox{(SEN{\middel}SH\={U})}  \\\hline
    \notyet{来週}  & come + week = \textit{next week}  & ライ{\middel}シュウ \mbox{(RAI{\middel}SH\={U})}  \\\hline
\end{somme}

\begin{sentence}
    \underline{\ruby{月曜日}{ゲツヨウび}}{\;}\underline{から}{\;}\underline{\ruby{二週間}{ニシュウカン}}の\underline{\ruby{夏休}{なつやす}み}です。\\\hline
    GETSU{\middel}Y\={O}{\middel}bi kara NI{\middel}SH\={U}{\middel}KAN no natsu yasu{\middel}mi desu.\\\hline
    (Monday) (from) (two weeks) (summer vacation)\\\hline
    =\textit{There's a two-week summer vacation from Monday.}\\\hline
\end{sentence}
\end{multicols}
\vhrulefill{5pt}
%%--------------------------------------------------------------------
%6--------------------------------------------------------------------
\begin{multicols}{2}
\begin{glyph}{Pull (引)}{pull}\label{glyph:pull}
Pull.\\\hline
    This character also conveys the idea of ``attraction'' (best seen in the
    fourth example), and the sense of ``receding'' and ``withdrawing''.\\\hline
    You will notice it frequently on signs; it indicates ``pull'' on doors, and
    ``discount'' when used with numbers.
\end{glyph}

\begin{comps}
\comprtwocone & 弓 + 丨 \\\hline
    \comprthreecone & Bow + Line/Pole/Vertical stroke \\\hline
\end{comps}

\begin{remglyph}
    \textit{Somewhat obvious from the glyph itself.}\\\hline
\end{remglyph}

\begin{prons}
\pronsrtwocone & \textbf{イン (IN)} & \textbf{ひ (hi)}\\\hline
\pronsrthreecone & --- &  ---\\\hline
\end{prons}

\begin{remprons}
    \hooglig{Pull} \comkun{heaps} of \comon{ink}.\\\hline
\end{remprons}

%{\middel}
% &  ()\\\hline
\begin{somme}
    \rye{3}{As seen in the third example, this is another kanji that can act
    like 買.  The 訓読み can also become voiced in the second position, as
    happens in the fifth compound.}
    引く  & \textit{to pull}  & ひ{\middel}く \mbox{(hi{\middel}ku)}  \\\hline
    引き出す  & pull + exit = \textit{to withdraw (something)}  & ひ{\middel}き　だ{\middel}す \mbox{(hi{\middel}ki da{\middel}su)}  \\\hline
    引出し  & pull + exit = \textit{a withdrawal; a drawer}  & ひ{\middel}き　だ{\middel}し \mbox{(hi{\middel}ki da{\middel}shi)}  \\\hline
    引力  & pull + strength = \textit{gravitation}  & イン{\middel}リョク \mbox{(IN{\middel}RYOKU)}  \\\hline
    万引き  & ten thousand + pull = \textit{shoplifter}  & マン　び{\middel}き \mbox{(MAN bi{\middel}ki)}  \\\hline
    \notyet{引き受ける}  & pull + receive = \textit{to undertake}  & ひ{\middel}き　う{\middel}ける \mbox{(hi{\middel}ki u{\middel}keru)}  \\\hline
    \notyet{引き上げる}  & pull + upper = \textit{to pull up}  & ひ{\middel}き　あ{\middel}げる \mbox{(hi{\middel}ki a{\middel}geru)}  \\\hline
\end{somme}

\begin{sentence}
    \underline{やっと}{\;}\underline{\ruby{牛}{うし}}\ruby{を}{お}\underline{\ruby{川}{かわ}}{\;}\underline{から}{\;}\underline{\ruby{引}{ひ}き\ruby{上}{あ}げました}。\\\hline
yatto ushi o kawa kara hi{\middel}ki a{\middel}gemashita.\\\hline
    (finally) (cow) (river) (from) (pulled up)\\\hline
    =\textit{(We) finally pulled the cow from the river.}\\\hline
\end{sentence}
\end{multicols}
\vspace*{\fill}
\pagebreak
%%--------------------------------------------------------------------
%7--------------------------------------------------------------------
\vspace*{\fill}
\begin{multicols}{2}
\begin{glyph}{Second (秒)}{second}\label{glyph:second}
    Second (of time or arc).
\end{glyph}

\begin{comps}
\comprtwocone & 禾 + 少 \\\hline
    \comprthreecone & Two-branch tree + Few (\ref{glyph:few}) \\\hline
\end{comps}

\begin{remglyph}
    \hooglig{Seconds} until a \remcom{few} \remcom{crops} remained.\\\hline
\end{remglyph}

\begin{prons}
    \pronsrtwocone & \textbf{ビョウ (BY\={O})} &---\\\hline
\pronsrthreecone & --- &  ---\\\hline
\end{prons}

\begin{remprons}
    \hooglig{Seconds} untill it will \comon{be yours}.\\\hline
\end{remprons}

%{\middel}
% &  ()\\\hline
\begin{somme}
    一秒  & one + second = \textit{one second}  & イチ{\middel}ビョウ \mbox{(ICHI{\middel}BY\={O})}  \\\hline
    五秒  & five + second = \textit{five seconds}  & ゴ{\middel}ビョウ \mbox{(GO{\middel}BY\={O})}  \\\hline
    秒読み  & second + read = \textit{countdown}  & ビョウ　よ{\middel}み \mbox{(BY\={O} yo{\middel}mi)}  \\\hline
\end{somme}

\begin{sentence}
Note that the sample sentence contains one of the less common
    on-yomi for ``分'' (Entry 124) in the compound ``十五分''. \\\hline
    \underline{\ruby{十五分}{ジュウゴフン}}{\;}\underline{\ruby{二十秒}{ニジュウビョウ}}{\;}\underline{かかりました}。\\\hline
    J\={U}{\middel}GO{\middel}FUN NI{\middel}J\={U}{\middel}BY\={O} kakarimashita.\\\hline
    (fifteen minutes) (twenty seconds) (took)\\\hline
    =\textit{It took fifteen minutes and twenty seconds.}\\\hline
\end{sentence}
\end{multicols}
\vhrulefill{5pt}
%%--------------------------------------------------------------------
%8--------------------------------------------------------------------
\begin{multicols}{2}
\begin{glyph}{Shell (甲)}{shell}\label{glyph:shell}
Shell.\\\hline
This character rarely appears on its own, but proves valuable as a component.
\end{glyph}

\begin{comps}
\comprtwocone & 田 \\\hline
    \comprthreecone & Rice field (\ref{glyph:rice_field}) \\\hline
\end{comps}

\begin{remglyph}
---\\\hline
\end{remglyph}

\begin{prons}
    \pronsrtwocone & \textbf{コウ (K\={O})} & ---\\\hline
    \pronsrthreecone & カン (KAN) & --- \\\hline
\end{prons}

\begin{remprons}
    \hooglig{Shells} \comon{cause} \ucomon{cunning} schemes.\\\hline
\end{remprons}

%{\middel}
% &  ()\\\hline
\begin{somme}
    手の甲  & hand + shell = \textit{back of the hand}  & て{\middel}の{\middel}コウ \mbox{(te{\middel}no{\middel}K\={O})}  \\\hline
\end{somme}

\begin{sentence}
    \underline{\ruby{上田}{うえだ}さん}は\underline{\ruby{手}{て}の\ruby{甲}{コウ}}が\underline{かゆい}。\\\hline
    Ue{\middel}da-san wa te{\middel}no{\middel}K\={O} ga kayui.\\\hline
    (Ueda-san) (back of the hand) (itchy)\\\hline
    =\textit{The back of Ueda-san's hand is itchy.}\\\hline
\end{sentence}
\end{multicols}
\vspace*{\fill}
\pagebreak
%%--------------------------------------------------------------------
%9--------------------------------------------------------------------
\vspace*{\fill}
\begin{multicols}{2}
\begin{glyph}{Spirit (気)}{spirit}\label{glyph:spirit}
``Spirit'', ``air'', ``feeling'', and ``mood''.\\\hline
This kanji is found in many common words and expressions.
\end{glyph}

\begin{comps}
\comprtwocone & 气 + メ \\\hline
\comprthreecone & Steam/Air + ME (katakana) \\\hline
\end{comps}

\begin{remglyph}
\hooglig{Spiritual} \remcom{men} in the \remcom{steam} room.\\\hline
\end{remglyph}

\begin{prons}
\pronsrtwocone & \textbf{キ (KI)} & ---\\\hline
    \pronsrthreecone &  ケ (KE) &  --- \\\hline
\end{prons}

\begin{remprons}
    A \hooglig{spirit}-\comon{killing} \ucomon{keg}.\\\hline
\end{remprons}

%{\middel}
% &  ()\\\hline
\begin{somme}
    気  & \textit{spirit}  & キ \mbox{(KI)}  \\\hline
    元気  & basis + spirit = \textit{high spirits}  & ゲン{\middel}キ \mbox{(GEN{\middel}KI)}  \\\hline
    人気  & person + spirit = \textit{popularity}  & ニン{\middel}キ \mbox{(NIN{\middel}KI)}  \\\hline
    電気  & electric + spirit = \textit{electricity}  & デン{\middel}キ \mbox{(DEN{\middel}KI)}  \\\hline
    天気  & heaven + spirit = \textit{the weather}  & テン{\middel}キ \mbox{(TEN{\middel}KI)}  \\\hline
    気持ち  & spirit + hold = \textit{feeling}  & キ　も{\middel}ち \mbox{(KI mo{\middel}chi)}  \\\hline
\end{somme}

\begin{sentence}
The sample sentence is an extremely common greeting/question in
    Japanese.\\\hline
    \underline{\ruby{元気}{ゲンキ}}ですか。\\\hline
GEN{\middel}KI desu ka.\\\hline
    (high spirits)\\\hline
    =\textit{How are you?}\\\hline
\end{sentence}
\end{multicols}
\vhrulefill{5pt}
%%--------------------------------------------------------------------
%10-------------------------------------------------------------------
\begin{multicols}{2}
\begin{glyph}{Time (時)}{time}\label{glyph:time}
Time.\\\hline
This character will become familiar to you through its use in the hours of the
    day (the third and fourth compounds below are examples).
\end{glyph}

\begin{comps}
\comprtwocone & 日 + 寺  \\\hline
    \comprthreecone & Sun + Temple (\ref{glyph:temple}) \\\hline
\end{comps}

\begin{remglyph}
    \hooglig{Time} is calibrated in \remcom{temples} by means of the
    \remcom{sun}.\\\hline
\end{remglyph}

\begin{prons}
\pronsrtwocone & \textbf{ジ (JI)} & \textbf{とき (toki)}\\\hline
\pronsrthreecone & --- & --- \\\hline
\end{prons}

\begin{remprons}
    \hooglig{Time} the \comkun{torque} to \comon{jiggle}.\\\hline
\end{remprons}

%{\middel}
% &  ()\\\hline
\begin{somme}
    \rye{3}{Note how the final half of the second compound becomes voiced; this
    is an occasional feature of repeated compounds.}
    時  & \textit{time}  & とき \mbox{(toki)} \\\hline
    時々  & time + time = \textit{sometimes}  & とき{\middel}どき \mbox{(toki{\middel}doki)}  \\\hline
    四時  & four + time = \textit{four o' clock}  & よ{\middel}ジ \mbox{(yo{\middel}JI)} \\\hline
    九時半  & nine + time + half = \textit{nine-thirty}  & ク{\middel}ジ{\middel}ハン \mbox{(KU{\middel}JI{\middel}HAN)} \\\hline
    \notyet{同時}  & same + time = \textit{simultaneous}  &  ドウ{\middel}ジ \mbox{(D\={O}{\middel}JI)} \\\hline
\end{somme}

\begin{irrr}
    \notyet{時計}  & time + measure = \textit{watch; clock}  & と{\middel}ケイ \mbox{(to{\middel}KEI)}  \\\hline
\end{irrr}

\begin{sentence}
    \underline{\ruby{八時半}{ハチジハン}}に\underline{\ruby{東口}{ひがしぐち}}の\underline{\ruby{近}{ちか}く}で\underline{\ruby{会}{あ}いましょう}。\\\hline
    HACHI{\middel}JI{\middel}HAN ni higashi{\middel}guchi no chika{\middel}ku de a{\middel}imash\={o}.\\\hline
    (eight-thirty) (east entrance) (near) (let's meet)\\\hline
    =\textit{Let's meet near the east entrance at eight-thirty.}\\\hline
\end{sentence}
\end{multicols}
\vspace*{\fill}
\pagebreak
%%--------------------------------------------------------------------
%11-------------------------------------------------------------------
\vspace*{\fill}
\begin{multicols}{2}
\begin{glyph}{Speak (話)}{speak}\label{glyph:speak}
Speak, alon with the ideas of ``conversations'' and ``stories''.
\end{glyph}

\begin{comps}
\comprtwocone & 言 + 舌 \\\hline
    \comprthreecone & Say + Tongue \\\hline
\end{comps}

\begin{remglyph}
    \hooglig{``Speak''}, \remcom{says} the \remcom{tongue}.\\\hline
\end{remglyph}

\begin{prons}
\pronsrtwocone & \textbf{ワ (WA)} & \textbf{はなし、はな (hanashi, hana)}\\\hline
    \rye{3}{As in the final example below, the kun-yomi for this kanji is always
	voiced in the second position.}
\pronsrthreecone & --- & --- \\\hline
\end{prons}

\begin{remprons}
    \hooglig{Speak} how \comkun{hung-up} \comkun{Hanah shielded} us from being
    \comon{washed}.\\\hline
\end{remprons}

%{\middel}
% &  ()\\\hline
\begin{somme}
    話  & \textit{story}  & はなし \mbox{(hanashi)}  \\\hline
    話す  & \textit{to speak}  & はな{\middel}す \mbox{(hana{\middel}su)}  \\\hline
    話好き  & speak + like = \textit{talkative}  & はな{\middel}し　ず{\middel}き \mbox{(hana{\middel}shi zu{\middel}ki)}  \\\hline
    電話  & electric + speak = \textit{telephone}  &  デン{\middel}ワ \mbox{(DEN{\middel}WA)} \\\hline
    会話  & meet + speak = \textit{conversation}  & カイ{\middel}ワ \mbox{(KAI{\middel}WA)}  \\\hline
    立ち話  & stand + speak = \textit{chatting while standing}  & た{\middel}ち　ばなし \mbox{(ta{\middel}chi banashi)}  \\\hline
\end{somme}

\begin{sentence}
    \underline{\ruby{田中}{たなか}さん}は\underline{フランス\ruby{語}{ゴ}}\ruby{を}{お}\underline{\ruby{上手}{ジョウず}}に\underline{\ruby{話}{はな}せます}。\\\hline
    Ta{\middel}naka-san wa Furanso{\middel}GO o J\={O}{\middel}zu ni hana{\middel}semasu.\\\hline
    (Tanaka-san) (French) (well) (can speak)\\\hline
    =\textit{Tanaka-san can speak French well.}\\\hline
\end{sentence}
\end{multicols}
\vhrulefill{5pt}
%%--------------------------------------------------------------------
%12-------------------------------------------------------------------
\begin{multicols}{2}
\begin{glyph}{English (英)}{english}\label{glyph:english}
English.\\\hline
A secondary meaning relates to ``brilliance''.
\end{glyph}

\begin{comps}
\comprtwocone & 艸 + 央 \\\hline
    \comprthreecone & Grass + Center (\ref{glyph:center}) \\\hline
\end{comps}

\begin{remglyph}
    \hooglig{English} \remcom{grass} is \remcom{central}.\\\hline
\end{remglyph}

\begin{prons}
\pronsrtwocone & \textbf{エイ (EI)} & ---\\\hline
\pronsrthreecone & --- & --- \\\hline
\end{prons}

\begin{remprons}
    \hooglig{English} \comon{aim} for world domination.\\\hline
\end{remprons}

%{\middel}
% &  ()\\\hline
\begin{somme}
    英語  & English + words = \textit{English language}  & エイ{\middel}ゴ \mbox{(EI{\middel}GO)}  \\\hline
    英国  & English + country = \textit{England}  & エイ{\middel}コク \mbox{(EI{\middel}KOKU)}  \\\hline
    英会話  & English + meet + speak = \textit{English conversation}  & エイ{\middel}カイ{\middel}ワ \mbox{(EI{\middel}KAI{\middel}WA)}  \\\hline
    日英  & sun + English = \textit{Japan and England}  & ニチ{\middel}エイ \mbox{(NICHI{\middel}EI)}  \\\hline
\end{somme}

\begin{sentence}
    \underline{\ruby{英会話}{エイカイワ}}は\underline{\ruby{日本}{ニホン}}で\underline{\ruby{人気}{ニンキ}}があります。\\\hline
EI{\middel}KAI{\middel}WA wa NI{\middel}HON de NIN{\middel}KI ga arimasu.\\\hline
    (English conversation) (Japan) (popular) (is)\\\hline
    =\textit{English conversation is popular in Japan.}\\\hline
\end{sentence}
\end{multicols}
\vspace*{\fill}
\pagebreak
%%--------------------------------------------------------------------
%13-------------------------------------------------------------------
\vspace*{\fill}
\begin{multicols}{2}
\begin{glyph}{Instruction (訓)}{instruction}\label{glyph:instruction}
Instruction.\\\hline
This kanji tends to give a more serious or spiritual sense of ``teaching'' to
    the compounds in which it appears.
\end{glyph}

\begin{comps}
\comprtwocone & 言 + 川 \\\hline
    \comprthreecone & Say/Speech + River\\\hline
\end{comps}

\begin{remglyph}
    \hooglig{Instruct} the \remcom{river} to \remcom{say} something.\\\hline
\end{remglyph}

\begin{prons}
\pronsrtwocone & \textbf{クン (KUN)} & ---\\\hline
\pronsrthreecone & --- & --- \\\hline
\end{prons}

\begin{remprons}
    \hooglig{Instruct} the \comon{coon}?!\\\hline
\end{remprons}

%{\middel}
% &  ()\\\hline
\begin{somme}
    訓読み  & instruction + read = \textit{the kun-yomi}  & クン　よ{\middel}み \mbox{(KUN yo{\middel}mi)}  \\\hline
    訓話  & instruction + speak = \textit{moral discourse}  & クン{\middel}ワ \mbox{(KUN{\middel}WA)}  \\\hline
\end{somme}

\begin{sentence}
    \underline{これ}\ruby{を}{お}\underline{\ruby{訓読}{クンよ}み}で\underline{\ruby{読}{よ}んで}{\;}\underline{\ruby{下}{くだ}さい}。\\\hline
kore o KUN{\middel}yomi de yo{\middel}nde kuda{\middel}sai.\\\hline
    (this) (kun-yomi) (read) (please)\\\hline
    =\textit{Please read this with kun-yomi.}\\\hline
\end{sentence}
\end{multicols}
\vhrulefill{5pt}
%%--------------------------------------------------------------------
%14-------------------------------------------------------------------
\begin{multicols}{2}
\begin{glyph}{Push (押)}{push}\label{glyph:push}
``Push'', together with the ideas of ``suppression'' and ``restraint''.\\\hline
    This kanji also serves as the partner to 引 (Entry \ref{glyph:pull}) on door signs.
\end{glyph}

\begin{comps}
\comprtwocone & 手 + 田 \\\hline
    \comprthreecone & Hand (\ref{glyph:hand}) + Rice field
    (\ref{glyph:rice_field}) \\\hline
\end{comps}

\begin{remglyph}
    \hooglig{Push} your \remcom{hand} into the \remcom{rice fields}.\\\hline
\end{remglyph}

\begin{prons}
\pronsrtwocone & --- & \textbf{お (o)}\\\hline
    \pronsrthreecone & オウ (OU) & --- \\\hline
\end{prons}

\begin{remprons}
    \hooglig{Pushing} can be \ucomon{awefully} \comkun{obligatory}.\\\hline
\end{remprons}

%{\middel}
% &  ()\\\hline
\begin{somme}
    \rye{3}{As indicated by the last two examples below, this is another kanji
    that often shows up in the manner of 買.}
    押す  & \textit{to push}  & お{\middel}す \mbox{(o{\middel}su)}  \\\hline
    押し出す  & push + exit = \textit{to push out}  & お{\middel}し　だ{\middel}す \mbox{(o{\middel}shi da{\middel}su)}  \\\hline
    押し上げる  & push + upper = \textit{to push up}  & お{\middel}し　あ{\middel}げる \mbox{(o{\middel}shi a{\middel}geru)}  \\\hline
    押し開ける  & push + open = \textit{to push open}  & お{\middel}し　あ{\middel}ける \mbox{(o{\middel}shi a{\middel}keru)}  \\\hline
    押入れ  & push + enter = \textit{closet}  & おし　い{\middel}れ \mbox{(oshi i{\middel}re)}  \\\hline
    押売  & push + sell = \textit{high-pressure sales}  & おし{\middel}うり \mbox{(oshi{\middel}uri)}  \\\hline
\end{somme}

\begin{sentence}
    \underline{\ruby{押入}{おしい}れ}の\underline{\ruby{中}{なか}}に\underline{\ruby{物}{もの}}が\underline{\ruby{少}{すく}ない}。\\\hline
oshi i{\middel}re no naka ni mono ga suku{\middel}nai.\\\hline
    (closet) (middle) (things) (few)\\\hline
    =\textit{There aren't many things in the closet.}\\\hline
\end{sentence}
\end{multicols}
\vspace*{\fill}
\pagebreak
%%--------------------------------------------------------------------
%15-------------------------------------------------------------------
\vspace*{\fill}
\begin{multicols}{2}
\begin{glyph}{Year (年)}{year}\label{glyph:year}
Year.
\end{glyph}

\begin{comps}
\comprtwocone & 干 + 丿 + 一 \\\hline
    \comprthreecone & Dry + Slash + One \\\hline
\end{comps}

\begin{remglyph}
    \hooglig{Year} \remcom{one} \remcom{dries} out in a \remcom{slash}.\\\hline
\end{remglyph}

\begin{prons}
\pronsrtwocone & \textbf{ネン (NEN)} & \textbf{とし (toshi)}\\\hline
\pronsrthreecone & --- & --- \\\hline
\end{prons}

\begin{remprons}
    The \hooglig{year} the \comon{Nenets} acquired \comkun{Toshiba}.\\\hline
\end{remprons}

%{\middel}
% &  ()\\\hline
\begin{somme}
    年  & \textit{year}  & とし \mbox{(toshi)}  \\\hline
    年上  & year + upper = \textit{older}  & とし{\middel}うえ \mbox{(toshi{\middel}ue)}  \\\hline
    年下  & year + lower = \textit{younger}  & とし{\middel}した \mbox{(toshi{\middel}shita)}  \\\hline
    新年  & new + year = \textit{the New year}  & シン{\middel}ネン \mbox{(SHIN{\middel}NEN)}  \\\hline
    中年  & middle + year = \textit{middle age}  & チュウ{\middel}ネン \mbox{(CH\={U}{\middel}NEN)}  \\\hline
    一年生  & one + year + life = \textit{first-year student}  & イチ{\middel}ネン{\middel}セイ \mbox{(ICHI{\middel}NEN{\middel}SEI)}  \\\hline
\end{somme}

\begin{sentence}
    \underline{\ruby{青木}{あおき}さん}は\underline{\ruby{三年生}{ワンネンセイ}}です。\\\hline
Ao{\middel}ki-san wa SAN{\middel}NEN{\middel}SEI desu.\\\hline
    (Aoki-san) (third-year student)\\\hline
    =\textit{Aoki-san is a third-year student.}\\\hline
\end{sentence}
\end{multicols}
\vhrulefill{5pt}
%%--------------------------------------------------------------------
%16-------------------------------------------------------------------
\begin{multicols}{2}
\begin{glyph}{Come (来)}{come}\label{glyph:come}
Come.\\\hline
As the final four compounds below make clear, the related meanings of ``next''
    is also included.
\end{glyph}

\begin{comps}
\comprtwocone & 一 + 米 \\\hline
    \comprthreecone & One + Rice \\\hline
\end{comps}

\begin{remglyph}
    \hooglig{Come} \remcom{one} \remcom{rice} grain at a time.\\\hline
\end{remglyph}

\begin{prons}
\pronsrtwocone & \textbf{アイ (RAI)} & \textbf{く (ku)}\\\hline
    \pronsrthreecone & --- & き (ki) \\\hline
    \rye{3}{This character is infamous for being one of the only two irregular
    verbs in Japanese.  The character can be read variously as く, き, or こ,
    depending on which tense of the verb is being used.  The sample sentence
    offers an example of the polite past tense.  Note, however, an important
    exception in the second compound below (an extremely common word); when
    following the kanji 出, 来 is always read き.}
\end{prons}

\begin{remprons}
    \hooglig{Come} up with a \comkun{coup} to \ucomkun{keep} on \comon{rhyming}.\\\hline
\end{remprons}

%{\middel}
% &  ()\\\hline
\begin{somme}
    来る  & \textit{to come}  & く{\middel}る \mbox{(ku{\middel}ru)}  \\\hline
    出来る  & exit + come = \textit{to be able to}  & で　き{\middel}る \mbox{(de ki{\middel}ru)}  \\\hline
    来週  & come + week = \textit{next week}  & ライ{\middel}シュウ \mbox{(RAI{\middel}SH\={U})}  \\\hline
    来月  & come + moon = \textit{next month}  & ライ{\middel}ゲツ \mbox{(RAI{\middel}GETSU)}  \\\hline
    来年  & come + year = \textit{next year}  & ライ{\middel}ネン \mbox{(RAI{\middel}NEN)}  \\\hline
    来春  & come + spring = \textit{next spring}  & ライ{\middel}シュン \mbox{(RAI{\middel}SHUN)}  \\\hline
\end{somme}

\begin{sentence}
    \underline{\ruby{先生}{センセイ}}は\underline{\ruby{一週間}{イッシュウカン}}{\;}\underline{\ruby{前}{まえ}}に\underline{\ruby{来}{き}ました}。\\\hline
    SEN{\middel}SEI wa IS{\middel}SH\={U}{\middel}KAN mae ni ki{\middel}mashita.\\\hline
    (teacher) (one week) (before) (came)\\\hline
    =\textit{The teacher came a week ago.}\\\hline
\end{sentence}
\end{multicols}
\vspace*{\fill}
\pagebreak
%%--------------------------------------------------------------------
%17-------------------------------------------------------------------
\vspace*{\fill}
\begin{multicols}{2}
\begin{glyph}{Study (学)}{study}\label{glyph:study}
Study and learning.\\\hline
When used as a suffix, it is often equivalent to the English ``-ology''.  The
    final compound offers such an example.
\end{glyph}

\begin{comps}
\comprtwocone & 子 + 小 + 冖 \\\hline
    \comprthreecone & Child + Small + Cover/Crown \\\hline
\end{comps}

\begin{remglyph}
\hooglig{Study} of \remcom{small} \remcom{children's} \remcom{crowns}.\\\hline
\end{remglyph}

\begin{prons}
\pronsrtwocone & \textbf{ガク (GAKU)} & \textbf{まな (mana)}\\\hline
\pronsrthreecone & --- & --- \\\hline
\end{prons}

\begin{remprons}
    \hooglig{Study} how to \comkun{manage} all this \comon{guck}.\\\hline
\end{remprons}

%{\middel}
% &  ()\\\hline
\begin{somme}
    学ぶ  & \textit{to study}  & まな{\middel}ぶ \mbox{(mana{\middel}bu)}  \\\hline
    大学  & large + study = \textit{university}  & ダイ{\middel}ガク \mbox{(DAI{\middel}GAKU)}  \\\hline
    学生  & study + life = \textit{student}  & ガク{\middel}セイ \mbox{(GAKU{\middel}SEI)}  \\\hline
    休学  & rest + study = \textit{absence from school}  & キュウ{\middel}ガク \mbox{(KY\={U}{\middel}GAKU)}  \\\hline
    語学  & words + study = \textit{linguistics}  & ゴ{\middel}ガク \mbox{(GO{\middel}GAKU)}  \\\hline
    生物学  & life + thing + study = \textit{biology}  & セイ{\middel}ブツ{\middel}ガク \mbox{(SEI{\middel}BUTSU{\middel}GAKU)}  \\\hline
\end{somme}

\begin{sentence}
    \underline{この}{\;}\underline{\ruby{家}{いえ}}に\ruby{は}{わ}\underline{\ruby{大学生}{ダイガクセイ}}が\underline{\ruby{多}{おお}い}。\\\hline
    kono ie ni wa DAI{\middel}GAKU{\middel}SEI ga \={o}{\middel}i.\\\hline
    (this) (house) (university students) (many)\\\hline
    =\textit{There are many university students in this house.}\\\hline
\end{sentence}
\end{multicols}
\vhrulefill{5pt}
%%--------------------------------------------------------------------
%18-------------------------------------------------------------------
\begin{multicols}{2}
\begin{glyph}{Mind (意)}{mind}\label{glyph:mind}
Mind/Thought.\\\hline
This abstract character appears in many compounds conveying the ideas of
    ``will'', ``intention'', and ``opinion'', etc.
\end{glyph}

\begin{comps}
\comprtwocone & 音 + 心 \\\hline
    \comprthreecone & Sound + Heart \\\hline
\end{comps}

\begin{remglyph}
    \hooglig{Mind} and \remcom{heart} don't always \remcom{resonate}.\\\hline
\end{remglyph}

\begin{prons}
\pronsrtwocone & \textbf{イ (I)} & ---\\\hline
\pronsrthreecone & --- & --- \\\hline
\end{prons}

\begin{remprons}
    \hooglig{Mind} the \comon{indicators}.\\\hline
\end{remprons}

%{\middel}
% &  ()\\\hline
\begin{somme}
    意思  & mind + think = \textit{(one's) intent}  & イ{\middel}シ \mbox{(I{\middel}SHI)}  \\\hline
    意図  & mind + diagram = \textit{(one's) aim}  & イ{\middel}ト \mbox{(I{\middel}TO)}  \\\hline
    意外  & mind + outside = \textit{unforseen}  & イ{\middel}ガイ \mbox{(I{\middel}GAI)}  \\\hline
    意見  & mind + see = \textit{opinion}  & イ{\middel}ケン \mbox{(I{\middel}KEN)}  \\\hline
    意気  & mind + spirit = \textit{morale}  & イ{\middel}キ \mbox{(I{\middel}KI)}  \\\hline
    \notyet{注意}  & pour + mind = \textit{attention}  & チュウ{\middel}イ \mbox{(CH\={U}{\middel}I)}  \\\hline
\end{somme}

\begin{sentence}
    \underline{あなた}の\underline{\ruby{意見}{イケン}}が\underline{\ruby{分}{わ}かりません}。\\\hline
anata no I{\middel}KEN ga wa{\middel}karimasen.\\\hline
    (you) (opinion) (don't understand)\\\hline
    =\textit{I don't understand your opinion.}\\\hline
\end{sentence}
\end{multicols}
\vspace*{\fill}
\pagebreak
%%--------------------------------------------------------------------
%19-------------------------------------------------------------------
\vspace*{\fill}
\begin{multicols}{2}
\begin{glyph}{Again (又)}{again}\label{glyph:again}
Agains/Also.  A character that rarely appears on its own.\\\hline
The suggested immediate visual connection is ``ironing board''.
\end{glyph}

\begin{comps}
\comprtwocone & 又 \\\hline
    \comprthreecone & Again \\\hline
\end{comps}

\begin{remglyph}
Part of the Kanji radicals (\#{29}).\\\hline
\end{remglyph}

\begin{prons}
\pronsrtwocone & --- & \textbf{また (mata)}\\\hline
\pronsrthreecone & --- & --- \\\hline
\end{prons}

\begin{remprons}
    \hooglig{Again} and again the \comkun{matador} irked the bull.\\\hline
\end{remprons}

%{\middel}
% &  ()\\\hline
\begin{somme}
    \rye{3}{Note how the second kanji in the final example below is voiced.}
    又  & \textit{again}  & また \mbox{(mata)}  \\\hline
    又々  & again + again = \textit{once again}  & また{\middel}また \mbox{(mata{\middel}mata)}  \\\hline
    又聞き  & again + hear = \textit{hearsay}  & また{\middel}ぎき \mbox{(mata{\middel}giki)}  \\\hline
\end{somme}

\begin{sentence}
    The pronunciation of 来なかった in the sample sentence, incidentially, is
    the informal past tense reading of 来.\\\hline
    \underline{\ruby{北川}{きたがわ}さん}は\underline{\ruby{又}{また}}{\;}\underline{\ruby{来}{こ}なかった}よ。\\\hline
Kita{\middel}gawa-san wa mata ko{\middel}nakatta yo.\\\hline
    (Kitagawa-san) (again) (didn't come)\\\hline
    \textit{Kitagawa-san didn't come again!}\\\hline
\end{sentence}
\end{multicols}
\vhrulefill{5pt}
%%--------------------------------------------------------------------
%20-------------------------------------------------------------------
\begin{multicols}{2}
\begin{glyph}{Red (赤)}{red}\label{glyph:red}
Red.
\end{glyph}

\begin{comps}
\comprtwocone & 赤 \\\hline
    \comprthreecone & Red \\\hline
\end{comps}

\begin{remglyph}
Part of the Kanji radicals (\#{155}).\\\hline
\end{remglyph}

\begin{prons}
\pronsrtwocone & \textbf{セキ (SEKI)} & \textbf{あか (aka)}\\\hline
    \pronsrthreecone & シャク (SHAKU) & --- \\\hline
\end{prons}

\begin{remprons}
    \hooglig{Red} is \comkun{also known as} \comon{sexy}.  \ucomon{Shucks}!\\\hline
\end{remprons}

%{\middel}
% &  ()\\\hline
\begin{somme}
    赤  & \textit{red (noun)}  & あか \mbox{(aka)}  \\\hline
    赤い  & \textit{red (adjective)}  & あか{\middel}い \mbox{(aka{\middel}i)}  \\\hline
    赤ちゃん  & \textit{baby}  & あか{\middel}ちゃん \mbox{(aka{\middel}chan)}  \\\hline
    赤道  & red + road = \textit{equator}  & セキ{\middel}ドウ \mbox{(SEKI{\middel}D\={O})}  \\\hline
\end{somme}

\begin{sentence}
    \underline{あの}{\;}\underline{\ruby{赤}{あか}い}{\;}\underline{\ruby{花}{はな}}\ruby{は}{わ}\underline{\ruby{美}{うつく}しい}{\;}\underline{ですね}。\\\hline
    ano aka{\middel}i hana wa utsuku{\middel}shii desu ne.\\\hline
    (that) (red) (flower) (beautiful) (isn't it)\\\hline
    =\textit{That red flower is beautiful, isn't it?}\\\hline
\end{sentence}
\end{multicols}
\vhrulefill{5pt}
%%--------------------------------------------------------------------
%%--------------------------------------------------------------------
\vspace*{\fill}
\pagebreak

%
%\begin{multicols}{2}
%\chapter{Kanji 161--180}
%%%--------------------------------------------------------------------
%1--------------------------------------------------------------------
\begin{glyph}{Public (公)}{public}\label{glyph:public}
Public.\\\hline
In many instances, this character serves as the opposite to 私 (p. \pageref{glyph:private}).
\end{glyph}
\gindex{Public}{公}\strindex{4}{公}

\begin{comps}
\comprtwocone&八 + ム\\\hline
\comprthreecone&Eight (Volcano) + Private\\\hline
\end{comps}

\begin{remglyph}
\hooglig{Public} \remcom{volcanoes} are being \remcom{privatized}.\\\hline
\end{remglyph}

\begin{prons}
\pronsrtwocone&\textbf{コウ (K\={O})}&---\\\hline
\pronsrthreecone&ク (KU)&おおやけ (\={o}yake)\\\hline
\end{prons}
\rindex{KOU@K\={O}}{公}
\rindex{KU}{公}
\rindex{ooyake@\={o}yake}{公}

\begin{remprons}
\hooglig{Publicly} \comon{cause} \ucomon{cooking} of \ucomkun{ore yuck}.\\\hline
\end{remprons}

%{\middel}
% &  ()\\\hline
\begin{somme}
公有&public + have = \textit{publicly owned}&コウ{\middel}ユウ \mbox{(K\={O}{\middel}Y\={U})}\\\hline
公安&public + ease = \textit{public peace}&コウ{\middel}アン \mbox{(K\={O}{\middel}AN)}\\\hline
公開&public + open = \textit{open to the public}&コウ{\middel}カイ \mbox{(K\={O}{\middel}KAI)}\\\hline
公金&public + gold = \textit{public funds}&コウ{\middel}キン \mbox{(K\={O}{\middel}KIN)}\\\hline
公道&public + road = \textit{public road}&コウ{\middel}ドウ \mbox{(K\={O}{\middel}D\={O})}\\\hline
公会&public + meet = \textit{public meeting}&コウ{\middel}カイ \mbox{(K\={O}{\middel}KAI)}\\\hline
公園&public + garden = \textit{public park}&コウ{\middel}エン \mbox{(K\={O}{\middel}EN)}\\\hline
\notyet{公平}&public + flat = \textit{fairness}&コウ{\middel}ヘイ \mbox{(K\={O}{\middel}HEI)}\\\hline
\end{somme}

\begin{decomp}{4}
ここ&は&公道&です。\\\hline
ここ&は&コウドウ&です\\\hline
this&は&public road&be/is\\\hline
\multicolumn{4}{|r|}{\textit{This is a public road.}}\\\hline
\end{decomp}
\end{multicols}
%%--------------------------------------------------------------------
%2--------------------------------------------------------------------
\begin{multicols}{2}
\begin{glyph}{Cloth (巾)}{cloth}\label{glyph:cloth}
Cloth.  Rag.\\\hline
This character frequently appears as a component in other 漢字.
\end{glyph}
\gindex{Cloth}{巾}\strindex{3}{巾}

\begin{comps}
\comprtwocone&巾\\\hline
\comprthreecone&Cloth\\\hline
\end{comps}

\begin{remglyph}
Part of the 漢字 radicals (\#{50}).\\\hline
\end{remglyph}

\begin{prons}
\pronsrtwocone&\textbf{キン (KIN)}&---\\\hline
\pronsrthreecone&はば (haba)&---\\\hline
\end{prons}
\rindex{KIN}{巾}
\rindex{haba}{巾}

\begin{remprons}
\hooglig{Clothe} the \comon{kindred} with \ucomkun{habaneros}.\\\hline
{\warn}\textit{Habaneros} is a hot variety of chilli peppers.\\\hline
\end{remprons}

%{\middel}
% &  ()\\\hline
\begin{somme}
巾&\textit{cloth, rag}&キン \mbox{(KIN)}\\\hline
手巾&hand + cloth = \textit{handkerchief}&シュ{\middel}キン \mbox{(SHU{\middel}KIN)}\\\hline
{\iscov}三角巾&triangle + cloth = \textit{sling}&サン{\middel}カク{\middel}キン \mbox{(SAN{\middel}KAKU{\middel}KIN)}\\\hline
\notyet{巾着}&cloth + wear = \textit{pouch}&キン{\middel}チャク \mbox{(KIN{\middel}CHAKU)}\\\hline
\notyet{雑巾}&miscellaneous + cloth = \textit{house-cloth}&ゾウ{\middel}キン \mbox{(Z\={O}{\middel}KIN)}\\\hline
\notyet{雑巾掛け}&dust cloth + hang = \textit{dirty work}&ゾウ{\middel}キン{\wsplit}が{\middel}け \mbox{(Z\={O}{\middel}KIN{\wsplit}ga{\middel}ke)}\\\hline
\notyet{頭巾}&head + cloth = \textit{hood}&ズ{\middel}キン \mbox{(ZU{\middel}KIN)}\\\hline
\notyet{黒頭巾}&black + hood = \textit{black hood}&くろ{\middel}ズ{\middel}キン \mbox{(kuro{\middel}ZU{\middel}KIN)}\\\hline
\end{somme}

\begin{decomp}{6}
あそこ&に&巾着&を&掛けて&下さい。\\\hline
あそこ&に&キンチャク&を&かけて&ください\\\hline
over there&に&pouch&を&hang&please\\\hline
\multicolumn{6}{|r|}{\textit{Please hang the pouch over there.}}\\\hline
\end{decomp}
\end{multicols}
\pagebreak
%%--------------------------------------------------------------------
%3--------------------------------------------------------------------
\begin{multicols}{2}
\begin{glyph}{City (市)}{city}\label{glyph:city}
City.\\\hline
A secondary meaning relates to \textit{market}.  In this sense the character can refer to anything from a \textit{simple fair} to the idea of a \textit{general business market}.\\\hline
This will be the first of four 漢字 in this chapter connected with \textit{groups of people living in a defined area}; moving from larger to smaller units will provide a sense of their relative sizes.
\end{glyph}
\gindex{City}{市}\strindex{5}{市}

\begin{comps}
\comprtwocone&亠 + 巾\\\hline
\comprthreecone&Top/Lid + Cloth\\\hline
\end{comps}

\begin{remglyph}
\hooglig{The city} on \remcom{top} has \remcom{clothes}.\\\hline
\end{remglyph}

\begin{prons}
\pronsrtwocone&\textbf{シ (SHI)}&---\\\hline
\pronsrthreecone&--- &いち (ichi)\\\hline
\end{prons}
\rindex{SHI}{市}
\rindex{ichi}{市}

\begin{remprons}
\hooglig{City} \comon{shielded} from \ucomkun{itchiness}.\\\hline
\end{remprons}

%{\middel}
% &  ()\\\hline
\begin{somme}
市内&city + inside = \textit{within the city}&シ{\middel}ナイ \mbox{(SHI{\middel}NAI)}\\\hline
市立&city + stand = \textit{municipal}&シ{\middel}リツ \mbox{(SHI{\middel}RITSU)}\\\hline
市有&city + have = \textit{city-owned}&シ{\middel}ユウ \mbox{(SHI{\middel}Y\={U})}\\\hline
\notyet{都市}&metropolis + city = \textit{town; city}&ト{\middel}シ \mbox{(TO{\middel}SHI)}\\\hline
\notyet{市民}&city + people = \textit{citizen}&シ{\middel}ミン \mbox{(SHI{\middel}MIN)}\\\hline
\notyet{市場}&city + place = \textit{marketplace}&いち{\middel}ば \mbox{(ichi{\middel}ba)}\\\hline
\notyet{都市計画}&city + plan = \textit{urban planning}&ト{\middel}シ{\middel}ケイ{\middel}カク \mbox{(TO{\middel}SHI{\middel}KEI{\middel}KAKU)}\\\hline
\end{somme}
\end{multicols}

\begin{decomp}{8}
夜&は&一人&で&市内&を&歩くないで&下さい。\\\hline
よる&は&ひとり&で&シナイ&を&あるくないで&ください\\\hline
night&は&1 person&で&within the city&を&don't walk&please\\\hline
\multicolumn{8}{|r|}{\textit{Please don't walk alone at night in the city.}}\\\hline
\end{decomp}
%%--------------------------------------------------------------------
%4--------------------------------------------------------------------
\begin{multicols}{2}
\begin{glyph}{Primary (主)}{primary}\label{glyph:primary}
Primary.\\\hline
This character conveys a sense of \textit{importance} and \textit{authority}; it is used with both concrete objects and abstract ideas.\\\hline
Note the difference between this 漢字 and 玉 (p. \pageref{glyph:jewel}).
\end{glyph}
\gindex{Primary}{主}\strindex{5}{主}

\begin{comps}
\comprtwocone&丶 + 王\\\hline
\comprthreecone&Dot + King\\\hline
\end{comps}

\begin{remglyph}
\hooglig{The primary} \remcom{jewel} of the \remcom{king}.\\\hline
\end{remglyph}

\begin{prons}
\pronsrtwocone&\textbf{シュ (SHU)}&\textbf{ぬし (nushi)}\\\hline
\pronsrthreecone&ス (SU)&おも (omo)\\\hline
\end{prons}
\rindex{SHU}{主}
\rindex{nushi}{主}
\rindex{SU}{主}
\rindex{omo}{主}

\begin{remprons}
\hooglig{Primarily} made from \comkun{new sheep}, this \comon{shooting} \ucomon{suit} can be placed in \ucomkun{omo}.\\\hline
{\warn}\textit{Omo} is a laundry detergent used for stain removal.\\\hline
\end{remprons}

%{\middel}
% &  ()\\\hline
\begin{somme}
主力&primary + strength = \textit{main force}&シュ{\middel}リョク \mbox{(SHU{\middel}RYOKU)}\\\hline
主人公&primary + person + public = \textit{main character}&シュ{\middel}ジン{\middel}コウ \mbox{(SHU{\middel}JIN{\middel}K\={O})}\\\hline
売主&sell + primary = \textit{seller}&うり{\middel}ぬし \mbox{(uri{\middel}nushi)}\\\hline
持主&hold + primary = \textit{owner}&もち{\middel}ぬし \mbox{(mochi{\middel}nushi)}\\\hline
\notyet{主食}&primary + food = \textit{staple food}&シュ{\middel}ショク \mbox{(SHU{\middel}SHOKU)}\\\hline
\notyet{主語}&primary + word = \textit{subject}&シュ{\middel}ゴ \mbox{(SHU{\middel}GO)}\\\hline
\notyet{主題}&primary + topic = \textit{theme; motif}&シュ{\middel}ダイ \mbox{(SHU{\middel}DAI)}\\\hline
\notyet{主治医}&primary + reign + doctor = \textit{attending physician}&シュ{\middel}ジ{\middel}イ \mbox{(SHU{\middel}JI{\middel}I)}\\\hline
\end{somme}

\begin{decomp}{6}
主人公&は&有名&な&美人&です。\\\hline
シュジンコウ&は&ユウメイ&な&ビジン&です\\\hline
main character&は&famous&な&beautiful woman&be/is\\\hline
\multicolumn{6}{|r|}{\textit{The main character is a famously beautiful person.}}\\\hline
\end{decomp}
\end{multicols}
\pagebreak
%%--------------------------------------------------------------------
%5--------------------------------------------------------------------
\begin{multicols}{2}
\begin{glyph}{Block (丁)}{block}\label{glyph:block}
The main sense is of a city \textit{block}.\\\hline
More obscure meanings can range from \textit{pages}, to \textit{polite}.\\\hline
This character often shows up as a component in other 漢字.
\end{glyph}
\gindex{Block}{丁}\strindex{2}{丁}

\begin{comps}
\comprtwocone&丁\\\hline
\comprthreecone&Block\\\hline
\end{comps}

\begin{remglyph}
\hooglig{Blocks} have \remcom{T-junctions}.\\\hline
\end{remglyph}

\begin{prons}
\pronsrtwocone&\textbf{チョウ (CH\={O})}&---\\\hline
\pronsrthreecone&テイ (TEI)&---\\\hline
\end{prons}
\rindex{CHOU@CH\={O}}{丁}
\rindex{TEI}{丁}

\begin{remprons}
\hooglig{Blocks} full of \ucomon{taping} \comon{chores}.\\\hline
\end{remprons}

%{\middel}
% &  ()\\\hline
\begin{somme}
丁目&block + eye = \textit{city block area}&チョウ{\middel}め \mbox{(CH\={O}{\middel}me)}\\\hline
二丁目&two + block + eye = \textit{City Block 2}&ニ{\middel}チョウ{\middel}め \mbox{(NI{\middel}CH\={O}{\middel}me)}\\\hline
四丁目&four + block + eye = \textit{City Block 4}&よん{\middel}チョウ{\middel}め \mbox{(yon{\middel}CH\={O}{\middel}me)}\\\hline
\notyet{丁重}&polite + heavy = \textit{polite; hospitable}&テイ{\middel}チョウ \mbox{(TEI{\middel}CH\={O})}\\\hline
\notyet{横丁}&sideways + block = \textit{side street}&よこ{\middel}チョウ \mbox{(yoko{\middel}CH\={O})}\\\hline
\end{somme}
\end{multicols}

\begin{decomp}{8}
私&の&新しい&家&は&四丁目&に&あります。\\\hline
わたし&の&あたらしい&いえ&は&よんチョウメ&に&あります\\\hline
I&の&new&house&は&City Block 4&に&is\\\hline
\multicolumn{8}{|r|}{\textit{My new house is in City Block 4.}}\\\hline
\end{decomp}
%%--------------------------------------------------------------------
%6--------------------------------------------------------------------
\begin{multicols}{2}
\begin{glyph}{Friend (友)}{friend}\label{glyph:friend}
Friend.
\end{glyph}
\gindex{Friend}{友}\strindex{4}{友}

\begin{comps}
\comprtwocone&ナ + 又\\\hline
\comprthreecone&Ten variant/(NA) + Again\\\hline
\end{comps}

\begin{remglyph}
\hooglig{Friends} will watch \remcom{Naruto} \remcom{again}.\\\hline
{\warn}\textit{Naruto} is an アニメ about a young ninja seeking recognition.\\\hline
\end{remglyph}

\begin{prons}
\pronsrtwocone&\textbf{ユウ (Y\={U})}&\textbf{とも (tomo)}\\\hline
\pronsrthreecone&---&---\\\hline
\end{prons}
\rindex{YUU@Y\={U}}{友}
\rindex{tomo}{友}

\begin{remprons}
\hooglig{Friend}, let's \comon{youtube} \comkun{tomorrow}.\\\hline
\end{remprons}

%{\middel}
% &  ()\\\hline
\begin{somme}
友達&friend + achieve = \textit{friend(s)}&とも{\middel}だち \mbox{(tomo{\middel}dachi)}\\\hline
友人&friend + person = \textit{friend}&ユウ{\middel}ジン \mbox{(Y\={U}{\middel}JIN)}\\\hline
友好&friend + like = \textit{friendship}&ユウ{\middel}コウ \mbox{(Y\={U}{\middel}K\={O})}\\\hline
交友&mix + friend = \textit{acquaintance}&コウ{\middel}ユウ \mbox{(K\={O}{\middel}Y\={U})}\\\hline
女友達&woman + friend = \textit{female friend}&おんなともだち \mbox{(onna{\middel}tomo{\middel}dachi)}\\\hline
親友&parent + friend = \textit{close friend}&シン{\middel}ユウ \mbox{(SHIN{\middel}Y\={U})}\\\hline
真の友&truth + friend = \textit{true friend}&シン{\middel}の{\middel}とも \mbox{(SHIN{\middel}no{\middel}tomo)}\\\hline
\notyet{悪友}&evil + friend = \textit{bad company}&アク{\middel}ユウ \mbox{(AKU{\middel}Y\={U})}\\\hline
\notyet{遊び友達}&play + friend = \textit{playmate}&あそ{\middel}び{\wsplit}とも{\middel}だち \mbox{aso{\middel}bi{\wsplit}tomo{\middel}dachi}\\\hline
\notyet{茶飲み友達}&tea drinking + friend = \textit{tea-drinking companion}&チャ{\wsplit}の{\middel}み{\wsplit}とも{\middel}だち CHA{\wsplit}no{\middel}mi{\wsplit}tomo{\middel}dachi\\\hline
\end{somme}
\end{multicols}

\begin{decomp}{7}
西山さん&の&友人&は&南米&から&来ました。\\\hline
にしやまさん&の&ユウジン&は&ナンベイ&から&きました\\\hline
Nishiyama-san&の&friend&は&South America&from&came\\\hline
\multicolumn{7}{|r|}{\textit{Nishiyama-san's friend came from South America.}}\\\hline
\end{decomp}
\pagebreak
%%--------------------------------------------------------------------
%7--------------------------------------------------------------------
\begin{multicols}{2}
\begin{glyph}{Group (団)}{group}\label{glyph:group}
Group.
\end{glyph}
\gindex{Group}{団}\strindex{6}{団}

\begin{comps}
\comprtwocone&囗 + 寸\\\hline
\comprthreecone&Enclosure/Box/Prison + Inch\\\hline
\end{comps}

\begin{remglyph}
\hooglig{The group} in the \remcom{prison} is \remcom{tiny}.\\\hline
\end{remglyph}

\begin{prons}
\pronsrtwocone&\textbf{ダン (DAN)}&---\\\hline
\pronsrthreecone&トン (TON)&---\\\hline
\end{prons}
\rindex{DAN}{団}
\rindex{TON}{団}

\begin{remprons}
\hooglig{Group} \comon{dunking} is a \ucomon{tonic}.\\\hline
\end{remprons}

%{\middel}
% &  ()\\\hline
\begin{somme}
一団&one + group = \textit{a group}&イチ{\middel}ダン \mbox{(ICHI{\middel}DAN)}\\\hline
団体&group + body = \textit{organisation}&ダン{\middel}タイ \mbox{(DAN{\middel}TAI)}\\\hline
団子&group + child = \textit{dumpling}&ダン{\middel}ご \mbox{(DAN{\middel}go)}\\\hline
入団&enter + group = \textit{joining (a group)}&ニュウ{\middel}ダン \mbox{(NY\={U}{\middel}DAN)}\\\hline
\notyet{集団}&gather + group = \textit{crowd; mass}&シュウ{\middel}ダン \mbox{(SH\={U}{\middel}DAN)}\\\hline
\notyet{団結}&group + bind = \textit{unity; solidarity}&ダン{\middel}ケツ \mbox{(DAN{\middel}KETSU)}\\\hline
\notyet{団地}&group + place = \textit{apartment complex}&ダン{\middel}チ \mbox{(DAN{\middel}CHI)}\\\hline
\notyet{記者団}&reporter + group = \textit{press organisation}&キ{\middel}シャ{\middel}ダン \mbox{(KI{\middel}SHA{\middel}DAN)}\\\hline
\notyet{代表団}&delegate + group = \textit{delegation}&ダイ{\middel}ヒョウ{\middel}ダン \mbox{(DAI{\middel}HY\={O}{\middel}DAN)}\\\hline
布団{\notcov}&spread/linen + group = \textit{bedding}&フ{\middel}トン \mbox{(FU{\middel}TON)}\\\hline
\end{somme}
\end{multicols}

\begin{decomp}{6}
内田さん&は&団子&が&好き&です。\\\hline
うちださん&は&ダンご&が&すき&です\\\hline
Uchida-san&は&dumplings&が&likes&be/is\\\hline
\multicolumn{6}{|r|}{\textit{Uchida-san likes dumplings.}}\\\hline
\end{decomp}
%%--------------------------------------------------------------------
%8--------------------------------------------------------------------
\begin{multicols}{2}
\begin{glyph}{Every (毎)}{every}\label{glyph:every}
Every.
\end{glyph}
\gindex{Every}{毎}\strindex{6}{毎}

\begin{comps}
\comprtwocone&𠂉 + 毋\\\hline
\comprthreecone&Person + Do Not/Mother\\\hline
\end{comps}

\begin{remglyph}
\hooglig{Every} \remcom{person} has a \remcom{mother}.\\\hline
\end{remglyph}

\begin{prons}
\pronsrtwocone&\textbf{マイ (MAI)}&---\\\hline
\pronsrthreecone&---&ごと (goto)\\\hline
\end{prons}
\rindex{MAI}{毎}
\rindex{goto}{毎}

\begin{remprons}
\hooglig{Every} \ucomkun{goto} \comon{migration}.\\\hline
\end{remprons}

%{\middel}
% &  ()\\\hline
\begin{somme}
毎日&every + sun = \textit{daily}&マイ{\middel}ニチ \mbox{(MAI{\middel}NICHI)}\\\hline
毎週&every + week = \textit{weekly}&マイ{\middel}シュウ \mbox{(MAI{\middel}SH\={U})}\\\hline
毎月&every + moon (month) = \textit{montly}&マイ{\middel}つき \mbox{(MAI{\middel}tsuki)}\\\hline
毎年&every + year = \textit{yearly}&マイ{\middel}とし \mbox{(MAI{\middel}toshi)}\\\hline
毎朝&every + morning = \textit{every morning}&マイ{\middel}あさ \mbox{(MAI{\middel}asa)}\\\hline
毎回&every + rotate = \textit{each round}&マイ{\middel}カイ \mbox{(MAI{\middel}KAI)}\\\hline
毎時&every + time = \textit{hourly}&マイ{\middel}ジ \mbox{(MAI{\middel}JI)}\\\hline
毎秒&every + second = \textit{every second}&マイ{\middel}ビョウ \mbox{(MAI{\middel}BY\={O})}\\\hline
毎夕&every + evening = \textit{every evening}&マイ{\middel}ユウ \mbox{(MAI{\middel}Y\={U})}\\\hline
\notyet{毎晩}&every + late = \textit{every night}&マイ{\middel}バン \mbox{(MAI{\middel}BAN)}\\\hline
\notyet{毎度}&every + degree = \textit{each time}&マイ{\middel}ド \mbox{(MAI{\middel}DO)}\\\hline
\end{somme}

\begin{decomp}{5}
毎日&何時&に&起きます&か。\\\hline
マイニチ&なんジ&に&おきます&か\\\hline
every day&what time&に&get up&か\\\hline
\multicolumn{5}{|r|}{\textit{What time do you get up every day?}}\\\hline
\end{decomp}
\end{multicols}

\pagebreak
%%--------------------------------------------------------------------
%9--------------------------------------------------------------------
\begin{multicols}{2}
\begin{glyph}{Empty (空)}{empty}\label{glyph:empty}
Empty.  Sky.\\\hline
As in English, \textit{empty} can incorporate suggestions of \textit{uselessness} and \textit{futility}.
\end{glyph}
\gindex{Empty}{空}\strindex{8}{空}

\begin{comps}
\comprtwocone & 穴 + 工 \\\hline
\comprthreecone & Cave + Craft/Work \\\hline
\end{comps}

\begin{remglyph}
\hooglig{Empty} \remcom{caves} are \remcom{crafted} out.\\\hline
\end{remglyph}

\begin{prons}
\pronsrtwocone & \textbf{クウ (K\={U})} & \textbf{そら、から (sora, kara)} \\\hline
\rye{3}{クウ is by far the most common reading in compounds.\\\hline
そら means \textit{sky}.\\\hline
から occasionally appears in the first position with a meaning of \textit{empty}.}
\pronsrthreecone&---&あ (a)\\\hline
\rye{3}{The あ reading is used to make the verbs 空く and 空ける.  They can be thought of as interchangeable with the more common 開く and 開ける from Entry \ref{glyph:open} (p. \pageref{glyph:open}).}
\end{prons}
\rindex{KUU@K\={U}}{空}
\rindex{sora}{空}
\rindex{kara}{空}
\rindex{a}{空}

\begin{remprons}
\hooglig{Empty} \comkun{saw lumps} were \ucomkun{updated} to be \comon{cool}
and \comkun{caramelized}.\\\hline
\end{remprons}

%{\middel}
% &  ()\\\hline
\begin{somme}
\splitter{空く}{intr}&to be open&あ{\middel}く \mbox{(a{\middel}ku)}\\\hline
\splitter{空ける}{tr}&to open something&あ{\middel}ける \mbox{(a{\middel}keru)}\\\hline
空&\textit{sky}&そら \mbox{(sora)}\\\hline
青空&blue + empty = \textit{blue sky}&あお{\middel}ぞら \mbox{(ao{\middel}zora)}\\\hline
空手&empty + hand = \textit{karate}&から{\middel}て \mbox{(kara{\middel}te)}\\\hline
空間&empty + interval = \textit{space}&クウ{\middel}カン \mbox{(K\={U}{\middel}KAN)}\\\hline
空気&empty + spirit = \textit{air}&クウ{\middel}キ \mbox{(K\={U}{\middel}KI)}\\\hline
空中&empty + middle = \textit{mid-air}&クウ{\middel}チュウ \mbox{(K\={U}{\middel}CH\={U})}\\\hline
大空&large + open = \textit{firmament}&おお{\middel}ぞら \mbox{(\={o}{\middel}zora)}\\\hline
上空&upper + empty = \textit{the skies}&ジョウ{\middel}クウ \mbox{(J\={O}{\middel}K\={U})}\\\hline
\notyet{空港}&empty + port = \textit{Airport}&クウ{\middel}コウ \mbox{(K\={U}{\middel}K\={O})}\\\hline
\notyet{真空}&true + empty = \textit{vacuum; hollow}&シン{\middel}クウ \mbox{(SHIN{\middel}K\={U})}\\\hline
{\iscov}空き缶&empty + can = \textit{empty can}&あき{\middel}カン \mbox{(aki{\middel}KAN)}\\\hline
空き瓶{\notcov}&empty + bin = \textit{empty bin}&あき{\middel}ビン \mbox{(aki{\middel}BIN)}\\\hline
\end{somme}

\begin{decomp}{7}
朝&の&空&は&きれい&です&よ。\\\hline
あさ&の&そら&は&きれい&です&よ\\\hline
morning&の&skry&は&pretty&be/is&よ\\\hline
\multicolumn{7}{|r|}{\textit{The morning sky is lovely.}}\\\hline
\end{decomp}
\end{multicols}
%%--------------------------------------------------------------------
%10-------------------------------------------------------------------
\begin{multicols}{2}
\begin{glyph}{Town (町)}{town}\label{glyph:town}
Town.  One size smaller than a ``City'' (p. \pageref{glyph:city}).\\\hline
On occasion it can also refer to a \textit{street} or an \textit{alley}.
\end{glyph}
\gindex{Town}{町}\strindex{7}{町}

\begin{comps}
\comprtwocone&田 + 丁\\\hline
\comprthreecone&Rice field + Block (p. \pageref{glyph:block})\\\hline
\end{comps}

\begin{remglyph}
A \hooglig{town} of \remcom{rice-field} \remcom{blocks}.\\\hline
\end{remglyph}

\begin{prons}
\pronsrtwocone&\textbf{チョウ (CH\={O})}&\textbf{まち (machi)}\\\hline
\pronsrthreecone&---&---\\\hline
\end{prons}
\rindex{CHOU@CH\={O}}{町}
\rindex{machi}{町}

\begin{remprons}
\hooglig{Town} \comkun{marching} is such a \comon{chore}.\\\hline
\end{remprons}

%{\middel}
% &  ()\\\hline
\begin{somme}
町&\textit{town}&まち \mbox{(machi)}\\\hline
小町&small + town = \textit{belle; town-beauty}&こ{\middel}まち \mbox{(ko{\middel}machi)}\\\hline
町内&town + inside = \textit{in the town}&チョウ{\middel}ナイ \mbox{(CH\={O}{\middel}NAI)}\\\hline
町立&town + stand = \textit{established by the town}&チョウ{\middel}リツ \mbox{(CH\={O}{\middel}RITSU)}\\\hline
町会&town + meet = \textit{town assembly}&チョウ{\middel}カイ \mbox{(CH\={O}{\middel}KAI)}\\\hline
町中&town + middle = \textit{downtown}&まち{\middel}なか \mbox{(machi{\middel}naka)}\\\hline
町民&town + people = \textit{townspeople}&チョウ{\middel}ミン \mbox{CH\={O}{\middel}MIN}\\\hline
\notyet{町村}&town + village = \textit{towns and villages}&チョウ{\middel}ソン \mbox{(CH\={O}{\middel}SON)}\\\hline
\end{somme}

\begin{decomp}{7}
此の&町&の&ラーメン&は&有名&です。\\\hline
この&まち&の&ラーメン&は&ユウメイ&です\\\hline
this&town&の&ramen&は&famous&be/is\\\hline
\multicolumn{7}{|r|}{\textit{This town's ramen is famous.}}\\\hline
\end{decomp}
\end{multicols}
\pagebreak
%%--------------------------------------------------------------------
%11-------------------------------------------------------------------
\begin{multicols}{2}
\begin{glyph}{Dog (犬)}{dog}\label{glyph:dog}
Dog.
\end{glyph}
\gindex{Dog}{犬}\strindex{4}{犬}

\begin{comps}
\comprtwocone&犬\\\hline
\comprthreecone&Dog\\\hline
\end{comps}

\begin{remglyph}
Part of the 漢字 radicals (\#{94}).\\\hline
\end{remglyph}

\begin{prons}
\pronsrtwocone&\textbf{ケン (KEN)}&\textbf{いぬ (inu)}\\\hline
\pronsrthreecone&---&---\\\hline
\end{prons}
\rindex{KEN}{犬}
\rindex{inu}{犬}

\begin{remprons}
\hooglig{Dogs}, like \comkun{inuyasha}, have \comon{kennels}.\\\hline
{\warn}\textit{Inuyasha} is a famous アニメ character who is half-demon and half-human.\\\hline
\end{remprons}

%{\middel}
% &  ()\\\hline
\begin{somme}
犬&\textit{dog}&いぬ \mbox{(inu)}\\\hline
子犬&child + dog = \textit{puppy}&こ{\middel}いぬ \mbox{(ko{\middel}inu)}\\\hline
秋田犬&autumn + rice field + dog = \textit{an Akita dog}&あき{\middel}た{\middel}ケン \mbox{(aki{\middel}ta{\middel}KEN)}\\\hline
\notyet{愛犬}&love + dog = \textit{pet dog}&アイ{\middel}ケン \mbox{(AI{\middel}KEN)}\\\hline
\notyet{愛犬家}&love + dog + house = \textit{dog lover}&アイ{\middel}ケン{\middel}カ \mbox{(AI{\middel}KEN{\middel}KA)}\\\hline
\notyet{番犬}&order + dog = \textit{watch-dog}&バン{\middel}ケン \mbox{(BAN{\middel}KEN)}\\\hline
\notyet{野犬}&field + dog = \textit{stray dog}&ヤ{\middel}ケン \mbox{(YA{\middel}KEN)}\\\hline
\notyet{犬小屋}&dog + hut = \textit{kennel}&いぬ{\middel}ご{\middel}や \mbox{(inu{\middel}go{\middel}ya)}\\\hline
\notyet{軍犬}&army + dog = \textit{wardog}&グン{\middel}ケン \mbox{(GUN{\middel}KEN)}\\\hline
\end{somme}

\begin{decomp}{5}
古川さん&が&子犬&を&買いました。\\\hline
ふるかわさん&が&こいぬ&を&かいました\\\hline
Furukawa-san&が&puppy&を&bought\\\hline
\multicolumn{5}{|r|}{\textit{Furukawa-san bought a puppy.}}\\\hline
\end{decomp}
\end{multicols}
%%--------------------------------------------------------------------
%12-------------------------------------------------------------------
\begin{multicols}{2}
\begin{glyph}{Death (死)}{death}\label{glyph:death}
Death.
\end{glyph}
\gindex{Death}{死}\strindex{6}{死}

\begin{comps}
\comprtwocone&歹 + ヒ\\\hline
\comprthreecone&Death + Spoon/Ladle\\\hline
\end{comps}

\begin{remglyph}
\hooglig{Death} \remcom{spoon}.\\\hline
\hooglig{Death} on \remcom{one} \remcom{evening} from \remcom{spooning}.\\\hline
\end{remglyph}

\begin{prons}
\pronsrtwocone&\textbf{シ (SHI)}&\textbf{し (shi)}\\\hline
\pronsrthreecone&---&---\\\hline
\end{prons}
\rindex{SHI}{死}
\rindex{shi}{死}

\begin{remprons}
\hooglig{Death}\ldots  There is no \comon{shield} against or \comkun{shifting} away from it.\\\hline
\end{remprons}

%{\middel}
% &  ()\\\hline
\begin{somme}
\splitter{死ぬ}{intr}&\textit{to die}&し{\middel}ぬ \mbox{(shi{\middel}nu)}\\\hline
水死&water + death = \textit{drowning}&スイ{\middel}シ \mbox{(SUI{\middel}SHI)}\\\hline
死体&death + body = \textit{dead body}&シ{\middel}タイ \mbox{(SHI{\middel}TAI)}\\\hline
死火山&death + fire + mountain = \textit{extinct volcano}&シ{\middel}カ{\middel}ザン \mbox{(SHI{\middel}KA{\middel}ZAN)}\\\hline
死力&death + strength = \textit{desperate effort}&シリ{\middel}ョク \mbox{(SHI{\middel}RYOKU)}\\\hline
死者&death + individual = \textit{dead person; casualties}&シ{\middel}シャ \mbox{(SHI{\middel}SHA)}\\\hline
死語&death + word = \textit{extinct language}&シ{\middel}ゴ \mbox{(SHI{\middel}GO)}\\\hline
\notyet{必死}&certain + death = \textit{frantic; inevitable death}&ヒッ{\middel}シ \mbox{(HIS{\middel}SHI)}\\\hline
\notyet{死亡}&death + perish = \textit{death; mortality}&シ{\middel}ボウ \mbox{(SHI{\middel}B\={O})}\\\hline
\notyet{死産}&death + produce = \textit{stillbirth}&シ{\middel}ザン \mbox{(SHI{\middel}ZAN)}\\\hline
\notyet{死去}&death + depart = \textit{passing away}&シ{\middel}キョ \mbox{(SHI{\middel}KYO)}\\\hline
\notyet{死神}&death + god = \textit{god of death}&しに{\middel}がみ \mbox{(shini{\middel}gami)}\\\hline
\end{somme}

\begin{decomp}{4}
人間&は&必ず&死ぬ。\\\hline
ニンゲン&は&かならず&しぬ\\\hline
human&は&certainly&to die\\\hline
\multicolumn{4}{|r|}{\textit{Man is bound to die.}}\\\hline
\end{decomp}
\end{multicols}

\pagebreak
%%--------------------------------------------------------------------
%13-------------------------------------------------------------------
\begin{multicols}{2}
\begin{glyph}{Go (行)}{go}\label{glyph:go}
\textit{Go} is the primary meaning of this character.\\\hline
Related ideas are \textit{carrying out} or \textit{doing}.\\\hline
Far less common is the secondary sense of a \textit{stroke} or \textit{line} of text.
\end{glyph}
\gindex{Go}{行}\strindex{6}{行}

\begin{comps}
\comprtwocone&彳 + 一 + 丁\\\hline
\comprthreecone&Step + One + Block (p. \pageref{glyph:block})\\\hline
\end{comps}

\begin{remglyph}
\hooglig{Go} \remcom{step} on \remcom{one} \remcom{block}.\\\hline
\end{remglyph}

\begin{prons}
\pronsrtwocone&\textbf{コウ (K\={O})}&\textbf{い、ゆ (i, yu)}\\\hline
\rye{3}{For the general sense of the verb \textit{to go}, いく is used.\\\hline
ゆき is used in a few specific compounds, or as a suffix meaning \textit{bound for}.}
\pronsrthreecone&ギョウ、アン (GY\={O}, AN)&おこな (okona)\\\hline
\rye{3}{ギョウ is found primarily in compounds relating to this 漢字's secondary meaning of \textit{line} or \textit{stroke}.}
\end{prons}
\rindex{KOU@K\={O}}{行}
\rindex{i}{行}
\rindex{yu}{行}
\rindex{GYOU@GY\={O}}{行}
\rindex{AN}{行}
\rindex{okona}{行}

\begin{remprons}
\hooglig{Go} \comon{cause} a \ucomon{long-gyone} \comkun{interesting} \comkun{yule} in \ucomkun{a corner} with your \ucomon{uncle}.\\\hline
{\warn}\textit{Yule} is a festival historically observed by the Germanic peoples.\\\hline
\end{remprons}

%{\middel}
% &  ()\\\hline
\begin{somme}
\rye{3}{行 can also act like 買, but on rare occasions it incorporates ゆく as opposed to ゆき.}
行く&\textit{to go}&い{\middel}く \mbox{(i{\middel}ku)}\\\hline
行き先&go + precede = \textit{destination}&ゆ{\middel}き{\wsplit}さき \mbox{(yu{\middel}ki{\wsplit}saki)}\\\hline
東京行&east + capital + go = \textit{bound for Tokyo}&トウ{\middel}キョウ{\wsplit}ゆ{\middel}き \mbox{(T\={O}{\middel}KY\={O}{\wsplit}yu{\middel}ki)}\\\hline
夜行&night + go = \textit{night travel}&ヤ{\middel}コウ \mbox{(YA{\middel}K\={O})}\\\hline
歩行&walk + go = \textit{walking}&ホ{\middel}コウ \mbox{(HO{\middel}K\={O})}\\\hline
歩行者&walk + go + individual = \textit{pedestrian}&ホ{\middel}コウ{\middel}シャ \mbox{(HO{\middel}K\={O}{\middel}SHA)}  \\\hline
\notyet{行方不明}&whereabouts + unclear = \textit{missing}&ゆく{\middel}え{\wsplit}フ{\middel}メイ \mbox{(yuku{\middel}e{\wsplit}FU{\middel}MEI)}\\\hline
\notyet{急行}&hasten + go = \textit{hastening}&キュウ{\middel}コウ \mbox{(KY\={U}{\middel}K\={O})}\\\hline
\end{somme}

\begin{decomp}{5}
今日&は&歩行者&が&多い。\\\hline
きょう&は&ホコウシャ&が&おおい\\\hline
today&は&pedestrian&が&many\\\hline
\multicolumn{5}{|r|}{\textit{There are many pedestrians today.}}\\\hline
\end{decomp}
\end{multicols}
%%--------------------------------------------------------------------
%14-------------------------------------------------------------------
\begin{multicols}{2}
\begin{glyph}{Take (取)}{take}\label{glyph:take}
Take.
\end{glyph}
\gindex{Take}{取}\strindex{8}{取}

\begin{comps}
\comprtwocone&耳 + 又\\\hline
\comprthreecone&Ear + Again\\\hline
\end{comps}

\begin{remglyph}
\hooglig{Take} by the \remcom{ear} \remcom{again}.\\\hline
\end{remglyph}

\begin{prons}
\pronsrtwocone&---&\textbf{と (to)}\\\hline
\pronsrthreecone&シュ (SHU)&---\\\hline
\end{prons}
\rindex{to}{取}
\rindex{SHU}{取}

\begin{remprons}
\hooglig{Take} on the \comkun{torrential} \ucomon{shooting}.\\\hline
\end{remprons}

%{\middel}
% &  ()\\\hline
\begin{somme}
取る&\textit{to take}&と{\middel}る \mbox{(to{\middel}ru)}\\\hline
取引&take + pull = \textit{transactions}&とり{\wsplit}ひ{\middel}き \mbox{(tori{\wsplit}hi{\middel}ki)}\\\hline
\splitter{取り上げる}{tr}&take + upper = \textit{to pick up}&と{\middel}り{\wsplit}あ{\middel}げる \mbox{(to{\middel}ri{\wsplit}a{\middel}geru)}\\\hline
\splitter{取り入れる}{tr}&take + enter = \textit{to take in}&と{\middel}り{\wsplit}い{\middel}れる \mbox{(to{\middel}ri{\wsplit}i{\middel}reru)}\\\hline
\splitter{取り出す}{tr}&take + exit = \textit{to take out}&と{\middel}り{\wsplit}だ{\middel}す \mbox{(tori{\wsplit}dasu)}\\\hline
\splitter{気取る}{intr}&spirit + take = \textit{to put on airs}&キ{\wsplit}ど{\middel}る \mbox{(KI{\wsplit}do{\middel}ru)}\\\hline
\splitterny{受取る}{tr}&receive + take = \textit{to receive}&うけ{\wsplit}と{\middel}る \mbox{(uke{\wsplit}to{\middel}ru)}\\\hline
\notyet{受取り}&receive + take = \textit{receiving; receipt}&うけ{\wsplit}と{\middel}り \mbox{(uke{\wsplit}to{\middel}ri)}\\\hline
\end{somme}
\end{multicols}

\begin{decomp}{6}
良し、&此れ&で&引取&は&まとまった。\\\hline
よし&これ&で&とりひき&は&まとまった\\\hline
fine&this&で&transactions&は&to be settled\\\hline
\multicolumn{6}{|r|}{\textit{All right.  It's a deal.}}\\\hline
\end{decomp}

\pagebreak
%%--------------------------------------------------------------------
%15-------------------------------------------------------------------
\begin{multicols}{2}
\begin{glyph}{Village (村)}{village}\label{glyph:village}
Village.\\\hline
This character implies something smaller than a city or town.
\end{glyph}
\gindex{Village}{村}\strindex{7}{村}

\begin{comps}
\comprtwocone&木 + 寸\\\hline
\comprthreecone&Tree + Tiny\\\hline
\end{comps}

\begin{remglyph}
\hooglig{Village} \remcom{trees} are \remcom{tiny}.\\\hline
\end{remglyph}

\begin{prons}
\pronsrtwocone&\textbf{ソン (SON)}&\textbf{むら (mura)}\\\hline
\pronsrthreecone&---&---\\\hline
\end{prons}
\rindex{SON}{村}
\rindex{mura}{村}

\begin{remprons}
\hooglig{Village} \comon{songs} depicted in the \comkun{mural}.\\\hline
\end{remprons}

%{\middel}
% &  ()\\\hline
\begin{somme}
村&\textit{village}&むら \mbox{(mura)}\\\hline
村々&village $(\times 2)$ = \textit{villages}&むら{\middel}むら \mbox{(mura{\middel}mura)}\\\hline
村人&village + person = \textit{villager}&むら{\middel}びと \mbox{(mura{\middel}bito)}\\\hline
山村&mountain + village = \textit{mountain village}&サン{\middel}ソン \mbox{(SAN{\middel}SON)}\\\hline
村会&village + meet = \textit{village assembly}&ソン{\middel}カイ \mbox{(SON{\middel}KAI)}\\\hline
入村&enter + village = \textit{moving to a village}&ニュウ{\middel}ソン \mbox{(NY\={U}{\middel}SON)}\\\hline
\notyet{村民}&village + people = \textit{villager}&ソン{\middel}ミン \mbox{(SON{\middel}MIN)}\\\hline
\notyet{無医村}&nil + doctor + village = \textit{doctorless village}&ム{\middel}イ{\middel}ソン \mbox{(MU{\middel}I{\middel}SON)}\\\hline
\notyet{選手村}&athlete + village = \textit{Olympic village}&セン{\middel}シュ{\middel}むら \mbox{(SEN{\middel}SHU{\middel}mura)}\\\hline
\end{somme}

\begin{decomp}{5}
村人&は&小道&に&立っています。\\\hline
むらびと&は&こみち&に&たっています\\\hline
villager&は&path&に&is standing\\\hline
\multicolumn{5}{|r|}{\textit{A villager is standing in the path.}}\\\hline
\end{decomp}
\end{multicols}
%%--------------------------------------------------------------------
%16-------------------------------------------------------------------
\begin{multicols}{2}
\begin{glyph}{Character (字)}{character}\label{glyph:character}
Character in the sense of a \textit{Chinese character} or \textit{letter of an alphabet}.
\end{glyph}
\gindex{Character}{字}\strindex{6}{字}

\begin{comps}
\comprtwocone&宀 + 子\\\hline
\comprthreecone&Roof + Child\\\hline
\end{comps}

\begin{remglyph}
\hooglig{Character} under a \remcom{roof} is \remcom{childish}.\\\hline
\end{remglyph}

\begin{prons}
\pronsrtwocone&\textbf{ジ (JI)}&---\\\hline
\pronsrthreecone&---&あざ (aza)\\\hline
\end{prons}
\rindex{JI}{字}
\rindex{aza}{字}

\begin{remprons}
\hooglig{Characters} in the \ucomkun{azalea} \comon{jig}.\\\hline
{\warn}\textit{Azaleas} are flowering shrubs in the genus \textit{Rhododendron}.\\\hline
\end{remprons}

%{\middel}
% &  ()\\\hline
\begin{somme}
十字&ten + character = \textit{a cross}&ジュウ{\middel}ジ \mbox{(J\={U}{\middel}JI)}\\\hline
赤十字&red + ten + character = \textit{The Red Cross}&セキ{\middel}ジュウ{\middel}ジ \mbox{(SEKI{\middel}J\={U}{\middel}JI)}\\\hline
字引&character + pull = \textit{dictionary}&ジ{\middel}びき \mbox{(JI{\middel}biki)}\\\hline
文字&literature + character = \textit{letter; character}&モ{\middel}ジ \mbox{(MO{\middel}JI)}\\\hline
赤字&red + character = \textit{being in the red}&あか{\middel}ジ \mbox{(aka{\middel}JI)}\\\hline
大文字&large + character = \textit{capital letter}&おお{\middel}モ{\middel}ジ \mbox{(\={o}{\middel}MO{\middel}JI)}\\\hline
\notyet{黒字}&black + character = \textit{being in the black}&くろ{\middel}ジ \mbox{(kuro{\middel}JI)}\\\hline
\notyet{漢字}&China + character = \textit{kanji}&カン{\middel}ジ \mbox{(KAN{\middel}JI)}\\\hline
\notyet{数字}&number + character = \textit{number; digit}&スウ{\middel}ジ \mbox{(S\={U}{\middel}JI)}\\\hline
\notyet{文字通り}&character + passage = \textit{literally; to the letter}&モ{\middel}ジ{\wsplit}とお{\middel}り \mbox{(MO{\middel}JI{\wsplit}t\={o}{\middel}ri)}\\\hline
\end{somme}
\end{multicols}

\begin{decomp}{9}
多く&の&国&で&は&赤十字&は&大切&です。\\\hline
おおく&の&くに&で&は&セイジュウジ&は&タイセツ&です\\\hline
many&の&countries&で&は&Red Cross&は&important&be/is\\\hline
\multicolumn{9}{|r|}{\textit{The Red Cross is important in many countries.}}\\\hline
\end{decomp}

\pagebreak
%%--------------------------------------------------------------------
%17-------------------------------------------------------------------
\begin{multicols}{2}
\begin{glyph}{Tool (具)}{tool}\label{glyph:tool}
Tool.  Equipment.
\end{glyph}
\gindex{Tool}{具}\strindex{8}{具}

\begin{comps}
\comprtwocone&目 + 丌\\\hline
\comprthreecone&Eye + Table\\\hline
\end{comps}

\begin{remglyph}
\hooglig{Tools} for your \remcom{eyes} are on the \remcom{table}.\\\hline
\end{remglyph}

\begin{prons}
\pronsrtwocone&\textbf{グ (GU)}&---\\\hline
\pronsrthreecone&---&---\\\hline
\end{prons}
\rindex{GU}{具}

\begin{remprons}
\hooglig{Tool} \comon{goofing}.\\\hline
\end{remprons}

%{\middel}
% &  ()\\\hline
\begin{somme}
道具&road + tool = \textit{tool}&ドウ{\middel}グ \mbox{(D\={O}{\middel}GU)}\\\hline
家具&house + tool = \textit{furniture}&カ{\middel}グ \mbox{(KA{\middel}GU)}\\\hline
具体&tool + body = \textit{definite/concrete}&グ{\middel}タイ \mbox{(GU{\middel}TAI)}\\\hline
具体化&concrete + change = \textit{embodiment}&グ{\middel}タイ{\middel}カ \mbox{(GU{\middel}TAI{\middel}KA)}\\\hline
馬具&horse + tool = \textit{harness}&バ{\middel}グ \mbox{(BA{\middel}GU)}\\\hline
雨具&rain + tool = \textit{rain gear}&あま{\middel}グ \mbox{(ama{\middel}GU)}\\\hline
夜具&night + tool = \textit{bedding; bed clothes}&ヤ{\middel}グ \mbox{(YA{\middel}GU)}\\\hline
工具&craft + tool = \textit{tool; implement}&コウ{\middel}グ \mbox{(K\={O}{\middel}GU)}\\\hline
\notyet{具合}&tool + match = \textit{condition}&グ{\wsplit}あ{\middel}い \mbox{(GU{\wsplit}a{\middel}i)}\\\hline
\notyet{絵の具}&picture + tool = \textit{paint}&エ{\middel}の{\middel}グ \mbox{(E{\middel}no{\middel}GU)}\\\hline
\end{somme}

\begin{decomp}{7}
その&家具&は&母&の&物&です。\\\hline
その&カグ&は&はは&の&もの&です\\\hline
that&furniture&は&mother&の&thing&be/is\\\hline
\multicolumn{7}{|r|}{\textit{The furniture belongs to my mother.}}\\\hline
\end{decomp}
\end{multicols}
%%--------------------------------------------------------------------
%18-------------------------------------------------------------------
\begin{multicols}{2}
\begin{glyph}{Pour (注)}{pour}\label{glyph:pour}
Pour.\\\hline
\textit{Directing} something toward an object (a gaze or one's attention, for instance).
\end{glyph}
\gindex{Pour}{注}\strindex{8}{注}

\begin{comps}
\comprtwocone&氵 + 主\\\hline
\comprthreecone&Water + Primary (p. \pageref{glyph:primary})\\\hline
\end{comps}

\begin{remglyph}
\hooglig{Pour} \remcom{water} \remcom{primarily}.\\\hline
\end{remglyph}

\begin{prons}
\pronsrtwocone&\textbf{チュウ (CH\={U})}&\textbf{そそ (soso)}\\\hline
\pronsrthreecone&---&---\\\hline
\end{prons}
\rindex{CHUU@CH\={U}}{注}
\rindex{soso}{注}

\begin{remprons}
\hooglig{Pour} \comkun{so-so} while \comon{chewing}.\\\hline
\end{remprons}

%{\middel}
% &  ()\\\hline
\begin{somme}
\splitter{注ぐ}{tr}&\textit{to pour}&そそ{\middel}ぐ \mbox{(soso{\middel}gu)}\\\hline
注意&pour + mind = \textit{attention}&チュウ{\middel}イ \mbox{(CH\={U}{\middel}I)}\\\hline
不注意&not + attention = \textit{carelessness}&フ{\middel}チュウ{\middel}イ \mbox{(FU{\middel}CH\={U}{\middel}I)}\\\hline
注文書&order + write = \textit{written order}&チュウ{\middel}モン{\middel}ショ \mbox{(CH\={U}{\middel}MON{\middel}SHO)}\\\hline
注意書き&attention + write = \textit{notes; instructions}&チュウ{\middel}イ{\wsplit}が{\middel}き \mbox{(CH\={U}{\middel}I{\wsplit}ga{\middel}ki)}\\\hline
要注意&essential + attention = \textit{need for caution}&ヨウ{\middel}チュウ{\middel}イ \mbox{(Y\={O}{\middel}CH\={U}{\middel}I)}\\\hline
注水&pour + water = \textit{watering}&チュウ{\middel}スイ \mbox{(CH\={U}{\middel}SUI)}\\\hline
注入&pour + enter = \textit{infusion}&チュウ{\middel}ニュウ \mbox{(CH\={U}{\middel}NY\={U})}\\\hline
注目&pour + eye = \textit{notice}&チュウ{\middel}モク \mbox{(CH\={U}{\middel}MOKU)}\\\hline
\notyet{発注}&discharge + pour = \textit{ordering}&ハッ{\middel}チュウ \mbox{(HAC{\middel}CH\={U})}\\\hline
\notyet{注文品}&order + goods = \textit{ordered goods}&チュウ{\middel}モン{\middel}ヒン \mbox{(CH\={U}{\middel}MON{\middel}HIN)}\\\hline
\end{somme}

\begin{decomp}{5}
犬&に&注意&して&下さい。\\\hline
いぬ&に&チュウイ&して&ください\\\hline
dog&に&attention&do&please\\\hline
\multicolumn{5}{|r|}{\textit{Watch out for the dog.}}\\\hline
\end{decomp}
\end{multicols}

\pagebreak
%%--------------------------------------------------------------------
%19-------------------------------------------------------------------
\begin{multicols}{2}
\begin{glyph}{Goods (品)}{goods}\label{glyph:goods}
This character expresses the idea of \textit{goods} or \textit{articles} in the majority of compounds in which it appears.\\\hline
It can occasionally have the more abstract notion of \textit{refinement}.
\end{glyph}
\gindex{Goods}{品}\strindex{9}{品}

\begin{comps}
\comprtwocone&品\\\hline
\comprthreecone&Goods\\\hline
\end{comps}

\begin{remglyph}
An uncommon 漢字 radical.\\\hline
\end{remglyph}

\begin{prons}
\pronsrtwocone&\textbf{ヒン (HIN)}&\textbf{しな (shina)}\\\hline
\pronsrthreecone&---&---\\\hline
\end{prons}
\rindex{HIN}{品}
\rindex{shina}{品}

\begin{remprons}
The \hooglig{goods} \comkun{she nudged} at was a \comon{hint}.\\\hline
\end{remprons}

%{\middel}
% &  ()\\\hline
\begin{somme}
品&\textit{goods}&しな \mbox{(shina)}\\\hline
品物&goods + thing = \textit{merchandise}&しな{\middel}もの \mbox{(shina{\middel}mono)}\\\hline
品切れ&goods + cut = \textit{``out of stock''}&しな{\wsplit}ぎ{\middel}れ \mbox{(shina{\wsplit}gi{\middel}re)}\\\hline
上品&upper + goods = \textit{elegant}&ジョウ{\middel}ヒン \mbox{(J\={O}{\middel}HIN)}\\\hline
下品&lower + goods = \textit{vulgar}&ゲ{\middel}ヒン \mbox{(GE{\middel}HIN)}\\\hline
中古品&middle + old + goods = \textit{secondhand goods}&チュウ{\middel}コ{\middel}ヒン \mbox{(CH\={U}{\middel}KO{\middel}HIN)}\\\hline
手品&hand + goods = \textit{magic illusion}&て{\middel}じな \mbox{(te{\middel}jina)}\\\hline
\notyet{作品}&make + goods = \textit{a work (book, film, \ldots)}&サク{\middel}ヒン \mbox{(SAKU{\middel}HIN)}\\\hline
\notyet{食品}&eat + goods = \textit{food products}&ショク{\middel}ヒン \mbox{(SHOKU{\middel}HIN)}\\\hline
\notyet{品質}&goods + quality = \textit{quality}&ヒン{\middel}シツ \mbox{(HIN{\middel}SHITSU)}\\\hline
\notyet{部品}&section + goods = \textit{parts}&ブ{\middel}ヒン \mbox{(BU{\middel}HIN)}\\\hline
\end{somme}

\begin{decomp}{4}
此の&品物&は&安い。\\\hline
この&しなもの&は&やすい\\\hline
this&article&は&cheap\\\hline
\multicolumn{4}{|r|}{\textit{This article is cheap.}}\\\hline
\end{decomp}
\end{multicols}
%%--------------------------------------------------------------------
%20-------------------------------------------------------------------
\begin{multicols}{2}
\begin{glyph}{Hamlet (里)}{hamlet}\label{glyph:hamlet}
Hamlet.  Here we have the fourth and smallest of our populated units.\\\hline
This character can also signify the \textit{countryside}.
\end{glyph}
\gindex{Hamlet}{里}\strindex{7}{里}

\begin{comps}
\comprtwocone&里\\\hline
\comprthreecone&Village/3.93 km/2.44 miles\\\hline
\end{comps}

\begin{remglyph}
Part of the 漢字 radicals (\#{166}).\\\hline
\end{remglyph}

\begin{prons}
\pronsrtwocone&---&\textbf{さと (sato)}\\\hline
\pronsrthreecone&リ (RI)&---\\\hline
\end{prons}
\rindex{sato}{里}
\rindex{RI}{里}

\begin{remprons}
\hooglig{Hamlet} caused a \comkun{subtle} \ucomon{leap} in my heart.\\\hline
{\warn}\textit{Hamlet} is a tragedy written by Shakespeare.\\\hline
\end{remprons}

%{\middel}
% &  ()\\\hline
\begin{somme}
里&\textit{hamlet}&さと \mbox{(sato)}\\\hline
山里&mountain + hamlet = \textit{mountain hamlet}&やま{\middel}さと \mbox{(yama{\middel}sato)}\\\hline
古里&old + hamlet = \textit{the old hometown}&ふる{\middel}さと \mbox{(furu{\middel}sato)}\\\hline
里人&hamlet + person = \textit{country folk}&さと{\middel}びと \mbox{(sato{\middel}bito)}\\\hline
里心&hamlet + heart = \textit{homesickness}&さと{\middel}ごころ \mbox{(sato{\middel}gokoro)}\\\hline
里親&hamlet + parents = \textit{foster parents}&さと{\middel}おや \mbox{(sato{\middel}oya)}\\\hline
里子&hamlet + child = \textit{foster child}&さと{\middel}ご \mbox{(sato{\middel}go)}\\\hline
\notyet{人里離れた}&human habitation + separate = \textit{lonely/remote place}&ひと{\middel}ざと{\wsplit}はな{\middel}れた \mbox{(hito{\middel}zato{\wsplit}hana{\middel}reta)}\\\hline
\notyet{里帰り}&hamlet + homeward = \textit{returning home}&さと{\wsplit}かえ{\middel}り \mbox{(sato{\wsplit}kae{\middel}ri)}\\\hline
\end{somme}

\begin{decomp}{5}
先週&は&古里&へ&行きました。\\\hline
センシュウ&は&ふるさと&へ&いきました\\\hline
last week&は&old hometown&へ&went\\\hline
\multicolumn{5}{|r|}{\textit{Last week I went to my old hometown.}}\\\hline
\end{decomp}
\end{multicols}
%%--------------------------------------------------------------------
%%--------------------------------------------------------------------
\pagebreak

%
%\begin{multicols}{2}
%\chapter{Kanji 181--200}
%%%--------------------------------------------------------------------
%1--------------------------------------------------------------------
\vspace*{\fill}
\begin{glyph}{Grove (林)}{grove}\label{glyph:grove}
Grove.
\end{glyph}

\begin{comps}
\comprtwocone &  木 + 木 \\\hline
    \comprthreecone & Tree (\ref{glyph:tree}) + Tree (\ref{glyph:tree}) \\\hline
\end{comps}

\begin{remglyph}
--- \\\hline
\end{remglyph}

\begin{prons}
\pronsrtwocone & \textbf{リン (RIN)} & \textbf{はやし (hayashi)}\\\hline
\pronsrthreecone & --- & --- \\\hline
\end{prons}

\begin{remprons}
    The \hooglig{grove} is \comon{leaning} towards \comkun{high-ash-investments}.\\\hline
\end{remprons}

%{\middel}
% &  ()\\\hline
\begin{somme}
    \rye{3}{As the common Japanese family name in the second example below
    shows, this reading becomes voiced when not in the first position.}
    林  & \textit{grove}  & はやし (hayashi)  \\\hline
    小林さん  & small + grove = \textit{Ko{\middel}bayashi-san}  & こ{\middel}ばやし{\middel}さん
    (ko{\middel}bayashi{\middel}san)  \\\hline
    林道  & grove + road = \textit{woodland trail; road}  & リン{\middel}ドウ (RIN{\middel}D\={O})  \\\hline
    人工林  & person + craft + grove = \textit{planted forest}  & ジン{\middel}コウ{\middel}リン
    (JIN{\middel}K\={O}{\middel}RIN)  \\\hline
    公有林  & public + have + grove = \textit{public woodland}  & コウ{\middel}ユウ{\middel}リン
    (K\={O}{\middel}Y\={U}{\middel}RIN)  \\\hline
\end{somme}

\begin{sentence}
    \underline{あの}{\;}\underline{\ruby{国}{くに}}に\ruby{は}{わ}\underline{\ruby{公有林}{コウユウリン}}が\underline{\ruby{少}{すく}ない}。\\\hline
    ano kuni ni wa K\={O}{\middel}Y\={U}{\middel}RIN ga suku{\middel}nai.\\\hline
    (that) (country) (public woodland) (few)\\\hline
    =\textit{That country has little public woodland.}\\\hline
\end{sentence}
\end{multicols}
\vspace*{-\baselineskip}
\vhrulefill{5pt}
%%--------------------------------------------------------------------
%2--------------------------------------------------------------------
\begin{multicols}{2}
\begin{glyph}{Forest (森)}{forest}\label{glyph:forest}
Forest.
\end{glyph}

\begin{comps}
\comprtwocone & 木 + 木 + 木 \\\hline
    \comprthreecone & Tree (\ref{glyph:tree}) + Tree (\ref{glyph:tree}) + Tree
    (\ref{glyph:tree}) \\\hline
\end{comps}

\begin{remglyph}
---\\\hline
\end{remglyph}

\begin{prons}
\pronsrtwocone & \textbf{シン (SHIN)} & \textbf{もり (mori)}\\\hline
\pronsrthreecone & --- & --- \\\hline
\end{prons}

\begin{remprons}
    A \hooglig{forest} with a \comkun{moribund} \comon{sheen}.\\\hline
\end{remprons}

%{\middel}
% &  ()\\\hline
\begin{somme}
    森  & \textit{forest}  & もり (mori)  \\\hline
    森林  & forest + grove = \textit{forests}  & シン{\middel}リン (SHIN{\middel}RIN)  \\\hline
    森林学  & forest + grove + study = \textit{forestry}  & シン{\middel}リン{\middel}ガク
    (SHIN{\middel}RIN{\middel}GAKU)  \\\hline
\end{somme}

\begin{sentence}
    \underline{\ruby{森}{もり}}に\underline{\ruby{入}{はい}る}{\;}\underline{\ruby{時}{とき}}\ruby{は}{わ}\underline{\ruby{注意}{チュウイ}}{\;}\underline{して}{\;}\underline{\ruby{下}{くだ}さい}。\\\hline
    mori ni hai{\middel}ru toki wa CH\={U}{\middel}I shite kuda{\middel}sai.\\\hline
    (forest) (enter) (time) (attention) (do please)\\\hline
    =\textit{Please be careful when going into the forest.}\\\hline
\end{sentence}
\end{multicols}
\vspace*{-\baselineskip}
\vhrulefill{5pt}
\vspace*{\fill}
\pagebreak
%%--------------------------------------------------------------------
%3--------------------------------------------------------------------
\vspace*{\fill}
\begin{multicols}{2}
\begin{glyph}{Dawn (旦)}{dawn}\label{glyph:dawn}
Dawn.\\\hline
This kanji features more often as a component in other kanji.
\end{glyph}

\begin{comps}
\comprtwocone & 日 + 一 \\\hline
    \comprthreecone & Sun (\ref{glyph:sun}) + One (\ref{glyph:one}) \\\hline
\end{comps}

\begin{remglyph}
    \hooglig{Dawn} is stage \remcom{one} of the \remcom{day}.\\\hline
\end{remglyph}

\begin{prons}
\pronsrtwocone & \textbf{タン (TAN)} & --- \\\hline
    \pronsrthreecone & ダン (DAN) & --- \\\hline
\end{prons}

\begin{remprons}
    \hooglig{Dawn} \ucomon{dunking} by the \comon{tons}.\\\hline
\end{remprons}

%{\middel}
% &  ()\\\hline
\begin{somme}
    一旦  & once + dawn = \textit{once}  & イッ{\middel}タン (IT{\middel}TAN)  \\\hline
    元旦  & basis + dawn = \textit{New Year's Day}  & ガン{\middel}タン (GAN{\middel}TAN)  \\\hline
\end{somme}

\begin{sentence}
    \underline{\ruby{元旦}{ガンタン}}に\underline{お\ruby{寺}{てら}}\ruby{へ}{え}\underline{\ruby{行}{い}きました}。\\\hline
GAN{\middel}TAN ni o{\middel}tera e i{\middel}kimashita.\\\hline
    (New Year's Day) (temple) (went)\\\hline
    =\textit{(I) went to a temple on New Year's Day.}\\\hline
\end{sentence}
\end{multicols}
\vspace*{-\baselineskip}
\vhrulefill{5pt}
%%--------------------------------------------------------------------
%4--------------------------------------------------------------------
\begin{multicols}{2}
\begin{glyph}{Black (黒)}{black}\label{glyph:black}
Black.
\end{glyph}

\begin{comps}
\comprtwocone & 里 + 灬 \\\hline
    \comprthreecone & Black (\ref{glyph:black}) + Fire (\ref{glyph:fire}) \\\hline
\end{comps}

\begin{remglyph}
    \hooglig{Black} \remcom{fire} consumed the \remcom{hamlet}.\\\hline
\end{remglyph}

\begin{prons}
\pronsrtwocone & \textbf{コク (KOKU)} & \textbf{くろ (kuro)}\\\hline
\pronsrthreecone & --- & --- \\\hline
\end{prons}

\begin{remprons}
    \hooglig{Black} \comkun{cool objects} \comon{cocked} my interests.\\\hline
\end{remprons}

%{\middel}
% &  ()\\\hline
\begin{somme}
    \rye{3}{With the exception of the third example, the kun-yomi is always
    voiced in the second position (as it appears in the fourth compound).}
    黒  & \textit{black (noun)}  & くろ (kuro)  \\\hline
    黒い  & \textit{black (adjective)}  & くろ{\middel}い (kuro{\middel}i)  \\\hline
    白黒  & white + black = \textit{black-and-white}  & しら{\middel}くろ (shira{\middel}kuro)  \\\hline
    赤黒い  & red + black = \textit{dark red}  & あか　ぐろ{\middel}い (aka guro{\middel}i)  \\\hline
    黒人  & black + person = \textit{a black person}  & コク{\middel}ジン (KOKU{\middel}JIN)  \\\hline
\end{somme}

\begin{sentence}
    \underline{この}{\;}\underline{\ruby{黒}{くろ}い}{\;}\underline{\ruby{馬}{うま}}\ruby{は}{わ}\underline{\ruby{美}{うつく}しい}{\;}\underline{ですね}。\\\hline
kono kuro{\middel}i uma wa utsuku{\middel}shii desu ne.\\\hline
    (this) (black) (horse) (beautiful) (isn't it)\\\hline
    =\textit{This black horse is beautiful, isn't it?}\\\hline
\end{sentence}
\end{multicols}
\vspace*{-\baselineskip}
\vhrulefill{5pt}
\vspace*{\fill}
\pagebreak
%%--------------------------------------------------------------------
%5--------------------------------------------------------------------
\vspace*{\fill}
\begin{multicols}{2}
\begin{glyph}{Mix (交)}{mix}\label{glyph:mix}
A general sense of things mixing in both physical and figurative ways.
    ``Associating'' with people or ``exchanging'' words are examples of the
    latter.
\end{glyph}

\begin{comps}
\comprtwocone & 亠 + 父 \\\hline
    \comprthreecone & head/above/lid + Father (\ref{glyph:father}) \\\hline
\end{comps}

\begin{remglyph}
    \hooglig{Mix} up \remcom{father's} \remcom{head}.\\\hline
\end{remglyph}

\begin{prons}
    \pronsrtwocone & \textbf{コウ (K\={O})} & \textbf{ま (ma)}\\\hline
    \pronsrthreecone & --- & まじ、か (maji, ka) \\\hline
\end{prons}

\begin{remprons}
    A \hooglig{mixup} \comon{caused} by a \ucomkun{customer} \comkun{muddying}
    \ucomkun{much information}.\\\hline
\end{remprons}

%{\middel}
% &  ()\\\hline
\begin{somme}
    交ざる (intr.)  & \textit{to mingle}  & ま{\middel}ざる (ma{\middel}zaru)  \\\hline
    交ぜる (tr.)  & \textit{to mix (something)}  & ま{\middel}ぜる (ma{\middel}zeru)  \\\hline
    外交  & outside + mix = \textit{foreign policy}  & ガイ{\middel}コウ (GAI{\middel}K\={O})  \\\hline
\end{somme}

\begin{sentence}
    \underline{あの}{\;}\underline{\ruby{国}{くに}}の\underline{\ruby{外交}{ガイコウ}}が\underline{\ruby{分}{わ}かりません}。\\\hline
    ano kuni no GAI{\middel}K\={O} ga wa{\middel}karimasen.\\\hline
    (that) (country) (foreign policy) (don't understand)\\\hline
    =\textit{(I) don't understand that country's foreign policy.}\\\hline
\end{sentence}
\end{multicols}
\vspace*{-\baselineskip}
\vhrulefill{5pt}
%%--------------------------------------------------------------------
%6--------------------------------------------------------------------
\begin{multicols}{2}
\begin{glyph}{Not (不)}{not}\label{glyph:not}
Expressing the idea of ``not'' or ``un-''.  This is an important negating prefix
in Japanese.
\end{glyph}

\begin{comps}
\comprtwocone & 一 + 亻 + 丶 \\\hline
    \comprthreecone & One (\ref{glyph:one}) + Person (\ref{glyph:person}) + dot \\\hline
\end{comps}

\begin{remglyph}
    \hooglig{Not} \remcom{one} \remcom{person} without a \remcom{dot} of sin.\\\hline
\end{remglyph}

\begin{prons}
\pronsrtwocone & \textbf{フ (FU)} & --- \\\hline
    \pronsrthreecone & ブ (BU) & --- \\\hline
\end{prons}

\begin{remprons}
    \hooglig{Not} a \comon{foofy} \ucomon{bully}.\\\hline
\end{remprons}

%{\middel}
% &  ()\\\hline
\begin{somme}
    \rye{3}{Note how this kanji can function as a prefix with complete words, as
    it does in the final example.}
    不意  & not + mind = \textit{unexpectedly}  & フ{\middel}イ (FU{\middel}I)  \\\hline
    不安  & not + ease = \textit{uneasy}  & フ{\middel}アン (FU{\middel}AN)  \\\hline
    不明  & not + bright = \textit{unclear}  & フ{\middel}メイ (FU{\middel}MEI)  \\\hline
    不死  & not + death = \textit{immortal}  & フ{\middel}シ (FU{\middel}SHI)  \\\hline
    不親切  & not + parent + cut = \textit{unkind}  & フ{\middel}シン{\middel}セツ (FU{\middel}SHIN{\middel}SETSU)  \\\hline
\end{somme}

\begin{sentence}
    \underline{\ruby{小林}{こばやし}さん}の\underline{\ruby{意図}{イト}}\ruby{は}{わ}\underline{\ruby{不明}{フメイ}}です。\\\hline
Ko{\middel}bayashi{\middel}san no I{\middel}TO wa FU{\middel}MEI desu.\\\hline
    (Kobayashi-san) (aim) (unclear)\\\hline
    =\textit{Kobayashi-san's aim is unclear.}\\\hline
\end{sentence}
\end{multicols}
\vspace*{-\baselineskip}
\vhrulefill{5pt}
\vspace*{\fill}
\pagebreak
%%--------------------------------------------------------------------
%7--------------------------------------------------------------------
\vspace*{\fill}
\begin{multicols}{2}
\begin{glyph}{After (後)}{after}\label{glyph:after}
After/Later/Behind.  These meanings apply to the ideas of both physical location
and time.
\end{glyph}

\begin{comps}
\comprtwocone & 彳 + 幺 + 夊 \\\hline
\comprthreecone & walking slowly + small/young + go \\\hline
\end{comps}

\begin{remglyph}
\hooglig{Afterwards} I \remcom{go} with the \remcom{young} children
    \remcom{walking slowly}.\\\hline
\end{remglyph}

\begin{prons}
    \pronsrtwocone & \textbf{コウ、ゴ (K\={O}, GO)} & \textbf{あと、うし (ato,
    ushi)}\\\hline
    \rye{3}{Look for コウ in the first position; あと and うしろ (this reading
    is always accompanied by ろ), appear less frequently here.\\\hline
    ゴ pops up in only a few words.  In the second position, moreover, the
    reading will invariably be ゴ.}
    \pronsrthreecone & --- & のち、おく (nochi, oku) \\\hline
\end{prons}

\begin{remprons}
    \hooglig{After} \ucomkun{notching} your way up, \comkun{you ship} the
    \comon{gobs} to \comon{cause} the \comkun{atoll} to be \ucomkun{occupied}.\\\hline
\end{remprons}

%{\middel}
% &  ()\\\hline
\begin{somme}
    \rye{3}{It is worth having another look at the compounds for ``前'' in
    Chapter 7 (Entry 133), to see how the fourth, fifth, and sixth examples
    given there ``pair up'' with the three below.}
    後  & \textit{after}  & あと (ato)  \\\hline
    後ろ  & \textit{behind}  & うし{\middel}ろ (ushi{\middel}ro)  \\\hline
    後年  & after + year = \textit{in later years}  & コウ{\middel}ネン (K\={O}{\middel}NEN)  \\\hline
    後半  & after + half = \textit{second half}  & コウ{\middel}ハン (K\={O}{\middel}HAN)  \\\hline
    後者  & after + individual = \textit{the latter}  & コウ{\middel}シャ (K\={O}{\middel}SHA)  \\\hline
    午後  & noon + after = \textit{P.M.; afternoon}  & ゴ{\middel}ゴ (GO{\middel}GO)  \\\hline
最後  & most + after = \textit{the last}  & サイ{\middel}ゴ (SAI{\middel}GO)  \\\hline
\end{somme}

\begin{sentence}
    \underline{\ruby{後}{あと}}で\underline{\ruby{友人}{ユウジン}}{\;}\underline{と}{\;}\underline{\ruby{東京}{トウキョウ}}\ruby{へ}{え}\underline{\ruby{行}{い}きます}。\\\hline
    ato de Y\={U}{\middel}JIN to T\={O}{\middel}KY\={O} e i{\middel}kimasu.\\\hline
    (after) (friend) (with) (Tokyo) (go).\\\hline
    =\textit{Afterwards, I'll go to Tokyo with a friend.}\\\hline
\end{sentence}
\end{multicols}
\vspace*{-\baselineskip}
\vhrulefill{5pt}
%%--------------------------------------------------------------------
%8--------------------------------------------------------------------
\begin{multicols}{2}
\begin{glyph}{Most (最)}{most}\label{glyph:most}
Most.\\\hline
Note how the first stroke of ``耳'' has been lengthened.
\end{glyph}

\begin{comps}
\comprtwocone & 日 + 最 \\\hline
    \comprthreecone & Sun (\ref{glyph:sun}) + Take (\ref{glyph:take}) \\\hline
\end{comps}

\begin{remglyph}
    \hooglig{Mostly} the \remcom{sun} \remcom{takes} it all.\\\hline
\end{remglyph}

\begin{prons}
\pronsrtwocone & \textbf{サイ (SAI)} & --- \\\hline
    \pronsrthreecone & --- & もっと、も (motto, mo) \\\hline
\end{prons}

\begin{remprons}
    The \hooglig{most} \comon{psychedellic} \ucomkun{moral} \ucomkun{motto}.\\\hline
\end{remprons}

%{\middel}
% &  ()\\\hline
\begin{somme}
    最小  & most + small = \textit{smallest}  & サイ{\middel}ショウ (SAI{\middel}SH\={O})  \\\hline
    最少  & most + few = \textit{fewest}  & サイ{\middel}ショウ (SAI{\middel}SH\={O})  \\\hline
    最高  & most + tall = \textit{highest; the best}  & サイ{\middel}コウ (SAI{\middel}K\={O})  \\\hline
    最古  & most + old = \textit{oldest}  & サイ{\middel}コ (SAI{\middel}KO)  \\\hline
    最新  & most + new = \textit{newest}  & サイ{\middel}シン (SAI{\middel}SHIN)  \\\hline
    最後  & most + after = \textit{the last}  & サイ{\middel}ゴ (SAI{\middel}GO)  \\\hline
\end{somme}

\begin{sentence}
    \underline{この}{\;}\underline{\ruby{天気}{テンキ}}\ruby{は}{わ}\underline{\ruby{最高}{サイコウ}}{\;}\underline{です}よ。\\\hline
    kono TEN{\middel}KI wa SAI{\middel}K\={O} desu yo.\\\hline
    (this) (weather) (the best) (is)\\\hline
    =\textit{This weather is fantastic.}\\\hline
\end{sentence}
\end{multicols}
\vspace*{-\baselineskip}
\vhrulefill{5pt}
\vspace*{\fill}
\pagebreak
%%--------------------------------------------------------------------
%9--------------------------------------------------------------------
\vspace*{\fill}
\begin{multicols}{2}
\begin{glyph}{Fat (太)}{fat}\label{glyph:fat}
    Fat/Thick.  Note the difference between this character and ``犬'' (Entry
    171).
\end{glyph}

\begin{comps}
\comprtwocone & 大 + 丶 \\\hline
    \comprthreecone & Large (\ref{glyph:large}) + dot \\\hline
\end{comps}

\begin{remglyph}
    \hooglig{Fat} \remcom{dots} are \remcom{large}.\\\hline
\end{remglyph}

\begin{prons}
\pronsrtwocone & \textbf{タイ (TAI)} & \textbf{ふと (futo)}\\\hline
    \pronsrthreecone & タ (TA) & --- \\\hline
\end{prons}

\begin{remprons}
    The \hooglig{fat} \comkun{futon} was \comon{tightly} packed in the
    \ucomon{tub}.\\\hline
\end{remprons}

%{\middel}
% &  ()\\\hline
\begin{somme}
    太い  & \textit{fat; thick}  & ふと{\middel}い (futo{\middel}i)  \\\hline
    太る  & \textit{to grow fat}  & ふと{\middel}る (futo{\middel}ru)  \\\hline
    太子  & fat + child = \textit{crown prince}  & タイ{\middel}シ (TAI{\middel}SHI)  \\\hline
\end{somme}

\begin{sentence}
    \underline{\ruby{私}{わたし}}\ruby{は}{わ}\underline{\ruby{冬休}{ふゆやす}み}に\underline{とても}{\;}\underline{\ruby{太}{ふと}った}。\\\hline
watashi wa fuyu yasu{\middel}mi ni totemo futo{\middel}tta.\\\hline
    (I) (winter vacation) (really) (grew fat)\\\hline
    =\textit{I really put on weight during the winter vacation.}\\\hline
\end{sentence}
\end{multicols}
\vspace*{-\baselineskip}
\vhrulefill{5pt}
%%--------------------------------------------------------------------
%10-------------------------------------------------------------------
\begin{multicols}{2}
\begin{glyph}{Car (車)}{car}\label{glyph:car}
Car.  As the compounds below indicate, this character can also refer to a wide
    range of other vehicles.
\end{glyph}

\begin{comps}
\comprtwocone & 車 \\\hline
    \comprthreecone & Car/Cart (\ref{glyph:car}) \\\hline
\end{comps}

\begin{remglyph}
--- \\\hline
\end{remglyph}

\begin{prons}
\pronsrtwocone & \textbf{シャ (SHA)} & \textbf{くるま (kuruma)}\\\hline
\pronsrthreecone & --- & --- \\\hline
\end{prons}

\begin{remprons}
    The \hooglig{car} was \comon{shut} to hear a \comkun{cool rumour}.\\\hline
\end{remprons}

%{\middel}
% &  ()\\\hline
\begin{somme}
    車  & \textit{car; vehicle}  & くるま (kuruma)  \\\hline
    電車  & electric + car = \textit{(electric) train}  & デン{\middel}シャ (DEN{\middel}SHA)  \\\hline
    空車  & empty + car = \textit{(taxi) ``for hire''}  & クウ{\middel}シャ (K\={U}{\middel}SHA)  \\\hline
    車内  & car + inside = \textit{inside the car}  & シャ{\middel}ナイ (SHA{\middel}NAI)  \\\hline
    中古車  & middle + old + car = \textit{used car}  & チュウ{\middel}コ{\middel}シャ
    (CH\={U}{\middel}KO{\middel}SHA)  \\\hline
\end{somme}

\begin{sentence}
    \underline{\ruby{中村}{なかむら}さん}\ruby{は}{お}\underline{\ruby{中古車}{チュウコシャ}}\ruby{を}{お}\underline{\ruby{買}{か}いました}。\\\hline
    Naka{\middel}mura-san wa CH\={U}{\middel}KO{\middel}SHA o ka{\middel}imashita.\\\hline
    (Nakamura-san) (used car) (bought)\\\hline
    =\textit{Nakamura-san bought a used car.}\\\hline
\end{sentence}
\end{multicols}
\vspace*{-\baselineskip}
\vhrulefill{5pt}
\vspace*{\fill}
\pagebreak
%%--------------------------------------------------------------------
%11-------------------------------------------------------------------
\vspace*{\fill}
\begin{multicols}{2}
\begin{glyph}{Receive (受)}{receive}\label{glyph:receive}
Receive/Accept.  This interesting character encompasses a wide range of meanings
having to do with the idea of reception: ``taking'' an exam, ``undergoing'' an
    operation, or ``suffering'' injuries are several examples.
\end{glyph}

\begin{comps}
\comprtwocone & 爪 + 冖 + 又 \\\hline
    \comprthreecone & claw + cover/roof + Again (\ref{glyph:again}) \\\hline
\end{comps}

\begin{remglyph}
    \hooglig{Receive} the \remcom{clawed} \remcom{cover} \remcom{again}.\\\hline
\end{remglyph}

\begin{prons}
\pronsrtwocone & \textbf{ジュ (JU)} & \textbf{う (u)}\\\hline
\pronsrthreecone & --- & --- \\\hline
\end{prons}

\begin{remprons}
    \hooglig{Receive} \comon{judicious} \comkun{oomph}.\\\hline
\end{remprons}

%{\middel}
% &  ()\\\hline
\begin{somme}
\rye{3}{The third and fourth examples below show this to be another kanji
    sharing the pronounciation characteristics of 買.}
    受ける  & \textit{to receive; to  accept}  & う{\middel}ける (u{\middel}keru)  \\\hline
    受け取る  & receive + take = \textit{to receive}  & う{\middel}け　と{\middel}る (u{\middel}ke to{\middel}ru)  \\\hline
    受取  & receive + take = \textit{recipient}  & うけ{\middel}とり (uke{\middel}tori)  \\\hline
    受取人  & receive + take + person = \textit{recipient}  & うけ{\middel}とり{\middel}ニン
    (uke{\middel}tori{\middel}NIN)  \\\hline
    引き受ける  & pull + receive = \textit{to undertake}  & ひ{\middel}き　う{\middel}ける (hi{\middel}ki
    u{\middel}keru)  \\\hline
    受注  & recieve + pour = \textit{receipt of an order}  & ジュ{\middel}チュウ
    (JU{\middel}CH\={U})  \\\hline
\end{somme}

\begin{sentence}
    \underline{\ruby{田中}{たなか}さん}\ruby{は}{わ}\underline{\ruby{新}{あたら}しい}{\;}\underline{\ruby{本}{ホン}}\ruby{を}{お}\underline{\ruby{受}{う}け\ruby{取}{と}りました}。\\\hline
Ta{\middel}naka-san wa atara{\middel}shii HON o u{\middel}ke to{\middel}rimashita.\\\hline
    (Tanaka-san) (new) (book) (received)\\\hline
    =\textit{Tanaka-san received a new book.}\\\hline
\end{sentence}
\end{multicols}
\vspace*{-\baselineskip}
\vhrulefill{5pt}
%%--------------------------------------------------------------------
%12-------------------------------------------------------------------
\begin{multicols}{2}
\begin{glyph}{Evil (凶)}{evil}\label{glyph:evil}
Evil/misfortune.
\end{glyph}

\begin{comps}
\comprtwocone & 丿 + 丶 + 凵  \\\hline
\comprthreecone & slash + dot + wide open mouth \\\hline
\end{comps}

\begin{remglyph}
    \hooglig{Evil} is \remcom{dotted} with \remcom{slashes}, it leaves your
    \remcom{mouth wide open}.\\\hline
\end{remglyph}

\begin{prons}
    \pronsrtwocone & \textbf{キョウ (KY\={O})} & --- \\\hline
\pronsrthreecone & ---  & --- \\\hline
\end{prons}

\begin{remprons}
    \hooglig{Evil} \comon{Kyoto}.\\\hline
\end{remprons}

%{\middel}
% &  ()\\\hline
\begin{somme}
    凶年  & evil + year = \textit{a bad year}  & キョウ{\middel}ネン (KY\={O}{\middel}NEN)  \\\hline
    凶行  & evil + go = \textit{violence}  & キョウ{\middel}コウ (KY\={O}{\middel}K\={O})  \\\hline
\end{somme}

\begin{sentence}
    \underline{\ruby{凶年}{キョウネン}}だ\underline{そう}です。\\\hline
    KY\={O}{\middel}NEN da s\={o} desu.\\\hline
    (bad year) (seems)\\\hline
    =\textit{Apparently it's a bad year.}\\\hline
\end{sentence}
\end{multicols}
\vspace*{-\baselineskip}
\vhrulefill{5pt}
\vspace*{\fill}
\pagebreak
%%--------------------------------------------------------------------
%13-------------------------------------------------------------------
\vspace*{\fill}
\begin{multicols}{2}
\begin{glyph}{Sea (海)}{sea}\label{glyph:sea}
Sea.
\end{glyph}

\begin{comps}
\comprtwocone & 氵 + 毎 \\\hline
    \comprthreecone & Water [variant] (\ref{glyph:water}) + Every (\ref{glyph:every}) \\\hline
\end{comps}

\begin{remglyph}
    The \hooglig{sea} has \remcom{water} \remcom{everywhere}.\\\hline
\end{remglyph}

\begin{prons}
\pronsrtwocone & \textbf{カイ (KAI)} & \textbf{うみ (umi)}\\\hline
\pronsrthreecone & --- & --- \\\hline
\end{prons}

\begin{remprons}
    \hooglig{Sea} \comon{kite-surfing} is \comkun{booming}.\\\hline
\end{remprons}

%{\middel}
% &  ()\\\hline
\begin{somme}
    海  & \textit{sea}  & うみ (umi)  \\\hline
    海水  & sea + water = \textit{sea water}  & カイ{\middel}スイ (KAI{\middel}SUI)  \\\hline
    海図  & sea + diagram = \textit{nautical chart}  & カイ{\middel}ズ (KAI{\middel}ZU)  \\\hline
    海王星  & sea + king + star = \textit{Neptune (planet)}  & カイ{\middel}オウ{\middel}セイ
    (KAI{\middel}\={O}{\middel}SEI)  \\\hline
    北海道  & north + sea + road = \textit{Hokkaido}  & ホッ{\middel}カイ{\middel}ドウ
    (HOK{\middel}KAI{\middel}D\={O})  \\\hline
\end{somme}

\begin{sentence}
    \underline{\ruby{夏休}{なつやす}み}に\underline{\ruby{北海道}{ホッカイドウ}}\ruby{へ}{え}\underline{\ruby{行}{い}きました}。\\\hline
    natsu yasu{\middel}mi ni HOK{\middel}KAI{\middel}D\={O} e i{\middel}kimashita.\\\hline
    (summer vacation) (Hokkaido) (went)\\\hline
    =\textit{(We) went to Hokkaido for summer vacation.}\\\hline
\end{sentence}
\end{multicols}
\vspace*{-\baselineskip}
\vhrulefill{5pt}
%%--------------------------------------------------------------------
%14-------------------------------------------------------------------
\begin{multicols}{2}
\begin{glyph}{School (校)}{school}\label{glyph:school}
School.\\\hline
A secondary meaning related to ``proof-reading'' shows up in only a few
    compounds.
\end{glyph}

\begin{comps}
\comprtwocone & 木 + 交 \\\hline
    \comprthreecone & Tree (\ref{glyph:tree}) + Mix (\ref{glyph:mix}) \\\hline
\end{comps}

\begin{remglyph}
    A \hooglig{school} is an environment where growing \remcom{trees}
    \remcom{mix}.\\\hline
\end{remglyph}

\begin{prons}
    \pronsrtwocone & \textbf{コウ (K\={O})} & --- \\\hline
\pronsrthreecone & --- & --- \\\hline
\end{prons}

\begin{remprons}
    The \hooglig{school} must have a noble \comon{cause}.\\\hline
\end{remprons}

%{\middel}
% &  ()\\\hline
\begin{somme}
    \rye{3}{Notice how ``学'' in the second example doubles up with 校, and how
    this remains constant in the other compounds in which this word occurs.}
    校外  & school + outside = \textit{off-campus}  & コウ{\middel}ガイ (K\={O}{\middel}GAI) \\\hline
    学校  & study + school = \textit{school}  & ガッ{\middel}コウ (GAK{\middel}K\={O})  \\\hline
    小学校  & small + study + school = \textit{elementary school}  &
    ショウ{\middel}ガッ{\middel}コウ (SH\={O}{\middel}GAK{\middel}K\={O})  \\\hline
    中学校  & middle + study + school = \textit{junior high school}  &
    チュウ{\middel}ガッ{\middel}コウ (CH\={U}{\middel}GAK{\middel}K\={O})  \\\hline
    高校  & tall + school = \textit{high school}  & コウ{\middel}コウ (K\={O}{\middel}K\={O})  \\\hline
    女学校  & woman + study + school = \textit{girls' school}  & ジョ{\middel}ガッ{\middel}コウ
    (JO{\middel}GAK{\middel}K\={O})  \\\hline
\end{somme}

\begin{sentence}
    \underline{あの}{\;}\underline{\ruby{女}{おんな}}の\underline{\ruby{子}{こ}}\ruby{は}{わ}\underline{\ruby{高校}{コウコウ}}{\;}\underline{\ruby{二年生}{ニネンセイ}}です。\\\hline
    ano onna no ko wa K\={O}{\middel}K\={O} NI{\middel}NEN{\middel}SEI desu.\\\hline
    (that) (woman) (child) (high school) (second year student)\\\hline
    =\textit{That girl is a second-year high school student.}\\\hline
\end{sentence}
\end{multicols}
\vspace*{-\baselineskip}
\vhrulefill{5pt}
\vspace*{\fill}
\pagebreak
%%--------------------------------------------------------------------
%15-------------------------------------------------------------------
\vspace*{\fill}
\begin{multicols}{2}
\begin{glyph}{Leg (足)}{leg}\label{glyph:leg}
Confusingly for English speakers, this character can signify both ``leg'' and
    ``foot''.\\\hline
    A secondary meaning relates to the sense of ``suffice''.
\end{glyph}

\begin{comps}
\comprtwocone & 口 + 止 \\\hline
    \comprthreecone & Mouth (\ref{glyph:mouth}) + Stop (\ref{glyph:stop}) \\\hline
\end{comps}

\begin{remglyph}
    The \hooglig{leg} \remcom{stopped} \remcom{mouthing} to me.\\\hline
\end{remglyph}

\begin{prons}
\pronsrtwocone & \textbf{ソク (SOKU)} & \textbf{あし (ashi)}\\\hline
    \pronsrthreecone & --- & た (tunnel) \\\hline
\end{prons}

\begin{remprons}
    The \hooglig{leg} was \comkun{ushered} in by the \ucomkun{tunnel}
    \comon{sock}.\\\hline
\end{remprons}

%{\middel}
% &  ()\\\hline
\begin{somme}
    \rye{3}{Note that ``取'' is voiced in the fourth example.}
    足  & \textit{leg; foot} & あし (ashi)  \\\hline
    足首  & leg + neck = \textit{ankle}  & あし{\middel}くび (ashi{\middel}kubi)  \\\hline
    足音& leg + sound = \textit{sound of footsteps}  & あし{\middel}おと (ashi{\middel}oto)  \\\hline
    足取り& leg + take = \textit{one's gait}  & あし　ど{\middel}り (ashi do{\middel}ri)  \\\hline
    土足& earth + leg = \textit{with shoes on}  & ド{\middel}ソク (DO{\middel}SOKU)  \\\hline
    足元注意 & leg + basis + pour + mind = \textit{``watch your step''}  &
    あし{\middel}もと{\middel}チュウ{\middel}イ (ashi{\middel}moto{\middel}CH\={U}{\middel}I)  \\\hline
\end{somme}

\begin{irrr}
裸足  & naked + leg = \textit{barefoot}  & はだし (hadashi) \\\hline
\end{irrr}

\begin{sentence}
    \underline{\ruby{足音}{あしおと}}\ruby{を}{お}\underline{\ruby{聞}{き}きました}か。\\\hline
ashi{\middel}oto o ki{\middel}kimashita ka.\\\hline
    (sound of footsteps) (heard)\\\hline
    =\textit{Did you hear footsteps?}\\\hline
\end{sentence}
\end{multicols}
\vspace*{-\baselineskip}
\vhrulefill{5pt}
%%--------------------------------------------------------------------
%16-------------------------------------------------------------------
\begin{multicols}{2}
\begin{glyph}{Thread (糸)}{thread}\label{glyph:thread}
Thread.  This character appears as a component in more than sixty other kanji.
\end{glyph}

\begin{comps}
\comprtwocone & 幺 + 小 \\\hline
    \comprthreecone & small/young/short thread + Small (\ref{glyph:small})  \\\hline
\end{comps}

\begin{remglyph}
    The \hooglig{thread} is \remcom{young} and \remcom{small}.\\\hline
\end{remglyph}

\begin{prons}
\pronsrtwocone & --- & \textbf{いと (ito)}\\\hline
    \pronsrthreecone & シ (SHI) & --- \\\hline
\end{prons}

\begin{remprons}
    \hooglig{Threaded} \ucomon{shields} against \comkun{e-tolling}.\\\hline
\end{remprons}

%{\middel}
% &  ()\\\hline
\begin{somme}
    \rye{3}{Note that ``口'' is voiced in the second example.}
    糸  & \textit{thread}  & いと (ito)  \\\hline
    糸口  & thread + mouth = \textit{beginning; clue}  & いとぐち (ito{\middel}guchi)  \\\hline
\end{somme}

\begin{sentence}
    \underline{\ruby{青}{あお}い}{\;}\underline{\ruby{糸}{いと}}が\underline{あります}か。\\\hline
ao{\middel}i ito ga arimasu ka.\\\hline
    (blue) (thread) (is)\\\hline
    =\textit{Is there any blue thread?}\\\hline
\end{sentence}
\end{multicols}
\vspace*{-\baselineskip}
\vhrulefill{5pt}
\vspace*{\fill}
\pagebreak
%%--------------------------------------------------------------------
%17-------------------------------------------------------------------
\vspace*{\fill}
\begin{multicols}{2}
\begin{glyph}{Same (同)}{same}\label{glyph:same}
Same.
\end{glyph}

\begin{comps}
\comprtwocone & 冂 + 一 + 口 \\\hline
    \comprthreecone & wilderness + One (\ref{glyph:one}) + Mouth
    (\ref{glyph:mouth}) \\\hline
\end{comps}

\begin{remglyph}
    The \hooglig{same} \remcom{one} \remcom{vampire} in the \remcom{wilderness}.\\\hline
\end{remglyph}

\begin{prons}
    \pronsrtwocone & \textbf{ドウ (D\={O})} & \textbf{おな (ona)}\\\hline
    \rye{3}{Remember that おな differs from おんな, a reading for ``女'' in Entry
    16.}
\pronsrthreecone & --- & --- \\\hline
\end{prons}

\begin{remprons}
    \hooglig{Same} \comkun{onanism} \comon{doorway}.\\\hline
\end{remprons}

%{\middel}
% &  ()\\\hline
\begin{somme}
    同じ  & \textit{same}  & おな{\middel}じ (ona{\middel}ji)  \\\hline
    同意  & same + mind = \textit{agreement}  & ドウ{\middel}イ (D\={O}{\middel}I)  \\\hline
    同一  & same + one = \textit{identical}  & ドウ{\middel}イツ (D\={O}{\middel}ITSU)  \\\hline
    同時  & same + time = \textit{simultaneous}  & ドウ{\middel}ジ (D\={O}{\middel}JI)  \\\hline
    同化  & same + change = \textit{assimilation}  & ドウ{\middel}カ (D\={O}{\middel}KA)  \\\hline
    同行者  & same + go + individual = \textit{traveling companion}  &
    ドウ{\middel}カ{\middel}シャ (D\={O}{\middel}K\={O}{\middel}SHA)  \\\hline
\end{somme}

\begin{sentence}
    \underline{\ruby{私}{わたし}}と\underline{\ruby{中村}{なかむら}さん}\ruby{は}{わ}\underline{\ruby{同時}{ドウジ}}に\underline{\ruby{来}{き}ました}。\\\hline
    watashi to Naka{\middel}mura-san wa D\={O}{\middel}JI ni ki{\middel}mashita.\\\hline
    (I) (Nakamura-san) (simultaneous) (came)\\\hline
    =\textit{Nakamura-san and I arrived at the same time.}\\\hline
\end{sentence}
\end{multicols}
\vspace*{-\baselineskip}
\vhrulefill{5pt}
%%--------------------------------------------------------------------
%18-------------------------------------------------------------------
\begin{multicols}{2}
\begin{glyph}{Noon (午)}{noon}\label{glyph:noon}
Noon.
\end{glyph}

\begin{comps}
\comprtwocone & ノ + 干 \\\hline
    \comprthreecone & NO/slash + Dry (\ref{glyph:dry}) \\\hline
\end{comps}

\begin{remglyph}
    \hooglig{Noon} \remcom{drying} is a \remcom{no-no}.\\\hline
\end{remglyph}

\begin{prons}
\pronsrtwocone & \textbf{ゴ (GO)} & --- \\\hline
\pronsrthreecone & --- & --- \\\hline
\end{prons}

\begin{remprons}
    \hooglig{Noon} \comon{gobbling}.\\\hline
\end{remprons}

%{\middel}
% &  ()\\\hline
\begin{somme}
    午前  & noon + before = \textit{A.M.; morning}  & ゴ{\middel}ゼン (GO{\middel}ZEN)  \\\hline
    午後  & noon + after = \textit{P.M.; afternoon}  & ゴ{\middel}ゴ (GO{\middel}GO)  \\\hline
    午前中  & noon + before + middle = \textit{all morning}  & ゴ{\middel}ゼン{\middel}チュウ
    (GO{\middel}ZEN{\middel}CH\={U})  \\\hline
\end{somme}

\begin{sentence}
    \underline{\ruby{午前}{ゴゼン}}{\;}\underline{\ruby{八時}{ハチジ}}に\underline{\ruby{高校}{コウコウ}}\ruby{へ}{え}\underline{\ruby{行}{い}きます}。\\\hline
    GO{\middel}ZEN HACHI{\middel}JI ni K\={O}{\middel}K\={O} e i{\middel}kimasu.\\\hline
    (A.M.) (eight o'clock) (high school) (go)\\\hline
    =\textit{I go to high school at 8 A.M.}\\\hline
\end{sentence}
\end{multicols}
\vspace*{-\baselineskip}
\vhrulefill{5pt}
\vspace*{\fill}
\pagebreak
%%--------------------------------------------------------------------
%19-------------------------------------------------------------------
\vspace*{\fill}
\begin{multicols}{2}
\begin{glyph}{Love (愛)}{love}\label{glyph:love}
Love.
\end{glyph}

\begin{comps}
\comprtwocone & 爪 + 冖 + 心 + 夊 \\\hline
    \comprthreecone & claw + cover/roof + Heart (\ref{glyph:heart}) + go \\\hline
\end{comps}

\begin{remglyph}
    \hooglig{Love} from the \remcom{heart} \remcom{goes} through the
    \remcom{clawed} \remcom{roof}.\\\hline
\end{remglyph}

\begin{prons}
\pronsrtwocone & \textbf{アイ (AI)} & --- \\\hline
\pronsrthreecone & --- & --- \\\hline
\end{prons}

\begin{remprons}
    \hooglig{Love} melts the \comon{ice} away.\\\hline
\end{remprons}

%{\middel}
% &  ()\\\hline
\begin{somme}
    愛好  & love + like = \textit{love; like}  & アイ{\middel}コウ (AI{\middel}K\={O})  \\\hline
    愛犬  & love + dog = \textit{pet dog}  & アイ{\middel}ケン (AI{\middel}KEN)  \\\hline
    愛犬家  & love + dog + house = \textit{dog lover}  & アイ{\middel}ケン{\middel}カ (AI{\middel}KEN{\middel}KA)  \\\hline
    愛鳥家  & love + bird + house = \textit{bird lover}  & アイ{\middel}チョウ{\middel}カ
    (AI{\middel}CH\={O}{\middel}KA)  \\\hline
    愛国心  & love + country + heart = \textit{patriotism}  & アイ{\middel}コク{\middel}シン
    (AI{\middel}KOKU{\middel}SHIN)  \\\hline
\end{somme}

\begin{irrr}
可愛い  & possible + love = \textit{cute}  & カ{\middel}ワイ{\middel}い (KA{\middel}WAI{\middel}i)  \\\hline
\end{irrr}

\begin{sentence}
    \underline{\ruby{本田}{ほんだ}さん}\ruby{は}{わ}\underline{\ruby{愛鳥家}{アイチョウカ}}{\;}\underline{として}{\;}\underline{\ruby{有名}{ユウメイ}}です。\\\hline
    Hon{\middel}da-san wa AI{\middel}CH\={O}{\middel}KA toshite Y\={U}{\middel}MEI desu.\\\hline
    (Honda-san) (bird lover) (as) (famouns)\\\hline
    =\textit{Honda-san is famous as a bird lover.}\\\hline
\end{sentence}
\end{multicols}
\vspace*{-\baselineskip}
\vhrulefill{5pt}
%%--------------------------------------------------------------------
%20-------------------------------------------------------------------
\begin{multicols}{2}
\begin{glyph}{Separate (離)}{separate}\label{glyph:separate}
Separate/Leave.
\end{glyph}

\begin{comps}
\comprtwocone & 亠 + 凶 + 冂 + 厶 + 隹 \\\hline
    \comprthreecone & head/above + Evil (\ref{glyph:evil}) + wilderness +
    self/private + small bird \\\hline
\end{comps}

\begin{remglyph}
    \hooglig{Separate} \remcom{evil} from the \remcom{heads} of \remcom{small
    birds} in the \remcom{private} \remcom{wilderness}.\\\hline
\end{remglyph}

\begin{prons}
\pronsrtwocone & \textbf{リ (RI)} & \textbf{はな (hana)}\\\hline
\pronsrthreecone & --- & --- \\\hline
\end{prons}

\begin{remprons}
    \hooglig{Separate} \comkun{hung-up} \comon{leaps}.\\\hline
\end{remprons}

%{\middel}
% &  ()\\\hline
\begin{somme}
    離れる (intr.)  & \textit{to separate; to leave}  & はな{\middel}れる (hana{\middel}reru)  \\\hline
    離す (tr.)  & \textit{to separate (something)}  & はな{\middel}す (hana{\middel}su)  \\\hline
    離れ離れ  & separate + separate = \textit{scattered; separated}  &
    はな{\middel}れ　ばな{\middel}れ (hana{\middel}re bana{\middel}re)  \\\hline
    切り離す (tr.)  & cut + separate = \textit{to cut off}  & き{\middel}り　はな{\middel}す (ki{\middel}ri
    hana{\middel}su)  \\\hline
    分離  & part + separate = \textit{separation}  & ブン{\middel}リ (BUN{\middel}RI)  \\\hline
    離村  & separate + village = \textit{rural exodus}  & リ{\middel}ソン (RI{\middel}SON)  \\\hline
\end{somme}

\begin{sentence}
    \underline{\ruby{山}{やま}}で\underline{その}{\;}\underline{\ruby{親子}{おやこ}}\ruby{は}{わ}\underline{\ruby{離}{はな}れ\ruby{離}{ばな}れ}に\underline{なってしまいました}。\\\hline
yama de sono oya{\middel}ko wa hana{\middel}re bana{\middel}re ni natte shimaimashita.\\\hline
    (mountain) (that) (parent and child) (separated) (ended up)\\\hline
    =\textit{That parent and child ended up getting separated on the
    mountain.}\\\hline
\end{sentence}
\end{multicols}
\vspace*{-\baselineskip}
\vhrulefill{5pt}
%%--------------------------------------------------------------------
%%--------------------------------------------------------------------
\vspace*{\fill}
\pagebreak

%
%\begin{multicols}{2}
%\chapter{Kanji 201--225}
%%%--------------------------------------------------------------------
%1--------------------------------------------------------------------
\vspace*{\fill}
\begin{glyph}{Trust (信)}{trust}\label{glyph:trust}
Trust.\\\hline
As can be seen in the final entry of the ``Common Words and Compounds'' section
below, a secondary meaning relates to the ideas of messages and correspondence.
\end{glyph}

\begin{comps}
\comprtwocone & 亻 + 言  \\\hline
    \comprthreecone & Person (\ref{glyph:person}) + Say (\ref{glyph:say}) \\\hline
\end{comps}

\begin{remglyph}
    \hooglig{Trust} \remcom{people's} \remcom{sayings}.\\\hline
\end{remglyph}

\begin{prons}
\pronsrtwocone & \textbf{シン (SHIN)} & --- \\\hline
\pronsrthreecone & --- & --- \\\hline
\end{prons}

\begin{remprons}
    \hooglig{Trust} has a beautiful \comon{sheen} to it.\\\hline
\end{remprons}

%{\middel}
% &  ()\\\hline
\begin{somme}
    信じる  & \textit{believe}  & シン{\middel}じる (SHIN{\middel}jiru)  \\\hline
    自信  & self + trust = \textit{self confidence}  & ジ{\middel}シン (JI{\middel}SHIN)  \\\hline
    信号  & trust + title = \textit{signal}  & シン{\middel}ゴウ (SHIN{\middel}G\={O})  \\\hline
    通信  & passage + trust = \textit{correspondence}  & ツウ{\middel}シン (TS\={U}{\middel}SHIN)  \\\hline
\end{somme}
\end{multicols}
\vspace*{-\baselineskip}
\vhrulefill{5pt}
%%--------------------------------------------------------------------
%2--------------------------------------------------------------------
\begin{multicols}{2}
\begin{glyph}{Literature (文)}{literature}\label{glyph:literature}
Literature.\\\hline
This character relates to writing and composition, from grand literary works to
sentences and individual letters.
\end{glyph}

\begin{comps}
\comprtwocone & 文 \\\hline
    \comprthreecone & Literature (\ref{glyph:literature}) \\\hline
\end{comps}

\begin{remglyph}
---\\\hline
\end{remglyph}

\begin{prons}
\pronsrtwocone & \textbf{ブン、モン (BUN, MON)} & --- \\\hline
    \pronsrthreecone & --- & ふみ (fumi) \\\hline
\end{prons}

\begin{remprons}
    The \hooglig{literature} in the \comon{monastery} is a \comon{boon} for
    \ucomkun{fumigation}.\\\hline
\end{remprons}

%{\middel}
% &  ()\\\hline
\begin{somme}
    文化  & literature + change = \textit{culture}  & ブン{\middel}カ (BUN{\middel}KA)  \\\hline
    文明  & literature + bright = \textit{civilisation}  & ブン{\middel}メイ (BUN{\middel}MEI)  \\\hline
    注文  & pour + literature = \textit{order; request}  & チュウ{\middel}モン
    (CH\={U}{\middel}MON)  \\\hline
    天文学  & heaven + literature + study = \textit{astronomy}  & テン{\middel}モン{\middel}ガク
    (TEN{\middel}MON{\middel}GAKU)  \\\hline
\end{somme}

\begin{irrr}
    文字 & literature + character = \textit{character/letter (of an alphabet)} &
    モジ (MO{\middel}JI)\\\hline
\end{irrr}
\end{multicols}
\vspace*{-\baselineskip}
\vhrulefill{5pt}
\vspace*{\fill}
\pagebreak
%%--------------------------------------------------------------------
%3--------------------------------------------------------------------
\vspace*{\fill}
\begin{multicols}{2}
\begin{glyph}{Advance (進)}{advance}\label{glyph:advance}

\end{glyph}

\begin{comps}
\comprtwocone &  \\\hline
\comprthreecone &  \\\hline
\end{comps}

\begin{remglyph}
\\\hline
\end{remglyph}

\begin{prons}
\pronsrtwocone & \textbf{ ()} & \textbf{ ()}\\\hline
\pronsrthreecone &  &  \\\hline
\end{prons}

\begin{remprons}
\\\hline
\end{remprons}

%{\middel}
% &  ()\\\hline
\begin{somme}

\end{somme}
\end{multicols}
\vspace*{-\baselineskip}
\vhrulefill{5pt}
%%--------------------------------------------------------------------
%4--------------------------------------------------------------------
\begin{multicols}{2}
\begin{glyph}{Sak\'e (酒)}{sake}\label{glyph:sake}

\end{glyph}

\begin{comps}
\comprtwocone &  \\\hline
\comprthreecone &  \\\hline
\end{comps}

\begin{remglyph}
\\\hline
\end{remglyph}

\begin{prons}
\pronsrtwocone & \textbf{ ()} & \textbf{ ()}\\\hline
\pronsrthreecone &  &  \\\hline
\end{prons}

\begin{remprons}
\\\hline
\end{remprons}

%{\middel}
% &  ()\\\hline
\begin{somme}

\end{somme}
\end{multicols}
\vspace*{-\baselineskip}
\vhrulefill{5pt}
\vspace*{\fill}
\pagebreak
%%--------------------------------------------------------------------
%5--------------------------------------------------------------------
\vspace*{\fill}
\begin{multicols}{2}
\begin{glyph}{Door (戸)}{door}\label{glyph:door}

\end{glyph}

\begin{comps}
\comprtwocone &  \\\hline
\comprthreecone &  \\\hline
\end{comps}

\begin{remglyph}
\\\hline
\end{remglyph}

\begin{prons}
\pronsrtwocone & \textbf{ ()} & \textbf{ ()}\\\hline
\pronsrthreecone &  &  \\\hline
\end{prons}

\begin{remprons}
\\\hline
\end{remprons}

%{\middel}
% &  ()\\\hline
\begin{somme}

\end{somme}
\end{multicols}
\vspace*{-\baselineskip}
\vhrulefill{5pt}
%%--------------------------------------------------------------------
%6--------------------------------------------------------------------
\begin{multicols}{2}
\begin{glyph}{Matter (事)}{matter}\label{glyph:matter}

\end{glyph}

\begin{comps}
\comprtwocone &  \\\hline
\comprthreecone &  \\\hline
\end{comps}

\begin{remglyph}
\\\hline
\end{remglyph}

\begin{prons}
\pronsrtwocone & \textbf{ ()} & \textbf{ ()}\\\hline
\pronsrthreecone &  &  \\\hline
\end{prons}

\begin{remprons}
\\\hline
\end{remprons}

%{\middel}
% &  ()\\\hline
\begin{somme}

\end{somme}
\end{multicols}
\vspace*{-\baselineskip}
\vhrulefill{5pt}
\vspace*{\fill}
\pagebreak
%%--------------------------------------------------------------------
%7--------------------------------------------------------------------
\vspace*{\fill}
\begin{multicols}{2}
\begin{glyph}{Wait (待)}{wait}\label{glyph:wait}

\end{glyph}

\begin{comps}
\comprtwocone &  \\\hline
\comprthreecone &  \\\hline
\end{comps}

\begin{remglyph}
\\\hline
\end{remglyph}

\begin{prons}
\pronsrtwocone & \textbf{ ()} & \textbf{ ()}\\\hline
\pronsrthreecone &  &  \\\hline
\end{prons}

\begin{remprons}
\\\hline
\end{remprons}

%{\middel}
% &  ()\\\hline
\begin{somme}

\end{somme}
\end{multicols}
\vspace*{-\baselineskip}
\vhrulefill{5pt}
%%--------------------------------------------------------------------
%8--------------------------------------------------------------------
\begin{multicols}{2}
\begin{glyph}{Elder brother (兄)}{elder_brother}\label{glyph:elder_brother}

\end{glyph}

\begin{comps}
\comprtwocone &  \\\hline
\comprthreecone &  \\\hline
\end{comps}

\begin{remglyph}
\\\hline
\end{remglyph}

\begin{prons}
\pronsrtwocone & \textbf{ ()} & \textbf{ ()}\\\hline
\pronsrthreecone &  &  \\\hline
\end{prons}

\begin{remprons}
\\\hline
\end{remprons}

%{\middel}
% &  ()\\\hline
\begin{somme}

\end{somme}
\end{multicols}
\vspace*{-\baselineskip}
\vhrulefill{5pt}
\vspace*{\fill}
\pagebreak
%%--------------------------------------------------------------------
%9--------------------------------------------------------------------
\vspace*{\fill}
\begin{multicols}{2}
\begin{glyph}{Prefecture (県)}{prefecture}\label{glyph:prefecture}

\end{glyph}

\begin{comps}
\comprtwocone &  \\\hline
\comprthreecone &  \\\hline
\end{comps}

\begin{remglyph}
\\\hline
\end{remglyph}

\begin{prons}
\pronsrtwocone & \textbf{ ()} & \textbf{ ()}\\\hline
\pronsrthreecone &  &  \\\hline
\end{prons}

\begin{remprons}
\\\hline
\end{remprons}

%{\middel}
% &  ()\\\hline
\begin{somme}

\end{somme}
\end{multicols}
\vspace*{-\baselineskip}
\vhrulefill{5pt}
%%--------------------------------------------------------------------
%10-------------------------------------------------------------------
\begin{multicols}{2}
\begin{glyph}{Back (背)}{back}\label{glyph:back}

\end{glyph}

\begin{comps}
\comprtwocone &  \\\hline
\comprthreecone &  \\\hline
\end{comps}

\begin{remglyph}
\\\hline
\end{remglyph}

\begin{prons}
\pronsrtwocone & \textbf{ ()} & \textbf{ ()}\\\hline
\pronsrthreecone &  &  \\\hline
\end{prons}

\begin{remprons}
\\\hline
\end{remprons}

%{\middel}
% &  ()\\\hline
\begin{somme}

\end{somme}
\end{multicols}
\vspace*{-\baselineskip}
\vhrulefill{5pt}
\vspace*{\fill}
\pagebreak
%%--------------------------------------------------------------------
%11-------------------------------------------------------------------
\vspace*{\fill}
\begin{multicols}{2}
\begin{glyph}{Clear (晴)}{clear}\label{glyph:clear}

\end{glyph}

\begin{comps}
\comprtwocone &  \\\hline
\comprthreecone &  \\\hline
\end{comps}

\begin{remglyph}
\\\hline
\end{remglyph}

\begin{prons}
\pronsrtwocone & \textbf{ ()} & \textbf{ ()}\\\hline
\pronsrthreecone &  &  \\\hline
\end{prons}

\begin{remprons}
\\\hline
\end{remprons}

%{\middel}
% &  ()\\\hline
\begin{somme}

\end{somme}
\end{multicols}
\vspace*{-\baselineskip}
\vhrulefill{5pt}
%%--------------------------------------------------------------------
%12-------------------------------------------------------------------
\begin{multicols}{2}
\begin{glyph}{Run (走)}{run}\label{glyph:run}

\end{glyph}

\begin{comps}
\comprtwocone &  \\\hline
\comprthreecone &  \\\hline
\end{comps}

\begin{remglyph}
\\\hline
\end{remglyph}

\begin{prons}
\pronsrtwocone & \textbf{ ()} & \textbf{ ()}\\\hline
\pronsrthreecone &  &  \\\hline
\end{prons}

\begin{remprons}
\\\hline
\end{remprons}

%{\middel}
% &  ()\\\hline
\begin{somme}

\end{somme}
\end{multicols}
\vspace*{-\baselineskip}
\vhrulefill{5pt}
\vspace*{\fill}
\pagebreak
%%--------------------------------------------------------------------
%13-------------------------------------------------------------------
\vspace*{\fill}
\begin{multicols}{2}
\begin{glyph}{Wash (洗)}{wash}\label{glyph:wash}

\end{glyph}

\begin{comps}
\comprtwocone &  \\\hline
\comprthreecone &  \\\hline
\end{comps}

\begin{remglyph}
\\\hline
\end{remglyph}

\begin{prons}
\pronsrtwocone & \textbf{ ()} & \textbf{ ()}\\\hline
\pronsrthreecone &  &  \\\hline
\end{prons}

\begin{remprons}
\\\hline
\end{remprons}

%{\middel}
% &  ()\\\hline
\begin{somme}

\end{somme}
\end{multicols}
\vspace*{-\baselineskip}
\vhrulefill{5pt}
%%--------------------------------------------------------------------
%14-------------------------------------------------------------------
\begin{multicols}{2}
\begin{glyph}{Oneself (己)}{oneself}\label{glyph:oneself}

\end{glyph}

\begin{comps}
\comprtwocone &  \\\hline
\comprthreecone &  \\\hline
\end{comps}

\begin{remglyph}
\\\hline
\end{remglyph}

\begin{prons}
\pronsrtwocone & \textbf{ ()} & \textbf{ ()}\\\hline
\pronsrthreecone &  &  \\\hline
\end{prons}

\begin{remprons}
\\\hline
\end{remprons}

%{\middel}
% &  ()\\\hline
\begin{somme}

\end{somme}
\end{multicols}
\vspace*{-\baselineskip}
\vhrulefill{5pt}
\vspace*{\fill}
\pagebreak
%%--------------------------------------------------------------------
%15-------------------------------------------------------------------
\vspace*{\fill}
\begin{multicols}{2}
\begin{glyph}{Roll (転)}{roll}\label{glyph:roll}

\end{glyph}

\begin{comps}
\comprtwocone &  \\\hline
\comprthreecone &  \\\hline
\end{comps}

\begin{remglyph}
\\\hline
\end{remglyph}

\begin{prons}
\pronsrtwocone & \textbf{ ()} & \textbf{ ()}\\\hline
\pronsrthreecone &  &  \\\hline
\end{prons}

\begin{remprons}
\\\hline
\end{remprons}

%{\middel}
% &  ()\\\hline
\begin{somme}

\end{somme}
\end{multicols}
\vspace*{-\baselineskip}
\vhrulefill{5pt}
%%--------------------------------------------------------------------
%16-------------------------------------------------------------------
\begin{multicols}{2}
\begin{glyph}{Music (楽)}{music}\label{glyph:music}

\end{glyph}

\begin{comps}
\comprtwocone &  \\\hline
\comprthreecone &  \\\hline
\end{comps}

\begin{remglyph}
\\\hline
\end{remglyph}

\begin{prons}
\pronsrtwocone & \textbf{ ()} & \textbf{ ()}\\\hline
\pronsrthreecone &  &  \\\hline
\end{prons}

\begin{remprons}
\\\hline
\end{remprons}

%{\middel}
% &  ()\\\hline
\begin{somme}

\end{somme}
\end{multicols}
\vspace*{-\baselineskip}
\vhrulefill{5pt}
\vspace*{\fill}
\pagebreak
%%--------------------------------------------------------------------
%17-------------------------------------------------------------------
\vspace*{\fill}
\begin{multicols}{2}
\begin{glyph}{Reason (理)}{reason}\label{glyph:reason}

\end{glyph}

\begin{comps}
\comprtwocone &  \\\hline
\comprthreecone &  \\\hline
\end{comps}

\begin{remglyph}
\\\hline
\end{remglyph}

\begin{prons}
\pronsrtwocone & \textbf{ ()} & \textbf{ ()}\\\hline
\pronsrthreecone &  &  \\\hline
\end{prons}

\begin{remprons}
\\\hline
\end{remprons}

%{\middel}
% &  ()\\\hline
\begin{somme}

\end{somme}
\end{multicols}
\vspace*{-\baselineskip}
\vhrulefill{5pt}
%%--------------------------------------------------------------------
%18-------------------------------------------------------------------
\begin{multicols}{2}
\begin{glyph}{Order (番)}{order}\label{glyph:order}

\end{glyph}

\begin{comps}
\comprtwocone &  \\\hline
\comprthreecone &  \\\hline
\end{comps}

\begin{remglyph}
\\\hline
\end{remglyph}

\begin{prons}
\pronsrtwocone & \textbf{ ()} & \textbf{ ()}\\\hline
\pronsrthreecone &  &  \\\hline
\end{prons}

\begin{remprons}
\\\hline
\end{remprons}

%{\middel}
% &  ()\\\hline
\begin{somme}

\end{somme}
\end{multicols}
\vspace*{-\baselineskip}
\vhrulefill{5pt}
\vspace*{\fill}
\pagebreak
%%--------------------------------------------------------------------
%19-------------------------------------------------------------------
\vspace*{\fill}
\begin{multicols}{2}
\begin{glyph}{Both (両)}{both}\label{glyph:both}

\end{glyph}

\begin{comps}
\comprtwocone &  \\\hline
\comprthreecone &  \\\hline
\end{comps}

\begin{remglyph}
\\\hline
\end{remglyph}

\begin{prons}
\pronsrtwocone & \textbf{ ()} & \textbf{ ()}\\\hline
\pronsrthreecone &  &  \\\hline
\end{prons}

\begin{remprons}
\\\hline
\end{remprons}

%{\middel}
% &  ()\\\hline
\begin{somme}

\end{somme}
\end{multicols}
\vspace*{-\baselineskip}
\vhrulefill{5pt}
%%--------------------------------------------------------------------
%20-------------------------------------------------------------------
\begin{multicols}{2}
\begin{glyph}{People (民)}{people}\label{glyph:people}

\end{glyph}

\begin{comps}
\comprtwocone &  \\\hline
\comprthreecone &  \\\hline
\end{comps}

\begin{remglyph}
\\\hline
\end{remglyph}

\begin{prons}
\pronsrtwocone & \textbf{ ()} & \textbf{ ()}\\\hline
\pronsrthreecone &  &  \\\hline
\end{prons}

\begin{remprons}
\\\hline
\end{remprons}

%{\middel}
% &  ()\\\hline
\begin{somme}

\end{somme}
\end{multicols}
\vspace*{-\baselineskip}
\vhrulefill{5pt}
\vspace*{\fill}
\pagebreak
%%--------------------------------------------------------------------
%21-------------------------------------------------------------------
\vspace*{\fill}
\begin{multicols}{2}
\begin{glyph}{Quality (質)}{quality}\label{glyph:quality}

\end{glyph}

\begin{comps}
\comprtwocone &  \\\hline
\comprthreecone &  \\\hline
\end{comps}

\begin{remglyph}
\\\hline
\end{remglyph}

\begin{prons}
\pronsrtwocone & \textbf{ ()} & \textbf{ ()}\\\hline
\pronsrthreecone &  &  \\\hline
\end{prons}

\begin{remprons}
\\\hline
\end{remprons}

%{\middel}
% &  ()\\\hline
\begin{somme}

\end{somme}
\end{multicols}
\vspace*{-\baselineskip}
\vhrulefill{5pt}
%%--------------------------------------------------------------------
%22-------------------------------------------------------------------
\begin{multicols}{2}
\begin{glyph}{Proper (正)}{proper}\label{glyph:proper}

\end{glyph}

\begin{comps}
\comprtwocone &  \\\hline
\comprthreecone &  \\\hline
\end{comps}

\begin{remglyph}
\\\hline
\end{remglyph}

\begin{prons}
\pronsrtwocone & \textbf{ ()} & \textbf{ ()}\\\hline
\pronsrthreecone &  &  \\\hline
\end{prons}

\begin{remprons}
\\\hline
\end{remprons}

%{\middel}
% &  ()\\\hline
\begin{somme}

\end{somme}
\end{multicols}
\vspace*{-\baselineskip}
\vhrulefill{5pt}
\vspace*{\fill}
\pagebreak
%%--------------------------------------------------------------------
%23-------------------------------------------------------------------
\vspace*{\fill}
\begin{multicols}{2}
\begin{glyph}{Direction (方)}{direction}\label{glyph:direction}

\end{glyph}

\begin{comps}
\comprtwocone &  \\\hline
\comprthreecone &  \\\hline
\end{comps}

\begin{remglyph}
\\\hline
\end{remglyph}

\begin{prons}
\pronsrtwocone & \textbf{ ()} & \textbf{ ()}\\\hline
\pronsrthreecone &  &  \\\hline
\end{prons}

\begin{remprons}
\\\hline
\end{remprons}

%{\middel}
% &  ()\\\hline
\begin{somme}

\end{somme}
\end{multicols}
\vspace*{-\baselineskip}
\vhrulefill{5pt}
%%--------------------------------------------------------------------
%24-------------------------------------------------------------------
\begin{multicols}{2}
\begin{glyph}{Voice (声)}{voice}\label{glyph:voice}

\end{glyph}

\begin{comps}
\comprtwocone &  \\\hline
\comprthreecone &  \\\hline
\end{comps}

\begin{remglyph}
\\\hline
\end{remglyph}

\begin{prons}
\pronsrtwocone & \textbf{ ()} & \textbf{ ()}\\\hline
\pronsrthreecone &  &  \\\hline
\end{prons}

\begin{remprons}
\\\hline
\end{remprons}

%{\middel}
% &  ()\\\hline
\begin{somme}

\end{somme}
\end{multicols}
\vspace*{-\baselineskip}
\vhrulefill{5pt}
\vspace*{\fill}
\pagebreak
%%--------------------------------------------------------------------
%25-------------------------------------------------------------------
\vspace*{\fill}
\begin{multicols}{2}
\begin{glyph}{Reside (住)}{reside}\label{glyph:reside}

\end{glyph}

\begin{comps}
\comprtwocone &  \\\hline
\comprthreecone &  \\\hline
\end{comps}

\begin{remglyph}
\\\hline
\end{remglyph}

\begin{prons}
\pronsrtwocone & \textbf{ ()} & \textbf{ ()}\\\hline
\pronsrthreecone &  &  \\\hline
\end{prons}

\begin{remprons}
\\\hline
\end{remprons}

%{\middel}
% &  ()\\\hline
\begin{somme}

\end{somme}
\end{multicols}
\vspace*{-\baselineskip}
\vhrulefill{5pt}
\vspace*{\fill}
\pagebreak
%%--------------------------------------------------------------------

%
\begin{multicols}{2}
\chapter{Kanji 226--250}
%%--------------------------------------------------------------------
%1--------------------------------------------------------------------
\begin{glyph}{Attach (付)}{attach}\label{glyph:attach}
Attach, in both physical and figurative senses.\\\hline
This 漢字 conveys the meanings of many English words---from \textit{setting} something on fire, to \textit{turning on} an electrical device.
\end{glyph}
\gindex{Attach}{付}\strindex{5}{付}

\begin{comps}
\comprtwocone&亻 + 寸\\\hline
\comprthreecone&Person + Tiny/Inch\\\hline
\end{comps}

\begin{remglyph}
\hooglig{Attach} the \remcom{human} to the \remcom{tiny} parasite.\\\hline
\end{remglyph}

\begin{prons}
\pronsrtwocone&\textbf{フ (FU)}&\textbf{つ (tsu)}\\\hline
\rye{3}{The reading in the second position is つ when 付 is preceded by another \mbox{漢字} that is read with its verb stem: \mbox{見付ける} (み{\wsplit}つ{\middel}ける: \textit{to find}) is one example.\\\hline
When the preceding reading does not form a verb with its 漢字, 付 will sometimes becomes voiced as づ: e.g. \mbox{気付く} (キ{\wsplit}づ{\middel}く: \textit{to notice}) as 気 does not form a standalone verb with キ.}
\pronsrthreecone&---&---\\\hline
\end{prons}

\begin{remprons}
Getting \hooglig{attached} to \comon{foosball} can stir up a \comkun{tsunami}.\\\hline
\end{remprons}

%{\middel}
% &  ()\\\hline
\begin{somme}
\splitter{付く}{intr}&\textit{to be attached}&つ{\middel}く \mbox{(tsu{\middel}ku)} \\\hline
\splitter{付ける}{tr}&\textit{to attach/put on}&つ{\middel}ける \mbox{(tsu{\middel}keru)} \\\hline
\splitter{見付ける}{tr}&see + attach = \textit{to find}&み{\wsplit}つける \mbox{(mi{\wsplit}tsu{\middel}keru)} \\\hline
\splitter{見せ付ける}{tr}&see + attach = \textit{to make a display of}&み{\middel}せ{\wsplit}つ{\middel}ける \mbox{(mi{\middel}se{\wsplit}tsu{\middel}keru)}\\\hline
\splitter{気付く}{tr}&spirit + attach = \textit{to notice (sth.)}&キ{\wsplit}づ{\middel}く \mbox{(KI{\wsplit}zu{\middel}ku)}\\\hline
\splitter{名付ける}{tr}&name + attach = \textit{to name}&な{\wsplit}づ{\middel}ける \mbox{(na{\wsplit}zu{\middel}keru)} \\\hline
\splitter{近付ける}{intr}&near + attach = \textit{to bring near}&ちか{\wsplit}づ{\middel}ける \mbox{(chika{\wsplit}zu{\middel}keru)}\\\hline
\splitter{受け付ける}{tr}&receive + attach = \textit{to accept/receive}&う{\middel}け{\wsplit}つ{\middel}ける \mbox{(u{\middel}ke{\wsplit}tsu{\middel}keru)}\\\hline
\splitter{言付ける}{tr}&say + attach = \textit{to order/direct}&いい{\wsplit}つ{\middel}ける \mbox{(\={i}{\wsplit}tsu{\middel}keru)}\\\hline
\splitter{引き付ける}{tr}&pull + attach = \textit{to draw near}&ひ{\middel}き{\wsplit}つ{\middel}ける \mbox{(hi{\middel}ki{\wsplit}tsu{\middel}keru)}\\\hline
\splitter{押し付ける}{tr}&push + attach = \textit{to press against}&お{\middel}し{\wsplit}つ{\middel}ける \mbox{(o{\middel}shi{\wsplit}tsu{\middel}keru)}\\\hline
\splitter{押さえ付ける}{tr}&push + attach = \textit{to hold down}&お{\middel}さえ{\wsplit}つ{\middel}ける \mbox{(o{\middel}sae{\wsplit}tsu{\middel}keru)}\\\hline
\splitter{売り付ける}{tr}&sell + attach = \textit{to force a sale}&う{\middel}り{\wsplit}つ{\middel}ける \mbox{(u{\middel}ri{\wsplit}tsu{\middel}keru)}\\\hline
\splitter{聞き付ける}{tr}&hear + attach = \textit{to hear; to catch the sound of}&き{\middel}き{\wsplit}つ{\middel}ける \mbox{(ki{\middel}ki{\wsplit}tsu{\middel}keru)}\\\hline
付近&attach + near = \textit{neighbourhood}&フ{\middel}キン \mbox{(FU{\middel}KIN)} \\\hline
受付&receive + attach = \textit{receipt; reception desk}&うけ{\middel}つけ \mbox{(uke{\middel}tsuke)} \\\hline
\splitterny{貸し付ける}{tr}&lend + attach = \textit{to lend/loan}&か{\middel}し{\wsplit}つ{\middel}ける \mbox{(ka{\middel}shi{\wsplit}tsu{\middel}keru)}\\\hline
\splitterny{結び付ける}{tr}&bind + attach = \textit{to combine}&むす{\middel}び{\wsplit}つ{\middel}ける \mbox{(musu{\middel}bi{\wsplit}tsu{\middel}keru)}\\\hline
\end{somme}
\end{multicols}

\begin{decomp}{6}
これ&を&受付&に&お出し&下さい。\\\hline
これ&を&うけつけ&に&おだし&ください\\\hline
this&を&reception desk&に&submit&please\\\hline
\multicolumn{6}{|r|}{\textit{Please hand this in at the front desk.}}\\\hline
\end{decomp}
\pagebreak

%%--------------------------------------------------------------------
%2--------------------------------------------------------------------

\begin{multicols}{2}
\begin{glyph}{Mourn (弔)}{mourn}\label{glyph:mourn}
Mourn.
\end{glyph}
\gindex{Mourn}{弔}\strindex{4}{弔}

\begin{comps}
\comprtwocone&弓 + 丨\\\hline
\comprthreecone&Bow + Pole\\\hline
\end{comps}

\begin{remglyph}
\hooglig{Mourn} the \remcom{bow} that turned into a \remcom{pole}.\\\hline
\end{remglyph}

\begin{prons}
\pronsrtwocone&\textbf{チョウ (CH\={O})}&---\\\hline
\pronsrthreecone&--- &とうら (tomura)\\\hline
\end{prons}

\begin{remprons}
I \hooglig{mourned} about the \comon{chores} such that \ucomkun{tomme luxury} could not cheer me up.\\\hline
\end{remprons}

%{\middel}
% &  ()\\\hline
\begin{somme}
弔意& mourn + mind = \textit{condolence}&チョウ{\middel}イ \mbox{(CH\={O}{\middel}I)}\\\hline
弔文& mourn + literature = \textit{funeral address}&チョウ{\middel}ブン \mbox{(CH\={O}{\middel}BUN)}\\\hline
\notyet{弔問}& mourn + question = \textit{sympathy call}&チョウ{\middel}モン \mbox{(CH\={O}{\middel}MON)}\\\hline
\end{somme}

\begin{decomp}{4}
弔い&は&今日&です。\\\hline
とむらい&は&きょう&です\\\hline
funeral&は&today&be/is\\\hline
\multicolumn{4}{|r|}{\textit{The funeral is today.}}\\\hline
\end{decomp}
\end{multicols}

%%--------------------------------------------------------------------
%3--------------------------------------------------------------------
\begin{multicols}{2}
\begin{glyph}{Cry (泣)}{cry}\label{glyph:cry}
Cry.
\end{glyph}
\gindex{Cry}{泣}\strindex{8}{泣}

\begin{comps}
\comprtwocone&氵 + 立\\\hline
\comprthreecone&Water + Stand\\\hline
\end{comps}

\begin{remglyph}
\hooglig{Cry} \remcom{water} as I \remcom{stand}.\\\hline
\end{remglyph}

\begin{prons}
\pronsrtwocone&---&\textbf{な (na)}\\\hline
\pronsrthreecone&キュウ (KY\={U})&---\\\hline
\end{prons}

\begin{remprons}
\hooglig{Crying} myself \comkun{numb} from all the \ucomon{`kill you'} threats.\\\hline
\end{remprons}

%{\middel}
% &  ()\\\hline
\begin{somme}
泣く&\textit{cry}&な{\middel}く \mbox{(na{\middel}ku)}\\\hline
泣き出す&cry + exit = \textit{burst out crying}&な{\middel}き{\wsplit}だ{\middel}す \mbox{(na{\middel}ki{\wsplit}da{\middel}su)}\\\hline
泣き声&cry + voice = \textit{tearful voice}&な{\middel}き{\wsplit}ごえ \mbox{(na{\middel}ki{\wsplit}goe)}\\\hline
\notyet{泣き虫}&cry + insect = \textit{crybaby}&な{\middel}き{\wsplit}むし \mbox{(na{\middel}ki{\wsplit}mushi)}\\\hline
\end{somme}

\begin{decomp}{3}
母&は&泣いていた。\\\hline
はは&は&ないていた\\\hline
mother&は&cry\\\hline
\multicolumn{3}{|r|}{\textit{My mother was crying.}}\\\hline
\end{decomp}
\end{multicols}

\pagebreak
%%--------------------------------------------------------------------
%4--------------------------------------------------------------------

\begin{multicols}{2}
\begin{glyph}{Silver (銀)}{silver}\label{glyph:silver}
Silver.\\\hline
Recall from Entry \ref{glyph:gold} (page \pageref{glyph:gold}) that when the character for gold functions as a component in other 漢字, it means `metal' in general.
\end{glyph}
\gindex{Silver}{銀}\strindex{14}{銀}

\begin{comps}
\comprtwocone & 金 + 艮\\\hline
\comprthreecone & Metal + Boundary\\\hline
\end{comps}

\begin{remglyph}
\hooglig{Silver} \remcom{metal} \remcom{boundaries}.\\\hline
\end{remglyph}

\begin{prons}
\pronsrtwocone & \textbf{ギン (GIN)} & --- \\\hline
\pronsrthreecone & --- & --- \\\hline
\end{prons}

\begin{remprons}
\hooglig{Silver} \comon{ginkgo} capsules.\\\hline
{\warn}\textit{Ginkgo}, also known as maidenhair tree, is a species of tree native to China.\\\hline
\end{remprons}

%{\middel}
% &  ()\\\hline
\begin{somme}
銀  & \textit{silver}  & ギン \mbox{(GIN)}  \\\hline
銀行  & silver + go = \textit{bank}  & ギン{\middel}コウ \mbox{(GIN{\middel}K\={O})}  \\\hline
\notyet{銀世界}  & silver + society + world = \textit{snow-covered scene}  & ギン{\middel}セ{\middel}カイ \mbox{(GIN{\middel}SE{\middel}KAI)}\\\hline
\end{somme}

\begin{decomp}{5}
金&は&銀&より&重い。\\\hline
キン&は&ギン&より&おもい\\\hline
gold&は&silver&more&heavy\\\hline
\multicolumn{5}{|r|}{\textit{Gold is heavier than silver.}}\\\hline
\end{decomp}
\end{multicols}

%%--------------------------------------------------------------------
%5--------------------------------------------------------------------
\begin{multicols}{2}
\begin{glyph}{Question (問)}{question}\label{glyph:question}
Question.  Inquiry.\\\hline
In some words (the second compound below is an example), this 漢字 can convey the sense of \textit{problem}.
\end{glyph}
\gindex{Question}{問}\strindex{11}{問}

\begin{comps}
\comprtwocone & 門 + 口 \\\hline
\comprthreecone & Gate + Mouth \\\hline
\end{comps}

\begin{remglyph}
\hooglig{Questions} have their \remcom{gateway} through the \remcom{mouth}.\\\hline
\end{remglyph}

\begin{prons}
\pronsrtwocone & \textbf{モン (MON)} & \textbf{と (to)}\\\hline
\pronsrthreecone & --- & --- \\\hline
\end{prons}

\begin{remprons}
\hooglig{Questionable} \comkun{torrents} on the \comon{monastery}.\\\hline
\end{remprons}

%{\middel}
% &  ()\\\hline
\begin{somme}
質問  & quality + question = \textit{question}  & シツ{\middel}モン \mbox{(SHITSU{\middel}MON)}  \\\hline
\notyet{問題}  & question + topic = \textit{problem; issue}  & モン{\middel}ダイ \mbox{(MON{\middel}DAI)}  \\\hline
\notyet{問い合わせる}  & question + match = \textit{inquire; apply to}  & と{\middel}い　あ{\middel}わせる \mbox{(to{\middel}i a{\middel}waseru)}  \\\hline
\end{somme}

\begin{irrr}
\notyet{問屋}  & question + roof = \textit{wholesale dealer}  & とん{\middel}や (ton{\middel}ya)  \\\hline
\end{irrr}

\begin{decomp}{5}
其れ&は&問題外&だ&よ。\\\hline
それ&は&モンダイガイ&だ&よ\\\hline
that&は&unthinkable&be/is&よ\\\hline
\multicolumn{5}{|r|}{\textit{That is out of the question.}}\\\hline
\end{decomp}
\end{multicols}

\pagebreak
%%--------------------------------------------------------------------
%6--------------------------------------------------------------------

\begin{multicols}{2}
\begin{glyph}{Dark (暗)}{dark}\label{glyph:dark}
Dark, with its various implications.
\end{glyph}
\gindex{Dark}{暗}\strindex{13}{暗}

\begin{comps}
\comprtwocone & 日 + 音 \\\hline
\comprthreecone & Sun/Day + Sound \\\hline
\end{comps}

\begin{remglyph}
\hooglig{Dark} and ominous \remcom{sun} \remcom{sounds}.\\\hline
\end{remglyph}

\begin{prons}
\pronsrtwocone & \textbf{アン (AN)} & \textbf{くら (kura)}\\\hline
\pronsrthreecone & --- & --- \\\hline
\end{prons}

\begin{remprons}
\hooglig{Dark} \comkun{curator} of an \comon{uncle}.\\\hline
\end{remprons}

%{\middel}
% &  ()\\\hline
\begin{somme}
暗い  & \textit{dark}  & くら{\middel}い \mbox{(kura{\middel}i)}  \\\hline
暗黒  & dark + black = \textit{darkness; blackness}  & アン{\middel}コク \mbox{(AN{\middel}KOKU)}  \\\hline
\notyet{暗記}  & dark + note = \textit{memorization}  & アン{\middel}キ \mbox{(AN{\middel}KI)}  \\\hline
\end{somme}

\begin{decomp}{4}
外&が&暗く&なってきた。\\\hline
そと&が&くらく&なってきた\\\hline
outside&が&dark&to grow\\\hline
\multicolumn{4}{|r|}{\textit{It's getting darker outside now.}}\\\hline
\end{decomp}
\end{multicols}

%%--------------------------------------------------------------------
%7--------------------------------------------------------------------
\begin{multicols}{2}
\begin{glyph}{Picture (絵)}{picture}\label{glyph:picture}
Picture.
\end{glyph}
\gindex{Picture}{絵}\strindex{12}{絵}

\begin{comps}
\comprtwocone & 糸 + 会 \\\hline
\comprthreecone & Thread + Meet (\ref{glyph:meet}) \\\hline
\end{comps}

\begin{remglyph}
\hooglig{Picturesque} \remcom{threads} \remcom{meet}.\\\hline
\hooglig{Picture} a place where \remcom{threads} \remcom{meet} artistically.\\\hline
\end{remglyph}

\begin{prons}
\pronsrtwocone & \textbf{エ (E)} & --- \\\hline
\pronsrthreecone & カイ (KAI) & --- \\\hline
\rye{3}{This reading appears in only one compound.  The word, however, is commonly used: 絵画 (カイ{\middel}ガ) \textit{pictures and paintings}, composed of, naturally enough, the 漢字 for `picture' and `painting' (Entry \ref{glyph:painting}, page \pageref{glyph:painting}).}
\end{prons}

\begin{remprons}
\hooglig{Picture} an \comon{egg} \ucomon{kite}.\\\hline
\end{remprons}

%{\middel}
% &  ()\\\hline
\begin{somme}
絵 & \textit{picture}  & エ \mbox{(E)}\\\hline
絵本 & picture + main (book) = \textit{picture book}  & エ{\middel}ホン \mbox{(E{\middel}HON)}\\\hline
絵画 & picture + painting = \textit{pictures and paintings} & カイ{\middel}ガ　\mbox{(KAI{\middel}GA)}\\\hline
\end{somme}

\begin{decomp}{5}
壁&の&絵&を&見て。\\\hline
かべ&の&エ&を&みて\\\hline
wall&の&painting&を&look at\\\hline
\multicolumn{5}{|r|}{\textit{Look at the picture on the wall.}}\\\hline
\end{decomp}
\end{multicols}

\pagebreak
%%--------------------------------------------------------------------
%8--------------------------------------------------------------------

\begin{multicols}{2}
\begin{glyph}{Surname (氏)}{surname}\label{glyph:surname}
Surname.\\\hline
This 漢字 can also mean \textit{Clan} in some instances, as well as convey a more official-sounding form of \textit{Mr.} when it is attached to a family name (as in the final compound below).
\end{glyph}
\gindex{Surname}{氏}\strindex{4}{氏}

\begin{comps}
\comprtwocone &  氏\\\hline
\comprthreecone & Surname/Clan\\\hline
\end{comps}

\begin{remglyph}
 Part of the 漢字 radicals (\#{83}).\\\hline
\end{remglyph}

\begin{prons}
\pronsrtwocone & \textbf{シ (SHI)} & --- \\\hline
\pronsrthreecone & --- & うじ (uji) \\\hline
\end{prons}

\begin{remprons}
I \hooglig{surnamed} my \comon{shield} according to the \ucomkun{uji} customs.\\\hline
{\warn}\textit{Uji} is a city in Ky\={o}to.\\\hline
\end{remprons}

%{\middel}
% &  ()\\\hline
\begin{somme}
氏名  & surname + name = \textit{one's full name}  & シ{\middel}メイ \mbox{(SHI{\middel}MEI)}  \\\hline
中村氏  & middle + village + surname = \textit{Mr. Nakamura}  & なか{\middel}むら{\middel}シ \mbox{(naka{\middel}mura{\middel}SHI)}  \\\hline
\end{somme}

\begin{decomp}{7}
あなた&の&氏名&は&何&です&か\\\hline
あなた&の&シメイ&は&なに&です&か\\\hline
you&の&full name&は&what&be/is&か\\\hline
\multicolumn{7}{|r|}{\textit{What's your name?}}\\\hline
\end{decomp}
\end{multicols}

%%--------------------------------------------------------------------
%9-------------------------------------------------------------------
\begin{multicols}{2}
\begin{glyph}{Young (若)}{young}\label{glyph:young}
Young.
\end{glyph}
\gindex{Young}{若}\strindex{8}{若}

\begin{comps}
\comprtwocone & 艹 + 右 \\\hline
\comprthreecone & Grass + Right\\\hline
\end{comps}

\begin{remglyph}
\hooglig{Young} \remcom{grass} on the \remcom{right}-hand side.\\\hline
\end{remglyph}

\begin{prons}
\pronsrtwocone & --- & \textbf{わか (waka)}\\\hline
\pronsrthreecone & ジャク、ニャク (JAKU, NYAKU) & --- \\\hline
\end{prons}

\begin{remprons}
\hooglig{Young} \comkun{wakaberry} enthusiasts \ucomon{jucked} away his cup from the \ucomon{`knee-yuck'} flavour.\\\hline
\end{remprons}

%{\middel}
% &  ()\\\hline
\begin{somme}
\rye{3}{Note that the second compound uses the less common pronunciation for 者.}
若い  & \textit{young}  & わか{\middel}い \mbox{(waka{\middel}i)}  \\\hline
若者  & young + individual = \textit{young person}  & わか{\middel}もの \mbox{(waka{\middel}mono)}  \\\hline
若々しい  & young + young = \textit{youthful}  &  わか{\middel}わか{\middel}しい \mbox{(waka{\middel}waka{\middel}shii)} \\\hline
\end{somme}

\begin{irrr}
若人  & young + person = \textit{young person}  & わこう{\middel}ど \mbox{(wak\={o}{\middel}do)} \\\hline
\end{irrr}

\begin{decomp}{3}
僕&は&若い。\\\hline
ボク&は&わかい\\\hline
I&は&young\\\hline
\multicolumn{3}{|r|}{\textit{I'm young.}}\\\hline
\end{decomp}
\end{multicols}

\begin{decomp}{7}
此れ&が&若者&特有&の&欠点&です。\\\hline
これ&が&わかもの&トクユウ&の&ケッテン&です\\\hline
this&が&youth&peculiar to&の&defect&be/is\\\hline
\multicolumn{7}{|r|}{\textit{This is a weakness peculiar to young people.}}\\\hline
\end{decomp}

\pagebreak
%%--------------------------------------------------------------------
%10-------------------------------------------------------------------

\begin{multicols}{2}
\begin{glyph}{Cause (由)}{cause}\label{glyph:cause}
Cause.  Reason.\\\hline
This is written in the same order as the 漢字 for ``Rice Field'' (Entry \ref{glyph:rice_field}, page \pageref{glyph:rice_field}), but with the third stroke lengthened.
\end{glyph}
\gindex{Cause}{由}\strindex{5}{由}

\begin{comps}
\comprtwocone & 田 \\\hline
\comprthreecone & Rice Field \\\hline
\end{comps}

\begin{remglyph}
What \hooglig{caused} the \remcom{sun} to get \remcom{polarized}?\\\hline
\end{remglyph}

\begin{prons}
\pronsrtwocone & \textbf{ユ、ユウ (YU, Y\={U})} & --- \\\hline
\pronsrthreecone & ユイ (YUI) & よし (yoshi) \\\hline
\end{prons}

\begin{remprons}
\hooglig{`Cause} \ucomon{you insulted} the \comon{youthful} \ucomkun{yoshi} on \comon{Youtube}.\\\hline
{\warn}\textit{Yoshi} is an アニメ　character.\\\hline
\end{remprons}

%{\middel}
% &  ()\\\hline
\begin{somme}
\rye{3}{This character only appears in about 10 Japanese compounds, which makes the large number of readings for it seem extreme.  The examples below, however, are all very common.}
由来  & cause + come = \textit{origin}  & ユ{\middel}ライ \mbox{(YU{\middel}RAI)}\\\hline
自由  & self + cause = \textit{freedom}  & ジ{\middel}ユウ \mbox{(JI{\middel}Y\={U})}\\\hline
理由  & reason + cause = \textit{reason}  & リ{\middel}ユウ \mbox{(RI{\middel}Y\={U})}\\\hline
\end{somme}
\end{multicols}

\begin{decomp}{7}
多く&の&英単語&は&ラテン語&に&由来する。\\\hline
おおく&の&エイタンゴ&は&ラテンゴ&に&ユライする\\\hline
many&の&English words&は&Latin&に&origin\\\hline
\multicolumn{7}{|r|}{\textit{Many English words derive from Latin.}}\\\hline
\end{decomp}

%%--------------------------------------------------------------------
%11-------------------------------------------------------------------
\begin{multicols}{2}
\begin{glyph}{Declare (申)}{declare}\label{glyph:declare}
\textit{Declare}, in the sense of a more humble way to say the verb \textit{say}.\\\hline
Ensure that you are able to differentiate between this \mbox{漢字} and the preceding
entry for 由, as well as between these two and the character for ``Shell''
(甲), from Entry \ref{glyph:shell} (page \pageref{glyph:shell}).
\end{glyph}
\gindex{Declare}{申}\strindex{5}{申}

\begin{comps}
\comprtwocone & 田 \\\hline
\comprthreecone & Rice Field \\\hline
\end{comps}

\begin{remglyph}
\hooglig{Declare} the \remcom{sun's} sovereignty as \remcom{poles} of sunlight impale your skin.\\\hline
\end{remglyph}

\begin{prons}
\pronsrtwocone & \textbf{シン (SHIN)} & \textbf{もう (m\={o})}\\\hline
\pronsrthreecone & --- & --- \\\hline
\end{prons}

\begin{remprons}
\hooglig{Declare} \comon{sheen} \comkun{morphing}.\\\hline
\end{remprons}

%{\middel}
% &  ()\\\hline
\begin{somme}
申す& \textit{say (humble form)}  & もう{\middel}す \mbox{(m\={o}{\middel}su)}\\\hline
申し上げる& declare + upper = \textit{inform; tell (humble form)}  & もう{\middel}し　あ{\middel}げる \mbox{(m\={o}{\middel}shi a{\middel}geru)}\\\hline
内申& inside + declare = \textit{confidential report}  & ナイ{\middel}シン \mbox{(NAI{\middel}SHIN)}\\\hline
\end{somme}

\begin{decomp}{5}
私&は&周&と&申します。\\\hline
わたし&は&シュウ&と&もうします\\\hline
I&は&Sh\={u}&と&have the honour to\\\hline
\multicolumn{5}{|r|}{\textit{My surname is Zhou.}}\\\hline
\end{decomp}
\end{multicols}

\pagebreak
%%--------------------------------------------------------------------
%12-------------------------------------------------------------------

\begin{multicols}{2}
\begin{glyph}{Essential (要)}{essential}\label{glyph:essential}
Essential.  Necessary.
\end{glyph}
\gindex{Essential}{要}\strindex{9}{要}

\begin{comps}
\comprtwocone & 襾覀 + 女\\\hline
\comprthreecone & West + Woman\\\hline
\end{comps}

\begin{remglyph}
\hooglig{Essential} \remcom{western} products for \remcom{women}.\\\hline
\end{remglyph}

\begin{prons}
\pronsrtwocone & \textbf{ヨウ (Y\={O})} & \textbf{い (i)}\\\hline
\pronsrthreecone & --- & --- \\\hline
\end{prons}

\begin{remprons}
\hooglig{Essentially} an \comkun{interesting} \comon{yawn}.\\\hline
\end{remprons}

%{\middel}
% &  ()\\\hline
\begin{somme}
要る  & \textit{need; be needed}  & い{\middel}る \mbox{(i{\middel}ru)}  \\\hline
{\diff}要する  & \textit{require; need}  & ヨウ{\middel}する \mbox{(Y\={O}{\middel}suru)} \\\hline
\notyet{必要}  & certain + essential = \textit{essential; necessary}  & ヒツ{\middel}ヨウ \mbox{(HITSU{\middel}Y\={O})}  \\\hline
\notyet{重要}  & heavy + essential = \textit{importance}  & ジュウ{\middel}ヨウ \mbox{(J\={U}{\middel}Y\={O})}  \\\hline
\end{somme}

\begin{decomp}{7}
その&家&は&大&修理&を&要する。\\\hline
その&いえ&は&ダイ&シュウリ&を&ヨウする\\\hline
that&house&は&big&repairs&を&require\\\hline
\multicolumn{7}{|r|}{\textit{This house requires repairs.}}\\\hline
\end{decomp}

\begin{decomp}{5}
もっと&時間&が&必要&だ。\\\hline
もっと&ジカン&が&ヒツヨウ&だ\\\hline
more&time&が&necessary&be/is\\\hline
\multicolumn{5}{|r|}{\textit{I need more time}}\\\hline
\end{decomp}
\end{multicols}

%%--------------------------------------------------------------------
%13-------------------------------------------------------------------
\begin{multicols}{2}
\begin{glyph}{Reach (至)}{reach}\label{glyph:reach}
\textit{Reach}, in the sense of \textit{arrive at} or \textit{extend to}.
\end{glyph}
\gindex{Reach}{至}\strindex{6}{至}

\begin{comps}
\comprtwocone & 至\\\hline
\comprthreecone & Arrive\\\hline
\end{comps}

\begin{remglyph}
Part of the 漢字 radicals (\#{133}).\\\hline
\end{remglyph}

\begin{prons}
\pronsrtwocone & \textbf{シ (SHI)} & \textbf{いた (ita)}\\\hline
\rye{3}{The second compound below is the only instance in which the 音読み becomes voiced.}
\pronsrthreecone & --- & --- \\\hline
\end{prons}

\begin{remprons}
\hooglig{Reach} for the \comkun{italic} \comon{shield}.\\\hline
\end{remprons}

%{\middel}
% &  ()\\\hline
\begin{somme}
至る & \textit{reach; extend}  & いた{\middel}る \mbox{(ita{\middel}ru)}\\\hline
冬至 & winter + reach = \textit{winter solstice}  & トウ{\middel}ジ \mbox{(T\={O}{\middel}JI)}\\\hline
\notyet{至る所}  & reach + location = \textit{everywhere}  & いた{\middel}る　ところ \mbox{(ita{\middel}ru tokoro)}\\\hline
\notyet{至急} & reach + hasten = \textit{urgent}  & シ{\middel}キュウ \mbox{(SHI{\middel}KY\={U})}\\\hline
\end{somme}

\begin{decomp}{6}
この&道&は&公園&に&至る。\\\hline
この&みち&は&コウエン&に&いたる\\\hline
this&road&は&public park&に&arrive\\\hline
\multicolumn{6}{|r|}{\textit{This road goes to the park.}}\\\hline
\end{decomp}
\end{multicols}

\begin{decomp}{6}
英語&は&世界中&至る所&で&使われている。\\\hline
エイゴ&は&セカイジュウ&いたる　ところ&で&つかわれている\\\hline
English&は&throughout the world&everywhere&で&to use\\\hline
\multicolumn{6}{|r|}{\textit{English is used in every part of the world.}}\\\hline
\end{decomp}

\pagebreak
%%--------------------------------------------------------------------
%14-------------------------------------------------------------------

\begin{multicols}{2}
\begin{glyph}{Compare (比)}{compare}\label{glyph:compare}
Compare.\\\hline
The sense of \textit{ratio} is also present in a few compounds.
\end{glyph}
\gindex{Compare}{比}\strindex{4}{比}

\begin{comps}
\comprtwocone & 比 \\\hline
\comprthreecone & Compare/Side-by-side \\\hline
\end{comps}

\begin{remglyph}
Part of the 漢字 radicals (\#{81}).\\\hline
\end{remglyph}

\begin{prons}
\pronsrtwocone & \textbf{ヒ (HI)} & \textbf{くら (kura)}\\\hline
\pronsrthreecone & --- & ---\\\hline
\end{prons}

\begin{remprons}
\hooglig{Compare} the \comon{healing} \comkun{curators}.\\\hline
\end{remprons}

%{\middel}
% &  ()\\\hline
\begin{somme}
比べる& \textit{compare}  & くら{\middel}べる \mbox{(kura{\middel}beru)}  \\\hline
背比べ& back + compare = \textit{compare heights}  & せい　くら{\middel}べ \mbox{(sei kura{\middel}be)}  \\\hline
見比べる& see + compare = \textit{comparison}  & み　くら{\middel}べる \mbox{(mi kura{\middel}beru)}  \\\hline
\notyet{対比}& versus + compare = \textit{contrast; compare}  & タイ{\middel}ヒ \mbox{(TAI{\middel}HI)}  \\\hline
\end{somme}
\end{multicols}

\begin{decomp}{6}
僕&と&兄&を&比べないで&下さい。\\\hline
ボク&と&あに&を&くらべないで&ください\\\hline
me&と&brother&を&don't compare&please\\\hline
\multicolumn{6}{|r|}{\textit{Please don't compare me with my brother.}}\\\hline
\end{decomp}

%%--------------------------------------------------------------------
%15-------------------------------------------------------------------
\begin{multicols}{2}
\begin{glyph}{Measure (計)}{measure}\label{glyph:measure}
Measure.  Estimate.
\end{glyph}
\gindex{Measure}{計}\strindex{9}{計}

\begin{comps}
\comprtwocone & 言 + 十 \\\hline
\comprthreecone & Say + Ten \\\hline
\end{comps}

\begin{remglyph}
\hooglig{Measure} up to the \remcom{speech} of the \remcom{scarecrow}.\\\hline
\end{remglyph}

\begin{prons}
\pronsrtwocone & \textbf{ケイ (KEI)} & \textbf{はか (haka)}\\\hline
\pronsrthreecone & --- & --- \\\hline
\end{prons}

\begin{remprons}
\hooglig{Measure} up to the \comkun{HK's} \comon{catering} requests.\\\hline
\end{remprons}

%{\middel}
% &  ()\\\hline
\begin{somme}
\rye{3}{Don't forget that the reading of 時 in the third compound is irregular.}
計る  & \textit{measure; estimate}  & はか{\middel}る \mbox{(haka{\middel}ru)}\\\hline
会計  & meet + measure = \textit{accounts; accounting}  & カイ{\middel}ケイ \mbox{(KAI{\middel}KEI)}\\\hline
時計  & time + measure = \textit{watch; clock}  & と{\middel}ケイ \mbox{(to{\middel}KEI)}\\\hline
\notyet{温度計}  & warm + degree + measure = \textit{thermometer}  & オン{\middel}ド{\middel}ケイ \mbox{(ON{\middel}DO{\middel}KEI)}  \\\hline
\end{somme}

\begin{decomp}{6}
僕&の&時計&は&どこ&だ。\\\hline
ボク&の&とケイ&は&どこ&だ\\\hline
I&の&watch&は&where&be/is\\\hline
\multicolumn{6}{|r|}{\textit{Where is my watch?}}\\\hline
\end{decomp}

\begin{decomp}{3}
血圧&を&計りましょう。\\\hline
ケツアツ&を&はかりましょう\\\hline
blood pressure&を&to measure\\\hline
\multicolumn{3}{|r|}{\textit{Let me take your blood pressure.}}\\\hline
\end{decomp}
\end{multicols}

\pagebreak
%%--------------------------------------------------------------------
%16-------------------------------------------------------------------
\begin{multicols}{2}
\begin{glyph}{Location (所)}{location}\label{glyph:location}
Location.
\end{glyph}
\gindex{Location}{所}\strindex{8}{所}

\begin{comps}
\comprtwocone & 戸 + 斤\\\hline
\comprthreecone & Door + Axe\\\hline
\end{comps}

\begin{remglyph}
\hooglig{Locate} the \remcom{axed} \remcom{door}.\\\hline
\hooglig{Locate} the \remcom{door} with an \remcom{axe} on it.\\\hline
\end{remglyph}

\begin{prons}
\pronsrtwocone & \textbf{ショ (SHO)} & \textbf{ところ (tokoro)}\\\hline
\rye{3}{In the first position, ショ is by far the more common reading.\\\hline
In the second position, the 訓読み is always voiced, with the exceptions of only a few words in which this 漢字 is preceded by a verb in its る form (such as 至る所, the third example in Entry \ref{glyph:reach}, page \pageref{glyph:reach}).\\\hline
The 音読み, by contrast is almost never voiced in second position; the only time you are likely to encounter it voiced here is in the second example below.\\\hline
In the third position and beyond, however, there is no way to predict on first encountering a compound whether or not the 音読み is voiced (the 訓読み never appears after the second position).}
\pronsrthreecone & --- & --- \\\hline
\end{prons}

\begin{remprons}
The \hooglig{location} of the \comkun{tokoloshe} \comon{shop}.\\\hline
{\warn}In African folklore, a \textit{tokoloshe} is a mischievous and lascivious hairy water sprite.\\\hline
\end{remprons}

%{\middel}
% &  ()\\\hline
\begin{somme}
所 & \textit{location; place}  & ところ \mbox{(tokoro)}  \\\hline
近所 & near + location = \textit{neighbourhood; vicinity}  & キン{\middel}ジョ \mbox{(KIN{\middel}JO)}  \\\hline
\notyet{台所} & platform + location = \textit{kitchen}  & ダイ{\middel}どころ \mbox{(DAI{\middel}dokoro)}  \\\hline
\notyet{場所} & place + location = \textit{place}  & ば{\middel}ショ \mbox{(ba{\middel}SHO)}  \\\hline
\end{somme}

\begin{decomp}{5}
台所&に&猫&が&いる。\\\hline
ダイどころ&に&ねこ&が&いる\\\hline
kitchen&に&cat&が&to be\\\hline
\multicolumn{5}{|r|}{\textit{There is a cat in the kitchen.}}\\\hline
\end{decomp}

\begin{decomp}{4}
場所&を&教えて&下さい。\\\hline
ばショ&を&おしえて&ください\\\hline
place&を&to inform&please\\\hline
\multicolumn{4}{|r|}{\textit{Please tell me your location.}}\\\hline
\end{decomp}
\end{multicols}
%%--------------------------------------------------------------------
%17-------------------------------------------------------------------
\begin{multicols}{2}
\begin{glyph}{No. (第)}{number}\label{glyph:number}
This 漢字 is usually used as a prefix for numbers, much as English uses \textit{No.}, or \textit{\#}.\\\hline
In a few compounds it can also mean \textit{exam}.
\end{glyph}
\gindex{Number}{第}\strindex{11}{第}

\begin{comps}
\comprtwocone & 竹 + 弔 + ノ \\\hline
\comprthreecone & Bamboo + Mourn (\ref{glyph:mourn}) + NO/slash \\\hline
\end{comps}

\begin{remglyph}
\hooglig{Numerous} people \remcom{mourned} the \remcom{slashed} \remcom{bamboo}.\\\hline
\end{remglyph}

\begin{prons}
\pronsrtwocone & \textbf{ダイ (DAI)} & --- \\\hline
\pronsrthreecone & --- & --- \\\hline
\end{prons}

\begin{remprons}
\hooglig{Numerous} \comon{dyes}\ldots\\\hline
\end{remprons}

%{\middel}
% &  ()\\\hline
\begin{somme}
\rye{3}{Recall that `ease' + `complete' in the second compound below gives us \textit{safety}.}
第一  & number + one = \textit{No. 1}  & ダイ{\middel}イチ \mbox{(DAI{\middel}ICHI)}\\\hline
安全第一  & safety + number + one = \textit{``safety first''}  & アン{\middel}ゼン{\middel}ダイ{\middel}イチ \mbox{(AN{\middel}ZEN{\middel}DAI{\middel}ICHI)}\\\hline
\notyet{次第}  & next + Number = \textit{order; circumstances}  & シ{\middel}ダイ \mbox{(SHI{\middel}DAI)}\\\hline
\end{somme}

\begin{decomp}{8}
日本&は&天下&第一&の&国&が&ある。\\\hline
ニホン&は&テンカ&ダイイチ&の&くに&が&ある\\\hline
Japan&は&whole world&best&の&country&が&ある\\\hline
\multicolumn{8}{|r|}{\textit{Japan is the best country under the sun.}}\\\hline
\end{decomp}
\end{multicols}
\pagebreak
%%--------------------------------------------------------------------
%18-------------------------------------------------------------------

\begin{multicols}{2}
\begin{glyph}{Gist (旨)}{gist}\label{glyph:gist}
Gist.\\\hline
Though rarely encountered on its own, this 漢字 figures in other more common characters.
\end{glyph}
\gindex{Gist}{旨}\strindex{6}{旨}

\begin{comps}
\comprtwocone & ヒ + 日 \\\hline
\comprthreecone & Spoon/Ladle + Sun \\\hline
\end{comps}

\begin{remglyph}
The \hooglig{gist} of the \remcom{spoon} is as clear as \remcom{day}.\\\hline
\end{remglyph}

\begin{prons}
\pronsrtwocone & --- & \textbf{むね (mune)}\\\hline
\pronsrthreecone & シ (SHI) & --- \\\hline
\end{prons}

\begin{remprons}
The \hooglig{gist} is that the \comkun{moon ebbs} away to \ucomon{shield} itself from the sun.\\\hline
\end{remprons}

%{\middel}
% &  ()\\\hline
\begin{somme}
旨&\textit{gist; meaning}&むね \mbox{(mune)}  \\\hline
\end{somme}

\begin{decomp}{5}
あなた&の&主旨&は&分かる。\\\hline
あなた&の&シュシ&は&わかる\\\hline
you&の&point&は&understand\\\hline
\multicolumn{5}{|r|}{\textit{I see your point.}}\\\hline
\end{decomp}
\end{multicols}

%%--------------------------------------------------------------------
%19-------------------------------------------------------------------
\begin{multicols}{2}
\begin{glyph}{Wear (着)}{wear}\label{glyph:wear}
This important 漢字 has two main meanings:  \textit{to wear}, and \textit{to arrive}.\\\hline
It may help to think that, in a figurative sense, a person \textit{wears} their surroundings when they arrive at a place.\\\hline
This 漢字, in fact, consistently has this idea of \textit{things coming together} in certain ways.
\end{glyph}
\gindex{Wear}{着}\strindex{12}{着}

\begin{comps}
\comprtwocone & 羊 + 目 + 丿 \\\hline
\comprthreecone & Sheep + Eye + Slash/NO\\\hline
\end{comps}

\begin{remglyph}
\hooglig{Arrive} at the \remcom{grass} \remcom{king} and \remcom{slash} his \remcom{eye}.\\\hline
\hooglig{Arrive} with \remcom{sheep} \remcom{eyes} that are \remcom{not} blemished.\\\hline
\end{remglyph}

\begin{prons}
\pronsrtwocone & \textbf{チャク (CHAKU)} & \textbf{き、つ (ki, tsu)}\\\hline
\rye{3}{As the famous word in the third example makes clear, き is used when this 漢字 refers to \textit{wearing} clothes; it is always voiced in the second and third position.\\\hline
つ is used in the sense of \textit{arrive}, and is never voiced in the second position.\\\hline
チャク is the most common of these readings and has the overarching meaning of \textit{things coming together}.}
\pronsrthreecone & ジャク (JAKU) & --- \\\hline
\end{prons}

\begin{remprons}
\hooglig{Weary} \comon{Chuck} \comkun{keeps} \ucomon{jucking} the \comkun{tsunami} away.\\\hline
\end{remprons}

%{\middel}
% &  ()\\\hline
\begin{somme}
着る  & \textit{wear}  & き{\middel}る \mbox{(ki{\middel}ru)}  \\\hline
着く  & \textit{arrive}  & つ{\middel}く \mbox{(tsu{\middel}ku)}  \\\hline
着物  & wear + thing = \textit{kimono}  & き{\middel}もの \mbox{(ki{\middel}mono)}  \\\hline
\notyet{接着}  & join + come together = \textit{adhesion}  & セッ{\middel}チャク \mbox{(SET{\middel}CHAKU)}  \\\hline
\end{somme}

\begin{decomp}{5}
私&は&駅&に&着いた。\\\hline
わたし&は&エキ&に&ついた\\\hline
I&は&station&に&arrive\\\hline
\multicolumn{5}{|r|}{\textit{I arrived at the station.}}\\\hline
\end{decomp}

\begin{decomp}{5}
何&と&素晴らしい&着物&でしょう。\\\hline
なに&と&すばらしい&きもの&でしょう\\\hline
what&と&wonderful&dress&でしょう\\\hline
\multicolumn{5}{|r|}{\textit{What a heavenly dress!}}\\\hline
\end{decomp}
\end{multicols}

\pagebreak
%%--------------------------------------------------------------------
%20-------------------------------------------------------------------

\begin{multicols}{2}
\begin{glyph}{Lack (欠)}{lack}\label{glyph:lack}
Lack.
\end{glyph}
\gindex{Lack}{欠}\strindex{4}{欠}

\begin{comps}
\comprtwocone & 欠 \\\hline
\comprthreecone & Gap/Lack/Yawn \\\hline
\end{comps}

\begin{remglyph}
Part of the 漢字 radicals (\#{76}).\\\hline
\end{remglyph}

\begin{prons}
\pronsrtwocone & \textbf{ケツ (KETSU)} & \textbf{か (ka)}\\\hline
\pronsrthreecone & --- & --- \\\hline
\end{prons}

\begin{remprons}
A \hooglig{lack} of \comon{kets} \comkun{customs}.\\\hline
{\warn}\textit{Ket} people represent an ethnic group in Siberia.\\\hline
\end{remprons}

%{\middel}
% &  ()\\\hline
\begin{somme}
欠く& \textit{lack}  & か{\middel}く \mbox{(ka{\middel}ku)}\\\hline
\notyet{欠点}& lack + point = \textit{defect; shortcoming}  & ケッ{\middel}テン \mbox{(KET{\middel}TEN)}\\\hline
\notyet{欠員}& lack + member = \textit{vacancy; post}  & ケツイン \mbox{(KETSU{\middel}IN)}\\\hline
\end{somme}

\begin{decomp}{5}
欠点&なき&人&は&なし。\\\hline
ケッテン&なき&ひと&は&なし\\\hline
fault&non existent&person&は&non existent\\\hline
\multicolumn{5}{|r|}{\textit{No man is without his faults.}}\\\hline
\end{decomp}

\begin{decomp}{6}
他人&の&欠点&を&探す&な。\\\hline
タニン&の&ケッテン&を&さがす&な\\\hline
others&の&fault&を&to find&な\\\hline
\multicolumn{6}{|r|}{\textit{Don't find fault with others.}}\\\hline
\end{decomp}
\end{multicols}

%%--------------------------------------------------------------------
%21-------------------------------------------------------------------
\begin{multicols}{2}
\begin{glyph}{Show (示)}{show}\label{glyph:show}
To show.
\end{glyph}
\gindex{Show}{示}\strindex{5}{示}

\begin{comps}
\comprtwocone & 示 \\\hline
\comprthreecone & Altar/Display \\\hline
\end{comps}

\begin{remglyph}
Part of the 漢字 radicals (\#{113}).\\\hline
\end{remglyph}

\begin{prons}
\pronsrtwocone & \textbf{ジ (JI)} & \textbf{しめ (shime)}\\\hline
\pronsrthreecone & シ (SHI)  & ---\\\hline
\end{prons}

\begin{remprons}
The \hooglig{show} \comkun{she met} was \comon{jigged} with \ucomon{shields}.\\\hline
\end{remprons}

%{\middel}
% &  ()\\\hline
\begin{somme}
示す& \textit{show}  & しめ{\middel}す \mbox{(shime{\middel}su)}\\\hline
暗示& dark + show = \textit{hint}  & アン{\middel}ジ \mbox{(AN{\middel}JI)}\\\hline
\notyet{指示}& finger + show = \textit{indication; direction}  & シ{\middel}ジ \mbox{(SHI{\middel}JI)}\\\hline
\notyet{展示}& display + show = \textit{display; exhibition}  & テン{\middel}ジ \mbox{(TEN{\middel}JI)}\\\hline
\end{somme}

\begin{decomp}{7}
雪&は&冬&の&到来&を&示す。\\\hline
ゆき&は&ふゆ&の&トウライ&を&しめす\\\hline
snow&は&winter&の&coming&を&show\\\hline
\multicolumn{7}{|r|}{\textit{Snow indicates the coming of winter.}}\\\hline
\end{decomp}
\end{multicols}

\begin{decomp}{9}
彼&の&言葉&は&何&を&暗示している&の&か。\\\hline
かれ&の&ことば&は&なに&を&アンジしている&の&か\\\hline
he&の&word&は&what&を&hint&の&か\\\hline
\multicolumn{9}{|r|}{\textit{What do his words imply?}}\\\hline
\end{decomp}

\pagebreak
%%--------------------------------------------------------------------
%22-------------------------------------------------------------------

\begin{multicols}{2}
\begin{glyph}{Serve (仕)}{serve}\label{glyph:serve}
Serve.  Work.
\end{glyph}
\gindex{Serve}{仕}\strindex{5}{仕}

\begin{comps}
\comprtwocone & 亻 + 士 \\\hline
\comprthreecone & Person + Gentleman \\\hline
\end{comps}

\begin{remglyph}
\hooglig{Serve} that \remcom{person} as if he is a \remcom{gentlemen}.\\\hline
\end{remglyph}

\begin{prons}
\pronsrtwocone & \textbf{シ (SHI)} & --- \\\hline
\pronsrthreecone & ジ (JI) & つか (tsuka) \\\hline
\end{prons}

\begin{remprons}
I \hooglig{serve} by \ucomon{jigging} \comon{shields} and \ucomkun{tsukas}.\\\hline
{\warn}The \textit{tsuka} is the hilt or handle of a Japanese sword.\\\hline
\end{remprons}

%{\middel}
% &  ()\\\hline
\begin{somme}
\rye{3}{{\diff}This 漢字 is unusual in that its 音読み almost always forms compounds with the 訓読み of other characters, as in the examples below.}
仕方 & serve + direction = \textit{way; method}  & シ{\middel}かた \mbox{(SHI{\middel}kata)}\\\hline
仕事 & serve + matter = \textit{job; work}  & シ{\middel}ごと \mbox{(SHI{\middel}goto)}\\\hline
仕切り & serve + cut = \textit{partition; dividing line}  & シ{\middel}きり \mbox{(SHI{\middel}kiri)}\\\hline
仕上げる & serve + upper = \textit{finish; complete}  & シ あ{\middel}げる \mbox{(SHI a{\middel}geru)}\\\hline
\end{somme}

\begin{decomp}{4}
仕方&が&無い&よ。\\\hline
シかた&が&ない&よ\\\hline
way&が&non-existent&よ\\\hline
\multicolumn{4}{|r|}{\textit{There's no choice.}}\\\hline
\end{decomp}

\begin{decomp}{6}
彼&は&仕事&の&鬼&だ。\\\hline
かれ&は&シごと&の&おに&だ\\\hline
he&は&work&の&demon&be/is\\\hline
\multicolumn{6}{|r|}{\textit{He is an eager beaver.}}\\\hline
\end{decomp}
\end{multicols}

%%--------------------------------------------------------------------
%23-------------------------------------------------------------------
\begin{multicols}{2}
\begin{glyph}{Heavy (重)}{heavy}\label{glyph:heavy}
Heavy, both in the literal and figurative sense.
\end{glyph}
\gindex{Heavy}{重}\strindex{9}{重}

\begin{comps}
\comprtwocone &  丿 + 里  \\\hline
\comprthreecone & Slash/NO + Hamlet \ref{glyph:hamlet} \\\hline
\end{comps}

\begin{remglyph}
\hooglig{Heavily} \remcom{slash} the \remcom{hamlet}.\\\hline
\end{remglyph}

\begin{prons}
\pronsrtwocone & \textbf{ジュウ (J\={U})} & \textbf{おも (omo)}\\\hline
\rye{3}{ジュウ is by far the more common reading.}
\pronsrthreecone & チョウ (CH\={O}) & かさ、え (kasa, e) \\\hline
\end{prons}

\begin{remprons}
\hooglig{Heavy} \ucomkun{edging} \ucomon{chores} for \ucomkun{casanova} due to his \comon{juice} stains that required \comkun{omo}.\\\hline
\end{remprons}

%{\middel}
% &  ()\\\hline
\begin{somme}
重い  & \textit{heavy}  & おも{\middel}い \mbox{(omo{\middel}i)}  \\\hline
重大  & heavy + large = \textit{important; grave}  & ジュウ{\middel}ダイ \mbox{(J\={U}{\middel}DAI)}  \\\hline
重要  & heavy + essential = \textit{important; essential}  & ジュウ{\middel}ヨウ \mbox{(J\={U}{\middel}Y\={O})}  \\\hline
\end{somme}

\begin{decomp}{5}
彼&は&口&が&重い。\\\hline
かれ&は&くち&が&おもい\\\hline
he&は&mouth&が&heavy\\\hline
\multicolumn{5}{|r|}{\textit{He is a man of few words.}}\\\hline
\end{decomp}

\begin{decomp}{6}
結婚&は&重大&な&問題&だ。\\\hline
ケッコン&は&ジュウダイ&な&モンダイ&だ\\\hline
marriage&は&serious&な&matter&be/is\\\hline
\multicolumn{6}{|r|}{\textit{Marriage is a serious matter.}}\\\hline
\end{decomp}
\end{multicols}

\pagebreak

%%--------------------------------------------------------------------
%24-------------------------------------------------------------------

\begin{multicols}{2}
\begin{glyph}{Light (軽)}{light}\label{glyph:light}
Light, in terms of weight or feeling.
\end{glyph}
\gindex{Light}{軽}\strindex{12}{軽}

\begin{comps}
\comprtwocone&車 + 又 + 土\\\hline
\comprthreecone & Car + Again + Earth\\\hline
\end{comps}

\begin{remglyph}
The \hooglig{light} \remcom{car} is \remcom{again} on \remcom{soil}.\\\hline
\end{remglyph}

\begin{prons}
\pronsrtwocone&\textbf{ケイ (KEI)}&\textbf{かる (karu)}\\\hline
\rye{3}{The 訓読み is voiced when not in the first position, as seen in
	the second and third compounds (the lattter a mix of 音読み and 訓読み).}
\pronsrthreecone&---&かろ (karo)\\\hline
\end{prons}

\begin{remprons}
\hooglig{Light} \comkun{caruncles} were \comon{catering} for the
\ucomkun{carousel}.\\\hline
{\warn}\textit{Caruncle} is a part of the human eye anatomy.\\\hline{\warn}\textit{Carousel} is a merry-go-round.\\\hline
\end{remprons}

%{\middel}
% &  ()\\\hline
\begin{somme}
軽い&\textit{light}&かる{\middel}い \mbox{(karu{\middel}i)}\\\hline
手軽&hand + light = \textit{simple; easy}&て{\middel}がる \mbox{(te{\middel}garu)}\\\hline
気軽&sprit + light = \textit{lighthearted}&キ{\middel}がる \mbox{(KI{\middel}garu)}\\\hline
\notyet{軽食}&light + eat = \textit{light meal; snack}& ケイ{\middel}ショク \mbox{(KEI{\middel}SHOKU)}\\\hline
\end{somme}

\begin{decomp}{7}
軽食&の&出来る&カフェ&が&あります&か。\\\hline
ケイショク&の&できる&カフェ&が&あります&か\\\hline
light snack&の&be able to&cafe&が&to be&か\\\hline
\multicolumn{7}{|r|}{\textit{Is there a cafe where I can have a light meal?}}\\\hline
\end{decomp}
\end{multicols}

%%--------------------------------------------------------------------
%25-------------------------------------------------------------------

\begin{multicols}{2}
\begin{glyph}{Match (合)}{match}\label{glyph:match}
Match.  Combine.\\\hline
This is a 漢字 you will encounter frequently, as it forms many compound verbs (such as in the third example) that suggest some kind of \textit{reciprocal action} between people or things.  These words are often translated into English with the phrase ``each other'': 押し合う (お{\middel}し　あ{\middel}う) \textit{to push each other}, is one example.\\\hline
This 漢字 can also act like 買 in absorbing its ひらがな accompaniment in certain compounds.
\end{glyph}
\gindex{Match}{合}\strindex{6}{合}

\begin{comps}
\comprtwocone & (亼 = 𠆢 + 一) + 口 \\\hline
\comprthreecone & (Assembly = Person + One) + Mouth \\\hline
\end{comps}

\begin{remglyph}
\hooglig{Match} that \remcom{assembly's} \remcom{vampire} speech.\\\hline
\end{remglyph}

\begin{prons}
\pronsrtwocone & \textbf{ゴウ (G\={O})} & \textbf{あ (a)}\\\hline
\rye{3}{ゴウ is the most common choice in the first position, but both it and あ occur frequently in second position.}
\pronsrthreecone & ガッ、カッ (GA-, KA-) & --- \\\hline
\rye{3}{These 2 readings only occur in a few words, and always double up with the 音読み of the 漢字 that follows them (as the small カタカナ ッ indicates).\\\hline
The compound 合宿 (ガッ{\middel}シュク) \textit{to lodge together}, from the 漢字 `match' and `lodging' is an example.}
\end{prons}

\begin{remprons}
The \hooglig{match} \comkun{update} was filled with \ucomon{cuss-worthy} \ucomon{gushing} \comon{gore}.\\\hline
\end{remprons}

%{\middel}
% &  ()\\\hline
\begin{somme}
合う (intr.)  & \textit{match; fit}  & あ{\middel}う \mbox{(a{\middel}u)}  \\\hline
合わせる (tr.)  & \textit{put together; combine}  & あ{\middel}わせる \mbox{(a{\middel}waseru)}  \\\hline
押し合う (tr.) & push + match = \textit{to push each other} & お{\middel}し　あ{\middel}う \mbox{(o{\middel}shi a{\middel}u)} \\\hline
話し合う  & speak + match = \textit{discuss; talk}  & はな{\middel}し　あ{\middel}う \mbox{(hana{\middel}shi a{\middel}u)}  \\\hline
合意  & match + mind = \textit{mutual agreement}  & ゴウ{\middel}イ \mbox{(G\={O}{\middel}I)}  \\\hline
合宿 & match + lodging = \textit{to lodge together} & ガッ{\middel}シュク \mbox{(GAS{\middel}SHUKU)}\\\hline
\end{somme}

\begin{decomp}{3}
合言葉&を&言う。\\\hline
あいことば&を&いう\\\hline
password&を&say\\\hline
\multicolumn{3}{|r|}{\textit{Give the password.}}\\\hline
\end{decomp}

\begin{decomp}{4}
会合&は&先月&あった。\\\hline
カイゴウ&は&センゲツ&あった\\\hline
meeting&は&last month&happened\\\hline
\multicolumn{4}{|r|}{\textit{The meeting was last month.}}\\\hline
\end{decomp}
\end{multicols}

\pagebreak
%%--------------------------------------------------------------------


\end{document}
